% ============================================================
%  R. Nielsen, Evangelietroen og den moderne Bevidsthed
%  Forelæsninger over Jesu Liv.  Første Deel.
%  Kjøbenhavn: C. A. Reitzel, 1849.  (X + 530 pp., Fraktur)
%  Published 19 May 1849. The title page marks it "Første Deel"; no
%  further part ever appeared.
%  TRANSCRIPTION (Danish)
%  Source: KB scan, bibliotek/Nielsen, Rasmus/1849-evangelietroen-modern-bevidsthed.pdf
%          (KB e-mat/dod/11030800268B.pdf; 551 PDF pages, ABBYY layer).
%
%  PAGE MAP:  see pagemap.py. UNIFORM throughout, unlike the 1850 companion
%             volume, whose scan has a duplicated leaf:
%
%               printed   1--530  =  PDF + 17          (PDF 18--547)
%               printed roman V--X (Fortale) = PDF 10--15
%               PDF  6  title page          [I]
%               PDF  8  dedication          [III]
%               PDF 16  Indhold             (unnumbered)
%               PDF 17  Rettelser           (unnumbered)
%               PDF 548--551  blank
%
%             Verified against the ABBYY text layer for all 551 pages: every
%             apparent change of offset is a single-page misreading of the
%             header numeral which the following page reverts (PDF 447 reads
%             "480" for 430; PDF 500 reads "482" for 483). There is no
%             duplicated leaf. PDF 16 only LOOKS like a second copy of printed
%             p. 1 in the text layer, because the Indhold reprints the Afsnit
%             heading verbatim above its list of chapter page ranges.
%
%  OCR:       tesseract with the Fraktur script model reads this scan cleanly;
%             the PDF's own ABBYY layer does not (it had no Fraktur model, and
%             renders æ as "cr", ø as "o", ſ as "j" — "Prcrst" for Præst,
%             "Begyndelje" for Begyndelse). The ABBYY layer is nevertheless kept
%             as a SECOND WITNESS, because it fails differently: where the two
%             agree the reading is safe, and where they differ the image is
%             consulted.
%
%  ERRATA:    The book prints a Rettelser page (PDF 17) with nine corrections.
%             EDITORIAL POLICY FOR THIS EDITION: each correction is APPLIED in
%             the body and the printed reading is logged in a % comment at the
%             site, in the form
%                 % Rettelser: printed „X“; corrected to „Y“
%             The Rettelser page itself is transcribed verbatim in the front
%             matter, so both states of the text are recoverable. This follows
%             the practice of texts/nielsen/evangelietroen-theologien.
%
%             The nine, read off the image at 600 dpi (the OCR of this leaf is
%             unusable — damaged numerals, "12^" for 124, "3t6" for 346):
%
%               Side 97  Linie  8 f. o.  Undergjerningens -> Undergjerningernes
%               Side 124 Linie 11 f. o.  at han tør       -> at han her før
%               Side 165 Linie  3 f. n.  Loø              -> Lov
%               Side 193 Linie 11 f. o.  Himmerige        -> Himmeriges Rige
%               Side 322 Linie 10 f. o.  le peuble        -> le peuple
%               Side 346 Linie  9 f. n.  hvor i henlægge  -> hvor I henlægge
%               Side 359 Linie  1 f. o.  Taabeligheden    -> Taaligheden
%               Side 368 Linie 13 f. o.  fortjene         -> fortjener
%               Side 406 Linie  9 f. o.  fordrage         -> foredrage
%
%             The Rettelser leaf itself prints `l:` without the point on its
%             last line; transcribed as printed.
%
%  INDHOLD vs. BODY: one disagreement is already recorded. The Indhold leaf
%             prints an antiqua „A.“ before the standing head of the second
%             Afsnit — „A. Forløseren.“ — and no „B.“ exists anywhere on the
%             leaf. This volume is marked „Første Deel“ on its title page and no
%             second part ever appeared, so the „A.“ is most likely the head of
%             a division whose sequel was never printed. Both readings stand;
%             they are not normalised.
%
%  DIPLOMATIC ELSEWHERE: apart from the nine Rettelser, printer's errors are
%             transcribed AS PRINTED and logged in a % comment at the site,
%             never silently corrected. Unmatched quotation marks are therefore
%             expected, and the running „ / “ balance is NOT expected to be
%             zero; the current total and the list of logged defects are kept in
%             RESUME-NOTES.md. Do not "fix" an imbalance without reading the
%             comment at its site.
%
%  ORTHOGRAPHY: 19th-century spelling exactly as printed — Christendommen,
%             Videnskaben, Theologie, høiere, aa, capitalised nouns. Danish
%             quotation marks „…“ low-high. Em-dash ---.
%
%  EMPHASIS — THREE DEVICES, not two:
%             (1) LETTERSPACING (Sperrsatz) -> \emph{}. Detected mechanically by
%                 spacing.py and confirmed on the page image.
%             (2) ANTIQUA against the surrounding Fraktur -> \textit{}. Latin,
%                 French, the dialogue speaker letters A./B. of pp. 92--95, and
%                 the „A.“ and „I.“ before the heads on p. 184.
%                 CAUTION: several German proper names are NOT antiqua but a
%                 LIGHTER FRAKTUR — „Strauß“ (p. 119), „Schleiermacher“
%                 (p. 133), „Lessing“ (p. 145). Those stay plain.
%             (3) BOLD FRAKTUR -> \textbf{}, used for certain scripture
%                 quotations and for the refrain „Varslerne komme og svinde…“
%                 (pp. 179 twice, 180, 181, 182). This is WEIGHT, not size, and
%                 must be distinguished from the larger DISPLAY FOUNT used for
%                 gospel passages (set with \large, or \begingroup\large …
%                 \par\endgroup at full measure). The test, at 600 dpi:
%                     body    x-height 37--39   stem 7
%                     bold    x-height 40--43   stem 8--9
%                     \large  x-height 45--55   stem 8
%                 Paragraphs that merely look enlarged are usually loose
%                 justification; measure before marking. Verse set SMALLER than
%                 the body (p. 129, p. 166) is set \small in a verse block.
%             Device (3) was not noticed until p. 63. Pages 1--55 were
%             transcribed before it was known and HAVE NOT BEEN RE-SWEPT for it;
%             see RESUME-NOTES.md, open review items.
%
%  THE FRAKTUR x/r PROBLEM — the single largest source of error in this book.
%             In this fount the x and the r are distinguishable only at 600 dpi
%             and above, by a stroke descending below the baseline on the x.
%             BOTH OCR witnesses confuse them, in both directions. Worse, the
%             compositor really does set a wrong sort here often: „vore“ for
%             voxe (pp. VIII, 79), „strar“ for strax (pp. 56, 86, 110, 118, 120,
%             126), „vorte“ (p. 127), „Grundtert“ for Grundtext (pp. 58, 132),
%             „forverle“ (pp. 73, 137), „Omverlinger“ (p. 84), „Reflerion“
%             (pp. 140, 142), „opvore“ (p. 157) — while a true „strax“ can stand
%             twelve lines below a „strar“ on the same page (p. 120), and a true
%             „voxer“ one line above a „vorer“ (p. 141). Every occurrence has
%             been judged on its own at 1200 dpi and each wrong sort is logged
%             at its site. Do not normalise them.
%
%  SCRIPTURE CITATIONS: the book cites the Gospels constantly, and is not
%             internally consistent between forms (Joh. / Ioh., spacing around
%             the comma). The inconsistency is transcribed, not normalised.
%
%  PAGE JOINTS: fragments are spliced into the markers below WITHOUT blank
%             lines between them, because a blank line is a paragraph break and
%             would invent one at every leaf. A word broken across a leaf is
%             joined with `\-%` (never a bare hyphen, which typesets as
%             "Ufuld- kommenhed" — hyphen AND space, inside the word). This is
%             the fault that had to be repaired retrospectively in both
%             texts/nielsen/evangelietroen-theologien and
%             texts/nielsen/speculative-methode; it is designed out here.
%
%  QUOTE BALANCE: +4 („ exceeds “ by four) as of p. 183, and this is CORRECT.
%             The compositor leaves speeches unclosed and once sets a closer
%             where an opener belongs. Sites logged in the text, all verified on
%             the image at 600--1200 dpi:
%               p.  53  objector's speech opened, never closed
%               p.  86  „hellig — opened, never closed
%               p.  89  „videnskabelig Iagttagelse — the only “ falls after
%                       „Undtagelse.“, where sense does not want it
%               p. 111  „Og du Barnlille — opened, never closed; and a closing
%                       „)“ with no opener
%               p. 148  a high closing “ set where the opening „ belongs, before
%                       „Nei, ingenlunde!“
%               p. 150  „Evangelietroen skyer… — ends at „upartisk.“ bare
%               p. 157  Maria's „Det var et saligt Øieblik… — ends at
%                       „Skikkelse?“ with an em-dash and no “
%               p. 158  „sindige Drøftelse — followed by a comma
%               p. 181  an unmatched closing “ after „alt Folket.“
%             Do not "fix" the balance. Running it back to zero would mean
%             editing the compositor.
%
%  TURNED SORTS: this printing turns n for u and u for n constantly, most
%             densely in the gatherings covering pp. 96--160 (Naturloveu, udeu,
%             Paastaud, uok, Universnm, kuu, Standpnnkt, Cyrenins, Frelsereus,
%             Tiugenes, and some fifty more). Every one is transcribed as
%             printed and logged at its site.
%
%  STATUS:    front matter (title, dedication, Fortale V--X, Indhold, Rettelser)
%             and FØRSTE AFSNIT COMPLETE — printed pp. 1--183, image-verified,
%             no gaps, no duplicate page markers, compiles clean.
%             ANDET AFSNIT (pp. 184--530) NOT YET BEGUN. Its head stands at the
%             top of p. 184, a fresh page, set in this order: bold Fraktur
%             „Andet Afsnit.“ / letterspaced „Den hellige Historie:“ / display
%             Fraktur „Jesu Liv i Verden.“ / the epigraph „Medens jeg er i
%             Verden, er jeg Verdens Lys (Joh. 9, 5).“ in smaller type / a short
%             centred rule / antiqua „A.“ / letterspaced „Forløseren.“ / a short
%             centred rule / antiqua „I.“ / bold Fraktur „Forgjængeren.“ / the
%             body, indented.
% ============================================================
\documentclass[12pt,a4paper]{book}
\usepackage[utf8]{inputenc}
\usepackage[T1]{fontenc}
\usepackage{libertinus}
\usepackage{libertinust1math}
\usepackage{textalpha}   % Greek typed directly
\usepackage{graphicx}    % for \reflectbox (the old Danish »ɔ:« id-est mark)
\DeclareUnicodeCharacter{0254}{\reflectbox{c}}  % ɔ (U+0254) = reversed c
\usepackage[danish]{babel}
\usepackage[protrusion=true,expansion=true]{microtype}
\usepackage{setspace}
\usepackage[margin=1.2in]{geometry}
\usepackage{enumitem}
\usepackage{fancyhdr}
\setlength{\headheight}{28pt}
\addtolength{\topmargin}{-13.5pt}

% Footnotes as the 1849 printing sets them: *) rather than 1, 2, 3.
% A constant \thefootnote rather than \fnsymbol with footmisc[perpage], both of
% which were tried on the 1850 volume and are wrong:
%   - \fnsymbol is math mode, so it sets U+2217 ASTERISK OPERATOR, not a text
%     asterisk;
%   - [perpage] resets on the TYPESET page, not the book's printed page, so a
%     second note landing on the same typeset page came out as a dagger.
% If a printed page of this book is found to carry two notes with distinct
% marks, log it and revisit this line — do not silently renumber.
\renewcommand{\thefootnote}{*)}

\usepackage{hyperref}

\hypersetup{
  pdftitle={Evangelietroen og den moderne Bevidsthed},
  pdfauthor={R. Nielsen},
  hidelinks
}

\pagestyle{fancy}
\fancyhf{}
\fancyhead[LE,RO]{\thepage}
\fancyhead[RE]{\textit{Evangelietroen og den moderne Bevidsthed}}
\fancyhead[LO]{\textit{\leftmark}}

\setstretch{1.3}

\title{\textbf{Evangelietroen og den moderne Bevidsthed}\\[6pt]
  \large Forelæsninger over Jesu Liv\\[4pt]
  \large Første Deel}
\author{R.\ Nielsen}
\date{Kjøbenhavn: C.\ A.\ Reitzels Forlag, 1849}

\begin{document}
\maketitle
\thispagestyle{empty}
\newpage

\tableofcontents
\newpage

%% ============================================================
%%  FRONT MATTER
%%  title page [I] = PDF 6; dedication [III] = PDF 8;
%%  Fortale V--X = PDF 10--15; Indhold = PDF 16; Rettelser = PDF 17.
%%  Transcribed from the images as one fragment: .parts/front.texfrag
%% ============================================================
%% ------------------------------------------------------------
%%  FRONT MATTER, transcribed from the page images.
%%    PDF  6  title page                       [I]
%%    PDF  7  blank
%%    PDF  8  dedication                       [III]
%%    PDF  9  blank
%%    PDF 10--15  Fortale, V--X (V is unnumbered; VI--X carry numerals)
%%    PDF 16  Indhold    (unnumbered leaf)
%%    PDF 17  Rettelser  (unnumbered leaf)
%%  Roman page markers are used throughout this fragment. check.py counts
%%  only arabic `% --- p. N ---` markers, so these are correctly ignored.
%% ------------------------------------------------------------

%% The document's own \maketitle above gives a modern title page; the block
%% below is the DIPLOMATIC record of the 1849 one, with the printed lineation
%% reproduced by \\ and the printed rules by \rule. The numerals [I] and [III]
%% are supplied in square brackets: they are not printed on the leaf.

% --- p. [I] ---   (PDF 6)
% "Lic. theol." is set in a bold ANTIQUA against the surrounding Fraktur
% (hence \textit{}); "R. Nielsen" beside it is Fraktur. The standing line
% "Professor i Philosophien." is very lightly inked in this copy but legible.
\begin{center}
  {\LARGE Evangelietroen}\\[10pt]
  og\\[10pt]
  {\LARGE den moderne Bevidsthed.}\\[12pt]
  \rule{0.28\textwidth}{0.4pt}\\[14pt]
  {\large Forelæsninger over Jesu Liv}\\[10pt]
  af\\[10pt]
  {\large \textit{Lic.\ theol.}\ R.\ Nielsen,}\\[2pt]
  {\footnotesize Professor i Philosophien.}\\[12pt]
  \rule{0.28\textwidth}{0.4pt}\\[16pt]
  {\large Første Deel.}\\[16pt]
  \rule{0.62\textwidth}{0.4pt}\\[16pt]
  {\large Kjøbenhavn.}\\[4pt]
  I Commission hos Universitetsboghandler C.\ A.\ Reitzel.\\[4pt]
  Trykt hos E.\ C.\ Løser.\\[6pt]
  {\large 1849.}
\end{center}

\newpage

% --- p. [III] ---   (PDF 8)
% The queen's name is set in a heavily ornamented Fraktur display face. The
% final letter of the last word is the same sort as the "a" of "Amalia" and of
% "Caroline", not the round "e" that ends "Caroline": read AMALIA, not Amalie.
% "denne Bog" and "allerunderdanigst" are letterspaced (hence \emph{}); the
% word "om" between "denne Bog" and "Evangelietroen" is set in a very small
% and very lightly inked type, but is legible at 900 dpi.
% The last three lines are set toward the RIGHT of the measure in the original;
% they are centred here.
\begin{center}
  {\large Deres Majestæt}\\[18pt]
  {\LARGE Enkedronning Caroline Amalia,}\\[20pt]
  hvis Trøst er en Tro,\\
  som troer Evangeliet,\\[20pt]
  helliges\\[14pt]
  \emph{denne Bog}\\[8pt]
  {\footnotesize om}\\[10pt]
  {\large Evangelietroen}\\[22pt]
  \emph{allerunderdanigst}\\[8pt]
  af\\[8pt]
  {\large Forfatteren.}
\end{center}

\newpage

%% ============================================================
%%  FORTALE  (printed pp. V--X = PDF 10--15)
%% ============================================================
\chapter*{Fortale}
\addcontentsline{toc}{chapter}{Fortale}
\label{ch:fortale}

% --- p. V ---
% The heading "Fortale." is letterspaced display type over a short rule; the
% first word "Om" opens with an enlarged initial O.
Om og hvorvidt den raisonnerende Bibelkritik kan beslutte sig til at erklære
vore fire kanoniske Evangelier for ægte eller uægte, overeensstemmende eller
uovereensstemmende, paalidelige eller upaalidelige, er et Skolespørgsmaal,
hvis Besvarelse vel allerhelst maa overlades de Lærde.

En Kjendsgjerning er det imidlertid, at der nu een Gang gives saadanne fire
Evangelier, i hvilke Jesus af Nazaret er os fremstillet som Guds og
Menneskens Søn, som Messias og Verdensfrelser.

At Forfatterne af disse Evangelier selv maae have troet paa den af dem
beskrevne Messias, er der vel heller ingen Grund til at betvivle; men
Spørgsmaalet er, hvorledes det nu lader sig tænke, at noget Menneske, der ret
er blevet deelagtigt i Culturens Goder og fortroligt med Nutidens høiere
Dannelse, fremdeles vil kunne agte det for det Høieste, at dele Tro med de
fire Evangelister.
% --- p. VI ---
Den Opgave, til hvis Løsning nærværende Skrift formeentligen yder nogle
Bidrag, kan da kortelig fremsættes i det Livsspørgsmaal: \emph{Hvad er det,
at troe paa Evangeliernes Christus?}

% The three answers are letterspaced entire, "Det er, at" included; the
% intervening dashes are not. "Livsspørgsmaal:" is NOT letterspaced.
En tredobbelt Besvarelse er mulig: \emph{Det er, at tilegne sig den
aabenbarede Sandhed}; --- \emph{Det er, at lade sig bedaare af en from
Indbildning}; --- \emph{Det er, at trøste sig ved en besynderlig Blanding af
Sandhed og Indbildning.}

Det første Svar egner sig for --- \emph{Evangelietroen}; det andet for ---
\emph{den moderne Bevidsthed}; det tredie viser hen til en mæglende
Anskuelse, der helst vil holde Middelveien.

Evangelietroen lever i --- \emph{den Troende}; den moderne Bevidsthed i ---
\emph{den frie Tænker, Fritænkeren}; Middelveislæren har anseelige Tilhængere
iblandt Skolens Theologer.

Den Hovedtanke, at lade Ordførerne for de her betegnede Anskuelser forene sig
om at gjennemgaae de fire Evangelier sammen og udtale sig hver i sin
Charakteer, blev allerede lagt til Grund for de Forelæsninger over Jesu Liv,
Forfatteren holdt i Universitetsbygningen, i Vinteren 1848, for et hæder\-%
% --- p. VII ---
ligt Antal Tilhørere af begge Kjøn; men en udførligere Fremstilling lod sig
kun meddele skriftlig.

I en mere udviklet Form anbefales da disse Forelæsninger til Læseverdenens
Opmærksomhed med det, maaskee vel dristige, Ønske, at det skrevne Ord nu
maatte finde samme velvillige Modtagelse, som forhen blev det mundtlige til
Deel.

For dog ikke at opvække nogen utidig Forventning eller foranledige nogen for
Referenter tidsspildende Feiltagelse, vil Forfatteren herved tillade sig at
gjøre det større Publicum opmærksomt paa, at den ivrige Strid, der i disse
Forhandlinger føres for og imod Evangelisternes Troværdighed, ikke just
egentlig faaer et saadant Udfald, at man med Bestemthed skal kunne sige, om
det nu er „Fritænkeren“ eller „den Troende“, der tilsidst gaaer af med
Seieren.

Derimod vil den opmærksomme Læser, der interesserer sig for
Grundspørgsmaalet i den Grad, at han foreløbigen tager til Takke med
Fremstillingen, og, om end ikke overseer, saa dog undskylder dens Mangler,
let overbevise sig om, at Striden ikke er uden Resultat. Thi ogsaa dette er
jo et Resultat, at de modsatte Anskuelser klare sig ud, at det sees, hvad der
boer i de Stridende. Striden er da som en Pro\-%
% --- p. VIII ---
ces; Processen gaaer sin Gang; Sagen indlades til Dom; Læserens
Samvittighed er den Jury, der afsiger Dommen.

For dog imidlertid heller ikke at lægge Dolgsmaal paa sin egen personlige
Overbeviisning og derved forvandle det Hele til et Slags philosophisk
Tanke-Experiment, har Forfatteren i flere Partier deels udtalt sig ganske
ligefrem, deels ogsaa taget Ordet paa de Stridendes Vegne; men, da det mere
og mere viste sig for ham, at den ene af Parterne umulig kunde gjendrive den
anden, har heller ikke han turdet indlade sig paa nogen endelig Gjendrivelse.

% Printer's defect: a small stray mark (a broken quad or a fragment of type)
% stands in the indent immediately before "Evangelietroen"; it is not a
% quotation mark and is not transcribed.
Evangelietroen kan ikke gjendrives; den har sit eget, overnaturlige Princip:
den Troende holder ud indtil Enden.

% Printer's defect: "vore" as printed (v-o-r-e, confirmed at 600 dpi), where
% the sense wants "voxe". Transcribed as printed; not in the Rettelser list.
Den moderne Bevidsthed kan ikke gjendrives; ogsaa den har sit eget Princip,
og frygter ikke dets Conseqvenser: den frie Tænker mattes ikke, hans Kræfter
vore under Stridens Gang.

Den eneste af de ideelle Magter, paa hvis Gjendrivelse her altsaa kunde
tænkes, maatte da være Middelveislæren, og dette saameget mere, som den
vistnok her befinder sig i en temmelig kneben Stilling, midt imellem den
gamle eenfoldige Evangelietro og den
% --- p. IX ---
moderne Videns negative Bevidsthed. Men ogsaa dette er umuligt:
Middelveislæren kan virkelig ikke gjendrives. Vel er den ikke meget stiv i
Principerne, men saa er den til Gjengjeld desto mere alsidig; vel har den
ikke Kraftens umiddelbare Overlegenhed, men saa har den desto mere Seighed;
vel kan den ikke opflamme med nogen synderlig Begeistring, men saa kan den
virke med desto større Sindighed og taktisk Beregning.

% "Favete lingvis" is set in a bold ANTIQUA against the Fraktur, and is
% spelled with v, not u, as printed.
Hvorledes skulde Middelveislæren vel kunne gjendrives, naar man betænker, at
ikke blot Maadeholdenhed er en Cardinaldyd, men Middelmaadigheden selv en
Magt i Verden? Hvorledes skulde Middelveislæren kunne gjendrives, naar man
betænker, at den jo kan beraabe sig deels paa det Gamle og deels paa det Nye,
deels paa Troen og deels paa Viden, deels paa Kritiken og deels paa
Dogmatiken? Hvorledes skulde Middelveislæren kunne gjendrives, naar man
betænker, at Den, der overtog sig et saa uhyre Arbeide, maaskee tilsidst
vilde komme til at gjendrive saa godt som hele den nuværende Skoletheologie?
--- \textit{Favete lingvis}: ingen Formastelighed!

Men, endskjøndt Forfatteren veed med sig selv, at han hverken kan eller vil
gjendrive Skoletheolo\-%
% --- p. X ---
gien, er han naturligviis dermed ingenlunde sikkret for, at Skoletheologien,
hvis den en Gang fandt det Umagen værd, ikke skulde kunne gjendrive ham.
Imidlertid er det dog ikke sagt, at den just vil, om den end sagtens kunde.

Dog, hvad enten nu den i dette Skrift fremsatte Grundtanke skulde foranledige
en stille Misbilligelse, eller en lydelig Gjendrivelse, eller begge Dele,
eller --- slet Intet; vil Undertegnede trøste sig med, at han ærlig og
redelig har arbeidet for Besvarelsen af sit Livsspørgsmaal, uden i nogen
Maade at ville krænke enten den dybere Lærdom eller den høiere Videnskab.

\medskip
\noindent\hspace*{2em}{\small Kjøbenhavn, den 1ste Mai 1849.}

\medskip
\noindent\hfill\textbf{R.\ Nielsen.}\hspace*{3em}

%% ============================================================
%%  INDHOLD  (unnumbered leaf = PDF 16)
%%  Transcribed from the image, not from the OCR layer.
%%  NOTE FOR THE BODY TRANSCRIBERS: the Andet Afsnit's standing head is
%%  printed here as "A. Forløseren." — with an antiqua "A." before the
%%  Fraktur word. The skeleton comment further down calls it merely
%%  „Forløseren“; the leaf reads A. Forløseren. No "B." appears anywhere on
%%  the leaf, which is consistent with the book being Første Deel only.
%%  The printed page ranges use an en-dash; set here as --.
%% ============================================================
\chapter*{Indhold}
\addcontentsline{toc}{chapter}{Indhold}
\label{ch:indhold}

\begin{center}
  {\large Første Afsnit.}\\[6pt]
  \emph{Det Nye Testamentes For-Evangelium:}\\[4pt]
  \textbf{\large Jesu Komme til Verden.}\\[6pt]
  {\footnotesize Jeg udgik fra Faderen og kom til Verden (Joh. 16, 28).}
\end{center}

\begingroup
\setlength{\parindent}{0pt}
\makebox[2.6em][l]{I.}Engelen Gabriel og Præsten Zacharias\dotfill Side 1--42\par
\makebox[2.6em][l]{II.}Engelen Gabriel og Jomfru Maria\dotfill --- 42--80\par
\makebox[2.6em][l]{III.}Den hellige Familie\dotfill --- 81--183\par
\hspace*{2.6em}\makebox[1.8em][l]{1.}Marias Besøg hos Elisabeth --- Johannes\\
\hspace*{4.4em}Døberens Fødsel\dotfill ---100--128\par
\hspace*{2.6em}\makebox[1.8em][l]{2.}Jesu Fødsel --- Hyrderne i Bethlehem\dotfill ---128--166\par
\hspace*{2.6em}\makebox[1.8em][l]{3.}De evangeliske Forvarsler\dotfill ---166--183\par
\endgroup

\begin{center}
  {\large Andet Afsnit.}\\[6pt]
  \emph{Den hellige Historie:}\\[4pt]
  \textbf{\large Jesu Liv i Verden.}\\[6pt]
  {\footnotesize Medens jeg er i Verden, er jeg Verdens Lys (Joh. 9, 5).}\\[8pt]
  \textit{A.}\ Forløseren.
\end{center}

\begingroup
\setlength{\parindent}{0pt}
\makebox[2.6em][l]{I.}Forgjængeren\dotfill Side 184--201\par
\makebox[2.6em][l]{II.}Christus iblandt sine Egne\dotfill --- 201--227\par
\makebox[2.6em][l]{III.}Christus er et Tegn, som imodsiges\dotfill --- 228--289\par
\makebox[2.6em][l]{IV.}Christus aabenbarer Hjerternes Tanker\dotfill --- 289--334\par
\makebox[2.6em][l]{V.}Christus og Mørkets Aand\dotfill --- 334--393\par
\makebox[2.6em][l]{VI.}Christus er mægtig i Ord og Gjerning\dotfill --- 393--458\par
\makebox[2.6em][l]{VII.}Jesus Christus er Verdens Frelser\dotfill --- 458--530\par
\endgroup

% The leaf closes with a typographic flourish (an engraved ornament), not
% reproduced here.

%% ============================================================
%%  RETTELSER  (unnumbered leaf = PDF 17)
%%  NINE corrections, read off the image at 900 dpi; every numeral confirmed.
%%  The head is "Rettelser." over the list; a short rule closes the list.
%%  The first line spells out "Side ... Linie ... fra oven: ... læs:";
%%  the rest abbreviate to "--- N --- N f. o.:" (or "f. n." / "f. neden")
%%  and "l.:". The LAST line prints "l:" without the point — as printed.
%%  EDITORIAL POLICY (see the file header): these nine are APPLIED in the
%%  body and the printed reading logged in a % comment at each site.
%% ============================================================
\chapter*{Rettelser}
\addcontentsline{toc}{chapter}{Rettelser}
\label{ch:rettelser}

\begingroup
\setlength{\parindent}{0pt}
\begin{tabular}{@{}l@{~}r@{~~}l@{~}r@{~~}l@{}}
Side & 97  & Linie & 8  & fra oven: Undergjerningens læs: Undergjerningernes\\
---  & 124 & ---   & 11 & f. o.: at han tør l.: at han her før\\
---  & 165 & ---   & 3  & f. neden: Loø l.: Lov\\
---  & 193 & ---   & 11 & f. o.: Himmerige l.: Himmeriges Rige\\
---  & 322 & ---   & 10 & f. o.: \textit{le peuble} l.: \textit{le peuple}\\
---  & 346 & ---   & 9  & f. n.: hvor i henlægge l.: hvor I henlægge\\
---  & 359 & ---   & 1  & f. o.: Taabeligheden l.: Taaligheden\\
---  & 368 & ---   & 13 & f. o.: fortjene l.: fortjener\\
---  & 406 & ---   & 9  & f. o.: fordrage l: foredrage\\
\end{tabular}
\endgroup

%% ============================================================
%%  FØRSTE AFSNIT  (printed pp. 1--183 = PDF 18--200)
%%
%%  The Afsnit head is printed at the top of p. 1 and is set here in the
%%  skeleton; the batch agent for pp. 1--14 must NOT repeat it, and begins its
%%  fragment with the epigraph line as printed.
%%
%%  CHAPTER HEADS ARE NOT IN THE SKELETON. Every chapter of this book except
%%  the first begins PART-WAY DOWN a page (the Indhold's ranges overlap: I
%%  1--42, II 42--80), so the head must be written by whichever batch owns the
%%  page it falls on, exactly where it falls, as
%%      \section*{II. Engelen Gabriel og Jomfru Maria}
%%      \addcontentsline{toc}{section}{II. Engelen Gabriel og Jomfru Maria}
%%  with the wording taken from the PRINTED HEAD, and any disagreement with the
%%  Indhold recorded rather than normalised.
%%
%%  Chapter ranges per the Indhold (PDF 16):
%%    I.   Engelen Gabriel og Præsten Zacharias .............  1--42
%%    II.  Engelen Gabriel og Jomfru Maria ..................  42--80
%%    III. Den hellige Familie .............................. 81--183
%%       1. Marias Besøg hos Elisabeth --- Johannes
%%          Døberens Fødsel ................................ 100--128
%%       2. Jesu Fødsel --- Hyrderne i Bethlehem ............ 128--166
%%       3. De evangeliske Forvarsler ...................... 166--183
%% ============================================================
\chapter*{Første Afsnit\\[4pt]
  \large Det Nye Testamentes For-Evangelium: Jesu Komme til Verden}
\addcontentsline{toc}{chapter}{Første Afsnit. Det Nye Testamentes For-Evangelium: Jesu Komme til Verden}
\markboth{Første Afsnit}{Første Afsnit}
\label{ch:afsnit1}

% --- p. 1 ---
% The Afsnit head is already set in the skeleton above. Below it p. 1 prints,
% in this order: the epigraph, a short centred rule, the chapter number "I." on
% its own line, the chapter head, and a second short rule.
% The epigraph is printed WITHOUT quotation marks (checked at 450 dpi); it
% agrees word for word with the Indhold, including "(Joh. 16, 28)".
\begin{center}
  {\footnotesize Jeg udgik fra Faderen og kom til Verden (Joh.\ 16, 28).}\\[10pt]
  \rule{0.22\textwidth}{0.4pt}
\end{center}

% The printed head sets the numeral on a line of its own above the title, and
% ends the title with a full stop. The full stop is kept in the \section*, and
% dropped in the \addcontentsline so that the running TOC matches the Indhold
% ("I. Engelen Gabriel og Præsten Zacharias ... Side 1--42"). Later batches
% setting chapter heads should follow the same rule.
\section*{I.\\[4pt] Engelen Gabriel og Præsten Zacharias.}
\addcontentsline{toc}{section}{I. Engelen Gabriel og Præsten Zacharias}

\begin{center}
  \rule{0.14\textwidth}{0.4pt}
\end{center}

% The two scripture passages on pp. 1--2 (Malach. 4, 5--6 and the Lukas
% pericope) are set at FULL MEASURE in a type one size LARGER than the body,
% not indented and not in quotation marks; Nielsen's connecting sentence
% between them is in the ordinary body size. Reproduced here with \large.
\begingroup\large
See, siger Jehova til Israels Folk, jeg sender Eder Elias
Propheten, forend den store og forfærdelige Herrens Dag
kommer. Og han skal omvende Fædrenes Hjerte til Børnene, og
Børnenes Hjerte til deres Fædre, at jeg ikke skal komme og slaae
Jorden med Band.
\par\endgroup

Saaledes lyde de sidste Ord af den sidste Prophet i det gamle
Testamente (Malach.\ 4, 5--6), og med Opfyldelsen af denne Spaadom
begynder Evangeliet i det nye Testamente. (Luk.\ 1, 5--25).

\begingroup\large
I Kong Herodes's Dage levede i Judæa en Præst af Abiæ Skifte,
ved Navn Zacharias, hvis Hustru Elisa\-%
% --- p. 2 ---
beth var af Aarons Døttre. Begge vare retfærdige for Gud, og
vandrede ustraffelige i alle Herrens Bud og Rette. Men de havde
intet Barn; thi Elisabeth var ufrugtsommelig, og de vare begge
bedagede. Men det begav sig, der Zacharias forrettede
Præste-Embedet i sin Skiftes Orden for Gud, og traadte ind i det
Hellige, for at antænde Røgelse-Offeret, at Engelen Gabriel
aabenbarede sig for ham. Og han saae Engelen staae ved den høire
Side af Røgelsens Alter, og han forfærdedes, og Frygt faldt over
ham. Men Engelen sagde til ham: frygt ikke Zacharias! thi din
Bøn er hørt, og din Hustru Elisabeth skal føde dig en Søn, og du
skal kalde hans Navn Johannes. Han skal gaae frem i Elias's Aand
og Kraft, at omvende Fædrenes Hjerter til Børnene og de Ulydige
til de Retfærdiges Forstand, at beskikke Herren et vel forberedt
Folk. Og Zacharias sagde til Engelen: hvorpaa skal jeg kjende
dette? thi jeg er gammel, og min Hustru er ogsaa bedaget. Og
Engelen svarede og sagde til ham: jeg er Gabriel, som staaer for
Gud, og jeg er udsendt at tale til dig, og at forkynde dig dette
til Glæde. --- Men efterdi Zacharias tvivlede, gjorde Engelen ham
stum, og gav ham i denne Straf et Tegn, at Forjættelsen skulde
vorde opfyldt. Og hele det forsamlede Folk, som var i Forgaarden
og bad i Røgelsens Time, ventede paa Zacharias, og undrede sig,
hvi han tøvede saalænge, og der han endelig traadte ud, mærkede
Forsamlingen, at han havde havt et Syn i Helligdommen. Men, da
hans Tjenestes Dage vare til Ende, gik Zacharias hjem i sit Huus,
og det blev opfyldt, hvad Engelen havde forkyndt ham.
\par\endgroup

Dette er i en Hovedsum Indholdet af den Fortælling, hvor\-%
% --- p. 3 ---
med Evangeliet om Verdens Frelser indledes. Forberedelsen ender,
at Fuldendelsen kan komme; Stykværket ophører, at det Fuldkomne
kan tage sin Begyndelse. En overordentlig, en vidunderlig
Begyndelse! Det Overordentlige viser sig ikke blot deri, at den
fromme Præst faaer et Varsel om, at Tidens Fylde er nær; men det
har en afgjørende Betydning, at dette Varsel kommer fra Himmelen,
at det er Guds Engel, der bringer ham det, at det er Prophetens
Spaadom, der lyder og gjenlyder i det glade Budskab. Det
Vidunderlige vil ikke lade sig indskrænke dertil, at Bebudelsen
virkelig gaaer i Opfyldelse; men det læres os udtrykkelig, at
denne Opfyldelse selv er et Under, en guddommelig Gjerning,
hvorved Forløsningens Værk forberedes. Evangelisten forkynder os,
at det er Israels Gud, der nu igjen besøger sit Folk, Propheternes
Gud, der opfylder sine Forjættelser, Frelsens Gud, der ihukommer
sin hellige Pagt efter den Ed, han svor Abraham.

Turde vi nu blot antage, at det i Virkeligheden er gaaet til med
Spaadommens Opfyldelse saaledes, som Fortællingen lyder, kunde vi
blot holde os forvissede om, at denne Israels Gud er den eneste
sande Gud, at denne Propheternes Gud dog virkelig er istand til at
opfylde sine Forjættelser, og at denne Frelsens Gud i Tidernes
Glemsomhed ikke har forglemt sig selv og den Ed, han svor Abraham;
ja var det os blot forundt --- at troe: da kunde vi jo endnu søge
Tilflugt i Aabenbaringens Helligdom, endnu, som i Evangeliets
Begyndelse, høre Engelens Røst: „jeg er udsendt at tale til dig,
og forkynde dig dette til Glæde.“

„Men,“ sige Nutidens Oplyste, „alt Dette har jo slet ingen sand
Virkelighed. Bebudelseshistorien opløser sig i et reent
Blændværk, naar Fortællingen sees igjennem Kritikens Brille.“
Skulde de maaskee have Ret? Skulde den kritiske Videnskab nu
virkelig være kommen saavidt, at den uden videre kan begynde sin
Opløs\-%
% --- p. 4 ---
ning af den evangeliske Historie med den Paastand, at Aabenbaring
og Undere høre blandt de umulige Ting? En saadan Paastand kunde
vist høilig trænge til et for eenfoldige Mennesker fatteligt
Beviis. Men et saadant Beviis kan dog vel ikke føres, uden at
Evangelietroen frit og uhindret først kommer til Orde, og det ikke
afsides for sig som en Partisag, men i Forhandlingernes egen Gang
som en Part i Sagen, ifølge hiin gyldne Regel: \textit{audiatur et
altera pars.} At angribe den foreliggende Fortælling med en
ubeviist Paastand, er jo ikke blot usømmeligt, uretfærdigt og
uvidenskabeligt, men synes endog at være saa ukritisk, at en
særlig Kritik af Fortællingen selv derved bliver aldeles
overflødig. Hvad der kan bringes ud af den eenfoldige Historie om
Zacharias og Elisabeth, naar Engelen vises bort, om den da maaskee
kan lade sig forklare som en overtroisk Forvanskning af en ganske
naturlig Begivenhed, eller endnu bedre omforklare til en Mythe,
beskjæftiger os aldeles ikke. Gaaer Engle-Aabenbarelsen tabt, saa
er Historien tabt for os. Den kritiske Undersøgelse dreier sig
derfor udelukkende om det Hovedspørgsmaal, hvorvidt den
foreliggende Beretning, saavel i sin Heelhed som i sine enkelte
Træk, er i fuldkommen Overeensstemmelse med Aabenbaringens
Grundcharakteer, eller ikke. Kun naar Opgaven stilles saaledes,
kan der indledes en videnskabelig Forhandling imellem Tvivlen og
Troen.

Lad den kritiske Tvivler saa kun fremkomme med sine Indsigelser;
disse Indsigelser ville netop vinde i videnskabelig Vægtfylde, jo
mere sympathetisk Tvivleren forøvrigt kan forholde sig til, hvad
han kalder Christendommens Grundidee.

„Ad den blot rationelle Vei,“ siger Tvivleren, „er det vistnok
umuligt, at opnaae nogen afgjørende religiøs Overbeviisning. Men
falder det nu paa den anden Side allerede haardt nok at skulle
troe paa Underets Realitet i Almindelighed, saa bliver Van\-%
% --- p. 5 ---
skeligheden endnu langt føleligere, naar man skal beslutte sig til
at antage dette eller hiint enkelte, bestemte Under i Form af
historisk Begivenhed. Hver Gang jeg tager det N.\ T. for mig, og
bereder mig til at gjennemgaae den evangeliske Historie,
indprænter jeg mig atter og atter den Grundsætning, at der jo dog,
saa sandt der gives en virkelig Aabenbaring, maa begyndes fra et
bestemt, overnaturligt Udgangspunkt, at altsaa guddommelige
Meddelelser, overordentlige Phænomener o.\ dsl. her maae være paa
rette Sted. Men, naar jeg saa skal til at gjøre mig fortroelig
med den allerførste Beretning hos Lukas, ja saa --- staaer
Forstanden stille. Jo mere jeg anstrenger mig for en Tilegnelse,
desto haardere bliver jeg tilbagestødt; jo længere jeg stirrer paa
Billedet, desto klarere indseer jeg, at Evangelisten, der forklarer
Alt for Theophilus, Intet forklarer for mig. For det Første
finder jeg nu ingen tilstrækkelig Grund til, at Døberens Fødsel
just skulde udkræve en saa overordentlig Foranstaltning af
Himmelen selv. Vel bør det indrømmes, at Johannes den Døber, som
Bodsprædikant, som Christi Forgjænger, som prophetisk Midler
imellem den gamle og den nye Pagt, maa have været en Mand af
udmærkede Aandsgaver. Men for at udrustes med guddommelig
Aandskraft, er det da aldeles nødvendigt, at Ens Fødsel betinges
ved et Mirakel? Ifølge denne Antagelse maatte jo baade Abraham og
Moses og alle Propheterne, ja Apostlene selv kunne godtgjøre, at
deres Fødsel skyldtes en Undergjerning. Men, om man nu ogsaa
kunde oversee denne Vanskelighed, hvorfor skulde da Johannes's
Fødsel ved et saa overordentligt Sendebud forud forkyndes
Forældrene? Hvad berettiger Zacharias og Elisabeth til en saa
extraordinær Begunstigelse? Denne Mislighed falder endnu stærkere
i Øinene, naar vi betragte den her beskrevne Engle-Aabenbarelse
lidt nærmere. Kun naar Menneske-Aanden arbeider i den allerhøieste
% --- p. 6 ---
Spænding, kan der tænkes paa overnaturlige Syner: og kun, naar
denne Spænding er bevirket ved en stor, afgjørende Verdenskrise,
kun naar Forviklingernes Knude er strammet saa haardt, at Ingen
uden Gud selv kan løse den, kun da tør vi vel antage
Verdensstyrerens umiddelbare Indgriben i Tingenes Gang for meer
end en from Illusion. Da Jerusalem beleiredes af Romerne, da
Jødefolket, bragt til det Yderste af Fortvivlelse, Raseri og
religiøs Fanatisme, kjæmpede den sidste Kamp; da Forhaanelsen mod
Jehovas Helligdom var saa oprørende, at Fædrene maatte vende sig i
deres Grave, og Elendigheden saa grændseløs, at en Moder fortærede
sit eget Barn, da syntes just et saadant storartet, Himlen
udfordrende Vendepunkt at være indtraadt. Havde Evangelisten nu i
vor Fortælling skildret et lignende Optrin, en Offerscene i det af
Luerne allerede antændte Tempel, ja da kunde jeg forstaae, at en
gammel Præst saae Syner i Helligdommen, og Folket i Forgaarden
troede paa Aabenbarelsen. Men denne Aabenbarelse mangler
Øieblikkets Pointe, Ideens uvilkaarlige Indskydelse. Dette Møde
imellem Zacharias og Engelen er holdt i en saa tør, man kunde
næsten sige, i en saa borgerlig Stiil, at den mirakuløse
Foranstaltning næsten tager sig ud, som en lille privat
Overraskelse midt i Hverdagslivet. Præsten er jo rigtignok
indtraadt i det Hellige; han bliver ogsaa noget forskrækket, det
første han seer Synet; men han fatter sig dog strax, saasnart
Engelen begynder at tale. Og hvad er saa det for et Indtryk denne
Englehilsen med de derpaa følgende tilforladelige Erklæringer
frembringer? „Vær ikke bange, Zacharias! jeg er hverken en ond
Aand, ei heller en Gjenganger; mit Navn er Gabriel. Jeg bringer
fra Himlen det Budskab, at Gud Jehova nu endelig har hørt din Bøn,
og besluttet at skjænke dig en Søn, der skal hedde Johannes, og
virke som Propheten Elias.“ Og nu den Koldblodighed, hvormed
Zacharias,
% --- p. 7 ---
efterat den første Forskrækkelse har sat sig, paa egne og Hustrues
Vegne betænker Sagens Vanskelighed, og for en Sikkerheds Skyld
begjærer et nyt Tegn, som om Engelens udtrykkelige Erklæring ikke
var ham nok. Hellige Gud, skulde Sligt være muligt? Skulde der
udenfor Indbildningens Enemærker nogetsteds kunne findes et
saadant Samqvem imellem de Himmelske og de Jordiske, et saa
fortroeligt, o hvad skal jeg kalde det, et saa kammeratligt
Forhold imellem Evighedens Majestæt og det endelige Menneskes
Ubetydelighed? Nei, Nei! det er umuligt! jeg kan, jeg tør ikke
troe denne Fortælling; thi jeg kan og tør ikke trøste mig selv ved
et Selvbedrag. Altsaa er det da denne lille Historie, der hver
Gang skal reise sig imod mig, og stille sig som en Slagbom imellem
mig og Evangeliet; thi jeg veed det, idet jeg forkaster den,
forkaster jeg jo --- „Guds Aabenbaring“. Saa er det da
Bebudelsen af dette Barns Fødsel, der skal udelukke mig fra
Barnlighedens Paradiis; thi jeg veed det, idet jeg bryder med dens
Engel, har jeg med det Samme brudt Staven over den hellige
Historie.“

Det er forgjæves at ville imødegaae disse og lignende Indvendinger
med almeenvidenskabelige Grunde, eftersom her jo ikke handles om
en relativ Forskjel imellem en høiere og lavere Grad af Viden; men
om den absolute Forskjel imellem Tro og Vantro. Det er kun et
moderne Hykleri, naar Mange nuomstunder lade, som om de gjerne
vilde troe paa Evangeliets Undere, hvis deres høiere Oplysning,
deres friere Dannelse, deres rene Videnskabelighed bare vilde
tillade dem det. Den, der ikke veed med sig selv, at
Modsætningen imellem Tro og Ikke-Tro er saa langt fra at være
afhængig af nogensomhelst Videnskab, at denne Modsætning meget
mere selv begrunder, gjennemtrænger og behersker vor væsentlige
Viden, kan maaskee gjerne rose sig af megen Oplysning og
mangesidig Dannelse, men er dog endnu ikke saa gjennemdannet, at
% --- p. 8 ---
han ret indseer, hvad det er, hvorom her egentlig handles.
Vistnok er det Tvivleren umuligt at fatte Troens overbevisende
Kraft; men det er derimod ingenlunde umuligt for den Troende, at
gjennemgaae Tvivlens Reflexioner. Med denne Selvbesindelse vende
vi da atter tilbage til vort Evangelium, hengive os i
Betragtningen, og bestræbe os for at udtyde Fortællingens Bogstav
i Troens Aand.

Hvor hemmelighedsfuld, og dog hvor aabenbar er ikke denne
Begyndelse! Historien dvæler, og Tanken hviler som i et
Vendepunkt. Tidens Aander have lagt sig til Ro; den udvortes
Verden er bleven saa stille: det er i Keiser August's, det er i
Kong Herodes's Dage. Endnu staaer Templet paa Zions Bjerg:
Israels Folk har endnu et Hjem i Patriarchernes Land. Ingen
udvortes Fiende tør ængste Judæa, thi den snilde Herodes er
Keiserens Ven; intet Oprørsskrig tør lyde i Jerusalems Gader, thi
den grumme Herodes --- straffer med Døden. Dagens Gjerning i
Husene, Offertjenesten i Templet, Alt gaaer sin stille, rolige
Gang efter Skik og Sædvane: Ingen venter, Ingen fordrer det
Overordentlige. Ogsaa Zacharias træder ind i det Hellige efter
Præstedommets Sædvane; det tilfaldt ham nu den Dag at „gjøre
Røgelse“ i Templet. Folkets Mangfoldighed er samlet i Forgaarden,
efter Skik og Sædvane; det burde sig jo, at bede i „Røgelsens
Time“. Ingen venter, Ingen fordrer det Vidunderlige. Da kommer
Aabenbarelsen. Pludselig som et Lyn sees Engelen i Helligdommen,
og Synet slaaer Zacharias med Forfærdelse, og Præsten tøver i det
% Printer's defect: the last sort of „Sædvane“ and the point after it are
% damaged in this copy. Read as „Sædvane.“ at 500 dpi; the ABBYY layer sees no
% punctuation at all here and tesseract reads „Sædvan?:“. A full stop is what
% the sense and the following capital „Og“ require.
Hellige, og Menigheden maa vente --- mod Skik og Sædvane. Og de
Ventende undre sig, hvi han dog tøver derinde; men Ingen aner,
Ingen fordrer det Overnaturlige. See, da træder den gamle Præst
ud i Forgaarden, stum og gaadefuld, og med en saa underlig Mine;
nu først mærker Forsamlingen, at det Vid\-%
% --- p. 9 ---
underlige maa være hændet ham, at han maa have seet „et Syn i
Templet“.

Saaledes er Evangeliets Begyndelse. Lad den verdslige Viisdom kun
forarge sig paa denne Begyndelse; det er jo naturligt, at man
fordrer et andet Udgangspunkt, naar man har en anden Gud. En
Guddom, hvis Væsen er umiddelbart Eet med Naturens Kræfter og
Historiens Aand, kan ikke begynde sin Aabenbarelse saa afgjørende
og dog saa roligt; thi med heftige Smerter, under stærke
Trækninger maa Verdensaanden føde sine nye Ideer. Den tvetydige
Verdensgud træffer ingen Afgjørelse under en Keiser Augustus, og
gjør ingen Underværker under en Kong Herodes. Gjærende i mørke
Muligheder, maa han først oppebie en eller anden fortvivlet
Katastrophe, som t.\ Ex. Jerusalems Forstyrrelse eller Roms
Ødelæggelse. Men Evangeliets Gud med de evige Raadslutninger har
slet ingen Trang til slige endelige Vækkelser. Eller skulde
Forløsningens Gud vel glemme Menneskets aandelige Nød, fordi
Slægten døser, og Folkene leve hen i Tryghed, sigende: „her er
Fred og ingen Fare?“ Eller skulde maaskee Evighedens og Forsynets
Gud ikke formaae at see 70 Aar fremad i Tiden, og fuldende
Aabenbaringens Kredsløb, før den gamle Pagts Helligdom falder i
Gruus, og det udvalgte Folk adspredes over Jorden?

Det er ganske naturligt, at den Ikke-Troende maa forarges paa
denne Begyndelse; thi en Verdensgud kan ikke begynde med saa
klare, saa ideale Kjendsgjerninger. Verdensguden maa begynde med
dunkle Rørelser, umodne Forsøg, tvetydige Tilløb. Først
efterhaanden, som Aandens nye Værk udklarer sig af sit Chaos og
fremviser de storartede Skikkelser i Udførelsens bestemtere
Omrids, træder den digtende Phantasie i Erindringens Tjeneste, og
forherliger den raae Begyndelse ved Sagn og Myther. Men
Aabenbaringens Gud begynder med Underet, og Underet, der selv
% --- p. 10 ---
er Eenheden af det Ideale og det Virkelige som hellig
Kjendsgjerning, vil ingen phantastisk Udsmykkelse af usikkre
Hændelser, men en eenfoldig og sanddru Beretning om den
overnaturlige Begyndelse.

Med Ringeagt seer den opløsende Kritik ned paa de Fromme, som
ifølge vort Evangelium benaades med Engelens Budskab. Med
indbildt Overlegenhed spørger man, hvad der berettiger Zacharias
og Elisabeth til en saadan Begunstigelse. Denne Indvending er nu
i Verdens Øine ligesaa træffende som bitter; thi dersom de
paagjældende Personer virkelig ere saa ringe og ubetydelige, som
Kritiken antager, da gaaer det jo til i Evangeliet, som i
Eventyret og Legenden, at en Gud, der faaer det Indfald at ville
overraske ved et Mirakel, endog i den Forstand ikke anseer
Personer, at han indlader sig med den Første den Bedste. Men er
det dog ikke høist ukritisk, at forvexle Geniets og
Handlekraftens glimrende Storhed, som kan forbause en Verden, med
den Lydighedens og Selvfornegtelsens stille, ydmyge Storhed, som
alene behager Gud? I Verdenshistoriens Øine ere Zacharias og
Elisabeth tilvisse ubetydelige; den kjender Intet til deres
Gjerninger, den ændser neppe engang deres Navne. En
% Printer's error, transcribed as printed: „phychologisk“ for
% „psychologisk“. Confirmed on the image at 500 dpi; both OCR witnesses read
% it the same way. Not among the nine Rettelser.
phychologisk Iagttager, der ellers med ufortrøden Flid opsporer
Gammelt og Nyt, at en dannet Læseverden altid kan være forsynet
med et passende Udvalg af interessante Skyggebilleder til
behagelig Underholdning, vil uden Tvivl gaae vore Gamle med
Ligegyldighed forbi; thi Zacharias's Huus er et stille Huus uden
pikante Scener og nerverystende Optrin. Tage vi derimod til Takke
med den Sjælestyrke, som, uden just udadtil at gjøre larmende
Opsigt, øves og prøves i al Inderlighed for den Alvidende, da har
Evangelisten ikke ladet Zacharias og hans Hustru uden Vidnesbyrd.
Han beretter jo om dem, at de begge vare retfærdige for Gud, og
vandrede ustraffelige i alle Herrens Bud og
% --- p. 11 ---
Rette. Vel sandt, dette Vidnesbyrd, affattet i faa Ord, tager sig
snarere ud som en flygtig Bemærkning, en forbigaaende Antydning,
end som en begeistret Lovtale. Men deels kommer det deraf, at der
slet ikke kan tænkes paa, at Zacharias og Elisabeth skulde ved
deres gode Gjerninger have fortjent en Aabenbarelse, der jo alene
skyldes den guddommelige Naade; deels ogsaa deraf, at Sproget i
faa og fyndige Ord er istand til at udtrykke det, som kræver et
heelt Livs fortsatte Anstrengelse, naar det skal udtrykkes i
Gjerningen. Det er dog, i Alvor talt, ikke saa let at vandre
ustraffelig i alle Herrens Bud og Rette. Den maatte vel kaldes
stor i Israel, som formaaede dette. Og dog kunde En maaskee synes
sig stor i Lydighed, uden derfor at prøves i megen Selvfornegtelse
(thi Israels Gud velsigner og belønner de Lydige); og en Anden
kunde maaskee synes sig stærk i Troen, uden dog at være forsøgt i
megen Taalmodighed, efterdi hans Gudsfrygt aldrig savnede sin
timelige Løn. Men, naar den forventede Løn udebliver, naar den
Gave, der først skal give det nærværende Liv sin rette og fulde
Betydning, den Gave, Patriarchernes Gud har lovet alle dem, der
finde Naade for hans Øine, negtes den Troende; naar det Dag for
Dag bliver mere og mere usandsynligt, at denne Gave nogensinde vil
vorde den Haabende til Deel, da skal der Sjælestyrke til, for at
bevare Troen og Haabet. Det er denne Sjælestyrke, vi ære og
beundre hos Zacharias og Elisabeth. --- Den, der forsmaaer at
sætte sig ind i Hebræerens religiøs-borgerlige Vilkaar, kan i fri
selvbehagelig Dannelse gjerne opholde sig over, at et Par gamle
Ægtefolk saaledes bekymre sig for en Livsarving, at de tilsidst
falde paa den Tanke, at Gud maaskee selv vil komme dem til Hjælp
ved et Mirakel. Men det er dog alligevel sandt, at Ideen om
Familielivets Fortsættelse i talrige og lykkelige Efterkommere var
saaledes sammenvoxet med det patriarchalske Gudsforhold, at
% --- p. 12 ---
den Israelit, der, uden at forplante sit Navn til Efterslægten,
sank ned i Glemselens Skyggerige, kunde ansees som forkastet af
Jehova og udslettet af sit Folk.

Saa kunne vi da ret forestille os, hvorledes Zacharias og
Elisabeth maae have sluttet deres Pagt i den barnlige Tro og det
ungdommelige Haab, at Abrahams, Isaaks og Jakobs Gud vilde
skjænke ogsaa dem sin Velsignelse, hvorledes de i denne Tro maae
have lovet hinanden at vandre ustraffelige for Gud i alle Herrens
Bud og Rette. Men --- Ungdommen svandt, og den glade Forventning
opfyldtes ikke. Naboer og Frænder fik Velsignelsen; men Zacharias
og Elisabeth maatte staae tilbage. Drog da Savnet med lønlig Sorg
ind at boe hos dem, maatte de end stundom see hinanden i Veemod,
saa fattige, ak, saa forladte, fordi Jehova ikke vilde see til
dem, ikke besøge dem: saa var den ydmyge Tro dog større end den
stille Kummer, Haabet stærkere end Frygten, og Forventningens Dage
endnu ikke omme. Forjættelsen staaer med gylden Skrift:
\emph{Den Retfærdige skal grønnes som et Palmetræ, han skal voxe
som en Ceder paa Libanon.} „Hvo der dog kunde findes retfærdig
for Gud, og saa --- faae sin Deel i Forjættelsen! Hvo der dog
kunde vandre ustraffelig i alle Herrens Bud og Rette, og saa finde
Naade for hans Øine, og --- vorde smykket med Velsignelsen!“ ---
Saaledes have vore Fromme da trøstet hinanden med Forjættelsens
Ord, og regnet Dagenes Mangfoldighed efter Lovens Opfyldelse og
Aarenes Tal efter Retfærdighedens Tjeneste, at de begge maatte
være „retfærdige for Gud og vandre ustraffelige i alle hans Bud og
Rette.“ Men Dagene svandt, og Aarene gik, og --- Forventningen
opfyldtes ikke. Palmetræet grønnedes, Cederen voxte, Viintræet
bar sin Frugt; men --- Forventningen opfyldtes ikke. Fromme
Naboer, lykkelige Frænder opløftede Røsten, og priste Jehovas
Miskundhed fra
% --- p. 13 ---
Slægt til Slægt; kun Zacharias og Elisabeth maatte tie og bie.
Fromme Naboer, lykkelige Frænder saae Fædrenes Velsignelse paa
Børn og Børnebørn, naar Festens Dag forsamlede dem til Offringens
Høitid og Gjæstebudets Glæde; men Zacharias og hans Hustru
forbleve alene med deres Offer, alene med deres Taksigelse, alene
i deres Huus. Alderdommen kom, og fandt Zacharias barnløs, og
Elisabeth fik et Tilnavn som den af Herren Forskudte, og hun følte
sin Skam iblandt Menneskene; men --- de forbleve alene. Saa var
det da en Prøvelse, hvori de skulde forsøges, og Zacharias og
Elisabeth forstode det, at det var en Prøvelse; thi de misundte
ikke de lykkelige Frænder, og de knurrede ikke imod Jehova, og de
tabte ikke Haabet i Mismodighed; men --- de fornyede deres
Ungdoms Pagt, sigende i deres Hjerte: „Vi ville dog vandre
ustraffelige i alle Herrens Bud og Rette; thi dette er jo en
Prøvelse. Hvorfor vi, just vi alene, skulle saaledes staae
tilbage i vor Slægt, og bære Menneskenes Ringeagt, som vare vi
ringeagtede af den Gud, vi ville tjene i al Retfærdighed, ak som
vare vi forkastede af ham, der dog ikkun forkaster de Ugudelige!
Hvorfor vi, just vi alene, skulle nedstige i Graven med graae
Haar, som vore Fædre, mætte af Dage, som vore Fædre, og dog ikke
faae vore Fædres Velsignelse, men efterlade vor Bolig øde, som
havde vort Liv kun været en Skygge, ja overlade vore Navne til
Forglemmelse, ak som vare vi forglemte af ham, der dog aldrig kan
forglemme sine Fromme! Hvorfor, o Gud, hvorfor? Men lader os tie
og bie og være stille for Gud; thi hvo har styret Jehovas Aand, og
hvilken Raadgiver har underviist ham! Lader os dog gaae i Rette
med vore egne Tanker, at vi ikke gaae i Rette med Herren; thi, hvo
veed, hvad den Almægtige endnu vil gjøre imod os? Ja store og
underlige Ting har Herren jo fordum gjort imod sine Udvalgte, som
imod Abraham og Sara, Manoah og hans Hustru, Elkana og Hanne;
% --- p. 14 ---
ja Herren er Vidskabers Gud, skulde hans Gjerninger ikke være
rette!“ Dette er da Prøvelsens Vendepunkt. De, som saalænge
havde troet med Haab, troede da nu tilsidst mod Haab, dog haabende
paa den Gud, der levendegjør det Døde, og kalder de Ting, der ikke
ere, som om de vare. See, det er Abrahams Tro, Abrahams Haab!
Med denne Tro og dette Haab møder Zacharias Engelen i det Hellige.

Hvis en ordrig Fremstilling imidlertid lagde an paa, at bevise den
moderne Kritik, hvor værdig og velforberedt Zacharias altsaa kunde
være paa at modtage Engelens Budskab, da vilde et saadant Forsøg
ikke blot i og for sig være temmelig ufrugtbart, men endogsaa
forraade en betænkelig Misforstaaelse af Evangeliets Øiemed. Lod
Sligt sig bevise, hvo var da nærmere til at føre Beviset, end
Evangelisten selv? Men det er Uforstand at kræve Beviis for det,
som ifølge sin Natur intet Beviis tilsteder. Om Fremstillingen
end var nok saa indtrængende, det skulde dog ikke hjælpe. Om
Sprogets hemmeligste Kræfter rørte sig, ja om der end taltes med
Engletunger til Zacharias's og Elisabeths Priis, hvad gavnede det
saa, naar den negative Kritik dog kun vil høre sit eget een Gang
vedtagne Sprog, og den moderne Bevidsthed veed med sig selv, at
den ikke kan lære Englenes Tungemaal? Derfor har Evangelisten
heller ikke indladt sig paa noget Forsøg i denne Retning; men, vel
vidende, at Evangeliet kun er for Troens Eenfoldige, fatter han
sig i al Korthed, og siger om vore Gamle: „Men de vare Begge
retfærdige for Gud, og vandrede ustraffelige i alle Herrens Bud og
Rette.“ Mere behøves ikke; altsaa --- til Engle-Aabenbarelsen.

Idet Kritiken bryder ind i Helligdommen og sigter paa Gabriel,
opstaaer en mærkelig Strid. Kritiken har nemlig udfundet, at,
eftersom den moderne Videnskab ikke anerkjender Englene, saa
% --- p. 15 ---
ere disse slet ikke til, og at, naar Engle ikke ere til, saa er en Engel
med det Navn Gabriel heller ikke til; endvidere har Kritiken udfundet,
at, naar Englene ikke ere til, saa kunne de heller ikke aabenbare sig,
samt at, naar Englene ikke kunne aabenbare sig, saa har en Engel ved
Navn Gabriel heller ikke kunnet aabenbare sig for Zacharias. Derfor
hedder det nu ogsaa i den moderne Videnskabs Sprog, at Gabriel ikke
kan bestaae for Kritikens Domstol, og at man desaarsag har fundet sig
beføiet til at anvise ham Plads iblandt de mythiske Figurer. Altsaa
dette har nu Kritiken udfundet, og det har Altsammen sin Rigtighed, og
der er aldeles Intet at indvende imod den kritiske Kjædeslutning uden
dette Ene, at dens Forudsætning er greben ud af Luften. Et virkeligt
Beviis for Umuligheden af Englenes Tilværelse skal den moderne
Videnskab vel lade være at føre, og naar Troen blot veed Muligheden
frelst, forlanger den ikke Mere af den hele Verdensviisdom. Imod den
Paastand, at det er Gabriel, der ikke kan bestaae for Kritikens
Domstol, kunne vi da paa Troens Vegne opstille den Paastand, at det nok
egentlig er Kritiken, som med alle sine Kritikere ikke kan staae sig
mod Gabriel, og derfor iler med jo før jo hellere at ønske ham ---
Mythen i Vold. Videnskaben gjør naturligviis Intet uden Overlæg. Har
Videnskaben altsaa efter et skarpt Forhør forviist Engelen fra
Virkelighedens Egne som en tvetydig og mistænkelig Person: saa kan man
da ogsaa være vis paa, at Videnskaben hertil har havt sine gode og
tilstrækkelige Grunde. Der er fra Fortolkernes Side gjort Alt, for at
bringe dette besynderlige Væsen til Raison, og forklare Historien paa en
rimelig og naturlig Maade; men alle Kunstens Bestræbelser have hidtil
været forgjæves. Sagen er denne: skal Engelen været Noget, saa maa han
% Printer's error, transcribed as printed: „skal Engelen været Noget“
% (checked at 600 dpi; „været“ is set as one word). Sense requires
% „skal Engelen være et Noget“ or „være Noget“; not corrected here, and
% the Rettelser page does not list p. 15.
jo enten være legemlig eller ulegemlig. Men nu har Evangelisten
fremstillet ham saaledes, at han baade er det Ene og det Andet, og dog
% --- p. 16 ---
hverken det Ene eller det Andet; og hvilket fornuftigt Menneske kan da
blive klog paa ham? Vil den naturlige Fortolkningskunst holde ham fast
som en udvortes Skikkelse, og lade ham svare til sit Navn Gabriel ɔ:
\emph{Guds Mand}; da maa denne Guds Mand naturligviis være en eller
anden høist besynderlig Mand, der har sneget sig ind i det Hellige, for
at overraske Præsten. Men dette Indfald falder strax igjennem.
Helligdommen var jo lukket, og man søger da forgjæves en Bagdør,
hvorigjennem Manden kan være sluppen ind, og hvorigjennem han igjen har
kunnet smutte ud. Vil det nu ikke lykkes at fastholde Engelen i
legemlig Udvorteshed, maa han naturligviis opfattes som en ulegemlig
Aand. Men er han en ulegemlig Aand, kan han jo heller ikke sees med
legemlige Øine: „saa maa han da være et Foster af Præstens egen
ophidsede Indbildningskraft“. Men imod denne Formodning protesterer
Fortællingen paa det Allerkraftigste. Fortolkeren, der vil begribe
Synet, griber i Luften, og see! nu staaer jo Engelen der, \emph{der}!
„ved den høire Side af Røgelsens Alter“, og seer paa den forfærdede
Zacharias, og taler saa lydeligt, saa myndigt, saa udførligt, saa
roligt, klart og besindigt, at ingen Guds Mand kan tale tydeligere, og
ingen Ekstatiker trylle sig ind i en saadan Vision. Var det nu ganske
afgjort og een Gang for alle beviist, at den moderne Videnskab udtømmer
Sandheden, og at Intet har Lov at være til, uden Det, hvis
Tilværelsesmaade Videnskaben begriber: ja, saa var det rigtignok forbi
med Engelen, og forbi med Underet, og forbi med Troen. Men sæt, at det
forholder sig omvendt; sæt, at den høieste Indsigt, hvortil al vor
Begriben fører, er den Indsigt, at det Ubegribelige er lagt til Grund
for det Begribelige, at Aabenbaringens Under hverken kan eller skal
begribes, saa faaer jo Troen Ret, og dermed faaer ogsaa Engelen Ret;
thi Engelen er Underets Talsmand:
% The re-quotation of Luc. 1, 19 is set in the same larger type as the
% pericope on pp. 1--2, but it is NOT a display block: it begins on the
% last line of p. 16 at full measure and ends in the middle of the first
% line of p. 17, where the body size resumes with „Saa bliver“. Hence
% {\large ...} inline rather than \begingroup\large ... \par\endgroup.
% Printer's defect: an inked quad (or foreign body) prints as a black
% blot between „Gabriel,“ and „som“; checked at 600 dpi. Not reproduced.
{\large jeg er Gabriel, som staaer for Gud, og er udsendt at
% --- p. 17 ---
tale til dig, og forkynde dig dette til Glæde.} Saa bliver det da vel
det Bedste, at lade Gabriel være, hvad han er, et overnaturligt Væsen,
en Bebudelsens Engel i den Helligdom, hvis Forhæng Kritiken ikke
formaaer at løfte.

Men selv et troende Øie er ikke i ethvert Øieblik lige forberedt paa at
møde Bebudelsens Engel. Den rolige, historiske Fortælling, der her
meddeler det Allervidunderligste i Form af umiddelbar Kjendsgjerning,
og uden alle Omsvøb indskyder et overnaturligt Led i de naturlige
Begivenheders Række, har for den opvaagnende Tænkning noget uendelig
Frastødende. Havde Evangelisten endda, før han udtrykkelig omtaler
Synet, forudskikket en indledende Bemærkning, som f.\ Ex., at
Zacharias, der han traadte ind i det Hellige, var i Aanden, eller: saae
Himlen aaben, --- eller forundt os et lignende Vink; da skulde det
maaskee endda ikke falde Fortolkeren saa vanskeligt, at udfinde en
Forklaring. For at oplyse dette Punkt lidt nærmere, ville vi hente et
Par Exempler hos Propheterne. Jesaias beskriver sin Kaldelse (Jes.\ 6,
1--8); hvilket Gjennembrud, hvilket Syn! Paa sin høie og ophøiede
Throne sidder Jehova; Fligene af hans Klædebon opfylde Templet.
Serapherne, der staae over ham, tilhylle deres Aasyn, og raabe, den ene
til den anden: „Hellig, hellig, hellig er den Herre Zebaoth, al Jorden
er fuld af hans Ære.“ Dørposterne skjælve ved den Raabendes Røst, og
Gudshuset fyldes af Røg. Forfærdet udbryder Seeren: „Vee mig, jeg maa
forgaae; thi jeg er en Mand med urene Læber, og jeg boer midt iblandt
et Folk, som har urene Læber, og dog have mine Øine seet Kongen, den
Herre Zebaoth.“ Da henter Seraphen en Glød fra Alteret, og berører med
den Prophetens Læber, „at hans Misgjerning maa vige fra ham, og hans
Synd udsones.“ I Sandhed, denne Indvielse virker paa os med
Aabenbaringens hele umiddelbare Kraft; det er,
% --- p. 18 ---
som fornam vi den hellige Ild, der gjennemgløder Prophetens Læbe. Men
Synet er os tillige fremstillet som et reent Aandesyn, og, just fordi
det er et Aandesyn, tilsteder Kjendsgjerningen en Forklaring.
Jehova-Skikkelsen forvandler sig til et Billede, og Billedet antager et
allegorisk Anstrøg; Underet gaaer op i Ideen, og Ideen tilfredsstiller
Reflexionen. Dette gjælder endnu i høiere Grad, naar de apokalyptiske
Billeder bevæge sig om Verdenshistoriens Kriser. Naar en Daniel i
Nattens Syner seer „den Gamle af Dage“ omgiven af Englenes mange
Tusinder, naar han seer Retten blive sat og Bøgerne opladte, og En i
Himmelens Skyer som et Menneskes Søn komme ind til „den Gamle af Dage“,
og af ham modtage Riget og Magten og Æren: da have disse natlige Syner
en saa aandig, luftig Form, at Endelighedens Tyngde forsvinder, og
Billederne drage forbi i store, men flygtige Omrids. Det er Historiens
Aander, der omsuse Prophetens Hoved; det er Alstyreren, der kalder
Folkeslagene for sin Throne, det er „Menneskesønnen“, der holder
Dommedag. Ja, saa paafaldende er Ligheden imellem disse apokalyptiske
Billeder og Historiephilosophiens aandrige Ideeformer, at en troende
Betragter saare let forvexler de overnaturlige Tegn fra Aabenbaringens
Forsyn med de naturlige Bestemmelser i Verdensstyrelsens Almeenbegreb,
og ender som philosophisk Tilskuer. Fornuften har vel Intet imod, at
Mennesket staaer i et almindeligt Forhold til Gud igjennem
Verdensideen; men dette tilfredsstiller ikke Troen. Troen er rodfæstet
i den Overbeviisning, at den frelsende Gud indgaaer i et umiddelbart
personligt Forhold med det faldne Menneske formedelst Underet. Staaer
det nu fast, at Underet er Aabenbaringens rette og egentlige Mærke, saa
staaer det ogsaa fast, at den Aabenbarelse, der nedlægger
Aandsbestemmelsernes hele Gehalt i Underets umiddelbare Kjendsgjerning,
er at ansee som Norm og Maalestok for
% --- p. 19 ---
Bedømmelsen af de Aabenbarelser, hvis overstrømmende Idee-Fylde mætter
Reflexionen og mildner det Anstødelige. Hermed have vi da vundet det
Synspunkt, hvorfra Bebudelsens Engel bliver at betragte. Det er nemlig
saa langt fra, at hine prophetiske Syner skulde have større aandelig
Gyldighed end vort Evangeliums Engle-Aabenbarelse, at Tilegnelsen af
denne meget mere er os den rette Borgen for, at vi ikke misforstaae
hine. Vel maa det indrømmes, at Engelens Komme her er holdt i en saa
naiv og udvortes Virkelighedsform, at Den, der slet ikke kan mærke det
Anstødelige, enten maa have den fuldkomne Tro i salig Eenfoldighed,
eller være hildet i Overtro, eller være en sløv og tankeløs Læser. Men
dette Kjendsgjerningens Haarde, paa hvilket Reflexionen støder sig, er
netop Tilegnelsens Prøvesteen. Det er Anstødet, der vækker Troen. Kan
jeg ikke troe Bebudelsens Under, fordi det nemlig er et Under, saa kan
jeg trods al Philosopheren over de prophetiske Fremtidssyner heller
aldrig troe Propheterne selv, det vil da sige: jeg har slet ikke Troen.
Den Troende derimod, der forstaaer Propheterne, saaledes som de,
egentlig talt, ville forstaaes, forstaaer ogsaa Engelen, saaledes som
Zacharias forstod ham; det er i --- „Frygt og Bæven“. Den Troende, der
fatter Aabenbaringens Aand, vil derfor ikke bestride Kjendsgjerningen,
men stride for Kjendsgjerningen, indtil Lyset opgaaer, og Engelen
vinder Seier. Den Troende, der fatter Aabenbaringens Aand, vil ikke
sætte sig hen at gruble over den Ulegemliges Legemlighed, at han ikke
skal tabe Forstanden; vil ikke fræk trænge sig ind i det Hellige, og
stirre paa den Guddommeliges Herlighed, at han ikke skal miste Synet;
men han vil i al Ydmyghed nærme sig Forhænget, og bøie sit Øre mod
Helligdommen, og lytte til Engelens Røst, som den lyder i
Helligdommen: i denne Røst vil han da
% --- p. 20 ---
gjenkjende den prophetiske Røst; thi Engelen er i Forstaaelse med
Propheterne.

Den moderne Bevidsthed kan her uden Tvivl ikke fornægte sin liberale
Anskuelse. „Skulde Himlen være saa naadig; skulde der hist bag
Forhænget virkelig findes en Engel, som endnu kunde hviske et Trøstens
Ord til en bedrøvet Sjæl, da var det sandelig Synd, om nogen Høirøstet
skulde komme ham for nær, og forstyrre Husvalelsen. Enhver maa i saa
Henseende have sin Frihed. Vil Troen blive staaende i al Eenfoldighed,
og præke det Ubegribelige: nu vel, lad den staae! Men den maa finde sig
i, at Verden gaaer fra den; thi Verden vil nu een Gang ikke blive
staaende; Verden gaaer med Tiden, og Tiden gaaer, og gaaer bestandig
fremad. Verden gaaer efter Lyset, og Oplysningen gaaer jo fremad,
bestandig fremad. Nutidens Livsspørgsmaal forudsætte derfor ganske
andre Idealer end Englene i Kong Herodes's Dage, og udkræve ganske
andre Forbilleder end Zacharias og Elisabeth.“ --- Men er Nutiden nu
virkelig saa frisindet, da er vort Evangelium tilvisse ogsaa frisindet,
og forsaavidt maa vel Striden snart kunne endes. Engelen er ikke
paatrængende; han træder ikke ud af Helligdommen; han vil jo ikke træde
i Veien for nogetsomhelst Verdensideal. Nei, Gabriel gaaer ikke omkring
at trætte Nogen med sin Nærværelse, eller at trættes med Nogen om sin
Bemyndigelse: som et Tilbud kommer han, naar han udtrykkelig sendes;
men han sendes kun til dem, der behøve en saadan Sendelse. Tro ham, han
har slet ingen Tid, at gaae frem med Tiden: han maa i al Eenfoldighed
blive staaende; thi --- han staaer for Gud. Tro ham, han er i Sandhed
ikke misundelig paa Oplysningens Fremskridt; han staaer for Gud, og
den, der staaer for Gud, skulde han vel savne Oplysning?

Ogsaa med vore troende Gamle, der i Forhold til Engle-%
% The hyphen at the foot of p. 20 is a real one: the compound is set
% „Engle-Aabenbarelse“ mid-line on p. 19, so the second element keeps its
% capital. Joined with \-%-style continuation so no space is introduced.
% --- p. 21 ---
Aabenbarelsen her ere stillede som evangeliske Forbilleder, vil Nutiden
saare let blive færdig. Thi disse fromme Ægtefolk ere jo heller ikke
paatrængende; de gjøre slet ingen Prætensioner, men holde sig ved sig
selv i al Eenfoldighed. Elisabeth er ikke saa forfængelig, at hun nu
skulde gaae hen at klæde sig om, for at være moderne: hun drager sig jo
tilbage i al Stilhed, og søger Skjul i sit eget Huus. Zacharias har
ikke den Ærgjærrighed at ville lede Tidens Bevægelser; Underet, der
slog ham med Forundring, standsede ham Tankernes Flugt: han maatte
forstumme. --- Uden at besværes af Evangeliets Under kan Verden altsaa,
som det sig bør, arbeide fort paa sine egne Opgavers Løsning, og
forsøge alle Timelighedens Midler til Opnaaelse af timelige Øiemed. Et
Spørgsmaal om, hvor det historiske Fremskridt saa dog tilsidst fører
hen, tør man imidlertid nok tillade sig. Skulde Tilværelsens dybere
Mening ufeilbarlig være denne, at enhver Tanke om det Evige snarest
muligt bør afskaffes, for at Aandens hele Stræben desto kraftigere kan
concentrere sig paa det Timelige, da vil Oplysningen jo nu med det
Allerførste være fuldkommen færdig med Aabenbaringen, og
Verdensklogskaben let affinde sig baade med Gud og hans Engle. Men naar
dette nu ikke kan være Meningen, naar Timeligheden dog vedbliver at
hungre midt i sin Overflod, og ikke selv har det, hvormed den kan
stille sin Hunger; naar Verden idelig maa skifte Farve, snart af Harme,
snart af Undseelse over, at Opgavernes forventede Løsning mod Alles
Forventning opløste sig i en Skuffelse: da kan Culturbevægelsen heller
ikke uden videre vedblive at løbe ud i uendelig Endelighed; men
Fremskridtet maa i hvert afgjørende Øieblik standse ved det
„tungsindige“ Spørgsmaal: er der saa en Gud, er der en Evighed? Deraf
kommer det, at Verden trods al sin Ustyrlighed dog aldrig ret kan
beqvemme sig til et aabenbart Brud med den høiere Styrelse. Kun
% --- p. 22 ---
i et Anfald af dæmonisk Fortvivlelse lader Oplysningen sig en Gang
imellem forlyde med, at den paa sin Fremskridens hurtige Fart er ilet
„Gud og Djævel, Himmel og Helvede“ saa langt forbi, at selv den
Feigeste ikke mere tør ængste sig for disse Skrækkebilleder; men det
er, som sagt, kun i et Anfald af opbrusende Lune, at Oplysningens
dumdristige Ordførere saaledes tale over sig. Verden selv bliver stedse
mere besindig, alt eftersom den bliver mere dannet; hvi skulde en
dannet Verden da ogsaa være saa ubesindig, at bryde med Gud i
Himmelen? Idet Menneskeheden gjør Fremskridt, gjør den ogsaa
Erfaringer, og, det forstaaer sig, af Erfaring bliver man jo klog.
Verdensklogskaben er allerede kommen saa vidt, at den ikke en Gang for
ramme Alvor tør rive sig løs fra Evangeliet. Hvorfor skulde den ogsaa
gjøre det? Sligt et Vovespil er usundt, og en indvortes Usundhed kunde
let foraarsage urolige Drømme. Sees nu Sagen fra denne Side, da er det
jo ingenlunde de evangeliske Forbilleder, der trænge sig ind paa
Oplysningen; men det er omvendt Oplysningen selv, der atter og atter
maa vende sig om, og give sig i Færd med Evangeliet. Ved første Øiekast
tager det sig da ud, som om Oplysningen var den overlegne Magt, der
skulde drage det evangeliske Indhold med ind i den høiere Udvikling,
rense Guldet fra de uædle Tilsætninger, og derved gjøre det aabenbart,
hvad Aabenbaring egentlig er. Men jo mere Verdensfornuften anstrenger
sig i denne Retning, desto klarere bliver det, at Overlegenheden er paa
Evangeliets Side, at det er Aabenbaringen selv, som, dybere forstaaet,
skal oplyse os om, hvad Oplysningen er i Forhold til Troen. Ere nu
Evangelietroen og Oplysningen, ifølge Tidens Vilkaar, ligesaa
uadskillelige, som uforenelige, da er der heller intet Andet for, end
at de uophørlig maae brydes med hinanden. Under denne Brydning viser
det sig, at Troen ikke blot bliver staaende i al Eenfoldighed,
% --- p. 23 ---
men tillige, een Gang berørt af Reflexionen, sætter stærke Tanker i
Bevægelse, for at hævde sin Eenfoldighed, og derved forvandler sig til
den meest kritiske af alle Kritikens Opgaver. Blandt Tidens mange
praktiske, uafviselige, paatrængende Spørgsmaal vil Spørgsmaalet om
Evangelietroens Væsen og Kjendemærke altid lyde med som et
Livsspørgsmaal. Skal altsaa vort foreliggende Evangelium rettelig
tilegnes, skal Tilegnelsen sikkres mod Oplysningens Misforstaaelse, og
den Troende gjøre sig Rede for, hvad det egentlig er, han troer, idet
han troer Evangeliet, da maae de hellige Forbilleder sætte sig i
Bevægelse, da maae Elisabeth og Zacharias, ja Engelen selv virke med
kritisk Kraft, og forhjælpe os til at udfinde Forklaringen.

Oplysningen bifalder gjerne, at Elisabeth betroede Gud sin Sorg,
efterdi hun troede paa Gud. Men hvis hun i sit Hjerte har næret det
Haab, at Alstyreren skulde træffe en overordentlig Foranstaltning, for
at aftørre hendes Taarer, og opfylde hendes Begjæring; da kan en oplyst
Forstand aldeles ikke bifalde dette. „Jeg, dette enkelte Menneske, er
ikke berettiget til at antage, at den Gud, der styrer de utallige
Verdener efter evige, uforanderlige Love, skulde indlade sig paa mine
smaalige Sorger, og tage sig af mine private Anliggender. Kun i samme
Grad, jeg bliver mig bevidst, at mit Livs Formaal har Betydning for det
store Hele, kun i samme Grad tør jeg ogsaa troe mine Tilskikkelser
optagne med i Forsynets særegne Styrelse. Lad os et Øieblik vende os
fra Elisabeth til et Forbillede, der staaer Nutiden nærmere. En
elskelig Fyrstinde deler Thronen med den Hersker, som med fast
overmenneskelig Kraft har grundlagt et nyt, et herligt Rige. Thronen
kræver en Arving, Herskeren en Søn; men Haabet --- svigter Fyrstinden.
Var en saadan Qvinde ikke oprindelig religiøs, hun maatte blive det ved
Omstændighedernes Magt. Saa lykkelig,
% --- p. 24 ---
og dog saa ulykkelig! Ja hun er jo den Lykkelige: ingen Underblomst er
vidunderlig som hendes vidunderlige Lykke; ikke kan Diademet funkle som
hendes Blik, naar det i stille Henrykkelse hviler paa Ham, Helten, hvis
Priis er Folkenes Jubel; paa Ham, Herskeren, for hvem de Mægtige bøie
sig i Beundring. Men --- „Thronen kræver en Arving, Herskeren en Søn“.
Hun føler det i sin Angst, hver Gang hun seer paa den tilbedede Herre,
--- at hans Øie stirrer; hun forstaaer det i hemmelig Qval, naar hun
ret betragter Ham, Herskeren, at en Skygge af Mismod mørkner hans
Aasyn. Saa er hun da blandt alle Ulykkelige den Allerulykkeligste, hvis
Haabet svigter. Denne Sorg kan en Qvinde ikke bære alene, og Ingen af
alle dem, der ville bære hende paa Hænderne, kan bære den for hende: da
vender hun sig i sin Sjæls Angst til Gud, at han dog for sin
Barmhjertigheds Skyld maa høre hendes Suk, og lade sig røre af hendes
Taarer. Gaves der nu en Himmelens Engel, der vilde tilsige dette
beklemte Hjerte et eneste, ak, blot et eneste Ønskes Opfyldelse: da
have vi jo her en Bedende, som i usigelig Lykke vilde offre sit hele
Liv til en stille Taksigelse. Gaves der nu en Forbarmelse for den, der
ellers maatte fortvivle over Tilskikkelserne: da var jo her en
Beængstet, som paa et eneste Tærningkast kunde vinde Alt, o ja, men
ogsaa tabe Alt i Lykkens forunderlige Spil. Gaves der nu en Styrelse,
som for det enkelte Menneskes Skyld vilde bøie Nødvendighedens Lov,
naar Individet havde forviklet sig i store Begivenheders Traade: da
havde vi jo her et Enkeltvæsen, hvis Skjæbne naaer langt ud over et
Familielivs snævre Grændse, og viser sig afgjørende for et Riges
Skjæbne, for Millioners Vee og Vel. Men det er en mislig Sag med denne
umiddelbare Tro paa Guds særegne Styrelse. Maaskee gaaer Forventningen
tilsidst i Opfyldelse: og da ville Tusinde lykønske den Lykkelige, og
Titusinde indbilde sig at prise det styrende
% --- p. 25 ---
Forsyn. Maaskee skal Haabet svigte: og da maa den Ulykkelige, forskudt
af Herskeren, forladt af Verden, og forglemt af Menneskene, see at
frelse sin Sjæl i den Forstaaelse, at Styrelsen gaaer sine egne Veie,
og kræver den Enkeltes Villie til Offer.“ --- I denne moderne
Hentydning maa det nu være let at udfinde den egentlige Mening. Den
gamle Elisabeth er jo ikke i nogen høi Stilling, og i det Hele taget
for ubetydelig til at vente det Overordentlige af Gud (thi at den
forventede Søn skulde blive en saa stor Mand, en Prophet som Elias,
kunde hun jo ikke vide, før Engelen havde aabenbaret det); „at hun
imidlertid saa trolig nærer det kjære Haab, og det i saa mange, lange
Aar, kan vel være et rørende Beviis paa Datidens patriarchalske
Eenfoldighed, men overtroisk er det alligevel; om de Vanskeligheder,
hvori den umiddelbare Forsynstro indvikler det tænkende Menneske, har
denne Aarons Datter endnu ingen Anelse, og kan følgelig heller ikke
opstilles som et Troens Forbillede for Nutiden“.

Her maae vi nu holde et vaagent Øie med Kritiken, at den ikke
kritiserer os ind i en bundløs Misforstaaelse. Vil Oplysningen være saa
resolut at forkaste Troen og dermed al religiøs Sandhed, da er det i
sin Orden, at den uden videre opløser og forkaster Forbilledet; men
saalænge der endnu i Verden er Noget til, der kaldes Religion, vil
Evangeliet hævde sine Forbilleder, og Aarons Datter vedblive indtil det
Sidste at vidne om Troen. Uden Tro kan Mennesket i det Høieste kun
fatte sig selv som Led i den fornuftige Verdensorden, og maa finde sig
i, at den lille Deel opoffres for „det store Hele“. Selvopholdelsen
strider en haard Strid med Nødvendigheden; men Loven er ubøielig: her
er ingen Redning for den Enkelte. Kaldet til at virke for almindelige
Formaal, vil den Begavede gjerne give sig hen i Ideens Magt, om den
maaskee kunde frelse ham ud af Forkrænkeligheden; men den rene Idee er
en troløs Ven: naar Individet mener sig sit
% --- p. 26 ---
høie Maal allernærmest, kaster Ideen det endelige Redskab bort, og
styrter sig i Verdensløbet, som Strømmen i Havet. „Er der da slet ingen
Frelse for Menneskets Selv?“ Jo vist er der Frelse, saa sandt Gud
lever, en Frelse i Troen. I Troen forandres Alt, Verdens Forhold til
Mennesket, og Menneskets Forhold til alle Ting; thi Mennesket forandrer
sin Stilling paa Jorden, idet han finder sin Gud i Himmelen. Hvad
Naturloven ikke formaaer, at redde Individet fra Undergang, efterdi den
kun bestaaer i Tingenes Omskiftelse; hvad Ideens Magt ikke formaaer, at
udrive det enkelte Menneske af de elementære Kræfters Vold, efterdi den
i sin Eiendommelighed dog kun er en relativ Magt iblandt de mange
Tilværelsens Magter: see, det formaaer de Troendes Gud; thi Gud er jo
den Almægtige. Hvorledes det iøvrigt lader sig tænke, at der trods
Endelighedens utallige Magter kan være en Almægtig, som har alle Ting i
sin Magt, derom veed Troen Intet; den indlader sig ikke paa nogen
Begriben, og tilbyder ingen forklarende Omskrivning, men siger til
hvert Menneske, der vil frelse sin Sjæl, kort og godt: dette skal du
troe, at Gud er den Almægtige. Er dette nu Troens Særkjende, at den saa
trolig bevarer det stærke, umiddelbare Indtryk af Guds Almagt; kan den
Troende nu een Gang ikke komme høiere end til den Almægtige, men maa
finde det saligt, at være under hans Varetægt, og boe under hans
Vingers Skygge; saa er Elisabeth jo endnu et Troens Forbillede.

Paa dette Punkt kunne vi maaskee bedst indlade os i nærmere
Underhandling med de Oplyste. --- I Verden gjøres der, og det med
Rette, en mangesidig Adskillelse imellem høie og lave, store og smaa,
stærke og svage, betydelige og ubetydelige Mennesker. Om de forskjellige
Bestræbelser og Øiemed gjælder det ligeledes, at nogle ere uvigtige,
andre vigtige, og den ene langt vigtigere end den anden. Men skulde den
samme Adskillelse nu ogsaa i samme Betydning kunne
% --- p. 27 ---
gjælde for Gud, den Almægtige? Skulde noget verdsligt Formaal virkelig
være saa storartet, at det kunde imponere Gud? eller skulde Nogen vel
have Behov at sætte sig paa en Throne, for at tildrage sig Forsynets
Opmærksomhed? For at tildrage sig et Publikums Opmærksomhed, er det
vistnok fordeelagtigt at indtage en høiere Plads i Samfundet; men med
Guds Opmærksomhed er det jo en ganske anden Sag, efterdi den Almægtige
tillige er den Alvidende. Hvad er den hele Verden for Gud? --- Et
Intet, saa sandt Gud er almægtig. (At dette ligger i Almagtens Begreb,
maae de Oplyste jo selv bedst vide, eftersom det just er dem, der
randsage Begreberne.) Enten gives der altsaa slet ingen Styrelse, fordi
Alt, selv det Allerstørste, er Forsynet for lidet, et Intet for den
Almægtige; eller ogsaa maa den styrende Villie antages at virke i det
Mindste, som i det Største, i den Enkeltes Liv, som paa Millioners
Skjæbne. For at forebygge enhver Mistydning, maa det udtrykkelig
fastholdes, at denne de verdslige Forskjelligheders Forsvinden
ingenlunde ophæver Begrebet om den særlige Styrelse, som om Alt i
Grunden havde lige Gyldighed, det vil sige: var i sig selv ligegyldigt
i Forsynets Øine; da Troen tvertimod forudsætter, at den Troende staaer
i et saadant Inderlighedsforhold til Gud, at hans Liv just derved faaer
en ganske eiendommelig og særegen Betydning. Den Troende lever nu
tillige i Endeligheden; og hermed er det da givet, at Haabet om
endelige Goder ogsaa kan optages med i den sande Tro. Men forsaavidt
den Troende sætter sit Haab til nogetsomhelst timeligt Gode, spørger
han slet ikke om, hvorvidt dette Gode i verdslig Forstand hører blandt
de betydelige eller ubetydelige Ting; men han spørger sig selv:
% Printer's defect: the colon after „selv“ prints only its upper dot
% (checked at 600 dpi); the lower dot failed to ink. Set as a colon.
„har det Gode, du attraaer, nu virkelig ogsaa en saadan Betydning for
din eiendommelige Tilværen, at du i Sandhed tør henføre Opnaaelsen
deraf til Gud selv, bede Gud derom, og, hvis det bliver dig til
% --- p. 28 ---
Deel, modtage det som en Guds Gave?“ --- Under denne Betingelse lader
det Timelige sig forene med det Evige; thi Opnaaelsen eller
Ikke-Opnaaelsen af det endelige Gode bliver da en Sag imellem Mennesket
og Gud; hvilket ogsaa kan udtrykkes saaledes: Mennesket henstiller sin
Sag til Gud. Hvormeget der end i Verden kan være afhængigt af
Omstændigheder og Hændelser, af Lykke og Ulykke, der er dog Eet, som er
unddraget Tilfældighedernes Spil, og det er den Sag, et Menneske i
Troen har henstillet til Gud. Er det nu ganske vist, at Oplysningen med
sine lange Fremskridt dog aldrig kan komme videre med det styrende
Forsyn, saa er Elisabeths Eenfoldighed jo den høieste Viisdom. Denne
Aarons Datter maa vel forstaae, hvad det vil sige, at henstille sin Sag
til Gud; hun har jo maattet offre et heelt Liv paa at lære det. Hun
ældes, og hører Menneskenes stridige Domme om Medgang og Modgang, Lykke
og Ulykke, hvorledes de snart i al Tryghed haabe det Bedste, snart
igjen forsage og befrygte det Værste; men hendes Tanker forvirre sig
ikke, hendes fromme Sind formørkes ikke: hun ældes i Jevnmodighed og
forynges i Mildhed; thi hun har henstillet sin Sag til Gud. Hun bliver
gammel, og hører Menneskenes ivrige Strid om Straf og Brøde, Løn og
Fortjeneste, hvorledes Den og Den fik sin fortjente Løn, fordi han
gjorde, hvad ondt er; men Den og Den fandt ingen Løn hos Gud, skjøndt
han dog meente sig saa retfærdig: hun veed det, at hun ikke er blandt
de Lykkelige; men der er ingen Bitterhed i hendes Sjæl. Hun veed det,
at hun ikke kan rose sig for Menneskene af sin Løn hos Gud; men hun har
ingen nagende Tvivl: hun tier og troer, og bærer sin Skam. Saa er hun
da nu paa Udvælgelsens Vei bleven gammel, og dog --- ung; forvandlet,
og dog --- uforandret: Ungdommens Forventning bærer sin Frugt i
Alderdommens Haab. Der er en inderlig Fred, en
% --- p. 29 ---
stille Forklarelse over den Bedagede: hun har henstillet sin Sag til Gud.

Hvad her er sagt om Elisabeth, gjælder i lige Maade om Zacharias.
Engelens Forkyndelse begynder med disse Ord: „Frygt ikke Zacharias, din
Bøn er hørt!“ Denne Hilsen vidner udtrykkelig om, at den fromme Præst
endnu i sin høie Alder (han siger jo selv V. 18: jeg er gammel, og min
Hustru er ogsaa bedaget) dog ikke har forglemt, at bede den samme Bøn,
han bad, da han stod i sin Manddoms Kraft; thi det er jo ikke tænkeligt,
at Engelen skulde minde ham om en Bøn, han for længst havde glemt. (Det
Guddommelige taaler ingen Glemsomhed; bønhører Gud forglemte Bønner, da
vorder det et Menneske til Straf, men ikke til Glæde.) Forsaavidt har
Zacharias altsaa ikke ændset Naturlovenes Orden, men i dette Punkt alene
adspurgt Guds Villie, henstillet Sagen til Gud, og er just derved kommen
til at overskride Sandsynlighedernes Grændse. Er Oplysningen nu ikke
tilfreds med Elisabeth, da kan den endnu mindre være det med Zacharias:
„denne eenfoldige, indskrænkede Mand med sine ubetydelige Familiesorger
har jo heller intet Krav paa Nutidens Sympathi.“ Forstaaer sig, den
historiske Bevidsthed kræver en ganske anden Tilfredsstillelse. Derfor er
det vel bedst, at vi (for ikke aldeles at oversee Nutidens Fordringer)
endnu en Gang betragte hiint moderne Forbillede, hvortil vi ovenfor
henvistes. „Verdensgeniet er altsaa kommet til sin Ret, den forgudede
Helt har vundet Diademet: for sin egen, for sin Thrones, for Millionernes
Skyld vender han sig da i en stille Tanke til Styrelsen, at den dog nu
maa sikkre hans Rige og skjænke ham den haabede Søn. Men det forstaaer
sig, Verdensloven maa her, som overalt, naturligviis staae ved Magt,
Lykken og Ulykken, Skjæbnen og dens Luner tages med i Beregning. Staaer
Sandsynligheden Ønsket imod, maa Herskeren see at finde en an\-%
% --- p. 30 ---
den Udvei: enten opgiver han da Forventningen og tager andre
Forholdsregler til Rigets Sikkerhed, eller ogsaa maa han endnu, før det
er for sildigt, dictere Skilsmisse, og indgaae en ny Forbindelse, at
Lykken dog maatte føie ham i Alt, og Sandsynligheden understøtte
Haabet.“ Hvorvidt Herskeren forøvrigt heri handler klogt eller uklogt,
derom er Talen ikke; men vist er det, at Sandsynlighedsberegninger af
denne Art Intet have med Evangelietroen at skaffe. Paa Troens Vegne gaaer
Zacharias da ogsaa en anden Vei: han opgiver ikke sit Haab, skjøndt Tiden
gaaer hen; han søger ikke Skilsmisse, skjøndt Elisabeth ældes; men han
holder sig til den Almægtige, henstiller Sagen til Gud, og --- beder.

Maaskee vilde det lyde som en ikke uvittig Spot, hvis en moderne Kritiker
nu kom paa det Indfald, at vi her opstillede „et ret træffende Forbillede
for alle Dem, der ved Bøn og Besværgelse søge at opnaae, hvad de i denne
syndige Verden ikke kunne opnaae ved Brugen af naturlige Kræfter, enten
fordi de overhovedet mangle al Kraft, eller fordi de ere for
ubehjælpsomme til at bruge de Kræfter, Vorherre har givet dem.“ Men, da
det ikke er for Spotterne og de Lattermilde, men for dem, der i Alvor
ville undervises om, hvad Troen formaaer, at Evangeliet forkyndes: saa
overveie vi nærmere det Spørgsmaal, om det ikke netop ligger i Sagens
Natur, at den Troende først kommer til Klarhed over sig selv, naar han
begynder at bryde med Sandsynligheden. Saalænge det Nærværende yder os
Tilfredsstillelse, uden at varsle om nogen Fare for det Tilkommende,
saalænge Sandsynligheden til alle Sider ligesom omfreder og betrygger vor
hele Tilværelse, er det ingen Sag at stole paa Gud; det er omtrent det
Samme, som at stole paa den naturlige Tingenes Orden. Men, saasnart
Sandsynligheden blot vil vende sig om, og vende den mørke Side til, kan
Prøven strax aflægges. --- Det er farligt, at gaae ind i et Pest\-%
% --- p. 31 ---
huus; der er al Sandsynlighed for, at man bliver smittet. Den Modige kan
imidlertid træde ind, og gjøre sin Pligt, skjøndt han mangler religiøs
Tro; den Troende kan ligeledes træde ind, og gjøre sin Pligt, skjøndt han
mangler naturligt Mod. Hiin bryder sig ikke om Sandsynligheden, fordi han
stoler paa sig selv. Denne derimod bryder med Sandsynligheden, fordi han
stoler paa Gud. Den Modige har den fuldeste Anerkjendelse af
Sandsynlighedens Betydning; han veed, at han her, fremfor paa noget andet
Sted, er i den største Fare; men han trodser Faren; thi han kan jo døe,
og dermed er Sagen endt. Den Troende derimod kan ikke trodse Faren,
altsaa byder han Sandsynligheden Trods, og giver sig hen i Guds Villie.
Vil Gud, at han skal gaae karsk derfra, saa kan ingen Smitte røre ham;
vil Gud, at han skal gjennemgaae Sygdommen, saa giver han ham Kraft til
at lide; vil Gud, at han skal døe, saa er det til hans Bedste. Men er det
nu et Troens Særkjende, at tillukke Øiet for de endelige
Naturbetingelser og trodse Sandsynlighederne, saa er Zacharias jo et
Troens Forbillede.

Oplysningen, der ikke kan undvære en almægtig Gud, har istedetfor
Evangelietroen efterhaanden faaet en anden Tro indført, som formeentlig
skal staae sig bedre med den voxende Erkjendelse af Naturlovene, nemlig
den saakaldte Fornufttro. Ifølge Fornufttroen tilkommer Gud vistnok den
Ære, at være Naturens Skaber og Verdens Lovgiver fra Begyndelsen af. Men
siden den Tid, Verden blev til (og det er vist langt over 6000 Aar siden,
at Alt blev indrettet paa det Bedste), maatte Gud paa en Maade suspendere
sin Almagt og lade Naturlovene raade sig selv; og hvorfor? fordi han er
altfor fornuftig til at ville forstyrre den een Gang indførte Orden.
Fastholdes nu Styrelsens Begreb i reen, ubestemt Almindelighed, da synes
Fornufttroen i alt Væsentligt at stemme med den evangeliske Tro. Fra
begge Sider indrømmes det, at
% --- p. 32 ---
Gud er almægtig, at Intet skeer uden hans Villie, at Mennesket altsaa bør
erkjende Forsynets Styrelse i alle Ting. Desuagtet er her i
Virkeligheden en uendelig Forskjel. Denne Forskjel viser sig
allertydeligst i Bestemmelsen af det Sandsynlige og det Mulige. Lad to
Mennesker enes om, at henføre Livets Tilskikkelser til Gud, men saaledes,
at den Første forstaaer det efter Fornufttroen, den Anden derimod holder
sig til Evangeliet: vi skulle da snart erfare, at der ligger en Verdens
Forskjel imellem dem. Med den gjensidige Forsikkring, at Menneskets Liv
er i Guds Haand, tiltræde de deres Reise. Underveis komme de i Havsnød,
og gribe begge i den samme Redningsplanke. Naar da den Første nu seer sig
om til alle Sider og Intet øiner uden den vilde Søe, naar
Sandsynligheden for at gaae under forholder sig til Sandsynligheden for
at frelses, som Hundrede til Een, saa udbryder han: „Redning er umulig,
vi maae forgaae!“ Men den Anden svarer: „Kjære! hvi skulde Redning her
være umulig? husk dog paa, at Menneskets Liv er i Guds Haand!“ Saa
frelses de da Begge, og forene sig om at takke Styrelsen. Men i det
fremmede Land, hvorhen de nu ere komne, vorde de af Barbarer dømte til at
kastes i en gloende Ovn. Da klager den Første i sin Sjæls Bitterhed:
„See, nu er al Redning umulig! Bedre havde det dog været for os, om vi
vare blevne paa den vilde Sø, end at vi skulde komme til dette barbariske
Land, for at lide en saa forsmædelig Død.“ Men den Anden svarer: „Jo vist
er Redning mulig; skulle vi end gaae igjennem Ilden, en Redning er dog
mulig.“ Dette Svar skal da være Fornuftmennesket en daarlig Tale; han vil
inderlig harmes over et saa taabeligt Spil med Muligheder: „Her, hvor
Flugt er umulig, hvor Barmhjertighed er umulig, hvor det er saa umuligt,
at Dødsstraffen forandres, som det er aldeles umuligt, at den brændende
Ild forandrer sin Natur: hvorledes kan dog Nogen
% --- p. 33 ---
her ville lege med Muligheden og drive Spot med Skjæbnen?“ Men den Anden
skal svare ham og sige: „Kjære, glem dog ikke, hvad du selv sagde, da vi
tiltraadte denne Reise: Menneskets Liv er i Guds Haand.“ Saa frelses de
da begge, og ere enige om \dots\ ja hvad er det nu, de egentlig talt ere
enige om? Begge ere vel enige om at bruge deres naturlige Forstand,
forsaavidt det gjælder at iagttage de naturlige Forhold og beregne
Sandsynlighederne. Heri behøver den Ene slet ikke at staae tilbage for
den Anden. Men i at anerkjende Styrelsen kunne de nok neppe være enige.
Den Første seer kun paa Naturloven, og forseer sig paa det Sandsynlige i
den Grad, at han i ethvert afgjørende Tilfælde forglemmer Gud; den Anden
har vel ogsaa et verdsligt Øie for de naturlige Betingelser; men i
Afgjørelsens Øieblik sammentrænger han, saa at sige, sin hele Seekraft i
det Øie, hvormed han seer efter Gud, saaledes, at han derved kommer til
at oversee Sandsynligheden. Hiin har et mat og taaget Begreb om en
guddommelig Almagt, som han i Gjerningen fornegter, hver Gang han just
skulde anvende det; denne har et saa levende Indtryk af den Almægtiges
Allestedsnærværelse, at ingen af Endelighedens Magter formaaer at
overvælde ham. Med andre Ord: Den Første har egentlig slet ingen Tro; den
Anden derimod er en Troende.

Denne Forskjel viser sig paa det Allerklareste, naar Opgaven er at
forvente et endeligt Gode, der skal modtages som en Guds Gave. Om de end
begge To nære det samme Ønske, og enes om, hver især, at ville henstille
Sagen til Gud; de kunne dog ikke længe følges ad: deres Veie skilles.
Fornufttroen gaaer ikke længere, end Sandsynligheden og de rimelige
Udsigter tillade det. „Er Ønskets Opfyldelse en plat Umulighed; kan
% Printer's defect: the „e“ of „Menneske“ (broken at the line end as
% „M|n-neske“) failed to ink; at 600 dpi only a faint trace of the sort
% prints, and both OCR witnesses read a blank. Set as „Mennesket“ requires,
% i.e. „Menneske“; the impression, not the setting, is at fault.
ethvert fornuftigt Menneske see, at det ikke bliver til Noget, saa er der
jo Intet videre ved den Sag at gjøre. Istedetfor at besvære Gud med
urimelige
% --- p. 34 ---
Bønner, gjør en fornuftig Mand da meget bedre i at opelske et andet
Ønske, og see sig om efter andre opnaaelige Goder.“ Evangelietroen
derimod kan ikke tillade en saa rig Afvexling i Ønsker og Forhaabninger.
Har den Troende først forstaaet sit Ønskes Betydning og henstillet Sagen
til Gud; da har han med det Samme bundet sig selv ved Ønsket, og maa nu
til Trods for alle Omstændighedernes og Sandsynlighedernes
Veirforandringer blive staaende, og vente paa, hvad Gud vil gjøre. Ved
Hjælp af Ønsket lærer Mennesket at forstaae sin Gud, idet han lærer at
forstaae sig selv; og derom er det jo ogsaa, han beder. Men vilde det nu
ikke være formasteligt, om en Bedende begyndte saaledes: „Almægtige Gud!
Jeg har hidtil haabet, at du skulde opfylde mit kjæreste Ønske; men nu,
da Omstændighederne have forandret sig, og da Tiden er skredet saa langt
frem, at du umulig \emph{kan} opfylde det, saa vil jeg nu ogsaa lade den
Sag fare, og bede dig en anden Bøn, som du maaskee endnu kan opfylde.“ At
en saadan Bøn er en Gudsbespottelse, kan dog vel ikke betvivles. Thi hvis
Ønsket var saa ubesindigt, at den Bedende strax maatte kunnet indsee dets
Urimelighed: hvorledes turde da Mennesket fordriste sig til, at
fremkomme med et saadant Ønske for Gud? Men er det kun verdslige
Usandsynligheder, der staae i Veien for Opfyldelsen: hvorledes tør da
Nogen vove at tale til den Almægtige om Umuligheden af at opfylde det?
Nei, det er noget ganske Andet, hvorpaa det her kommer an. Den Bedende
holder sig forvisset om, at Gud baade kan og vil bønhøre ham, saa sandt
Bønhørelsen er til hans Bedste; men hermed er da tillige givet, at
Ikke-Bønhørelsen ogsaa kan være til hans Bedste. Dette er det just, som
gjør Opgaven vanskelig. Hvilket af de tvende Tilfælde her muligviis
finder Sted, kan Tiden vel aldrig lære den Troende; men denne kan og skal
lære det med Tiden, idet han lærer at bie paa Gud og forstaae sig selv i
Forhold
% --- p. 35 ---
til den Almægtige. Dersom en saadan Venten nu hører med blandt Troens
væsentlige Fordringer, saa kunne vi jo nok forstaae, hvad der opholdt
Zacharias ved det ene Ønskes Opfyldelse, saaledes at han i en vis
Forstand slet ikke mærkede, at Omstændighederne forandrede sig, at Aarene
gik, og Alderdommen kom; men blev ved med det samme Ønske og den samme
Bøn, endnu da en heel Slægt var gaaet ham forbi, endnu da han maatte sige
sig selv: „jeg er gammel, og min Hustru er ogsaa bedaget.“

Det er just ikke saa let, som den vrængende Kritik foregiver, at
henstille sin Sag til Gud, og bie efter den Almægtige. Nei, da er det
meget lettere at gaae frem med Tiden. Thi Den, der gaaer frem med Tiden,
bæres af Omstændighederne, og drives af Sandsynlighederne. Men den, der
skal bie paa Gud, kan i Endelighedens Forstand end ikke staae stille (thi
da vilde Endeligheden rive ham med sig); men for at blive paa sin Post og
bie, maa han uophørlig brydes med de paatrængende Omstændigheder og
arbeide imod Sandsynlighedernes Strøm. Det er saa betryggende at gaae
Sandsynlighedernes brede Vei: finder jeg ikke et Fodfæste her, saa finder
jeg det jo der; selv om jeg skulde falde, blev Skaden vel ikke større,
end at jeg kunde staae op igjen: der er ikke saa langt til Jorden.
Derimod er det meget ængsteligt, at gaae ad en smal Stie hen over en
Afgrund; det mindste Feiltrin kan bringe Døden. Sandsynligheden for at
falde er saa stor, at jeg slet ikke tør ændse den, for ikke at svimle;
jeg tør slet ikke see til nogen af Siderne, men maa uafbrudt see
ligefrem, og trænger høilig til En, der kunde gaae foran og sige: see paa
mig!

Fornufttroen, der holder sig paa Sandsynlighedernes vidtstrakte
Overflade, kan sagtens skride fremad med sikkre Fjed, thi den brede Jord
er jo fast nok til at bære et Menneske. Men, netop fordi den bærer saa
godt, og Alt gaaer som af sig selv efter Tyng\-%
% --- p. 36 ---
dens og Ligevægtens Love, er Almagten her jo næsten overflødig.
Evangelietroen derimod har i Endeligheden aldrig mere Plads tilovers, end
et Menneske netop hver Gang behøver til at støtte Foden paa. Derfor maa
den Troende vel vogte sig for at see for dybt i Mulighederne og vende sig
for ofte efter Sandsynlighederne; i al Varsomhed maa han gaae sin Vei
ligefrem, og vedblive at see paa den Almægtige, for ikke at svimle. Det
er dog kun en saare lille Deel af Jordens Flade, et Menneske behøver til
et Fodfæste. Det Sted, hvorpaa Foden hviler, kan Øiet ikke see; det Sted,
Øiet seer, har Foden endnu ikke betraadt. Den, der gaaer, maa dog
alligevel have Øinene med sig, og see sig for, om det Sted, han vil træde
paa, nu ogsaa kan bære. Væsentlig talt, kan Øiet vel ikke med Bestemthed
forudsee, hvor paalideligt Stedet er; men saa føler man sig for og
forlader sig paa sit Skjøn: den naturlige Gang gaaer efter
Sandsynligheden. Vil Nogen derimod danne sig en Forestilling om,
hvorledes det vel kunde være at gaae lige imod Sandsynligheden; da er det
vel saare vanskeligt at beskrive en indvortes Bevægelse ved et udvortes
Billede: vi ville imidlertid anstille følgende Forsøg. Antag, at et
Menneske pludselig forandredes saaledes, at Jorden blev gjennemsigtig for
hans Øie, endnu mere gjennemsigtig end den stille, klare Sø, der i Dybet
afspeiler Skyerne og den blaae Himmel. Antag, at det samme Menneske
endvidere blev forvandlet saaledes, at den gjennemsigtige Jord ikkun da
vilde fastne under hans Fod og bære ham, naar han med besluttet Sind
traadte dristig til, men derimod opløste sig og veg tilbage, flygtig og
troløs som den rene Luft, hver Gang han i ængstelig Vaklen famlede efter
et Fodfæste. Saa havde vi da her en Vandringsmand, der ikke kunde gjøre
et eneste Skridt uden at komme i Strid med Sandsynligheden. Ja, det er en
% Printer's error, transcribed as printed: a comma stands where the sense
% requires a full stop, before the capitalised „See“ that opens the next
% sentence. Checked at 600 dpi on the last line of the page: the mark has
% the same descending tail as the comma after „See“ two words later, and
% is quite unlike the full stop after „Sagen endt.“ on p. 31. Both OCR
% witnesses read a comma.
underlig Stilling, See, han staaer paa den faste Jord, som alle de An\-%
% --- p. 37 ---
dre; og dog synes det ham, som han stod paa en Luftsøile imellem Himmel
og Himmel. Jorden kan virkelig bære ham, ligesaavel som alle de Andre,
naar han kun træder dristig til; og dog seer det ud i hans Øine, som
skulde han ved et eneste Skridt styrte dybt ned i --- Skyerne forneden,
ja igjennem Skyernes Aabning ned i det blaae uendelige Dyb. Men just som
Kraften forlader ham, og det vil sortne for hans Øine, lyder en sagte
Røst: „Frygt ikke, træd kun dristig til, sæt Foden her; men vakler du, da
vee dig, thi da er du Dødsens!“ Saa gjør han da ethvert Skridt lige imod
Sandsynligheden; men det er en langsom Gang, det er, som om Aandedraget
vilde standse, hver Gang han flytter en Fod.

De Andre kunne nu sagtens træde dristig til; thi for dem har den mørke
Jord jo ingen Gjennemsigtighed, og giver intet Speilbillede af den dybe
Himmel; de Andre kunne sagtens gaae rask frem; thi for dem har det ingen
Fare, at den faste Jord skulde opløse sig i Luft: hvorhen de see, finde
de overalt et sikkert Fodfæste, ja alle de Andre kunne jo sagtens ile vor
Vandrer langt forbi; thi det er i endelig Forstand langt lettere at gjøre
hundrede Skridt efter Sandsynligheden uden Almagtens udtrykkelige
Bistand, end et eneste Skridt imod det Sandsynliges Lov i ubetinget
Tillid til den Almægtige. Men er det desuagtet Evangelietroens
uafviselige Fordring, at et Menneske gaaer sin Vei, selv imod al
Sandsynlighed, da er Zacharias et Troens Forbillede; thi en saadan Gang
gik just den gamle Præst, der han gik ind i det Hellige, at staae for
Guds Engel.
% The last word of the paragraph is „at“, not „al“: both OCR witnesses read
% „al“, but at 600 dpi the second sort is plainly the short ascender of „t“
% (lightly inked crossbar), well below the height of the „ll“ of „Hellige“
% earlier in the same line, and unlike the full ascender of „al“ in „imod al
% Sand-“ on the line above.

Oplysningen vil imidlertid vide bedre Beskeed om det Sandsynlige, og
kaster sig derfor med alle sine Indsigelser over Gabriel og hans
underfulde Bebudelse. „Hvad der i nærværende Fortælling synes at
retfærdiggjøre det fromme Pars ubetingede Tro, og giver deres Forventning
et saa opbyggeligt Anstrøg, er dog i
% --- p. 38 ---
Grunden intet Andet end Engle-Aabenbarelsen selv med det derpaa følgende
høist glædelige Udfald. Efter denne Anviisning maa det da vel heller ikke
være saa besværligt for en Troende, at indlade sig i Strid med den
naturlige Omverden, da han jo i Nødsfald altid kan gjøre Regning paa en
Engels Bistand. Imod en saadan Bistand have vi, hvad den evangeliske
Forsynslære selv angaaer, Intet at indvende. Tvertimod, Englene ere jo
„tjenstagtige Aander, der udsendes til Tjeneste for dem, som skulle arve
Saliggjørelsen“. Da Alt nu desuden er saa viselig indrettet, at Englenes
Tal, der staae for Guds Throne, er titusinde Gange Titusinde og langt
derover, medens de sande Troendes Antal paa Jorden altid maa anslaaes til
„en saare liden Hjord“: saa vil der i enhver Tidsalder altid være langt
flere Engle, end der ere Troende, og „Tjenesten“ kan følgelig aldrig
falde de Himmelske besværlig. Forholder dette sig nu saaledes, saa er
Troen unegtelig den bedste Livsassurance, noget Menneske kan vælge sig,
og en troende Sjæl kan da sagtens byde Sandsynlighederne Trods. Men hvad
der paa den anden Side gjør Sagen noget betænkelig, er den besynderlige
Omstændighed, at man i vore Dage, da selv den mindste Mærkværdighed
sjeldent undgaaer den almindelige Opmærksomhed, slet ikke mærker det
Allermindste, enten til Englene eller deres Mirakler. Hænder det sig en
Gang, at et Menneske træder frem og beraaber sig paa en umiddelbar
Aabenbarelse eller et nyt Mirakel, saa ere alle fornuftige Folk jo strax
enige om at sigte ham for Sværmeri. Det er ikke blot de Letsindige, der
spotte ham, og de Vantroe, der holde sig op over ham; men selv de
Rettroende ryste paa Hovedet, og see paa hinanden med høist betænkelige
Miner: Herregud, hvad skal man da tænke? Tvende Tilfælde ere mulige; men
hvilket af disse man end vil antage, saa gaaer det alligevel ud over
Mirakeltroen. Enten har Gud siden hine Dage aldeles forandret sin
% --- p. 39 ---
Styrelse, og givet Englene en ganske anden Tjeneste, og da følger det
naturligviis af sig selv, at Mirakeltroen ikke længere er berettiget;
eller ogsaa maae Miraklerne være ophørte, og Englene selv have trukket
sig tilbage, fordi den rette Tro er forsvunden af Jorden. Men skulde det
nu ikke være et stort Beviis imod Mirakeltroens guddommelige Oprindelse,
at den saadan reent kan forsvinde af Jorden?“ --- Saavidt den moderne
Oplysning.

I Aabenbaringens Navn er Evangelietroen sikkret imod enhver formastelig
Klogt. Evangelisten har ikke fortiet os, at Zacharias forfærdedes, og
Frygt faldt paa ham, der han saae Engelen staae ved den høire Side af
Røgelsens Alter. Ja saa lidt har den fromme Mand været beredt paa det
Under, der forkyndes ham, at han endog i Overraskelsens Øieblik forsynder
sig imod den Hellige.
% Luc. 1, 20 is set in the same larger type as the pericope on pp. 1--2,
% without quotation marks, but it is not a display block: it begins a fresh
% line at full measure and ends in the middle of a line, where the body size
% resumes with „Saaledes“. The citation „(Luc. 1, 20).“ is in the larger
% type too. Hence inline {\large ...}.
{\large See, du skal vorde stum, og ikke kunne tale, indtil den Dag, da
dette skeer, fordi du ikke troede mine Ord, som skulle fuldkommes i deres
Tid (Luc.\ 1, 20).} Saaledes straffer Engelen den Tvivlende, og med denne
Straf maa vel Zacharias frifindes for den overtroiske Letsindighed, som
mener, at der hvert Øieblik vil skee et Mirakel. Om Bønnen end ikke var
bleven opfyldt paa denne Maade, skulde den Troende da maaskee have
forbrudt sit Haab og anseet sit Liv for spildt, og opgivet Troen paa det
Sidste? O nei, dertil var den Bedagede altfor inderlig forvisset om sin
Fortrøstnings Sandhed og Guds ubeskrivelige Godhed. Med denne Tro vilde
den modige Patriarch være gangen til Hvile, med denne Tro vilde han have
takket og priset den Almægtige i sin sidste Stund, ja med denne salige
Tro vilde han da velsignet sin Hustru, og smilet i Døden, naar Englene
kom at bære ham hen i Abrahams Skjød. Men det forstaaer sig, saa var
Opfyldelsen bleven en anden, saa var Døden kommen før Englene; nu derimod
kom Engelen før Døden, og overraskede ham med det glade
% --- p. 40 ---
Budskab, og forfærdede ham med sin Tale. Det er ganske rigtigt, hvad
Oplysningen paastaaer, at Troen forudsætter Underet; thi den drager sin
Næring deraf, og er selv et sjæleligt Under. I denne Forstand kan
Evangelietroen med Rette kaldes en Mirakeltro. Det er ligeledes vist, at
et Menneske, der lever oprigtig i denne Tro, maa i mange Henseender
forekomme Verden noget underligt. Hvi skulde ikke Den synes underlig, der
har sit egentlige Liv i et Under! Men heraf følger dog ingenlunde, at den
Troende skulde grunde paa vilkaarlige Mirakler, og gjennemstrøife
Tilværelsen i den fade Indbildning, at han maaskee etsteds her eller der
kunde træffe paa en Engel eller en ny Aabenbarelse. Higer Nogen
utaalmodig efter nye Aabenbarelser, da er en slig Higen saa langt fra at
bevise hans Troes Styrke, at den endog beviser det Modsatte, eftersom den
jo vidner om, at han endnu ikke har lært at tilegne sig den een Gang
givne Aabenbaring. Forsaavidt kan Evangelietroen meget vel bestaae med
Naturloven, og vil bestaae indtil Verdens Ende. Eller skulde
Verdensaanden vel have Grund til at føre Klage over Aabenbaringens
forstyrrende Mirakler? See, Riger opstaae --- uden Mirakler; de forgaae
igjen --- uden alle Mirakler. Throner reise sig --- uden Mirakler; de
styrte sammen igjen --- uden alle Mirakler. Saaledes gaaer det i
Nutiden; men saaledes gik det jo ogsaa i Keiser Augusts og Kong
Herodes's Dage. Verden er i det Hele stærk bygget, og Menneskeheden har
store Kræfter: om ogsaa ti Riger faldt sammen ti Gange; Slægten vilde
være stærk nok til at bygge dem op alle ti den ellevte Gang. Kun om eet
Rige gjælder det, at den hele Verden ikke har Kræfter nok til at bygge
det op fra Grunden af; om eet Rige gjælder det, at dersom det nogensinde
kunde styrte sammen, da vilde alle Timelighedens Magter være for svage
til at reise det af dets Fald, og det er Guds Rige, Sandhedens, Ret\-%
% --- p. 41 ---
færdighedens og Hellighedens Rige. Naar dette Rige skal grundlægges paa
Jorden, da maa der skee et Under, et stort Under, da maae Himmelens
Kræfter røre sig, og ---
% The re-quotation of Hebr. 1, 14 is letterspaced (body size), not enlarged;
% the sentence that follows it, Gabriel's word from Luc. 1, 19, IS enlarged
% and not letterspaced. Both readings confirmed on the image at 200 dpi;
% spacing.py flags the first as a RUN and not the second.
\emph{Englene udsendes til Tjeneste for dem, som skulle arve
Saliggjørelsen.} Derfor lyder Gabriels Røst i det Hellige: {\large Jeg er
udsendt at tale til dig, og forkynde dig dette til Glæde.}

Store Forandringer, mærkelige Fremskridt kunne foregaae i Folkenes
Historie, uden alle Mirakler. Tidens Aand har Arbeidere nok: hvor den Ene
ender, begynder den Anden; hvad den Ene lover, opfylder vel den Anden;
hvad den Ene forbereder, kan maaskee fuldføres af den Anden. Saaledes
gaaer Verden jo fremad. Forgjængere fødes i Hobetal; de vække, de lede,
de lære Folkene, og pege mod „den store Fremtid“ og „det store Maal“. Og
Folkene reise sig, og --- føle sig kaldede, og Helte opstaae, og føle sig
kaldede; og Bevægelsen begynder, og Omvæltningen skeer, men uden alle
Mirakler. Og der strides og seires til Lands og til Vands, men --- uden
alle Mirakler; og der strides og seires med Ord og med Sværd, men ---
uden alle Mirakler. Og Opvækkelsen ender, og Seiren skuffer, og
Forgjængerne modsige og anklage hinanden, og Heltene bekjæmpe og styrte
hinanden, og Folkene kives og mistrøste hinanden, og Jubelen forstummer,
og Bevægelsen standser, og det hele „Underværk“ er igjen som et tomt,
møisommeligt Intet, --- uden alle Mirakler. Alt dette tør Verdensaanden
vove; thi Slægtens moderlige Natur har Kræfter nok, og der ødsles med
naturlige Kræfter. Men Eet er der dog, som den moderlige Natur ikke
formaaer, om den end ødslede sine herligste Gaver til nok saa mange
Begavede: den kan ikke udruste nogen Mand saa stærk, at han ved sit Ords
Kraft skal kunne vække Menneskenes sovende Samvittighed, og føre en
vanartig Slægt tilbage til Lydighed under Guds Lov. Skal en saadan Mand
udrustes, maa Kraften gives ham herovenfra. Men
% --- p. 42 ---
med Kraften fra oven holder Viisdommen til Raade: der ødsles ikke med de
overnaturlige Kræfter. Da derfor Guds Søns Komme skulde forberedes, var der
kun een Mand, der blev kaldet til at udføre Forberedelsens store Værk. „Kun
een Mand“? Ja, men denne ene Mand blev ogsaa tilstrækkelig udrustet fra
Begyndelsen af; der blev sørget for, at hans Kald og Sendelse rettelig maatte
fattes og forstaaes fra Begyndelsen af; derfor udsendes Bebudelsens Engel,
derfor lyder Gabriels Røst i Helligdommen, som det himmelske Sendebuds
menneskelige Røst:
% Luc. 1, 15--17 is set in the same larger type as the pericopes on pp. 1--2,
% but it BEGINS mid-line, immediately after the colon, and ends the paragraph
% mid-line. Hence inline {\large ...}, as at p. 41 (Luc. 1, 19), not
% \begingroup\large ... \par\endgroup.
{\large Han skal være stor for Herren; Viin og stærk Drik skal han ei drikke,
og han skal fyldes med den Hellig-Aand alt fra Moders Liv, og omvende mange
af Israels Børn til Herren deres Gud. Og han skal gaae frem for ham i Elias'
Aand og Kraft, at vende Fædrenes Hjerter til Børnene og de Gjenstridige til
de Retfærdiges Sind, at beskikke Herren et vel forberedt Folk.}
% „Elias'“ takes a bare apostrophe here (verified at 600 dpi); the parallel
% passage on p. 2 prints „Elias's“. Both readings stand; not normalised.

% The head of chapter II falls part-way down p. 42, in this order: a DOUBLE
% rule (two hairlines, one above the other), the numeral „II.“ on its own
% line, the head, the scripture reference, and a second, much shorter single
% rule. Printed widths measured at 600 dpi against a 90 mm measure: the upper
% double rule c. 34 mm (0.38 of the measure), the lower c. 11 mm (0.12). The
% double rule is reproduced here as a single \rule, as elsewhere in this file.
\begin{center}
  \rule{0.38\textwidth}{0.4pt}
\end{center}

% The printed head agrees with the Indhold („II. Engelen Gabriel og Jomfru
% Maria ... 42--80“). As in chapter I, the numeral stands on a line of its own
% above the title, the printed full stop is kept in the \section* and dropped
% in the \addcontentsline.
\section*{II.\\[4pt] Engelen Gabriel og Jomfru Maria.}
\addcontentsline{toc}{section}{II. Engelen Gabriel og Jomfru Maria}

% The reference line is centred under the head in a type smaller than the
% body, like the epigraph on p. 1; hence \footnotesize. The book prints „Luk.“
% here and „Luc.“ on p. 53; the inconsistency is transcribed, not normalised.
\begin{center}
  {\footnotesize (Luk.\ 1, 26--38.\ \ Matth.\ 1, 18--25.)}\\[10pt]
  \rule{0.12\textwidth}{0.4pt}
\end{center}

\begingroup\large
I den sjette Maaned af Elisabeths Tid blev Engelen Gabriel sendt af Gud til
Jomfru Maria i Nazaret med det underfulde Budskab, at hun var den Benaadede,
den iblandt Qvinderne Velsignede, der skulde føde Verdens Frelser. Og
Jomfruen forfærdedes over hans Tale, og
% --- p. 43 ---
forundredes over hans Hilsen. Men, der Engelen med guddommelig Røst
forkyndte hende: Den Hellig Aand skal komme over dig, og den Høiestes Kraft
skal overskygge dig; derfor skal ogsaa det Hellige, som fødes af dig, kaldes
Guds Søn: da gav Maria sig hen som Herrens Tjenerinde, og Engelen skiltes
fra hende. --- Men Jomfru Maria var trolovet en Mand ved Navn Joseph af
Davids Huus, og Joseph miskjendte sin trolovede Hustru. Men, efterdi han var
en skaansom, om end streng og retfærdig, Mand, vilde han ikke beskjæmme
hende aabenbart, men besluttede, at skille sig hemmelig fra hende. Men, idet
han tænkte derpaa, see, da aabenbaredes Herrens Engel for ham i en Drøm, og
sagde: Joseph, Davids Søn! frygt ikke for at annamme din Hustru Maria; thi
hendes Foster er af den Hellig Aand. Men hun skal føde en Søn, og du skal
kalde hans Navn Jesus; thi han skal frelse sit Folk fra deres Synder. Og
Joseph vaagnede af sin Drøm, og fornegtede sig selv, og gjorde, som Herrens
Engel havde befalet ham.
\par\endgroup

Den inderlige Sammenhæng imellem denne og den foregaaende Fortælling indlyser
af sig selv. Den samme Engel, der aabenbarede sig i Templet, udsendes ogsaa
til Nazaret; Tiden, da han viser sig for Jomfruen, bestemmes efter det
Tidspunkt, da Zacharias fik sin Aabenbarelse; ja selv det Under, der foregik
med Elisabeth i hendes Alderdom, bliver endog ligefrem udlagt for Maria som
et Tegn paa, at \emph{ingen Ting skal være umulig hos Gud}. Den, der føler
sig tiltalt af den første Evangelieberetning, kunde da maaskee tænke ved sig
selv, at denne Bebudelseshistorie uden videre var hævdet og forklaret i hiin,
saa at man sagtens kunde troe denne, naar man først havde gjort sig fortrolig
% --- p. 44 ---
med hiin. Imidlertid turde dog en saa godtroende Forudsætning let underligge
en noget betænkelig Misforstaaelse. Vel have vi i den foregaaende Udvikling
allerede fjernet adskillige Misligheder, der ellers her vilde stille sig i
Veien for en troende Tilegnelse af dette Evangelium; vi have allerede
forstaaeliggjort os Troens Selvbekræftelse imod Forargelsens Anløb og
Tvivlens Underfundighed; men det gjælder i Forhold til den hellige Historie,
hvad der viser sig i det religiøse Liv, at Troens Strid aldrig standser.
Underets Gud er altid den samme; Aabenbaringens Formaal er væsentlig altid
det samme; men i enhver ny og overordentlig Begivenhed antager det
Vidunderlige en ny overraskende Form, og Tvivlen, der ikke blot som en
udvortes Fjende angriber den Troende udenfra, men tillige som en indvortes,
hemmelig Fjende lister sig frem i hans egen Sjæl, vil stedse gribe enhver ny
Anledning til at pukke paa Underets Urimelighed, og ved en uformodet Vending
forraske den Trygge. Skjøndt den hellige Historie i allerhøieste Forstand er
Aandens Historie, er, for at tale i Philosophiens Sprog, et organisk Hele, en
sammenhængende Kjæde af evige Ideer: er det dog saa langt fra, at en følgende
Undergjerning skulde kunne afledes af en foregaaende, at Troens Tilegnelse af
enhver overordentlig Begivenhed meget mere maa begynde forfra, for at stride
med Tvivlen om \emph{Alt og Intet}. Hver Gang en enkelt Evangelieberetning
trøster, opbygger og glæder os, er det jo, som om der gik et Lys op over hele
Evangeliet; thi vi ere i Troen, og besidde det Hele i det Enkelte. Men
omvendt, naar en enkelt Fortælling frastøder og forarger os, da er det jo
ogsaa, som om den hele Aabenbaring skulde opløse sig i Intet: vi tabe det
Hele med det Enkelte; thi vi tabe --- Troen. Dette finder nu en særegen
Anvendelse paa Forholdet imellem de tvende Bebudelser. Just fordi Underet i
„Maria Bebudelse“ er større
% --- p. 45 ---
og, saa at sige, mere underfuldt, end i den foregaaende Fortælling, just
derfor skjærper ogsaa Tvivlen sin Braad, saa at den indtrængende Tilegnelse
udkræver desto større Anstrengelse. For at overbevise os herom, behøve vi nu
kun et Øieblik at lytte til Fritænkernes Spot, og lade den negative Kritik
komme til Orde.

„Fortællingen om Englebesøget hos Maria er,“ for at tale i Fritænkerens Aand,
„dog unegtelig en af de allertroskyldigste Fortællinger i det N. Tst. Hvor de
Sjæle maae være stærke i Troen, der, uden at ane Vanskeligheden, kunne
sammensætte saa ulige Forestillinger, som: \emph{Gud i Himlen og Provindsen
Galilæa, Erkeengelen Gabriel og Flækken Nazaret, den Hellig Aand og Jomfru
Maria}. (Slige forbausende Sammensætninger af det Himmelske og det Jordiske
% Printer's defect: in „Sammensætninger“ the t of -sætninger printed only the
% tip of its top stroke --- the sort is broken or failed to take ink (checked
% 600 dpi). Both OCR witnesses read a gap. Transcribed as the word intended.
ere dog just ikke altid Beviis paa overvættes Styrke i Troen; de kunne
undertiden ogsaa forklares af en overvættes Svaghed i Tænkningen.) Navnene
\emph{Galilæa, Nazaret, Maria} etc. have nu i Aarhundreders Løb erholdt
Privilegium paa at sættes i umiddelbar Forbindelse med Gud og hans Engle. Men
slige Egennavne kunne jo let tænkes forandrede, uden at Sagen selv i mindste
Maade forvanskes. Det gjælder blot om, at vælge saadanne Benævnelser, som
kalde os det oprindelige Indtryk tilbage, derved, at de uvilkaarlig stemple
Tingen og Personen med Virkelighedens rette Hverdagspræg. Foretag da blot en
ganske simpel Navneforandring. Vælg kun Navnene paa Provindsen og Byen af
dine nærmeste Omgivelser; giv Jomfruen et mere jordisk, mindre æsthetisk
Navn: oplæs saa med præstelig Pathos de to første Vers af dit Evangelium, og
du skal fornemme, at det endda ikke er saa let at bringe det Overnaturlige i
en saa udvortes og umiddelbar Berøring med det Naturlige.“

„Men,“ siger man med Gabriel, „hos Gud skal ingen Ting
% --- p. 46 ---
være umulig. Troen er forberedt paa Alt, selv det Usandsynligste. Velan da!
Siden den gamle Mirakeltro holder sig saa godt, at man endnu i vort 19de
Aarhundrede offentlig tør fremsætte den Lære, at Maria har „undfanget“ sin
Søn af den Hellig Aand; maae vi saa bede om lidt nærmere Oplysning angaaende
denne Jomfru og hendes mythiske Forhold til sin Trolovede. Evangelisterne
Lukas og Matthæus ere, hvad disse delicate Punkter angaaer, temmelig
ordknappe; det træffer sig da ret heldigt, at vi have udførligere Traditioner
omtrent fra samme Tid og, efter Kirkefædrenes Opfattelse at dømme, omtrent i
samme Aand, som vore tvende kanoniske. Lader os da med Omhyggelighed
conferere de to ældgamle Kildeskrifter, nemlig: Evangeliet om Marias
vidunderlige Fødsel (\textit{Ev.\ de nativitate Mariæ}) og Apostelen Jacobs
Protevangelium; maaskee kunde der ved disse Bestræbelser opgaae et klarere
Lys over Sagens sande, mirakuløse Sammenhæng.“

„Jomfru Maria, der havde den overordentlige Bestemmelse, at blive Moder til
Guds Søn, kan naturligviis ikke have været nogen almindelig Jomfru. Man er
altsaa berettiget til at vente, at hun selv maa være kommen til Verden ifølge
en særegen guddommelig Foranstaltning, og fra Barndommen af have levet et til
Gud og hans Engle indviet Liv. Denne Formodning stadfæster sig da ogsaa, naar
vi kun raadføre os med begge de anførte Smaa-Evangelier. Joachim og Anna, et
Par gamle, meget gudfrygtige Folk, boede i Nazaret, og bleve omsider
miskjendte og forhaanede paa Grund af deres lange, barnløse Ægteskab. Men i
denne deres kummerfulde Stilling nøde de nu den store Naade, at Herrens Engel
aabenbarede sig for dem begge, hver især, og lovede dem en Datter, der skulde
kaldes Maria, og vorde udseet til at føde Verdens Frelser. Da Maria var 6
Maaneder gammel,
% --- p. 47 ---
satte Moderen hende en Gang ned paa Jorden; men Barnet gik kun 7 Skridt, og
søgte igjen tilbage til sin Moders Barm. Dette udlagde saa Anna meget viselig
som et Tegn paa, at hendes gudindviede Datter ikke vilde træde paa Jorden,
før hun var fremstillet i Herrens Tempel. Fra sit 3die til sit 14de Aar
levede Maria da i Templet, hvor hun hver Dag nød guddommelig Aabenbarelse;
thi Englene besøgte hende, ja hun modtog endog sin Føde af en Engels Haand.
% Printer's defect: a raised space took ink between „af“ and „en“, so that the
% line reads „af-en“. Both OCR witnesses record the mark; at 600 dpi it is a
% blind quad, not the „=“ this book uses for a hyphen. Set here as „af en“.
--- Dette maa nu være os nok, hvad det første Punkt angaaer.“

„Dernæst ønskede vi paa Troens Vegne, at kjende de Vilkaar, under hvilke
Trolovelsen fandt Sted, lidt nærmere; da den Maade, hvorpaa Engelen, ifølge
de kanoniske Beretninger, forvandler den virkelige Forbindelse til en blot
tilsyneladende, ophøier Maria til en \textit{Noli me tangere} og degraderer
Joseph til titulær Ægtemand med Pligter uden Rettigheder, ellers let kunde
støde de Svagere i Troen som en vilkaarlig og Gud lidet værdig
Fremgangsmaade. Ogsaa i denne Henseende gjøre de to anførte Skrifter god
Tjeneste. Det indlyser nemlig saavel af Evangeliet om Marias Fødsel, som af
det tilsvarende Protevangelium, at det aldrig har været enten Josephs eller
Marias Hensigt, at indlade sig i nogen ordinær Ægteforening med hinanden.
Først efterat Ypperstepræsten Zacharias havde adspurgt Gud i det
Allerhelligste, maatte Maria, ifølge det guddommelige Orakel (Jes.\ 11, 1)
indvillige i, at der blandt Mændene af Davids Huus valgtes hende en
tilsyneladende Ægtefælle eller idetmindste en Formynder. Det blev da bestemt,
at de giftefærdige Mænd af den davidiske Familie eller, nøiagtigere efter
Protevangeliet, alle Enkemænd i Folket skulde bringe deres Stave til Alteret,
at det ved et Jertegn (4 Moseb.\ 17) kunde vise sig, hvo der var den Værdige.
Joseph var den Gang allerede Olding, og vilde i Begyndelsen ikke komme
% --- p. 48 ---
frem med sin Stav, hvisaarsag heller intet Tegn viste sig. Ypperstepræsten
maatte da atter adspørge Gud, og da Josephs Stav var overleveret, skete
Jertegnet, idet en Blomst udskjød af Staven, og en Due, der kom ned fra
Himlen, satte sig paa Enden af den. --- See, dette er da en Forklaring
ligesaa opbyggelig, som troværdig. Skal man troe det store Under, hvorfor da
ikke ogsaa de mindre? Hos Gud skal, som Gabriel siger, ingen Ting være
umulig. Troen er jo forberedt paa Alt, selv paa det Usandsynligste.“

Det er ingenlunde saa let for de høilærde Theologer, at optage den
Stridshandske, den negative Kritik med dristigt, udfordrende Blik tilkaster
dem paa Videnskabens Kampplads. At tage sin Tilflugt til den naturlige
Fortolknings Kunstgreb, hjælper ikke; thi, naar Engelen forvandles til en
naturlig Elsker, da kan den Forargelse ikke afværges, at Maria enten med
Forsæt eller i Sværmeriets Vildelse har bedraget Joseph, og at Christus, den
Usyndige, der skulde frelse Verden fra Synd, altsaa selv er undfangen i Synd
og født i Misgjerning. Men oprøres man nu paa Religiøsitetens og
Sædelighedens Vegne over en saa lav Mishandling af den ærværdige
Overlevering, uden dog at kunne beqvemme sig til at antage Underet, saaledes
som vore kanoniske Evangelier fremstille det; da er der kun een Udvei
tilbage, den nemlig, at erklære hele Fortællingen for en Mythe. --- At den
opløsende Kritik især udvælger sig denne Fortælling som et flaaende Exempel
% Printer's defect: „flaaende“ as printed, with a full-crossbar f (verified at
% 600 dpi against the long s of „Misforstaaelser“ on p. 53); the sense requires
% „slaaende“. Transcribed as printed, not corrected.
paa, at Mythen ogsaa hører hjemme i det N. Tst., er ganske i sin Orden; men
selv saadanne Theologer, som med al deres videnskabelige Kritik dog søge at
bevare, hvad de kalde Troens Kjerne, have i en Tidsalder, der saa let
forvexler Aabenbaringens hellige Aand med den rene Verdensidee, og den
troende Tilegnelse med Philosophiens Viden, ikke kunnet modstaae Fristelsen.
Medens Underet nemlig, som det eneste i sit Slags, har al Analogie imod sig,
har Mythen
% --- p. 49 ---
derimod i nærværende Tilfælde en rig Mangfoldighed af Oldtids Analogier for
sig. Helte, som Herakles, Romulus og Alexander; Viismænd, som Zoroaster,
Pythagoras og Plato, forestilles jo i Mythedigtningen som jomfrufødte
Gudesønner. Mythens Begreb er desuden saa bøieligt, tvetydigt og svævende, at
det tillader de allerforskjelligste Nuanceringer: Ideal og Virkelighed, Fabel
og Kjendsgjerning kunne her blande sig, og indgaae Forbindelser efter meget
ulige Maal. Mytheformen er egentlig den Paaklædning, hvormed Phantasten
udstyrer Ideen, for at indføre den med overvirkelig Pragt i Virkelighedens
tarvelige Verden. Saaledes kommer den mythiske Idee da til at glimre, som den
Art Skjønhed, Kjenderne i Salonsproget kalde en „Distanceblænder“; saaledes
hviler Ideen blødagtig svulmende i Mytheus Former, og Reflexionen, der vil
% Printer's error: „Mytheus“ as printed (M-y-t-h-e-u-s, the u carrying its
% breve; verified at 600 dpi, and read so by both OCR witnesses); the sense
% requires „Mythens“. Transcribed as printed, not corrected.
begribe det Eviges Sammensmelten med det Timelige, tyer nu først derhen, hvor
den finder sig mindst generet af Tilværelsens Skranker. Men, naar den
ungdommelige Reflexion saa en Tidlang har holdt sig forelsket ved det
mythiske Phænomen, begynder den i al Ædruelighed at støde sig paa
Virkelighedens skarpere Kanter, og vaagner da omsider med den kritiske
Erkjendelse, at Mythens ideale Herlighed kun er et Blændværk.

„Men,“ indvender en theologisk Skole, „det er jo slet ikke nødvendigt at lade
sig omtumle mellem slige Yderligheder, saa at man enten med de gammeldags
Troende holder Underet fast i Fremstillingens umiddelbare Bogstavform, eller
ogsaa med de moderne Aabenbaringsfjender fornegter al Tro paa en guddommelig
Medvirkning, fordi Ideen af denne Medvirkning er os overleveret i Sagnets
mythiske Indklædning. Den theologiske Sands vil vide at finde Middelveien,
saaledes at Noget af Fortællingens væsentlige Indhold bevares for Troen,
medens den Deel derimod, der nærmest beroer paa Formen, opoffres for den
videnskabelige Kritik.“
% --- p. 50 ---
Altsaa et Forslag til en fornuftig og anstændig Ligevægt imellem Troen og
Kritiken! Skulde det være muligt? Skulde det virkelig kunne lykkes en kløgtig
Besindighed at tilveiebringe en saa lemfældig Overeenskomst imellem de
stridende Magter? Nu vel da! Hvi skulde vi være saa daarlige at flye
Middelveien, dersom det virkelig er den, der fører til Maalet? Det være langt
fra os at ville overdøve deres Røst, der vandre paa Middelveien, hvis denne
Røst kun er Sandhedens Røst. Det være langt fra os at ville negte
Middelveislæren tilbørlig Agtelse. Tvertimod! Vi beundre Kløgten, vi ære
Besindigheden, hvis den virkelig udretter saa store Ting; vi ville gjerne
lytte til et sindrigt Forslag. „Methoden er denne: Man antage, at Jesus,
ifølge den guddommelige Aands underfulde Medvirkning, er født til Verden uden
Synd: ved denne dogmatiske Forudsætning fyldestgjør man den christelige Tro;
man indrømme dernæst, at Joseph i sit lovlige Ægteskab med Maria er Jesu
naturlige Fader, og at, som Følge heraf, Engle-Aabenbarelsen med Alt, hvad
derunder henhører, alene skyldes den sagnmæssige Digtning: ved denne
Indrømmelse tilfredsstiller man Kritiken.“ --- At der nu blot ikke skulde
indløbe en lille Ubehagelighed! „Hvilken Ubehagelighed?“ Den nemlig, at baade
Troen og Kritiken føle sig saa krænkede ved den foreslaaede Deling, at de i
allerhøieste Uenighed enes om at protestere imod Forliget. „Hvad mener du?
Skulde Troen, den sande christelige Tro, virkelig lide af en saadan
Mirakelsyge, at den med Hensyn paa dette Under slet ikke tillod nogen
Forandring i Fremstillingens enkelte Træk?“ Ingenlunde! Den sande Tro er ikke
mirakelsyg (Mirakelsygen kan søges i Overtroen). Den sande Tro fordrister sig
ikke til at foreskrive Gud Forløsningens Vei; den paastaaer ikke, at Underet
jo ogsaa kunde være skeet paa en anden Maade. Den sande Tro benegter ikke, at
jo Markus og Johannes have, hver for sig, overleveret os
% --- p. 51 ---
et i sit Slags fuldstændigt Evangelium; den sande Tro veed med sig selv, at
den ikke vilde gaae til Grunde, om det endog behagede Gud at tage dette
Bebudelsens Evangelium bort, naar han kun i Naade lod Troesbekjendelsens Ord
blive tilbage. Den sande Tro forsøger derfor heller ikke paa at aflede
Bebudelsens Evangelium, som en nødvendig Følge af Aabenbaringsideens
nødvendige Grund. Nei, saa vilkaarlig, saa anmassende og selvraadig er den
sande Tro ikke. Naar Troen desuagtet protesterer imod det anførte
Mæglings-Forsøg, da fremgaae dens Indsigelser af en ganske anden Grund. Den
Fortælling, der i nærværende Tilfælde skulde falde som Offer for Kritiken, er
os jo overleveret som integrerende Led i den kanoniske Evangelie-Cyklus; den
samme Lov, der gjælder for Evangeliefortolkningen i Almindelighed, maa da vel
ogsaa gjælde for Fortolkningen af dette Evangelium i Særdeleshed.

Men Evangelietroens ufravigelige Grundlov er denne, at intet Under kan være
Gjenstand for en troende Opfattelse, hvis det ikke bliver os fremstillet i en
bestemt Virkelighedsform; thi Underet er en Kjendsgjerning, og enhver
Kjendsgjerning maa altid forestilles paa en vis Maade. Lærer jeg nu af
Kritiken, at vor foreliggende Evangelieberetning er en forfeilet og forfalsket
Fremstilling af Ideen om Jesu syndefrie Fødsel, da maa jeg naturligviis opgive
Troen paa dette Evangelium. Min Tanke kan maaskee beskjæftige sig med Ideen;
men Troen, der har tabt den evangeliske Forestilling, famler iblinde mellem
løse Formodninger om \emph{saaledes}? eller \emph{saaledes}? eller maaskee
\emph{ganske anderledes}? For ikke at misforstaaes, maae vi her udtrykkelig
bemærke, at det bestemte \emph{Hvorledes}, hvorom Troen spørger, ingenlunde er
et Begribningens, men et Forestillingens \emph{Hvorledes}. Ja kunde jeg
begribe, hvorledes Underet nødvendigviis maatte være foregaaet, ja da behøvede
jeg just ikke at tage det saa nøie med den evangeliske Forestillings\-%
% --- p. 52 ---
maade; men, just fordi Troen ikke har Begrebet, maa den nødvendigviis holde
sig til Forestillingen. Den, der hverken har nogen Indsigt i Tingens Væsen, ei
heller nogen Forestilling om dens ydre Fremtræden, hvad har han? Løse Tanker,
ustadige Meninger, tomme Formodninger. Den kritiske Fortolker, der slaaer ind
paa den anpriste Middelvei, maa nu vende og dreie sig, hvordan han vil, saa
kommer han dog i Forlegenhed med Troen. Siger han: „Vi betragte den
evangeliske Fremstilling som forfeilet, fordi det paagjældende Under
overhovedet slet ikke lader sig betegne ved nogensomhelst Forestilling;“ da
ligger jo heri, at Underet selv er utroligt, eftersom det modsiger enhver
mulig Antagelse. Evangelisten har da villet give os en Forestilling om Det,
der ifølge sin Natur strider imod al Forestilling; han har villet gjøre det
Umulige muligt, idet han har sat sig en aldeles uevangelisk Opgave. Men imod
en saadan Paastand maa Troen paa det Bestemteste protestere; thi en
Evangelieforfatter, der i den Grad misforstaaer sit Kald, at han begynder med
umulige Opgavers Løsning, kan dog vel i det Hele taget ansees for lidet
troværdig. Svarer Kritikeren derimod, at Opgavens Løsning ingenlunde er
umulig, kun at Evangelisten af let forklarlige Grunde har ladet sig vildlede
af det mythiske Sagn; paatager han sig desaarsag at indsætte os en mere
adæqvat Forestillingsform istedetfor den evangeliske: da --- ja hvad gjør
Kritikeren da? \emph{Han dicterer et nyt Evangelium!} Kaste vi nu alle
undvigende Tvetydigheder bort, saa lyder det nye Evangelium kort og bestemt
saaledes: „\emph{Jesus er vel en naturlig Frugt af Josephs og Marias
Ægteskab; men formedelst den guddommelige Aands Medvirkning kom han dog til
Verden, et usyndigt Barn af syndige Forældre.}“ Man maa dog i Sandhed være
meget vis i sin Sag for at vove et saa afgjørende Skridt. Troen er sandelig i
sin gode Ret, naar den her spørger, hvil\-%
% --- p. 53 ---
ken Nødvendighed der tvinger Fortolkeren til at forkaste det gamle
Evangelium, og dernæst: hvilken Hjemmel der haves for Paalideligheden af det
nye? Det sidste Punkt er let afgjort; det nye Evangelium er et selvgjort
Evangelium, en Formodning, der staaer idag og falder imorgen. Det første
Spørgsmaal kræver derimod en omhyggeligere Drøftelse. Man indvender mod det
gamle Evangelium, at det dog alligevel ikke forklarer den syndefrie Fødsel,
efterdi det lader den menneskelige Moder blive tilbage som syndigt Led. Men
denne Vanskelighed kan dog vel ikke gjøre det nødvendigt, at digte et nyt
Evangelium, der, istedetfor at fjerne hiint Led, indfører tvende syndige Led?

„Det er da,“ sige vore Middelveistheologer, „just heller ikke af den Grund,
at vi forkaste den gamle Evangelieform; men det skeer ifølge vægtige,
kritiske Hensyn. Hvilke Indvendinger kunne f.\ Ex.\ ikke hentes paa den ene
Side fra den Opfattelse af Jesu Fødselsforhold, der gjaldt saavel blandt hans
Samtidige, som blandt hans Allernærmeste (Luc.\ 4, 22; Matth.\ 13, 55;
Joh.\ 6, 42 o.\ fl.\ St.), og paa den anden Side fra den fuldkomne Taushed,
hvormed
dette Punkt forbigaaes saavel i den evangeliske som den apostoliske
Beviisførelse for Jesu gudmenneskelige Natur? ---
% Printer's defect: the objector's speech, opened with „ before „just heller
% ikke“, is never closed. The em-dash after „Natur?“ marks the turn back to
% the author's own voice, and the closing “ is simply wanting; verified at 600
% dpi both here and after „Misforstaaelser?“ below, where the 300 dpi OCR had
% reported a quotation mark that turns out to be a speck. Left unbalanced.
Men hvis disse Indvendinger nu alle tilhobe beroe paa lutter Misforstaaelser?
Hvis en troende Opfattelse af den hellige Historie nu godtgjør, at det laae i
Sagens Natur, at hine Samtidige maatte ansee Jesus for Josephs Søn, efterdi
Joseph jo, ifølge Engelens Paamindelse, annammede Maria som sin Hustru; og at
Jesu Nærmeste ligesaa lidt som de øvrige Samtidige kunde kjende Noget til en
Hemmelighed, der for de Paagjældende selv jo var saa vanskelig at meddele, at
Maria end ikke formaaede at betroe den til Joseph? Hvis endelig dette Under
ikke bliver anført som Beviis for Jesu guddommelige Natur, nemlig i
Evangelierne
% --- p. 54 ---
ikke, fordi det ikke er ved at fortælle om Undere, men ved selv at gjøre
Undergjerninger, at Jesus vækker Erkjendelsen af sin guddommelige Sendelse;
og i de apostoliske Foredrag ikke, fordi disse have en anden Form end de
evangeliske Fortællinger: hvis nu, sige vi, dette forholder sig saaledes, har
da Evangelietroen ikke Grund til at klage over en saadan Fremfærd?

„Dog, hvad hjælper det, at Troen klager, naar Kritiken alligevel har
Sandheden og Videnskaben paa sin Side?“ --- Ja! Sandheden er en hellig Magt,
og den Tro, der strider imod Sandheden, maa falde ved sin egen Afmagt. Men
for at finde Sandheden, maa man kjende Veien. Almindelige Sandheder paanøde
sig Tænkeren med Nødvendighed; men Evangeliets Sandhed er en Guds Gave, som
den Troende skal tilegne sig med Frihed. Forvexles nu Nødvendighedsprøven med
den frie Tilegnelse, saa er Forvirringen grændseløs. Det er denne Forvirring,
der forraader Middelveislærens videnskabelige Svaghed. Halvkritiken begynder
nemlig med at afsondre Alt, hvad den i Forhold til den almindelige
Aabenbaringsidee mener at kunne betragte som ligegyldigt og overflødigt.
Derfor tilsidesætter den ethvert Under, der ikke med Nødvendighed kan afledes
af Ideen. Saaledes bliver Engle-Aabenbarelsen for Joseph i Drømme forkastet
som umotiveret, eftersom Maria jo har havt et Englebesøg i vaagen Tilstand; ja
selv dette Englebesøg er jo nu, ifølge Halvkritikens „nye Evangelium“, ligesaa
overflødigt, som hiint. Slige oekonomiske Beregninger kunne ved første Øiekast
tage sig ret smukt ud; men man lægge vel Mærke til, at Middelveistheoriens
afsondrende Kritik med det Samme er bleven en opløsende Kritik. Er
Opløsningen først kommen i Gang, saa kan dens videnskabelige Fart ikke
standses ved at sige: holdt! og møde for den med en Levning af Troen. Har jeg
først tilladt mig den Vilkaarlighed, at bortvise Bebudelsens Engel, da
% --- p. 55 ---
er det kun en ny Vilkaarlighed, hvis jeg i Troens Navn vil paa min egen Viis
begynde at hævde den syndefrie Fødsel. For Ideen, som saadan, er Intet
nødvendigt uden Ideen selv. Lad Menneskeheden efter sin Idee være syndefri,
Syndefrihedens Idee kan jeg da maaskee kræve med Nødvendighed; men hvor er nu
den Middelveistheolog, der beviser mig, at denne Idee vilde opløse sig i
Intet, hvis Jesus af Nazaret ikke var født syndefri? Troen paa Jesu
Usyndighed staaer og falder med den barnlige Tro paa den hellige Historie;
men ingen historisk Kjendsgjerning lader sig bevise i Kraft af den rene
Nødvendighed. Den negative Kritik er derfor i sin gode, videnskabelige Ret,
naar den med ubøielig Conseqvens fuldbyrder den Opløsning, som
Middelveislæren allerede forlængst har paabegyndt.

Vi have fulgt Undersøgelsen igjennem dens forskjellige Instanser; vi have
hørt Nok af Vantroens Spot, af Tvivlens Omsvøb og Halvtroens kritiske
Forlegenheder, Nok til at overbevise os om, at der kun er eet afgjørende
Skridt, enten at vælge Troesfornegtelsen eller at blive i Troen.

Den Troende, der saaledes har sikkret sig mod ethvert videnskabeligt Angreb,
vender sig da nu bort fra Stridens Larm, for at hvile i Evangeliet med
Inderlighedens stille Glæde. Spotterens Overmod skal ikke saare hans Følelse,
og Kritikerens Gjøglebilleder ikke blænde hans Øie. Kun Modsigelsens sagte
Gjenlyd vil endnu en enkelt Gang lade sig høre i den Overveielse, hvorunder
han stræber at blive fortrolig, ret fortrolig med Evangeliets Under.

% The quotation opened here (Joh. 1, 1) closes on p. 56, after „iblandt
% Qvinderne.“; the opening „ is therefore unmatched inside this fragment.
% Only the words of Joh. 1, 1 are letterspaced; what follows is ordinary.
„\emph{I Begyndelsen var Ordet, og Ordet var hos Gud, og Ordet var Gud.} Og
Gud talede Almagtens Ord, og see! der blev en Verden formedelst Ordet; han
talede Udvælgelsens Ord, og see! der blev et Guds Folk i Verden formedelst
Ordet.
% --- p. 56 ---
Abraham hørte den talende Gud, og vandrede bort fra de stumme Guder; Moses
hørte den talende Gud, og --- Israel blev Guds Eiendoms Folk formedelst Ordet.
Men Folkets Hjerte var for Herren i Israels Huus, som en Qvindes kjødelige
Hjerte, der ikke kjender Bestandighed, men i Afmagt og Styrke, i Svaghed og
Trods bevæges ustadigt hid og did. Da hørtes det hellige Ord, som Vredens
bittre Ord, der revser den Gjenstridige, som Straffens sønderknusende Ord, der
tugter den Troløse. Og Folkets Hjerte var igjen i de Udvalgte som en elskelig
Qvindes Hjerte, der angrer sin Brøde, ydmyger sig i Lydighed, og bøier sig i
Tro og Hengivenhed. Da hørtes det hellige Ord som et Barmhjertighedens Ord, der
opreiser den Faldne, som et Naadens Ord, der husvaler den Sørgende, som et
Kjærlighedens Ord, der fylder Sjælen med Glæde. Det var en Lignelse, ja en mørk
Tale, at Folket var viet til Jehova, som Qvinden er viet til Manden, sin Herre.
Det var en Hellighedens Tugt, der skulde frelse de Udvalgte; en Kjærlighedens
Gaade, der skulde løses paa det Sidste. Aldrig lød et Verdens-Ord saa tungt, saa
strengt, som Lovens Ord til Jehovas Folk; aldrig blev en Verdens-Gaade løst, som
Pagtens Gaade i Israel, da Qvinden „fandt Naade hos Gud“. Og Gaaden løstes, og
Lignelsen blev en Sandhed, og Loven et Evangelium i Engelens Hilsen:
% The angel's greeting is set in the type one size LARGER than the body, but it
% begins mid-line after the colon; hence inline {\large ...}, as at pp. 41--42.
% The closing “ inside the group closes the quotation opened on p. 55 („I
% Begyndelsen var Ordet ...); it is the only closer in this fragment without a
% mate here.
{\large Hil være dig, du Benaadede; Herren er med dig, du velsignede iblandt
Qvinderne.“}%
% Printed mark *), which is what \thefootnote is set to. The note is very long:
% it runs across the feet of printed pp. 56--61, ending on p. 61. It is given
% here whole at its anchor; the page markers below therefore mark where each
% printed page's BODY begins, not where the note breaks.
\footnote{„Men Jomfru Maria var trolovet en Mand ved Navn Joseph af Davids
Huus, og Joseph miskjendte sin trolovede Hustru.“ Hvorledes Forholdet imellem
Maria og Joseph nærmere bliver at opfatte, vil i Hovedsagen beroe paa, hvorledes
man forklarer sig Marias Udvælgelse. Naar Evangelielegenden, ifølge den gamle
Kirkes asketiske Retning og luxuriøse Mirakeltro, fremstiller Maria som født
Helgeninde, som Englenes Legesøster, da kunne vore Evangelister ingenlunde siges
at bære Skylden for en saadan Misforstaaelse. Lukas omtaler Trolovelsen som en
virkelig Kjendsgjerning, og Matthæus, der forudsætter, at Underet allerede er
foregaaet med Maria, lægger en særegen Vægt paa Josephs naturlige Overraskelse.
At Fritænkeren tager sig de apokryphiske Opdigtelser til Indtægt kan desaarsag
ikke overraske Nogen, der veed, hvor han med Bestemthed vil hen.
Grundforskjellen imellem den kanoniske og apokryphiske Fremstilling beroer
imidlertid ingenlunde derpaa, at hiin optager færre, denne flere mirakuløse
Tildragelser; men hvad der her gjør Udslaget, er just dette, at Legendens
Mirakelfortællinger forraade deres uægte Herkomst, forsaavidt deres
Mangfoldighed kun tjener til at udstyre „den Benaadede“ som et phantastisk
Væsen. Med samme Ret som den asketisk fromme Phantasie har tilintetgjort
Trolovelsens naturlige Betydning ved overnaturlige Tilberedelser, og forvansket
det oprindelige Indtryk af Marias Billede ved en falsk Forgyldning, med samme
Ret kunde en moderne Phantasie nu ogsaa falde paa, at ville omgive de Trolovede
med en erotisk Nimbus, og gjennemgløde Joseph og Maria med Elskovsgudens stærke
Lidenskab. Istedetfor det aldeles usandselige Baand imellem en 90aarig Olding og
en ung Klosterjomfru, der sidder og venter paa et overordentligt
% The last line of the note on p. 57 is badly under-inked. „Englebesøg“: the e
% between „Engl“ and „besøg“ survives only as a grey ghost at 1200 dpi, and both
% OCR witnesses read „Engl besog“. The word-space and the sense settle it.
Englebesøg, fik vi altsaa et romantisk Kjærlighedsbaand, sandseligt straalende,
som det, der omslynger en Romeo og Julie, en
% Printer's error: „Arel“ as printed. The second sort is a plain r --- no tail
% below the baseline --- checked at 600 dpi against the x of „Reflexion“ later
% in this same note and against „luxuriøse“ above, both of which carry the
% tail clearly. Both OCR witnesses also read r („Arel“, „Artl“). The sense
% requires „Axel og Valborg“. Transcribed as printed.
Arel og Valborg; kun med den Forskjel, at det her ikke er en Familietvist, ei
heller nogen Sortebroder, men en Erkeengel, der som en
% „Deus ex machina“ is set in a bold ANTIQUA against the surrounding Fraktur,
% as „Lic. theol.“ on the title page and „Favete lingvis“ in the Fortale.
\textit{Deus ex machina} hugger Knuden over. Kan den sidste Digtning nu ikke med
mindste Skin af Ret retfærdiggjøres af Evangeliets
% Printer's error: „Grundtert“ as printed --- again a plain r, without the
% descending tail of the x, verified at 600 dpi; the sense requires
% „Grundtext“. Transcribed as printed.
Grundtert, da kan den første det ligesaa lidt; og der er saaledes ingen Grund
til at miskjende den evangeliske Fremstilling af ængsteligt Hensyn til Legendens
Udsvævelser. Blive vi altsaa, som det sig bør, ved den kanoniske Overlevering,
da beholder Trolovelsen sin fulde sædelige Gyldighed, og da vil ogsaa den
overnaturlige Meddelelse, hvorved Josephs Misforstaaelse hæves, vise sig
begrundet i Sagens Natur. At Maria har rakt Joseph sin Haand, beviser klart, at
hun, før Bebudelsen indtraf, ikke har havt nogen Forestilling om, at hun selv
var „den Benaadede“. Denne Omstændighed er imidlertid saa langt fra at svække
vor Overbeviisning om hendes sjælelige Forberedelse, at den meget mere tjener
til at sætte hendes rene Hjertes eenfoldige Tro i et endnu klarere Lys. Det er
dog unegtelig et af de allerskjønneste Beviser paa sand Ydmyghed, at Den, der i
sin dybe Syndserkjendelse, sin ubetingede Tillid til Guds Naade og uendelig
Hengivenhed i hans Villie fuldkommer Troens indre Gjerning, selv ingen Anelse
har om sit Fortrin fremfor Andre. Havde Maria tænkt ved sig selv: Jeg tør ikke
trolove mig med Joseph; jeg maa være forberedt paa noget langt Høiere; thi, hvis
Herren skulde behøve en Tjenerinde, for at udrette det Overordentlige, vil der
neppe findes nogen Værdigere, end jeg; da skulde en saadan selvbehagelig
Reflexion vel have været nok til at gjøre hende aldeles uværdig.

Stemmer det nu saavel med Ordet som med Aanden i Evangeliet, at de Trolovede
virkelig stode i Begreb med at fuldbyrde deres Ægteskab, da Bebudelsen indtraf;
saa bliver Marias Taushed angaaende det skete Under forstaaelig for enhver
Troende. Om hendes nærmere Forhold til Joseph og sammes Betydning i den
borgerlig-sædelige Verden havde Engelen jo Intet aabenbaret hende. Skulde hun
altsaa afbryde Forbindelsen? Skulde hun, som Moder til „den Høiestes Søn“, staae
ene i Verden, og udsætte ham og sig for Menneskenes Spot? Vilde Joseph forstaae
hende, hvis hun betroede ham sin Hemmelighed? Hvad skulde hun forlange af ham?
Vilde han indvillige i hendes Forlangende, hvis hun virkelig kunde udfinde, hvad
hun skulde bede ham om? At Kritikerne hidtil ikke ret have villet ændse disse og
lignende Vanskeligheder, er forklarligt nok. Naar man nemlig i Kraft af den
moderne Bevidsthed forud veed, at Underet kun er en Indbildning, saa farer man
slige Ubetydeligheder over. Det er jo ikke Umagen værd, at sætte sig levende ind
i en saa urimelig Fortælling. Derimod fremmer det Kritikens videnskabelige Seier
langt lettere, naar man lader, som om en Engle-Aabenbarelse var en
Hverdagsbegivenhed; thi derved aabner man sig jo den allerbeqvemmeste Udvei til
at slippe bort fra Underet.

Skulde det derimod ikke være Videnskaben aldeles uværdigt at sætte sig
tilbørligt ind i sin Gjenstand, skulde den moderne Kritik virkelig kunne
beqvemme sig til at gjøre et Indblik i de Trolovedes daværende Stilling, da vil
man sikkerlig blues ved at henføre Marias Taushed med Strauß til ---
„Dovenskabens Grundsætning“.

Er Aabenbarelsen for Joseph allerede tilstrækkelig motiveret ved Hensynet til
Marias Tilstand, saa erholder Meddelelsen endnu desuden sit eiendommelige Præg
derved, at Engelen viser sig for ham i Drømme. Denne Meddelelsesform svarer
aldeles til sit Øiemed. Josephs naturlige Smerte over Marias Tab stiller os
Manden i det rette Lys. Han føler sig krænket, og denne Krænkelse er saa meget
bittrere, som han tillige føler, hvad han taber, idet han taber Maria.

Han nærer Mistanke mod hendes Uskyldighed, thi hun er jo som
% Printer's error: „salden“ as printed. The sort is a long s (U+017F) --- a nub on the
% left of the stem only --- identical with the long s of „som“ two words earlier and
% unlike the full-crossbar f of „Derfor“ three lines below, all checked at
% 1200 dpi; both OCR witnesses also read s. The sense requires „falden“. This is
% the mirror of the „flaaende“ for „slaaende“ defect logged on p. 48.
% Transcribed as printed, not corrected.
salden; men han qvæler igjen denne Mistanke: hun har jo hidtil været ham et
Mønster paa den reneste Qvindelighed. Derfor tør han end ikke følge Iversygens
Indskydelse, at beskjæmme hende offentlig, men beslutter at skille sig hemmelig
fra hende. Ikke til noget Menneske kan han betroe sin Kummer, thi hun er ham
dyrebar fremfor Alle; skal han opgive Troen paa hende, maa han opgive Troen paa
Menneskeheden. Ingen Ven vilde forstaae hans Tale, om han end prøvede paa en
Forklaring; thi der er Forvirring og Modsigelse i hans egne Tanker. Endnu mindre
kan han nærme sig Maria selv: I hvilket Sprog skulde han vel tale til hende?
Bebreide hende sit Fald? Nei! hun kan jo ikke være skyldig! Forlange hendes
Retfærdiggjørelse? Nei! Hun kan umuligt retfærdiggjøre sig; hun maa jo dog være
skyldig! Saa har han da kun een Udvei tilbage, ɔ: at skjule sin Smerte, og
„skille sig hemmelig fra hende“. Man behøver i Sandhed ikke at ophøie
Tømmermanden Joseph af Davids Huus til Elskovsridder af den moderne Reflexion,
for at forklare sig hans hemmelige Smerte. Man indrømme blot Manden, hvad der,
efter Evangeliet, dog vel idetmindste tilkommer ham, et ædelt Hjerte, en sund
Forstand, og Striden i hans Indre forstaaes af sig selv. Kan Joseph med al sin
Grublen ikke selv løse den Gaade, der er sat ham ved Underet, da er det jo
ingenlunde overflødigt, men aldeles overeensstemmende med det tilsigtede Øiemed,
at en Herrens Engel løser den for ham i Drømme. Joseph er altsaa forberedt paa
en Aabenbarelse. Hans Hjerte fordrer intet Mindre, end hvad Drømmens Engel
forkynder ham: han kræver Underet, og derfor troer han ogsaa paa Underet. Denne
Josephs Tro er imidlertid uendelig forskjellig fra Drømmetyderens løse Overtro;
thi den beviser
% Printer's error: „strar“ as printed --- a plain r without the x's descending
% tail (600 dpi), against the x of „Reflexion“ five lines above in this same
% note. The sense requires „strax“. Transcribed as printed.
strar sin Sandhed i en stor sædelig Beslutning: han vaagner af sin Drøm, for at
gjøre, som Herrens Engel har befalet ham. Derved bliver Joseph et evangelisk
Forbillede for alle Dem, der med reen og mandig Selvfornegtelse fyldestgjøre
Pligtens hellige Lov, selv da, naar Pligtopfyldelsen, hvis den vorder kundbar,
kun fører Verdens Spot og Miskjendelse over dem.}

% --- p. 57 ---
% As on p. 55, only the words of Joh. 1, 1 are letterspaced; what follows is
% ordinary. This „ is not closed until p. 66.
„\emph{I Begyndelsen var Ordet, og Ordet var hos Gud, og Ordet var Gud.} Og Gud
talede Almagtens Ord, og see! der blev en Verden formedelst Ordet; han talede
Udvælgelsens Ord, og see! Israel blev Guds Eiendoms Folk formedelst Ordet. ---
Og efterdi Aanden, den Hellig Aand kom over Israel ved Ordet, og den Høiestes
Kraft overskyggede Guds Udvalgte formedelst Ordet, derfor blev Israels Folk, som
i et Billede
% --- p. 58 ---
kaldet den Høiestes Søn efter Herrens Ord: \emph{Israel er min førstefødte Søn}
(2 Mos.\ 4, 22). Og dog var Israels Hjerte haardt, og han vandrede efter sine
egne Tanker paa en Vei, som ikke var god. Og Jehova straffede sin Udvalgte, ja
Herren revsede sit Folk, som en Fader revser sin gjenstridige Søn. Men efterdi
% --- p. 59 ---
Aanden, den Hellig Aand, dog var over Israel ved Ordet, og den Høiestes Kraft
overskyggede Guds Udvalgte formedelst Ordet, saa bleve da de Ypperste, som
fødtes af Herrens Folk, hilsede som \emph{den Høiestes Børn} (Ps.\ 82, 6). Og
hvo var vel ypperligere iblandt de Udvalgte, end Hyrden fra Bethlehem,
Sangerkongen med den
% --- p. 60 ---
% The citation is printed „(2 Ps. 7)“ --- the numeral of the psalm BEFORE
% the abbreviation, unlike „(Ps. 82, 6)“ on p. 59. Ps. 2, 7 is meant. As printed.
gyldne Harpe, Manden efter Guds Hjerte? Saa lød da Herrens Ord til David:
\emph{Du er min Søn; jeg fødte dig idag} (2 Ps.\ 7). --- Og Gud Jehova var med
David, og Davids Navn var stort, „som de Stores Navn, der ere paa Jorden,“ og
han var den Høiestes Søn formedelst Aanden. Men han var
% --- p. 61 ---
ogsaa en Syndens og Misgjerningens Søn formedelst Kjødet. Og Manden efter Guds
Hjerte maa grue for sin Misgjerning, som ingen Anden, Herrens Udvalgte fornemmer
Guds Vrede, som ingen Anden; Sangerkongen med den gyldne Harpe er jo bedrøvet,
som ingen Anden: han er saa bleg af Sorg, „han er
% --- p. 62 ---
saa træt af Sukke i den lange Nat, han har vædet sit Leie med Taarer.“ Det var
da kun som i et Billede, i en Lignelse, at David blev kaldet \emph{den Høiestes
Søn}. Han fik en Evighedens Forjættelse i Ordet, men Aanden straffede ham i
Kjødet; --- „han gik bort ad al Jordens Vei.“ Men, efterdi Aanden dog var over
Israel, og den Høiestes Kraft overskyggede Guds Udvalgte, virkede Sandhedsordet
i Lignelsen. Nathan stod for Davids Throne, og talte med prophetisk Røst om den
Høiestes Søn, der skulde bygge Herren et Huus, og have sit Riges Stol stadfæstet
evindeligen. \emph{Naar dine Dage ere opfyldte, og du hviler med dine Fædre, da
vil jeg opreise din Sæd efter dig, som skal udkomme af dit Liv, og jeg vil
stadfæste hans Rige} (2 Sam.\ 7, 12). --- Og hvo var nu ypperligere blandt de
Udvalgte, end Salomo med Lykkens Skatte og Aandens herlige Gaver? Han bad om
Forstand, og Gud Herren gav ham et viist og forstandigt Hjerte, ja gav ham
endogsaa Det, han ikke havde begjæret, ɔ: baade Rigdom og Ære. Al Israel hørte
Kong Salomos Dom; og de frygtede for Kongens Ansigt, thi de saae, at Guds
Viisdom var inden i ham til at holde Dom. Og Salomo blev større end alle Konger
paa Jorden med Rigdom og Viisdom. Dronningen fra Saba kom til Jerusalem, at
friste ham med mørke Taler; og der var ikke et Ord skjult for Kongen, som han jo
kundgjorde hende. Folket priste ham i sit Hjerte, naar han vendte sit Ansigt
omkring, og velsignede al Israels Forsamling; ja Folket priste hans Navn, thi
han var en Fredens Fyrste: Juda og Israel boede tryggeligen, hver under sit
Vintræ og under sit Figentræ fra Dan indtil Beerschaba alle Salomos Dage. Og dog
var det kun som i et Billede, i en Lignelse, at denne Davids Søn kunde kaldes
\emph{den Høiestes Søn}. Viisdommens Søn tog Viisdommen forfængelig, og for For\-%
% --- p. 63 ---
% „Forfængelighedens Forfængelighed. Alt Forfængelighed!“ is set in BOLD
% Fraktur (fette Fraktur), tight, not letterspaced; the sentence that follows,
% „Jeg Prædikeren ... Aandsfortærelse.“, is letterspaced in the ordinary weight.
% Both checked at 600 dpi: cap heights are identical, so this is weight, not
% size. Bold is a third emphasis device in this book, distinct from the
% letterspacing (\emph) and the antiqua (\textit) recorded in the header.
fængelighedens Viisdom ere alle Ting møisommelige, at Ingen kan udsige det: Øiet
mættes ei af at see, og Øret fyldes ei af at høre; og Prædikeren sukker saa
tungsindig: \textbf{Forfængelighedens Forfængelighed. Alt Forfængelighed!}
\emph{Jeg Prædikeren, jeg var Konge over Israel i Jerusalem; jeg saae alle de
Gjerninger, der ere gjorte under Solen, og see! det var altsammen Forfængelighed
og Aandsfortærelse.} Nei, ak nei! Salomo i al sin Herlighed var ikke den sande
Guds Søn. Aandens Kraft svigtede ham, og hans Tjenere spottede ham; Kjødets
Sands bedaarede ham, og hans Hustruer bøiede hans Hjerte. --- Israel er da sig
selv en Gaade: Aanden er kraftig; men Kjødets Sands er Fjendskab mod Guds Aand,
og Ordet kan ikke blive Kjød. See Ephraim og Juda, hvor de hade hinanden med et
dødeligt Had. Falske Aander og falske Propheter, fremmede Fiender og fremmede
Guder besnære Israel, fortære hans Kraft, og true ham med Fald og Fordærvelse.
Hvor er nu \emph{den Høiestes Søn}, der skulde sidde paa Davids Throne, og
regjere over Jacobs Huus evindelig? Ak! Paa Thronen sidder en Achas, ængstet af
Phekah og Rezin. Det er kundgjort for Davids Huus, at Syrerne forlade sig trygt
paa Ephraim; og Kongens Hjerte bæver og hans Folks Hjerte, „som Træerne i Skoven
bæve for Vinden.“ --- Men, efterdi Aanden, den Hellig Aand, dog var over Israel,
og den Høiestes Kraft overskyggede Guds Udvalgte, var Ordet mægtigt formedelst
Aanden, og Herrens Propheter vældige ved Ordets Kraft; thi Aanden var i
Propheternes Ord, og Ordet var som „en Hammer, der kløver Fjelde“. --- Propheten
gaaer Kongen imøde; Jesaias taler til Achas; saa taler en Troens Mand i Farens
Stund til den forsagte Tvivler: \emph{Tag dig iagt, hold dig rolig, ei frygte
du, og ei vorde dit Hjerte forsagt} (Jesaj.\ 7, 4). --- Jehovas Folk er i den
Almægtiges Haand. Kongen i Damaskus kan
% --- p. 64 ---
ikke betvinge os, Ephraim og Remaljas Søn kunne ikke ødelægge os; kun naar
Jehova fløiter ad Fluen, som er ved Ægyptens Floders Grændse, og ad Bien, som er
i Assurs Land, kommer Fordærvelsen over os. Assyrien kan ikke redde os; et
Menneskes Arm kan ikke beskytte os mod den Helliges Vrede. Vi ere i Guds Haand,
vi ydmyges under hans Vælde, vi opreises ved hans Arms Kraft, vi skulle frelses
ved Almagtens Under. --- \emph{Derfor vil Jehova give Eder et Tegn:}
% Jes. 7, 14 and Jes. 9, 6 are both in BOLD Fraktur, like the Ecclesiastes verse
% on p. 63; the words that introduce them are letterspaced. The citations
% themselves are in the ordinary weight.
\textbf{See, Jomfruen vorder frugtsommelig, og føder en Søn, og du skal kalde
hans Navn Immanuel} (Jes.\ 7, 14). Men Tvivleren paa Thronen forstod det ikke,
hans Tjenere fattede det ikke. Han, der skulde være som \emph{den Høiestes Søn},
brændte sine Sønner i Ilden efter Hedningernes Vederstyggelighed, og plyndrede
Herrens Huus, for at kjøbe sin Frelse af Assyriens Konge. Og dog --- drømmer
Propheten ikke, han taber aldrig Nærværelsens Vilkaar af Syne; men midt under
Tidens Trængsler drager et Fremtids Syn forbi hans Øie, og han fryder sig ved
Synet, som Den, der famler i Natten, fryder sig, naar Dagen gryer; ja han fryder
sig ved Synet, thi Folkets Frelse, den sande Frelse viser sig for ham i Synet:
% „udbryder han,“ stands OUTSIDE the letterspacing, in the ordinary tracking
% (600 dpi); the prophecy resumes with „skuer“.
\emph{Det Folk, der vandrede i Mørke,} udbryder han, \emph{skuer et stort Lys;
Lyset opgaaer for dem, der boe i Dødens Skyggers Land.} Jesaias drømmer ikke,
han taber ikke det Nærværendes Vilkaar af Syne; men af Nærværelsens Syn udklarer
sig et Evighedssyn, og den Høiestes Søn staaer for ham i Synet: \textbf{Et Barn
er født os, en Søn er given os. Og Fyrstendommet skal være paa hans Skulder; og
hans Navn skal kaldes Under, Raadgiver, stærke Gud, Evighedens Fader og Fredens
Fyrste} (Jes.\ 9, 6).

Hvor er dette udvalgte Folk med sine Konger og Propheter dog ikke et
vidunderligt Folk. Betragter nu Israel! Fjernt i
% --- p. 65 ---
Historien kommer han op, en eenlig Gestalt, i Mythernes Taage. Hans Bane krydser
Folkenes Veie trindt omkring: skulde han nu blande sig med de Fremmede? Nei! han
hader dem, de foragte ham: han vilde knuse alle Hedningene, om han kunde.
Betragter Israel, betragter ham nøie! Han er intet Taagebillede i den graae
Oldtid, men en virkelig Skikkelse med skarpe, udprægede Træk. Der er Lys i hans
Øie, naar han seer paa Virkeligheden; der er Forstand i hans Hjerne, naar han
dømmer om Tingenes Gang, og beregner Vilkaar og Øiemed. Du blænder ham ikke med
tomme Ord: han kjender Naturlighedens strenge Lov, han fornemmer Virkelighedens
Skranke; han kræver Kjendsgjerning. Og dog, hvor forunderligt! Han, just han er
en Troende, ja den eneste Troende midt iblandt alle Folkene, fordi han er den
Eneste, der ret kjender Forskjel mellem Skaber og Skabning. Han drømmer sig ikke
ind i Troen. Nei! han nødes til Troen: thi kun i Troen er hans Frelse. Almagtens
Gud frelser ham, Underets Gud løser hans Lænker, og bærer ham paa Ørnevinger
bort fra Trældommens Land. Ja, saa maa han vel troe; thi han staaer jo for Sinai
Bjerg, omgiven af Tegn og underlige Gjerninger. Betragter Israel, betragter ham
nøie! Han er ingen marvløs Skygge, men en kraftig Natur med Kjød og Blod og en
Sjæl, fuld af Lyster og Lidenskaber. Han seer Afgudsbilledet: Lysten vaagner,
Nydelsen vinker. Han seer Guldkalven, dens Glands fortryller ham: Øiet funkler,
Haanden zittrer, Læberne tilbede, thi han er næsten vild af Begjærlighed. Da
tordner Jehovas Stemme, og Synderen farer sammen i Forfærdelse. Afgudsalteret
falder, Glædessangen tier, Dandsen ender. Den Ulydige maa nu lide sin Straf, maa
fornegte sin Lyst, angre sin Overtrædelse, og forbande de falske Guder.
Betragter Israel, naar den Vældiges Arm hviler tungt paa ham, naar han strider
Fortvivlelsens Strid med de bittre Fiender, naar
% --- p. 66 ---
han raaber til Himlen om Hevn, men Herren svarer ham ikke, fordi Vredens Dag er
kommen: da synker hans Mod, da svinder hans Kraft; da komme Fienderne over ham,
og sønderrive ham, og sprede hans Lemmer ved Babels Floder. Da er han som en
Død; hvo skulde nu troe, at han atter kunde reise sig. Men betragter ham saa
igjen, naar Aanden kommer over ham, naar den Høiestes Kraft overskygger ham! Det
er et Guds Under. „De døde Been“ samle sig med Kjød og Blod; de spredte Lemmer
forene sig, Dødningen aander, Israel staaer op, finder atter Veien til
Forjættelsens Land og priser Jehova paa Zions Bjerg. Det udvalgte Folk er en
verdenshistorisk Gaade. Om alle øvrige, aandeligt begavede Folkeslag gjælder
det, at de begynde med mythiske Naturguder, med glimrende Ungdomsdrømme, gaae
dernæst efterhaanden over fra Mythetroens umiddelbare Forestillinger til en
reflecteret Opfattelse af Verdensideen, og derigjennem til Tvivl og Vantro.
Israel derimod har fra Begyndelsen af den Modsigelse i sig, at han paa een Gang
er en Troende og dog en Tvivler, en ivrig Gudstilbeder og dog en Spotter. Han
gaaer derfor ikke fremad fra Tro til Tvivl, som fra Umyndighedens til
Myndighedens Trin, for at ende i en Opløsning; men han gaaer idelig fra Tro
gjennem Vantro tilbage til Tro. Derfor staaer han ved Slutningen af sin Løbebane
i Grunden paa samme Trin, som ved Begyndelsen. Med alle de Prøvelser, det
udvalgte Folk har gjennemgaaet, alle de Forjættelser, det har modtaget, alle de
Prophetier, det har hørt, og alle de Bibestemmelser, hvormed Loven efterhaanden
er bleven forøget, begriber det ligesaalidt sig selv og sin Gud i Slutningen,
som i Begyndelsen. Kun dette forstaaer den troende Israelit, at han skal frygte
Gud og holde hans Bud, at Folket skal frelses ved Underet, og forvente den Tid
og Time, da Frelseren, den Høiestes Søn, kommer.“

% --- p. 67 ---
% „Sjel“, not „Sjæl“, as printed here (600 dpi); the book has „Sjæl“ elsewhere
% on this same page. Not normalised.
„Seer jeg nu paa alt Dette med et alvorligt prøvende Øie, da samler min Sjel sig
i den ene Afgjørelse, at \emph{enten} er Udvælgelsens Folk et forunderligt
Vrængebillede i Historiens eget verdslige Drama, Underet et Blændværk, Aandens
Drøm om Ordets Almagt en dristig Digtning, og Verdensgaaden et naturligt
Knudepunkt i Menneskehedens eget Væsen (og da har Tænkeren Ret, naar han opløser
alle overnaturlige Kjendsgjerninger som naturlige Symboler i den kosmiske
Idee-Udvikling, og da har den negative Kritiker fuldkommen Ret, naar han igjen
opløser Ideens mythiske Selvhed i gemytlige Indbildninger, og Halvtroens
Forsvarer kan i Videnskabens Sprog ikke svare ham Eet til Tusinde); ---
\emph{eller} ogsaa maa Verdensgaaden være en Guds Hemmelighed, Udvælgelsens Folk
Underets Folk, og Propheternes Syner Aandens Forvarsler for det evige Ords
Aabenbarelse i Kjødet: og da har Evangelietroen ubetinget Ret, naar den med
Uvillie vender sig bort fra alle tjenstvillige Halvforklaringer, og bringer alle
selvkloge Indsigelser til Taushed, for at kunne høre i Fred, hvorledes
Bebudelsens Engel løser Gaaden for Jomfruen i Nazaret. Halvtroens Mænd, der her
ane noget Overnaturligt, og dog see sig ængstelig om efter de naturlige
Betingelser, modsige sig selv og vide ikke, hvad de gjøre: kun Bebudelsens
Engel, den Engel, der „staaer for Guds Aasyn“ (Luc.\ 1, 19.), veed Gaadens
Opløsning. Engelen veed, at som Jord og Støv var den naturlige Betingelse for
Skabelsen af den første Adam, saa er nu en troende Qvindes rene Sjæl den
tilstrækkelige Natur-Betingelse for Undfangelsen af den anden Adam; Engelen
veed, at Frelsens Tegn er „Jomfruens Søn Immanuel“; ja Engelen, der staaer for
Guds Ansigt, Bebudelsens Engel, veed, at Ordet fra Himmelen bliver Kjød i
Israel, naar Aanden sænker sig med Skaberens Kraft i det jomfruelige Hjerte, og
Kjærlighedens Gud fuldbyrder sit Under,
% --- p. 68 ---
Derfor er Røsten saa klar, Stemmen saa høitidelig, saa mildt bevæget, og dog saa
myndig, naar det lyder til Maria: \textbf{Den Hellig Aand skal komme over Dig,
og den Høiestes Kraft skal overskygge Dig; derfor skal ogsaa det Hellige, som
fødes af Dig, kaldes Guds Søn.“}

% This „ is not closed on p. 68; it closes on p. 69, after „det overgik Sands og
% Forstand.“ (checked on the image of p. 69, which is not transcribed here).
„\emph{I Begyndelsen var Ordet, og Ordet var hos Gud, og Ordet var Gud.} Og
Ordets Gud talede til de Udvalgte, og Israel hørte Guds Ord, og forstod, at det
var et Troens Ord, der skulde bevares i Hjertet, og bevises i Gjerningen. Ogsaa
Verdensordet vil høres i Verden, ogsaa Naturaanden vil forherliges i
Menneskeheden efter sine Forbilleder. Gjælder det en Væddestrid om Skjønhedens,
Tankeviisdommens, Verdensherredømmets Priis: da er Kanaan ikke det forjættede
Land; da kan Israel hverken maale sig med Grækeren eller Romeren eller de Stærke
i Norden. Men spørges der om den Eviges Herlighed, om Gudfrygtighedens Viisdom
og Troens Kamp for Hellighedens Rige paa Jorden, da har Israel Ordet. Skal der
tales om et Menneskes Sind, der offrer sit Kjød og Blod, ja sin naturlige
Forstand for Opfyldelsen af et guddommeligt Bud; skal der tales om de Prøvelser,
hvori et Menneske forsøges, naar han strider Afgjørelsens Strid med Gud og sit
eget Hjerte; skal Talen vise, hvorledes en Aandens Mand kan trodse en heel
Verden med en guddommelig Forjættelse, og midt i Virkelighedens Farer bygge alt
sit Haab paa et Guds Under: da staaer Israel alene uden Sammenligning, da har
Jehovas Søn Ordet, og alle Folkerøster maae tie, naar han taler. Og dog er
Israels Folk med det aabenbarede Ord sig selv en Gaade: det hører, og forstaaer
dog ikke; det seer, og fatter dog ikke --- Troens Hemmelighed. At Troen har en
Evighedskraft i sig formedelst Ordet, maa indlyse for alle Dem, der betragte
dens Virkninger i Kræfternes Rige. Men, naar det kjøde\-%
% --- p. 69 ---
lige Sind antændes af Aanden, da bruser Lidenskaben, da maa den virkelige
Forskjel imellem Lydighed og Ulydighed træde ud i synlig Afgjørelse: Loven
kræver den udvortes Gjerning. Kun under denne Betingelse kan Folket forstaae den
Troendes Sjælestorhed. Selvfornegtelsens Kamp, Angerens Smerte, Opoffrelsens
ubegribelige Spændkraft: Alt maa komme tilsyne i saadanne Kjendsgjerninger, at
det kan blive Gjenstand for Beundring, og gaae over i Folkehistorien. At
Troeslivet derimod kan være skjult i reen Inderlighed, og dog opløfte et
Menneske over den ydre Verdens Forkrænkelighed, at den fuldkomne Gjerning kan
være en sagtmodig Sjæls ubetingede Hengivenhed i Guds Villie, at den Høiestes
Kraft kan være virksom i en stille Aands uforkrænkelige Væsen, det vilde aldrig
ret blive klart for Jehovas Folk. I denne Forstand gjælder det, at Israels
stærke Tro er mere udvortes-historisk end udvortes-sjælelig, mere
heroisk-mandig end gemytlig-qvindelig. Derfor kan Lovens Folk vel fatte sin
Troes Begyndelse, men ikke dens Fuldendelse. At Udvælgelsen begyndte med den
Mand, der fuldkommedes i Lydighed, da han beredte sig til at offre Gud sin
eneste Søn, det kunde det udvalgte Folk forstaae, og derfor priste Israel Troens
Fader; men, at Troens Fuldendelse skulde skee ved den Qvinde, der i Hjertets
Eenfoldighed hengav sig i den Helliges Villie, for at føde Guds eenbaarne Søn,
% This “ closes the „ opened on p. 68 („I Begyndelsen var Ordet …“). Verified on
% the image of p. 69: the closing mark stands after „Forstand.“, and the next
% line opens a fresh, indented paragraph with „Ikke, som vilde Israel …
det overgik Sands og Forstand.“

„Ikke, som vilde Israel stille Qvinden i Skygge: nei, langtfra! Jehovas Øie
hviler jo ogsaa paa hende; han hører hendes Suk, naar hun beder ydmygelig; han
trøster hende, naar hun taler af sin „Klages Mangfoldighed“; han velsigner
hende, naar hun „udøser sit Hjerte for hans Aasyn“. Hvor Troen udøver den
høieste Magt, tør Qvinden vel i Afgjørelsens Øieblik træde op, og maale sig med
Manden. Israel veed det, at Kraftens Aand kan
% --- p. 70 ---
i Farens Stund hvile over en Qvinde; Israel veed det, at stundom, naar de
Dristige tabte Modet, naar de Stærke bleve svage, da er det skeet, at Herren har
gjort den Svage stærk, og frelst sit Folk ved en Qvindes Haand. Tael blot et Ord
om Deborah under Palmetræet, og du skal see, hvor Hebræerens Ansigt straaler:
„Deborahs Raad blev Baraks Seier, Lapidoths Hustru var en Moder i Israel“. Nævn
Judiths Navn, det tænder en Gnist i Jødens Øie: „Judith drog Hævnens Sværd i
% Printer's error: „Jernsalems“ as printed. At 600 dpi the fourth letter is a
% Fraktur n (shoulder joined at the top), not the u of „Jerusalems“ printed on
% p. 76 of this same batch, with which it was compared side by side. Both OCR
% witnesses read „Jernsalems“ as well. Transcribed as printed.
Assyrerens Telt, og frelste Jernsalems Ære“. Ja det kan al Verden forstaae, at
Israels Folk maa lovsynge Deborah, maa prise en Esther og Judith.“

„Hvor mægtig er ikke Deborah i hendes høie, prophetiske Ro! --- Folket træller
under Jabins Haand, i tyve Aar har Forsmædelsen varet. Kun hos Prophetinden paa
Ephraims Bjerg finde de Fortrykte en Tilflugt: hun, en Qvinde, er den Eneste,
der skifter Ret i Israel. Sisera truer, og Folket raaber i Angst til Herren,
„ak, Jabins Høvidsmand har 900 Jernvogne“; men Deborah skrækkes ikke. Rolig
sidder hun under sit Palmetræ, og dømmer Israels Børn. Aanden rører sig i hende;
hun seer jo Folkets Nød, hun kjender Fiendens Anslag. Og dog sidder Qvinden med
Lovens Ord i Hjertet paa sit Dommersæde, saa fast og rolig, som den forfarne
Styrmand, naar Havet bruser. Da slaaer Befrielsens Time: hun kalder Abinoams
Søn. Hun betegner Helten sin Vei, hun hæver hans Mod, hun følger ham i Striden.
Der er en Vished i hendes Blik, en Tilforladelighed i hendes Ord, en
Overlegenhed i hendes Daadskraft, som kun Troen kan give Den, hvem Herren
udvælger at frelse sit Folk.“

% Printer's error: „Deiligbed“ as printed --- a b, not the h of „Deilighed“; at
% 600 dpi the letter has no descender, where the h of „hendes“ two words earlier
% clearly descends below the baseline. Both OCR witnesses read „Deiligbed“.
% Transcribed as printed. The comma in „bævede, for hendes Dristighed“ is also
% as printed (verified at 600 dpi); the parallel clauses on either side have none.
„Stor og herlig er Judiths Kraft; hun er Slægtens Roes: alle Folketroens
Lidenskaber brænde i Jødindens Hjerte. Assyrerens Sjæl fangedes af hendes
Deiligbed, Perseren bævede, for
% --- p. 71 ---
hendes Dristighed, Mederen forfærdedes saare for hendes Modighed. --- I Judith
blotter den stridende Jødedom sin allerhaardeste Modsigelse. Skal det udvalgte
Folk bestaae i Verden, saa maa ethvert Frelsens Middel jo være Jehova
velbehageligt; og dog kan List og Svig umuligt fremme det hellige Formaal. Der
behøves en Qvindes høieste Lidenskab, for at magte denne Modsigelse. Derfor er
Judiths Beslutning en Hemmelighed, hun end ikke kan betroe sine Egne. Ingen af
Bethuluas Mænd aner, hvorledes hun strider i Bønnen, og besværger
„Krigshærenes Gud, at han skal slaae Ødelæggelsens Tjener ved hendes Læbers
Besvigelse“. Ingen fatter, hvi den bedrøvede Enke dog saa pludselig har villet
aflægge sin Sørgedragt, og forynge sin Skjønhed ved duftende Salve og kostelige
Prydelser. De unge Mænd betragtede hende med taus Beundring, der hun gik ud af
Stadens Port; og de saae efter hende, „indtil hun kom ned af Bjerget, indtil hun
kom igjennem Dalen,“ og de ikke kunde see hende mere. Ingen Qvinde vover Mere
for sit Folk, end Judith. Hendes Opoffrelse udspringer af den uovervindelige
Jehovatro. Hun troer, at den Gud, der frelser hendes Folk, ogsaa vil frelse
hendes qvindelige Ære; og i det lykkelige Udfald finder hun sin Tro bekræftet.
Med den fuldeste Forvisning, at den Almægtiges Villie har aabenbaret sig i
Udførelsen af hendes Daad, bestiger hun igjen Bethuluas Bjerg, og raaber til
% Letterspaced, not bold and not in a larger fount: the words of Judith are set
% off by Sperrsatz alone, and carry no quotation marks. Verified at 600 dpi ---
% „Og der hun“ on the same line reverts to ordinary spacing, and so does „da
% forstod alt Folket, at hun var velsignet fremfor alle Qvinder, som ere paa
% Jorden“ two lines below, although it too paraphrases Judith 13.
dem, der holde Vagt: \emph{Oplader, oplader dog Porten! Gud, vor Gud er med os,
og gjør Styrke endnu i Israel!} Og der hun opløftede Krigsøverstens Hoved, og
raabte: \emph{Lover Gud! Lover, lover, lover Gud!} da forstod alt Folket, at hun
var velsignet fremfor alle Qvinder, som ere paa Jorden. Og Mænd og Qvinder
velsignede Judith, og de droge hende imøde med Krandse og med Sange; og hun gik
foran alt Folket, og førte Qvinderne
% --- p. 72 ---
% The printed mark is „*)“, standing between „Dandsen“ and the full stop.
i Dandsen\footnote{At Judith er apocryphisk, kan ikke være den Troende
ubekjendt.}. Ja det kan al Verden forstaae, at Israels Folk maa lovsynge
Deborah, maa prise en Esther og Judith.“

„Men Jomfruen i Nazaret! Hvo kjender vel hende? Hvo taler et Ord til hendes
Priis? Historiens Øie kan ikke opdage hende: hun dølger sig blufærdig i
Familielivets Skygge. Intet prophetisk Kald er udgaaet til hende: hun sætter sig
ikke paa Dommersædet, at skifte Ret i Israel. Ingen stor og mægtig
Befrielsestanke bruser i hendes Sind: man hører ikke hendes Røst paa Gaderne.
Den Trolovede sysler med sit eget Hjertes stille Tanker, stilfærdig, ubemærket,
for Verden ubetydelig, som Livet i Nazaret. Og dog! dog vidner Evangeliet: Her
er Mere end Judith! Her er Mere end Deborah! --- Naar Qvinden blusser i Kampens
Hede, og griber handlende ind i Historien, er hun for et Øieblik stærkere end
Manden, fordi hun er sat i en høiere Spænding. Uden at vakle, uden at snuble,
farer hun fremad mod sit Maal, thi hun handler efter Indskydelse; eller rettere,
hun handler ikke selv, men hun er et Redskab for en uimodstaaelig Magt, der
handler igjennem hende. Den Alt opoffrende Lidenskab, der har forvandlet hendes
Væsen, den Alt formaaende Beslutning, der har sammentrængt hendes Kraft, kort
sagt: den Indvielse, der forlener hende overnaturlig Styrke, er betinget ved
Troens Forudsætning, og Jehovatroen har tilstrækkelig viist, hvad Aanden formaaer
ved en Qvinde. Men i alle slige Optrin bliver Qvinden dog egentlig bortført fra
sig selv --- ud over sin væsentlige Naturgrændse, istedetfor at bevares for sin
oprindelige Bestemmelse og i sjælfuld Eiendommelighed oplade sit Hjertes skjulte
Rigdom. Er nu Qvindelighedens Blomst herved hæmmet i sin Udfoldning, da er
Troens Væsen det ikke mindre. Den religiøse Tro yttrer sig nemlig eensidig,
forsaavidt den kun yttrer sig som heroisk Handlekraft; thi en saadan
% --- p. 73 ---
Handlekraft kan jo ogsaa udspringe af en ganske anden Kilde. Modsætningen
imellem det Overnaturlige og det Naturlige, det Hellige og det Verdslige, den
reneste Selvfornegtelse og den kraftigste Egenvillie savner endnu sit bestemte
Udtryk. Troeslivet kræver derfor en ganske anden og renere Aabenbarelse, at
Sjælen i sand Jomfruelighed kan føle sig gjennemtrængt af det hellige Ord, og
see sig selv i Guds eget Lys. I denne Forstand vidner Evangeliet om Maria: her
er Mere end Judith, her er Mere end Deborah!“

„Kun Den, der forstaaer Troens Fuldendelse i Israel, forstaaer Maria; thi det,
hun gjemmer i sit Hjerte, er --- Troens egen Hemmelighed. Sindrige Theologer
vække en dunkel Anelse om den Benaadedes Naturbestemmelse, idet de lade hendes
Gemyt udtrykke Qvindens rene Modtagelighed, og see i hende den ædleste Blomst
paa Hebræerfolkets Livstræ. Herimod er forsaavidt Intet at indvende: Maria er
tilvisse en folkelig Blomst, den skjønneste Blomst i Palæstinas Have, det rene
Hjerteskud af Qvindelighedens Rod. Men med dette mystiske Billedsprog kommer man
dog ikke videre end til den Antydning, at Guds Søns Moder maa fremfor alle
Qvinder have været en reen, en elskelig, en livsalig Natur. Men hvilken Gaade!
Er det en hellig Natur, hvorom Talen er, saa er Maria jo syndefri, saa
% Printer's error: „forverle“ as printed. At 600 dpi the sixth letter is a plain
% r, with no trace of the tail by which this fount's x descends below the
% baseline --- it is indistinguishable from the r of „for“ in the same word. The
% sense requires „forvexle“. Same defect class as „strar“ for „strax“ on p. 56;
% both OCR witnesses read „forverle“. Transcribed as printed.
forverle vi jo Moderen med hendes Søn, og falde tilbage i Legendens barnagtige
Tankeleg. Skal hun være en syndig, men dog alligevel reen og skyldfri Natur, saa
støde vi jo her paa en Modsigelse, som ingen Menneskeforstand synes at kunne
udholde. Forlade vi derimod det mislige Spil med Naturbestemmelser, og sige det
med rene Ord: Maria er hverken en usyndig Natur, ei heller et i Syndighed bundet
Væsen, men hun er --- \emph{den troende Sjæl}; saa forandres Opgaven. Thi heri
ligger jo ikke blot, at den Benaadede er mere troende, end fromme Qvinder i
Almindelighed,
% --- p. 74 ---
eller at hun overhovedet har en usædvanlig stærk Tro; men det fremhæves
udtrykkelig, at hun har Troen selv som Tro, at hendes sande Væsen er Eet med
Troens Væsen, at altsaa enhver af hendes Sjælebestemmelser falder umiddelbart
sammen med en Troesbestemmelse. Nu først bliver det en religiøs Opgave at
forstaae, hvo den Benaadede er; og for denne Forstaaelse vil den Troende stride
indtil det Yderste, thi Troen vil forstaae sig selv. „Men,“ indvender den
raisonnerende Forstand, „er det ikke Andet, saa skulle vi snart blive færdige
med den Sag. Er Maria hverken Meer eller Mindre end den troende Sjæl i reen
Eenfoldighed og eenfoldig Reenhed, saa maa Enhver, der har nogen Tro, vel
sagtens kunne forstaae hende; den Troende er troende, det følger jo af sig selv:
det Lige forstaaes af det Lige.“ --- Men hvilken Gaade! I det første Tilfælde
blev Forstaaelsen saa vanskelig, at den overgik menneskelige Kræfter; i dette
sidste Tilfælde falder den saa let, at her, egentlig talt, Intet bliver
forstaaet, efterdi her slet Intet er at forstaae. Maaskee kan en psychologisk
Undersøgelse forhjælpe os til en dybere Opfattelse; lad os dog idetmindste
forsøge det.“

% The letterspacing takes in „den“, which stands at the end of the line: verified
% at 600 dpi, where „da“ before it keeps ordinary spacing.
„I qvindelig Umiddelbarhed er Troen en Følelse, en inderlig stærk, men
uforklarlig, uudsigelig Følelse. Den ubestemte Følelse maa imidlertid nærmere
bestemmes ved Det, der føles; og, da \emph{den troende Sjæl} veed sig som Led af
en syndig Verden, maa hun i Forhold til Hellighedens Gud vel ogsaa føle sin
Skyld. Lad da Psychologen forsøge det, om han kan beskrive os denne
ubeskrivelige Følelse. „Ja, Maria føler sin Skyld. Et Barn, der i al sin Uskyld
ikke ændser nogen Fare, ikke tynges af nogen Brøde, fornemmer hun dog en
Sjæleangst, en Frygt og Smerte, som kun tilfulde kan kjendes af den bodfærdige
Synderinde. Det er Lovens Trudsler, det er Propheternes Straffedomme, det er
Davids Klageraab, der forfærde hende; hun føler det i sit Hjerte, at
% --- p. 75 ---
\emph{Alle ere afvegne, at der er Ingen, som gjør Godt, end ikke Een;} ja hun
føler det i uforstaaet Angst, at Israels Gud er en nidkjær Gud, og hans Vrede en
fortærende Ild; hun føler det dybt (thi hendes Væsen er jo Følelse); hun føler
de Lidendes Trang til den Eviges Barmhjertighed; hun føler de Frommes Forlængsel
efter Israels Befrielse; der gaaer et Suk igjennem hendes Sjæl, et stille Suk,
men saa dybt og tungt, at Den, der forstod dette uforstaaede Suk, maatte sige
ved sig selv: o den Tungsindige, hun er ulykkelig!“ --- Men nu er „den troende
Sjæl“ jo ikke blot i et umiddelbart Forhold til den straffende Retfærdighed, men
i et ligesaa umiddelbart Forhold til den forbarmende Kjærlighed. Ogsaa denne
Følelse maa Psychologen kunne beskrive.
% „i Løn“ (in secret): the last letter is n or u --- the two are separated in this
% fount only by the breve over the u, and at 600 dpi the mark above this sort is
% no darker than the show-through around it. Both OCR witnesses read u, which
% would give the non-word „Løu“; „Løn“ is the reading adopted.
„Nei, Maria er ikke tungsindig, skjøndt hun sukker i Løn over Menneskenes Synd.
Hendes Suk henaandes som Barnets Suk, Sorgens Skygge gaaer den Uskyldige forbi:
hun aander igjen saa sundt og let som Den, der eier et saligt Haab, en unævnelig
Glæde. Forjættelsens Ord staaer skrevet i Loven, i Propheterne, i de trøsterige
Psalmer; saa skriver hun da dette Ord dybt i sit Hjerte, og Hjertet forstaaer
det aabenbarede Ord: her er da det usvigelige Pant paa den frelsende Naade. Ja
hun føler det (hendes Væsen er jo Følelse), at Frelsen kommer fra
\emph{Abrahams Gud}, fra den Gud, der \emph{bruger Magt med sin Arm, at neddrage
de Mægtige af deres Throner og ophøie de Ringe.} Hun føler det (thi hendes hele
Væsen er jo Følelse), at \emph{Guds Barmhjertighed varer fra Slægt til Slægt mod
dem, som frygte ham.} Ogsaa derfor kaldes hun jo den Benaadede, fordi hun
fremfor alle Sjæle saa reent og dybt, saa stærk og inderlig har følt den
guddommelige Naade.“ --- Saaledes har da Psychologen angivet de to modsatte
Følelser, hvori den religiøse Grundfølelse nødvendig maa skille sig; og Maria
maa vel fra Grunden af kjende dem begge, efterdi hun er --- \emph{den troende
% --- p. 76 ---
Sjæl.} Den, der ikke kan forstaae hendes Angst og Smerte, vil heller aldrig
forstaae hendes Haab og Glæde. Tag Smerten bort, og den tillidsfulde Tryghed er
kun en ungdommelig Freidighed, kun en kjødelig Sikkerhed uden religiøs Dybde og
hellig Alvor: alle hendes Trøstegrunde fare da hen som krafteløse Talemaader.
Tag Glæden bort, og al hendes Sorg over Verdens Synd og Folkets Elendighed er
kun et tungsindigt Sværmeri; en sygelig Grublen formørker hendes Væsen. --- Men
hvilken Gaade! Skulde disse modsatte Følelser maaskee svække hinanden gjensidig,
saa at Sjælen kun tildeels fornam Frygtens Smerte, fordi den nu ogsaa tildeels
fornam en Lindring i Haabet; saa maatte den Benaadede jo sygne i Troen, og leve
hen i halve Følelsers afmattende Ubestemthed? Eller skulde det være muligt, at
„den troende Sjæl“ kunde i eet og samme Nu føle en saadan Smerte og dog en
saadan Glæde, en saadan Frygt og dog et saadant Haab? Vilde ikke den ene af
disse stærke Følelser qvæle den anden, og den Benaadede finde sig fortabt i en
uhyggelig Stemning, en dump, beklæmmende Rugen, der gjør det umuligt at aande?
--- Lad da de modsatte Følelser skille sig ad i øieblikkelige Gjennembrud; lad
„den troende Sjæl“ bevæges i lidenskabelig Uro imellem begge Yderligheder:
Frygtens ubeskrivelige Smerte, Haabets unævnelige Glæde; lad hende synke brat i
Sorgens Mørke, stige brat i Glædens Henrykkelse: hvad er hun da med sine stærke,
overvældende Følelser? Hun er en gjærende Mulighed, istand til at vove Alt efter
Tilfældets Vilkaar og Øieblikkets Indskydelse; men saa er hun jo ogsaa i selviske
Lidenskabers Vold, og kan ikke være --- Herrens Tjenerinde. Hun er en heroisk
Mulighed, der kun venter paa et Skjæbnens Vink, for at føle sig istand til Alt,
istand til at offre sig paa Baalet for Israels Frelse, istand til en Judiths
List for Jerusalems Ophøielse; men saa er hun ogsaa en heftig,
% --- p. 77 ---
urolig Natur, og kan da ikke være den i Troen Forklarede, ikke den af Himlen
Benaadede, ikke den iblandt Qvinderne Velsignede.“

„Den religiøse Smerte og den fromme Fortrøstnings inderlige Glæde maae, eftersom
enhver af disse Stemninger hersker umiddelbart, vistnok yttre sig i aldeles
modstridende Tilstande: Smerten nedtrykker, Glæden opløfter; Smerten forkynder
den straffende Vrede, Glæden den milde, forbarmende Naade. „Den troende Sjæl“ er
altsaa i Strid med sig selv. Ogsaa denne Strid lader sig fremstille ved Hjælp af
den beskrivende Psychologie. Smertens første umiddelbare Gjennembrud kan saare
let bringe et ungdommeligt Hjerte til at forsage og mistvivle om Guds
Barmhjertighed, just fordi Indtrykket er saa stærkt, og det uerfarne Gemyt
overvældes af Følelsen. (At den i Angeren nedtrykte Sjæl ikke kan anklage sig
selv for nogen enkelt Brøde, kommer slet ikke her i Betragtning; det Afgjørende
er, at den ikke usyndige Natur føler sig i sin Syndighed ligeoverfor den
Hellige). Denne Følelse virker sønderknusende, efterdi den virker med
Umiddelbarhedens hele usvækkede Kraft. Men idet Fortvivlelsen nærmer sig,
vaagner Selvbesindelsen; Villien strider for at vinde et Haab, og det blødende
Hjerte tilegner sig Forbarmelsen. Da glimter Lyset, en Straale af Guds
Kjærlighed adspreder Mørket, og den nedtrykte Sjæl hæver sig --- i salig
Fortrøstning. Opløftelsen er imidlertid ogsaa umiddelbar, og Fortrøstningens
Fred en forhen ukjendt Følelse; derfor er Tilliden saa ubegrændset, og Glæden
overvættes. Ifølge den ungdommelige Forventning skulde den haardeste Strid altsaa
nu være forbi, og Faren overstaaet. Men dette Haab vil skuffe, og denne
Skuffelse er saa meget desto smerteligere, som den synes at forraade et Affald
fra Saligheden selv. „Hvad kan dog dette betyde, at Himmelens Fred besøger et
Menneske, for igjen at drage bort og lade Afmagt og Tomhed og Tvivl og Uro
tilbage?
% --- p. 78 ---
Er da det Himmelske ustadigt som det Jordiske? Var det ikke Gud selv, der
aandede paa dette Hjerte og forkyndte sin Tilgivelse? Har en ny Skyld maaskee
fortørnet ham, at han nu straffer i Vredens Kulde, eller er det en formastelig
Indbildning, at Kjød og Blod vil leve sammen med det Hellige? Kan den Skrøbelige
gjøre Retfærdigheden Fyldest? Kan Guds Kjærlighed ville dræbe en Sjæl? Umuligt,
Umuligt!“ --- Her have vi atter Smerten. Det er vanskeligt at afgjøre, hvorvidt
den følgende Overgang fra den fornyede Smerte til den gjenvundne Glæde
forudsætter en endnu haardere Strid end den, der begyndte med det første
Gjennembrud; men det er vist, at Troen kun lever i den fortsatte Strid. Det er
vist, at Troen forudsætter og virker den Strid, hvori Angerens Smerte og
Forsoningens Glæde idelig adskilles som modstridende Stemninger, efterdi
Nedtrykkelsen altid betinger Opløftelsen; men det er ogsaa vist, at den igjen
strider for at ende Striden.“

„Saa maa den troende Sjæl da lide mangen Anfægtelse og gjennemgaae mange Farer,
inden den rette Forstaaelse kan komme; saa maa Maria vel ogsaa, fremfor alle
Andre, være forsøgt i haarde Prøvelser og have stridt de indre Erfaringers lange
møisommelige Strid med Gud og sig selv, før hun kunde vinde Forstaaelsen, og
naae Fuldendelsens Trin, og vorde hilset af Engelen som den Benaadede? Det
skulde altsaa i denne Forstand være Alderdommens skjønne Forret, at eie Troens
rette Forklaring. Men hvilken Gaade! Det er jo ikke Maria, der er os fremstillet
som den Bedagede; men det er Elisabeth. Om Elisabeth kan det gjælde, at hun
igjennem Striden trinviis arbeider sig opad mod Fuldendelsen. Den beskrivende
Psychologie, der ikke kan naae Maria, kan desaarsag vel naae Elisabeth. Da nu
Elisabeth forstaaer Maria, saa vende vi os igjen til den Bedagede, om vi maaskee
under hendes Veiledning kunne gjøre et nyt Indblik i
% --- p. 79 ---
Troens Væsen. „Hvad er det altsaa for en vidunderlig Sjælekraft, der forsoner
% Printer's error: „vore“ as printed, the third instance of this book's r-for-x
% (cf. „strar“ p. 56, „forverle“ p. 73). At 600 dpi the r ends on the baseline
% with a plain flag and shows none of the x's descending tail. The sense
% requires „voxe“. Both OCR witnesses read „vore“. Transcribed as printed.
Angerens Smerte med Fortrøstningens Glæde, at de begge vore tilsammen, saa at
Glæden selv nærer sig af Smerten, og Smerten forvandler sig til en inderlig
Veemods stille Glæde?“ Spørg den Bedagede, der har stridt den gode Strid, og
Svaret vil lyde kort og bestemt: „det er den sande Ydmyghed.“ --- Stridens
umiddelbare Overgang i Forsoning, Sindets Frigjørelse fra det Selviske i
ubetinget Underkastelse under Hellighedens Lov, er given med Ydmygheden, og
Ydmygheden er jo, som Jomfrueligheden selv, paa een Gang baade en Troens og
Qvindelighedens Idealbestemmelse. Elisabeth, der af Erfaring har lært, hvad
Ydmyghed er, veed ogsaa, at de mange endelige Ydmygelser, Menneskene pleie at
gjennemgaae, for at vinde den sande Ydmyghed, ikke kunne undskyldes med vor
Naturs almindelige Syndighed, men ere en Følge af Hjerternes Haardhed. Derfor
roser hun sig heller ikke, hverken af sine mange Prøvelser, eller af sin lange
Strid, eller af sin høie Alder; men hun seer med ydmyg Glæde op til den
Benaadede, og vidner for hende, og udbryder: \emph{salig hun, som troede.}
Staaer det nu fast, at Ydmygheden selv kan være ligesaa umiddelbar som Troen, og
at Troen derved faaer den vidunderlige Gave, at kunne være baade Striden selv og
Stridens Forsoning, saa har Elisabeth løst os Gaaden, idet hun vidner om Maria:
\emph{salig hun, som troede.} --- Det er ikke de overordentlige Naturgaver, ikke
det rige Indbegreb af alle qvindelige Fuldkommenheder, vi her skulle beundre;
Maria er jo ikke kaldet til at glimre i Historien, men til at være Herrens
Tjenerinde. Hun er ikke bestemt til at gjøre Lykke i Verden, men til at finde
Naade hos Gud. Det er da kun den ene Fuldkommenheds Gave, den Gave, hvorved en
Qvindes Sjæl finder Naade hos Gud, den livsalige Gave at være reen for den
Hellige i Hjertets Ydmyghed,
% --- p. 80 ---
hvorom Elisabeth vidner, idet hun hilser den Benaadede: \emph{salig hun, som
troede.} Ja, salig hun, som troede! --- Store Gjerninger har Israels stærke Tro
i Nødens og Befrielsens Stund fuldbragt til Jehovas Ære; men kun i Marias Hjerte
virker Troen den fuldkomne Gjerning med Inderlighedens underfulde Kraft. Thi
dette er jo den fuldkomne Gjerning, at en jomfruelig Qvinde, hvis Hjerte
fornemmer en Verdens Nød, dog lever i Verden eenfoldig tryg, som Den, der ikke
kjender til Fare; at hun, hvis Sjæl kan føle en Verdens Smerte, dog lever i
Verden eenfoldig glad, som Den, der ikke veed, hvad Smerte betyder. Dristige,
herlige Ord har Israel talet i Troens kraftige Sprog; det lyder til os i
Smertens Skrig, i Salighedens Priis, i Forfærdelsens Vee-Raab, i Fortrøstningens
overtalende Styrke; i de Lidendes Suk, i de Seirendes Jubel, i de Nedtryktes
Klage, i de Frelstes høie Taksigelse, som intet andet Sprog i Verden. Men det
sande Troens Ord, det Ord, der fødes dybt i Menneskehedens Hjerte, og hæver sig
op af Sjælens stille Grund, det rene Troens Ord i Gudhengivenhedens inderlige
% Bold Fraktur, not a larger fount: at 600 dpi the caps of „See!“ and „Herrens“
% stand exactly as high as the body caps on the same line („Engelen“, „Marias“),
% but every stroke is markedly heavier. The closing “ --- which shuts the „
% opened on p. 78 („Saa maa den troende Sjæl …“) --- is inside the bold.
Sprog, er dog Marias Ord til Engelen: \textbf{See! jeg er Herrens Tjenerinde,
mig skee, som du haver sagt.“}

% Centred ornament: three asterisks in a triangle (two above, one below),
% verified at 600 dpi under autocontrast --- the upper pair is as black as the
% body type and is not show-through. Text-mode asterisks, never $*$: math mode
% would set U+2217 ASTERISK OPERATOR, as the preamble note on \thefootnote warns.
\begin{center}
  *\quad*\\[-4pt]
  *
\end{center}

% Joh. 1, 1--5 closes the chapter in the larger display fount, at full measure
% and with the ordinary paragraph indent, exactly like the pericopes of pp. 1--2
% and Luc. 1, 15--17 on p. 42; hence \begingroup\large … \par\endgroup rather
% than the inline {\large …} used for quotations that end mid-line.
\begingroup\large
„I Begyndelsen var Ordet, og Ordet var hos Gud, og Ordet var Gud. Det var i
Begyndelsen hos Gud. Alle Ting ere blevne ved det, og uden det er ikke een Ting
bleven til af det, som er. I det var Livet, og Livet var Menneskens Lys, og
Lyset skinnede i Mørket, og Mørket begreb det ikke (Joh. 1, 1--5).“
\par\endgroup

% The chapter-break rule falls at the FOOT OF P. 80, not on p. 81: a double rule
% (two hairlines), centred, measured at 600 dpi as 800 px against a 2100 px
% measure = 0.38, i.e. identical to the double rule that precedes the head of
% chapter II on p. 42. It is written here because it is printed here. The head
% „III. Den hellige Familie.“ and the short rule beneath it stand at the top of
% p. 81 and belong to the next batch, which must NOT repeat this rule.
\begin{center}
  \rule{0.38\textwidth}{0.4pt}
\end{center}
% --- p. 81 ---
% The head of chapter III stands at the VERY TOP of p. 81, immediately under the
% running number: „III.“ on a line of its own, then „Den hellige Familie.“, then a
% short single rule. There is NO scripture reference line under this head (unlike
% chapter II on p. 42, which prints „(Luk. 1, 26--38. Matth. 1, 18--25.)“ between
% head and rule). The double chapter rule belongs to the FOOT OF P. 80 and is
% written in the previous fragment; it is not repeated here.
%
% The printed head agrees with the Indhold („III. Den hellige Familie ... 81--183“).
% As in chapters I and II the numeral stands on a line of its own above the title,
% the printed full stop is kept in the \section* and dropped in the
% \addcontentsline.
\section*{III.\\[4pt] Den hellige Familie.}
\addcontentsline{toc}{section}{III. Den hellige Familie}

% The rule under this head is SHORTER than the one under the head of chapter II.
% Measured at 600 dpi on PDF p. 98: rule 176 px (x 1078--1254, y 1141--1146)
% against a body measure of 2104 px (left 141, right 2245) = 0.084, i.e. c. 7.5 mm;
% the corresponding rule on p. 42 was measured at 0.12 (c. 11 mm). The difference is
% real and is reproduced.
\begin{center}
  \rule{0.08\textwidth}{0.4pt}
\end{center}

Det er en Virkelighedens Lov, at Alt, hvad der gjør Fordring paa
Selvstændighed, maa have sit Sted. Derfor sige vi jo om enhver Yttring, der ved
nærmere Prøvelse viser sig uberettiget, at den \emph{intetsteds} hører hjemme.
Vel er den evige Aand allestedsnærværende, og følgelig ikke bunden til noget
enkelt Sted; men i sin virkelige Forskjel fra Verden maa Gud dog alligevel siges
at have sit særegne Sted; hvorfor Religionen jo ogsaa, og det med Rette, lærer
os, at Gud har sit Sted i Himmelen, at han boer deroppe i Evighedens og
Salighedens Hjem. Menneske-Aanden, skjøndt den i sin frie Stræben og uendelige
Higen ligeledes overskrider enhver udvortes Stedbegrændsning, har naturligviis
ogsaa sit Sted; thi den boer jo med Menneskene hernede paa Jorden. Der kan da
ikke blot tales om et sandseligt Stedforhold imellem den synlige Himmel og den
synlige Jord; men der kan i dybere Forstand tales om et oversandseligt
Stedforhold imellem den himmelske Aand og Aanderne, der leve paa Jorden.
Ligeledes kan der, hvad ethvert Folkeslag angaaer, ikke blot spørges om Stedet,
hvor det boer eller har boet i den udvortes Verden; men der kan og maa ogsaa
spørges om dets aandelige Sted ɔ: om den bestemte Plads, det indtager i
Folke-Aandernes Række. Hvor det israelitiske Folk i sandselig Forstand havde sit
Sted og sit Hjem, veed Enhver, der har nogen Kundskab til Landenes Beliggenhed
paa Jordkloden; men, naar dette Folks aandelige Sted skal nærmere bestemmes,
begynder Troen at tale et Sprog, hvori Verdensforstanden ikke ret kan finde
nogen Mening.
% --- p. 82 ---
Aabenbaringstroen paastaaer nemlig, „at den hebræiske Folke-Aand boede lidt
høiere end alle de øvrige i Aandernes Rige; at den just boede paa det
vidunderlige Sted, hvor Timeligheden grændser saa nær op til Evigheden, at et
dødeligt Øre tydeligt kan høre Guds egen Røst, og fornemme en Susen af
Seraphens Vinge; at den hørte hjemme i den Verdensegn, hvor Solen studser, og
Maanen standser for en Bedendes Ord; i det Land, hvor Skyerne bære Cherubens
Vogn, og en Jakobsstige naaer fra Jorden til Himlen.“ Det er nu dette aandelige
Sted, der har givet det sandselige sit Navn, og derfor er der et Land paa
Jorden, vi kalde --- \emph{det hellige Land}.

Det er en Virkelighedens Lov, at enhver selvstændig Virken maa have sin Tid.
Naar det nu desuagtet hedder, at Gud er ophøiet over enhver Tid, at tusinde Aar
ere for ham, som den Dag igaar, der er forbigangen; da udsiges hermed en
væsentlig Forskjel imellem den evige og den timelige Tid; men ingenlunde, at
Tidsbestemmelsen selv skulde falde udenfor Guds hele Væren og Virken. En
Guddom, der er ligegyldig mod Tiden, kan ikke være Andet end Tilværelsens
tidløse Grund; men en saadan Grund maa jo netop, efterdi den i sig selv er
tidløs, laane alle Tidsbestemmelser fra Endeligheden. Heraf følger da igjen, at
en tidløs Gud umiddelbart synker ned i den endelige Tid, og bliver en
Timelighedens Tjener, istedetfor at magte alle Tidsbestemmelser, og være
Tidernes Herre til evig Tid. Maa nu Alt, selv Evighedens Gud, have sin Tid, da
gjælder det Samme om den virkelige Aabenbaring. „Der var en Tid, da Gud steg ned
fra Himmelen og paatog sig synlig Skikkelse; der var en Tid, da den Evige selv
talte Timelighedens Sprog, og Guds Engle besøgte Menneskenes Børn.“ Denne
Begyndelse er tilstrækkelig til at reise en vidtløftig Strid imellem Tro og
Vantro. --- Hvilke Skjæbner og Omskiftel\-%
% --- p. 83 ---
ser det israelitiske Folk, verdslig talt, har gjennemgaaet i Tidernes Løb;
hvilken Kraft og Svaghed det har lagt for Dagen i de lykkelige og ulykkelige
Tider, dette Spørgsmaal kan meget vel afgjøres for Videnskabens almindelige
Domstol; og Enhver, der kjender Historien tilgavns, er berettiget til at tale
med. Men, naar det derimod gjælder en Undersøgelse om, hvorvidt disse store
Tidsforandringer ere at tilskrive de timelige Verdensbetingelser alene, eller
væsentligen at henføre til Styrelsens forudbeskikkede Raad i den evige Tid, da
bliver den menneskelige Viden splidagtig med sig selv, og uddeler Vaaben til
begge Sider. Det er jo muligt, at de Syner, om hvilke den religiøse Overlevering
beretter, kun have været saadanne Billeder, som henkastes af en from, men
levende og stærk bevæget Indbildningskraft; men det er jo ogsaa muligt, at hine
Guddomssyner have til visse Tider fremstillet sig for Sjælens Øie med en saa
ubetinget Nødvendighed og et saa afgjørende Virkelighedspræg, at de i Kraft af
deres overnaturlige Aarsag gjøre Krav paa samme Gjenstændighed, som den, vi
tillægge de udvortes Ting, fordi disse uvilkaarligt paanøde sig vore Sandser.
Det er jo muligt, at det Guddomsord, der taltes i Israel, dog i Grunden ikke har
havt anden Kilde, end den, der opvælder i ethvert Folkegemyt, naar Religionens
Idee rører sig kraftig derinde. Men det er jo ogsaa muligt, at den Alt virkende
Gud selv engang i fordums Dage har formet sine Evighedstanker i Menneskenes
timelige Sprog; at dette menneskelige Guddomsord har til sine Tider lydt saa
klart og reent, saa tydeligt og bestemt for Sjælens Øre, at den Hørende ligesaa
lidt turde kalde det sit eget Ord, som den Lærling, der har hørt og forstaaet
sin Mester, med Sandhed tør sige om sig selv, at han er Lærens Opfinder.
Gjælder det nu om Lærlingen, at han kun beskjæmmer sig selv, naar han negter
Mesteren den ham tilkommende Ære, naar han istedetfor at vise
% --- p. 84 ---
tilbage til Grundtankens Oprindelse, sigende: „saaledes taler den vise Mester“,
kun forvansker det oprindelige Ord ved at ville fremsætte det i sit eget Navn,
og forvirrer de urkraftige Tanker ved at omskrive dem i sit eget matte,
farveløse Sprog: da gjælder dette endnu i en langt anden og høiere Forstand,
naar Gud selv har paataget sig at være Menneskenes Lærer. Det er da ogsaa
muligt, at Moses og Propheterne ikke kunde, og i Sandhedens hellige Navn heller
ikke turde oplade deres Mund til Folket, uden at begynde med de Ord: „Saa taler
Jehova, Israels Gud.“ Saaledes staaer Mulighed mod Mulighed til begge Sider.
Tvivleren halter og sygner; men Troesfornegteren og den Troende vide, hver paa
sin Viis, at træffe et Valg. Aabenbaringen har sin Virkelighed og sine bestemte
Tider; denne Troens Grundtanke kan sprænge Naturlærens Skranker, som Gnisten i
Krudt kan sprænge en Klippe. „Aabenbaringen har sine egne Tider. Der var en Tid,
da Jehova gjæstede Abraham i Mamre Lunde; der var en Tid, da Aandens Himmel
aabnede sig, og det himmelske Tabernakel blev Forbillede for det jordiske; der
var en Tid, da Propheterne i Israel dreves af Guds hellige Aand, og kunde, midt
under Endelighedens mangehaande Misviisninger, saa sikkert stole paa dens
ufeilbare Indskydelser; der var en Tid, da den hebræiske Folke-Aand trods alle
Synder og Overtrædelser blev indviet til at bære det evige Ord gjennem Tidernes
% Printed „Omverlinger“ for „Omvexlinger“ --- a genuine r, not the Fraktur x.
% Verified at 600 dpi (PDF p. 101): the glyph after „Omve“ is identical to the
% final r of the same word and has no stroke descending below the baseline, unlike
% the true Fraktur x of „f. Ex.“ (p. 87, p. 93) and „(Extremer)“ (p. 90), each of
% which drops a diagonal below the line. Both OCR witnesses read r. This joins the
% book's documented series: „vore“ (pp. VIII, 79), „strar“ (pp. 56 and 86),
% „forverle“ (p. 73). Transcribed as printed.
Omverlinger indtil Tidens Fylde.“ Det er nu dette overnaturlige Ord, der har
givet den naturlige Folkecharakteer sit Navn, og derfor blev det sagt til det
udvalgte Folk, at det skulde være et Kongerige af Præster, --- \emph{et helligt
Folk}.

Det er en Virkelighedens Lov, at Alt, hvad der vil frem i Tilværelsen, maa have
Kraft til at bryde sig igjennem. Naar Sædekornet bryder, og Løvet springer ud,
har Naturen jo
% --- p. 85 ---
Foraar; naar Følelserne bryde i Sindets Grund, og stærke Tanker berede sig til
Gjennembrud, bebudes en aandelig Vaar. Naturens Elementer, Aanderigets
forskjellige Magter brydes med hinanden overalt, saa langt som Virkeligheden
naaer. Det nye Liv begynder med et Brud, og Døden brydes igjen med Livet. Og
hvad er Aabenbaringen med dens Udvælgelse vel Andet end et Brud? Det er et
Gjennembrud af det hellige Lys, der bringer Abraham til at bryde de naturlige
Slægtskabsbaand, og forlade sine Fædres Land. Og som det gaaer med
Stammefaderen, saaledes ogsaa med Slægten. Jehova brydes med sit Folk. Israel
er i mere end een Henseende den Mand, der --- brydes med Gud. Ved idelige Brud
gjentages Udvælgelsen indenfor de Udvalgtes Kreds, saa at denne, trods Slægtens
Formerelse, bliver snævrere og snævrere. Først falder Ismael bort, som et
Brudstykke, som en Green, der brydes af et Træ, for at plantes i en anden
Jordbund; siden skeer der et nyt Brud, hvorved Esau skydes til Side, at
% The Fraktur capital I/J is a single sort; „Ismael“, „Israel“ and „Isai“ are
% settled from the sense, as „Iagttager“ was on p. 22.
Velsignelsen kan hvile desto rigere over Jakob. Men ogsaa indenfor Jakobs Huus
vedbliver Udsondringens Brud; thi i Isai Stamme tilsigter Aanden et høiere
Gjennembrud, og tilsidst bryder det afgudiske Israel med Juda, der under
mangehaande Afbrydelser vedbliver at haabe paa Forjættelsen. Hvad der saaledes
skeer med det hele Folk, gjentager sig i visse Maader med de enkelte Udvalgte:
de forklares ikke i sædelig Fuldkommenhed; men Guddomsstraalerne bryde igjennem
den naturlige Sjæl, som Solens Straaler igjennem den mørke Sky. Det er denne
Aabenbaringens Straalebrydning i Naturaandens dunkle Medium, der giver Folket
med dets Udvalgte sin eiendommelige Charakteer. „Lyset skinner i Mørket, og
Mørket begriber det ikke.“ Jehovas Lys opgaaer over Israels Børn, som Solens Lys
over Judæas Bjerge: kun enkelte fremragende Punkter tindre i Dagens Glands; men
% --- p. 86 ---
mellem disse vise sig --- dybe Kløfter, natlige Skygger. --- Naar det indsees,
at Udvælgelsen skeer ved et Brud, da forstaaer den Troende, at dette just er
Betingelsen for, at Ordet kan blive Kjød, og fremtræde som selvstændigt Led i
Slægten. „Havde det været Formaalet, at hævde den hellige Lovs almindelige
Overmagt over Menneskenaturens Drifter og opsætsige Egenvillie, da skulde intet
Brud i Slægtlivet selv have været fornødent; da kunde den Almægtige jo
% Printed „strar“ for „strax“: a genuine r, verified at 600 dpi (PDF p. 103) ---
% no descender below the baseline, and the glyph is identical to the r two letters
% earlier in the same word. The same reading was verified on p. 56. Both OCR
% witnesses agree. Transcribed as printed.
strar have aabenbaret sig i Skyerne med Kraft og Herlighed, og som
Menneskeslægtens Gud forkyndt sine Trudsler og Forjættelser for hele
Menneskeheden, istedetfor at nedlade sig til at blive en Familiegud og begynde
med i al Stilhed at kalde Abraham til en Side, og udvælge ham alene af Massen.
Hvis det havde været Formaalet, at blande det Hellige sammen med det Verdslige i
tvetydig Eenhed, da skulde ethvert Udvælgelsens Brud været overflødigt; da kunde
den skabende Aand jo have omspredt sine Funker i alle Folke-Aanderne, udstrøet
sin Sæd ligeligen i alle Sjæle, og derved bevirket en Gjæring uden Maal og
Grændse. Men, naar det derimod er Udvælgelsens Endemaal, at Guds egen Søn skal
fødes i Slægten som et Menneskes Søn, da maa der beredes ham et Sted i Slægtens
egen Midte. Alle Udsondringer og Gjennembrud tilsigte derfor en sidste
Udsondring, forat den lille udvalgte Kreds kan danne sig, der omslutter den
Reneste i Slægten, den Benaadede, i hvem det hellige Ord selv kommer til
Gjennembrud.“ Det udvalgte Folk har saaledes den væsentlige Bestemmelse at danne
Grundlaget for --- \emph{den hellige Familie}.

% PRINTER'S DEFECT: the closing “ of „hellig“ is not printed. The line ends
% „... paa Tillægsordet „hellig“ and breaks there, well short of the right margin,
% with clear white space after the g; verified at 600 dpi under autocontrast
% (PDF p. 103). The parallel „Familie“ on p. 87 is properly closed. Transcribed as
% printed: this is the one unclosed „ in pp. 81--94, and the whole batch therefore
% runs one „ ahead of its “.
Den hellige Familie! Dette Begreb er sammensat af tvende ueensartede
Bestanddele, og kræver desaarsag en omhyggeligere Betragtning. Lægger man
Eftertrykket paa Tillægsordet „hellig alene, da faaer Aandsbestemmelsen derved
en saa eensidig Overvægt, at det naturlige Moment, der ligger i Hovedordet,
aldeles
% --- p. 87 ---
taber sin Betydning. Udtrykket: „den hellige Familie“ kommer da væsentlig til at
betegne et Samfund af hellige Medlemmer. At disse Medlemmer forøvrigt staae i
Familieforhold til hinanden indbyrdes, kan ansees for tilfældigt og ligegyldigt,
eftersom det jo ikke er Slægtskabsforbindelsen, der her helliger de Enkelte, men
omvendt de Enkeltes Hellighed, der giver Familien sit Tilnavn. Men er
Familiebegrebet paa denne Maade forflygtiget, saa bliver det med det Samme
umuligt, at forbinde noget klart Begreb med Benævnelsen: „hellig“. Medlemmerne
af den Familie, hvorom vi her tale, nemlig: Zacharias og Elisabeth, Maria og
% Only „hellige“ is letterspaced in this three-limbed sentence; the parallel
% „særegne“ and „allerhøieste“ are set plain (checked at 600 dpi, PDF p. 104).
% Not normalised into agreement.
Joseph, ere jo hverken \emph{hellige} i den lavere Forstand, hvori man til
dagligt Brug skjødesløst og halvt spottende siger om en Opvakt, at han er af de
% „f. Ex.“: a true Fraktur x --- at 600 dpi the glyph after the capital E throws a
% diagonal below the baseline, which the r of „naar“ and „taler“ on the same line
% does not. Contrast „strar“ (p. 86) and „Omverlinger“ (p. 84), where the glyph has
% no such descender. Both OCR witnesses read „Er.“ at every one of these sites and
% are therefore useless here; the reading is settled on the image alone.
Hellige; ei heller i den særegne Forstand, hvori Catholiken bruger Ordet, naar
han f. Ex. taler om „de Helliges Forbøn“; endnu mindre i den allerhøieste
Forstand, som naar vi tænke paa den hellige Gud. Lægger man derimod Eftertrykket
alene paa Hovedordet „Familie“, da bliver Naturbestemmelsen saa fremherskende,
at Tilnavnet „hellig“ maa betragtes som lydende paa en af Slægtskabsforholdet
umiddelbart afledet Egenskab. I en saadan Forudsætning ligger da den Mening, at
Medlemmerne af den hellige Familie skulde, hver især, kunne gjøre Fordring paa
at æres som Hellige paa Grund af deres overnaturlige Herkomst. Man kunde altsaa
tale om en hellig Familie omtrent i samme Forstand, som naar der i Mythologien
tales om Slægter, der umiddelbart nedstamme fra de høie Guder. Men, at en saadan
Opfattelse af Begrebet Hellighed staaer i afgjørende Modsigelse til
Aabenbaringens Idee, trænger ikke til yderligere Beviis. --- Den hellige Skrift
har vel ikke udtrykkelig givet den Familiekreds, i hvilken Johannes og Jesus
fødes og opdrages, Navn af hellig (Benævnelsen er fra den catholske Kirke gaaet
over i det almindeligere
% --- p. 88 ---
Kunstsprog); men Evangeliet stiller dog denne Familie i et saadant Lys, at vi
paa Troens Vegne meget vel kunne kalde den med dette betydningsfulde Navn, idet
vi af Aabenbaringen selv lære, hvorledes Benævnelsen maa forklares. Den strenge
Modsætning imellem Naturbestemmelsen og Aandsbestemmelsen har nemlig i sig
optaget det Overnaturlige som væsentligt Mellemled. Naturgaven, der viser hen
til Slægtforholdet, er tillige en Naadegave formedelst Udvælgelsen. Ikke med
enhver af Israels fromme Qvinder foregaaer det Vidunderlige, som med en
Elisabeth; og hvorledes kunde det guddommelige Ord blive Kjød i Marias Hjerte,
naar ikke det overnaturlige Princip formedelst Udkaarelsen havde helliget og
viet den naturlige Grund? I den hellige Familie ere altsaa Individerne af den
ædleste Slægt, ikke blot med Hensyn paa det naturlige Slægtforhold, men fordi de
staae i det allernærmeste Slægtskab med det Hellige. Ligeledes har det
Overnaturlige gjennemtrængt Aandsbestemmelsen; thi vel ere Haabet, Troen og
Selvfornegtelsen Frihedens indre Gjerninger i Mennesket selv; men et saadant
Haab, som det, der lever i denne Familiekreds, en saadan Tro og Troskab, at
disse Individer, skjøndt skrøbelige Mennesker, dog kunne betroes det
Allerhelligste, er kun tænkelig formedelst en ganske særegen overnaturlig
Bistand. De helliges da ved Guds Aand, idet de, hver især, hellige og indvie sig
til at modtage og bevogte det Hellige. Det skjulte Under gaaer saaledes forud
for det aabenbare, Forberedelsen forud for det egentlige Gjennembrud. --- De
Gamle forestillede sig, at Kjæmpen Atlas bar Himlen paa sine brede Skuldre.
Dette er kun en Fabel. Den sandselige Himmel tynger ikke Jorden paa noget Punkt;
ja den har, egentlig talt, ingen Tyngde. Derimod kunde man, billedligt forstaaet
snarere sige, at det er Himlen, der favner og bærer Jorden, eftersom jo Kloden
svæver i det vide Himmelrum.
% --- p. 89 ---
Men endnu mere vidunderlig end den gamle Fabel om Kjæmpen Atlas er dog den
evangeliske Fortælling om Jomfruen Maria, der bærer Verdens Frelser, den
Høiestes Søn, under sit Hjerte. Thi dette forudsætter jo, at Evighedens
aandelige, sande Himmel har paa sit Sted og til sin Tid sænket sig ned til
Jorden og hvilet sig paa en skrøbelig Qvinde; og dog fastholder Evangelietroen
denne Sandhed, fordi den veed, at Frelsen ligger i Almagtens Under. De Gamle
meente, at denne Jord var Universets sandselige Middelpunkt; videnskabelige
Iagttagelser have siden lært, at dette ikke er saaledes. Men Evangelietroen
gaaer endnu langt videre. Den antager, at en ringe, qvindelig Skabning har i
Tidens Fylde været hele Universets Middelpunkt, at alle Straaler af det evige
Lys have samlet sig i hendes Sjæl, som i deres naturlige Brændpunkt; og dog skal
% The „ opened before „videnskabelig Iagttagelse“ is never closed where the sense
% wants it: the line ends flush at the right margin with „Iagttagelse“ and no “
% (600 dpi, PDF p. 106). The only closing “ in the sentence is the faint one after
% „uden Undtagelse.“, which is printed and is reproduced. The quotation therefore
% balances numerically but is mis-paired; not normalised.
ingen „videnskabelig Iagttagelse kunne fravriste Troen denne Overbeviisning,
fordi den veed, at Frelsen er --- Kjærlighedens Under. --- Ved disse og lignende
„sære“ Paastande er Evangelietroen naturligviis nu efterhaanden igjen kommen til
at staae paa en meget spændt Fod, ikke blot med den moderne Bevidsthed og den
saakaldte sunde Menneskeforstand, men med alle videnskabelige Stormagter uden
Undtagelse.“ Kun een Videnskab er funden beredvillig til at antage sig Troens
Sag, nemlig den speculative Philosophie; dog ogsaa denne Coalition har opløst
sig, idet Speculationen efterhaanden er bleven Troen utro, og denne, der snart
mærkede Bedraget, har frasagt sig enhver tvetydig Alliance, og besluttet at føre
Krigen med egne Vaaben.

Hvorledes Misforstaaelsen imellem Philosophien og Troen er opkommen, kan, hvad
den hellige Historie angaaer, kortelig oplyses i denne Sammenhæng. Alt beroer
nemlig paa, hvorledes man opfatter Forholdet imellem Aabenbaringens væsentlige
For\-%
% --- p. 90 ---
beredelse og dens virkelige Gjennembrud. Den speculative Philosophie accentuerer
Forberedelsen, fordi den eensidig holder sig i Ideen; Troen derimod accentuerer
Brudet, efterdi den udelukkende holder sig til Kjendsgjerningen. At accentuere
Forberedelsen er, væsentlig forstaaet, det samme som at drive paa Formidlingen,
at opløse Underet som nødvendigt Led i Ideens indre Selvudvikling, og skjule det
Overnaturlige i Naturlivets høiere Former. Naar da Evangeliet i al Eenfoldighed
beretter om det store, ubegribelige Vidunder, at en Qvinde i Nazaret er bleven
Moder til Guds eenbaarne Søn, som selv er Gud, og Verden udbryder i Forargelsens
% „(Extremer)“ is Fraktur, not antiqua, and the glyph after the E carries the
% Fraktur x's descending diagonal (600 dpi, PDF p. 107); likewise „Extremerne“ on
% p. 92. Both OCR witnesses read „Ertremer“.
tusindstemmige Umuligt! da træder Speculationen frem, for at overbevise de
Vantroe, ved at indvie dem i sin gudmenneskelige Viden. --- De Yderled
(Extremer), hvis Formidling Philosophien overtager sig, ere paa den ene Side:
„den evige Gud, ophøiet over Tiden“, paa den anden Side: „det endelige Menneske
i Tiden“. „Den naturlige Forstand maa vistnok antage det for en afgjort
Umulighed, at Guds evige Ord, som selv er Gud, skulde kunne lade sig føde af et
Menneske, og fremtræde paa Jorden i menneskelig Begrændsning; men den
speculative Ideelære seer sig dog istand til at opløse den haarde Anstødssteen i
reen Fornuftighed. Det Uendelige og det Endelige, det Evige og det Timelige, det
Guddommelige og det Verdslige vise sig nemlig som Sider og Momenter i det ene og
samme Grundbegreb: Mennesket her paa Jorden er just det væsentlige Mellemled,
hvorigjennem hine Yderled slutte sig sammen til sand og fuldstændig Eenhed. Naar
altsaa den hellige Skrift taler om Guds Menneskevordelse i Tidens Fylde, da er
dette saa langt fra at være en overtroisk, vilkaarlig og urimelig Tale, at det i
dybere Forstand er den høieste Viisdom. Ideen om Gudmennesket er den Idee,
hvorom hele Universet dreier sig, den Idee, paa hvis Virke\-%
% --- p. 91 ---
liggjørelse Altilværelsen er anlagt lige fra Tingenes allerførste Begyndelse.
Guds Søn maa altsaa nødvendigviis aabenbare sig, thi hans Komme til Verden er
Skaberværkets egen Fuldendelse; han maa nødvendigviis fødes af en Qvinde, thi
kun under denne Betingelse kan han udskyde som Blomsten paa Slægtlivets Træ.“

Havde Philosophien nu kunnet løse sin Opgave, at omsætte alle evangeliske
Forestillinger i den rene Videns nødvendige Former, standse alle
Sandsynlighedsberegninger, og bringe alle Kritikens Modsigelser til Taushed: da
--- ja da skulde Evangelietroen selv vel ikke have vundet en Hvid ved den hele
Omsætning, men Speculationen vilde derimod have tilkjæmpet sig en
overmenneskelig Seier, og stødt Religionen fra Thronen. Opgavens Løsning var
imidlertid umulig. Philosophien, der begyndte med at bortkaste alle
Forudsætninger, for ved Tankens dialektiske Selvudvikling at fremtvinge
Erkjendelsen af den absolute Idee, saae sig efterhaanden nødsaget til at slaae
Noget af i sine videnskabelige Prætensioner. Evangeliets overnaturlige
Kjendsgjerninger vare altfor talrige, og de enkelte Mirakelfortællinger altfor
anstødelige, til at kunne retfærdiggjøres for den rene Tænkning. Man maatte
altsaa overlade slige „lavere Bestanddele“ til Kritikens Behandling, og anvende
al sin Kraft paa at begribeliggjøre Christendommens store Grundunder, nemlig
Sønnens Menneskevordelse. Da den evangeliske Incarnationsidee imidlertid vedblev
at trykke og tynge paa den rene Viden med det Overnaturliges hele Vægt, begyndte
Speculationen saa smaat at see sig om efter positive Støttepunkter, og vilde nu
% „Akkorden“, not „Afforden“: the two glyphs after the A are Fraktur k, with the
% k's arm and no descender (600 dpi, PDF p. 108). Tesseract's f/k confusion reads
% „Afforden“; the ABBYY layer reads „Akkorden“ and the image settles it.
kalde Troens almindelige Forudsætninger til Hjælp. Men dette var lige mod
Akkorden. Opgaven lød jo paa at formidle og retfærdiggjøre Alt i Kraft af den
rene Viden; hvorledes kunde den absolute Videnskab da være bekjendt at betjene
sig af en Nødhjælp? Lidt af Troen og lidt af den rene Viden giver en uklar
% --- p. 92 ---
Sammenblanding, der kun tager sig maadeligt ud, saavel i Livet som i
Videnskaben. Under slige Omstændigheder var det derfor ganske naturligt, at den
formummede Speculation tog Masken af, og viste Verden, hvad det var for en
Opgave, den kunde og vilde løse. Denne aabne Erklæring gjorde en forbausende
Virkning. Mirakelproblemet løstes; men Løsningen faldt lidt anderledes ud, end
de Godtroende havde ventet. Formidlingens aabenbarede Hemmelighed blev nemlig
denne, at Evangeliets Christus kun i meget uegentlig Forstand kan kaldes den
sande Messias. „Alle de Fuldkommenheder, hvormed Troen har udstafferet dette
enkelte Individ, ere i aabenbar Modsigelse med det indskrænkede Begreb om en
historisk Personlighed, og maae, i Ideens Navn, fordeles paa hele Slægten. Det
er Menneskeslægten selv, der er det sande, syndefrie Gudmenneske; det er
Menneskeslægtens egen Aand, der i Aarhundredernes stadige Fremskridt selv gjør
de store Undergjerninger; det er Menneskeslægten, der er sin egen Forløser.“

Forladt af enhver anden Hjælp, maa Evangelietroen altsaa nu see, at tage sig
sammen, om den maaskee ved egne Kræfter, eller rettere ved den overnaturlige
Bistand, hvorpaa den i Afgjørelsens Stund altid gjør Regning, kan forsvare sig
selv og sit Evangelium. Dette Forsvar lader sig kun gjennemføre ved at
accentuere Gjennembrudet, indskyde Underet som Mellemled imellem Extremerne, og
aflede Formidlingen af de overnaturlige Kjendsgjerninger. --- For imidlertid
nærmere at prøve, hvilke Kræfter der i denne Henseende staae til Troens
Raadighed, ville vi for et Øieblik betjene os af Samtalens Form, og lade den
% The dialogue letters are ANTIQUA against the surrounding Fraktur --- upright
% roman capitals, and heavier than the body (600 dpi, PDF pp. 109--111). Rendered
% \textit{} per this edition's antiqua convention (as the antiqua „A.“ of the
% Andet Afsnit head in the Indhold), not \textbf{}, which is reserved for the bold
% FRAKTUR of the scripture quotations on pp. 63--64, 68 and 79.
Troende \textit{(A)} møde adskillige Indvendinger, der kunne fremsættes af en
Ikke-Troende \textit{(B)}. Altsaa: \textit{B.} Naar vi fra et
naturvidenskabeligt Standpunkt betragte Stjernehimlen, og besinde os paa, at vor
Klode kun er som et Punkt, et Fnug i Universet, maa det da
% --- p. 93 ---
ikke forekomme os høist usandsynligt, at den uendelige Fornuftgud, der ellers
aldrig afbryder sine Love, just skulde have gjort denne Jord til Skuepladsen for
saa specielle og mirakuløse Aabenbarelser, som de, der beskrives i det Gl. og
det Nye Testamente? \textit{A.} Jo tilvisse! \textit{B.} Naar vi fremdeles fra
Kritikens almeenvidenskabelige Synspunkt anstille en Sammenligning imellem
Bibelens Mirakelfortællinger og de hedenske Myther, opdage vi da ikke saamange
træffende Analogier, at vi, for at fælde en rolig og upartisk Dom, maae finde
det i høieste Grad usandsynligt, at de bibelske Myther skulde have større Krav
paa historisk Troværdighed end de hedenske? \textit{A.} Ogsaa dette faaer jeg
vel at indrømme. \textit{B.} Nu vel da! Vi gaae endnu et Skridt videre. Lad os
f. Ex. antage, at det virkelig er skeet Maria, som der siges i Bebudelsen; jeg
% „Luk.“ here; the book prints „Luc.“ elsewhere (e.g. p. 53). Transcribed as
% printed at each site.
skal da vise dig, hvad der ligger i denne Forestilling. Der fortælles (Luk.\ 1,
39), at „Maria stod op i de Dage, og gik hasteligen paa Bjergene til \emph{Judæ
Stad}“. Aner du ikke, hvilken Rædsel der skjuler sig i disse faa Ord?
\textit{A.} Hvad mener du? \textit{B.} Giv nu Agt! Maaskee kan det lykkes mig,
at aabne dine Øine. --- Umiddelbart foran det anførte Sted (nemlig V.\ 38)
hedder det jo udtrykkelig, at „Engelen skiltes fra hende“. Maria vandrer altsaa
ganske alene paa Veien over Bjergene. Sæt dig nu levende ind i Begivenheden.
Forestil dig, at du selv tilligemed en Ledsager mødte hende paa Veien. Naar hun
er kommen forbi Eder, standser Ledsageren, peger paa hende, og siger til dig med
sagte, men høitidelig Stemme: „Betragt engang denne Qvinde; paa hendes Liv
beroer al Verdens Frelse!“ Vilde du da ikke blive maalløs af Forbauselse, ikke
over Qvinden, der iler opad Bjergstien (thi hun er som saa mangen anden Qvinde),
men over din Ledsagers Afsindighed? Jo nærmere Billedet sees, desto tydeligere
opdager man det Urimelige. Maria vandrer alene paa Bjergene! Sæt,
% --- p. 94 ---
% The book punctuates the two parallel clauses differently: „sæt, at et glubende
% Dyr“ with a comma, „sæt at hendes Fod snublede“ without. Both OCR witnesses
% agree. Transcribed as printed, not normalised.
at en Røver overfaldt og myrdede hende, saa vilde denne Røver jo med det Samme
myrde hele Verdens Frelse; sæt, at et glubende Dyr sønderrev hende, saa vilde
det glubende Dyr jo gjøre Verdens Frelse til Intet; sæt at hendes Fod snublede,
og hun styrtede ned ad Fjeldet i Afgrunden, saa vilde dette Fald jo gjøre
% PRINTER'S ERROR, transcribed as printed: „Gjøalespil“ for „Gjøglespil“. At
% 600 dpi (PDF p. 111) the glyph between ø and l is an unmistakable a --- round,
% flat-topped, with no descender, where Fraktur g descends well below the line.
% Both OCR witnesses read „Gjoalespil“. Rejected reading: „Gjøglespil“.
Verdens Opreisning umulig. Det er ikke min Agt at forvirre din Indbildning med
et eventyrligt Gjøalespil; jeg taler i Alvor den rene Sandhed. Er Maria et
virkeligt Menneske, og den Verden, hvori hun bevæger sig, en virkelig Verden: da
er hun ogsaa, selv i sin benaadede Tilstand, udsat for tilfældige Ulykker; thi
Tilfældigheder og Ulykker ere nu een Gang uadskillelige fra den virkelige
Tilværelse. Har hun altsaa, da hun gik paa Bjergene, baaret Ham, paa hvem hele
Verdens Frelse beroer, saa har der jo ogsaa været en Tid, da Gud gav hele
Verdens Frelse til Priis for Tilfældigheden. Men er en saadan Forestilling nu
ikke at ansee for en aldeles fornuftstridig og urimelig Forestilling?
\textit{A.} Jo i Sandhed, en høist fornuftstridig og urimelig Forestilling!
\textit{B.} Saa har du da hermed fældet Fordømmelsesdommen over den gamle
Evangelietro, og jeg holder dig fangen i dine egne Ord. \textit{A.} Mener du
virkelig? \textit{B.} Den Sag er soleklar. Har du ikke selv indrømmet mig, at
Aabenbaringsmiraklerne maatte saavel for Naturlærens som Mythekritikens
almeenvidenskabelige Domstol bortfalde som --- høist usandsynlige, urimelige og
fornuftstridige? \textit{A.} Jo! \textit{B.} Bliver du da nu ikke ogsaa nødsaget
til at tilstaae, at den Tro, hvis Indhold ifølge de foregaaende Indrømmelser, er
i allerhøieste Grad usandsynligt, urimeligt og fornuftstridigt, derved
nødvendigviis stempler sig selv som en i allerhøieste Grad urimelig,
fornuftstridig og forkastelig Tro? \textit{A.} Nei! Det har jeg ikke isinde at
indrømme dig. \textit{B.} Det er vel en Caprice? \textit{A.} Det er min ramme
% p. 94 ends mid-dialogue with a full stop; p. 95 opens flush with the next
% speech, „A. Evangelietroen hviler i Underet.“, in the same paragraph. No word is
% broken across the leaf. The final sort is thick and sits a little high, but it is
% a period, not a comma (900 dpi, PDF p. 111); the ABBYY layer reads a period too.
Alvor. \textit{B.} Maae vi saa bede om en Forklaring.
% --- p. 95 ---
% p. 95 opens FLUSH at the body margin, in the middle of the dialogue paragraph
% that runs from p. 92; verified on the image (PDF p. 112). No word is broken
% across the leaf. The antiqua speaker letters „A.“ / „B.“ continue to be set
% against the Fraktur and are transcribed \textit{} as in the previous batch;
% they are in fact BOLD antiqua, as is „Ergo:“ on p. 96 and (in a lighter
% weight) „Laterna magiea“ on p. 104. All three are rendered \textit{}.
\textit{A.} Evangelietroen hviler i Underet. \textit{B.} En maadelig Forklaring,
i samme Øieblik vi ere blevne enige om at indrømme Undergjerningens Umulighed.
\textit{A.} Undergjerningens Umulighed? Gud bevare mig for en saadan
Indrømmelse! Jeg har kun, for at vise dig, at jeg kunde følge din videnskabelige
Tankegang, tilstaaet, at Aabenbaringens Undere maae forekomme den almindelige
Verdensbevidsthed høist usandsynlige, urimelige og anstødelige. Mere har jeg
slet ikke tilstaaet. \textit{B.} Men hvortil nu dette Ordkløveri? Siden du saa
let kan „følge min Tankegang“, vil du vel ogsaa kunne indsee, at Underet, som
historisk Kjendsgjerning, maa bedømmes fra et historisk-kritisk Standpunkt. Men
nu veed Enhver, at Grændselinien imellem det Allerusandsynligste og det Usande,
imellem det Allerurimeligste og det Umulige er saa fiin, at det historiske Blik
ikke kan opdage den, og derfor med Rette heller ikke ændser den. Enten jeg
altsaa siger om den historiske Undergjerning, at den er i allerhøieste Grad
usandsynlig og urimelig, eller jeg ligefrem paastaaer dens Umulighed, løber da
vel omtrent ud paa Eet og det Selvsamme? \textit{A.} Hermed udtaler du en
Grundvildfarelse, der, lige fra den nyere Bibelkritiks Begyndelse i forrige
Aarhundrede og indtil denne Dag, har afstedkommet de allersørgeligste
Forvirringer saavel i Troen som i Videnskaben. \textit{B.} En dristig Paastand,
der nok kunde trænge til et Beviis. \textit{A.} Lad os overveie Sagen i al
Besindighed. Det ligger i Underets Begreb, at det i Sammenligning med de
sædvanlige Naturphænomener maa fremtræde som det Allerusandsynligste, man kan
forestille sig; tabte det sig i Sandsynligheden, lod det sig jo forklare paa en
ganske naturlig Maade, og var følgelig intet virkeligt Under. Men vil dette nu
med andre Ord sige, at Underet i sig selv er en Urimelighed, og desaarsag at
sætte i Klasse med den Art meningsløse Modsigelser, man i Logiken pleier at
sam\-%
% --- p. 96 ---
menligne med en rund Fiirkant? \textit{B.} Modsigelsen er rigtignok ikke slet
saa paafaldende, thi Aabenbaringslæren kalder jo det Overnaturlige til Hjælp;
men desuagtet er Miraklet dog en virkelig Umulighed, eftersom det jo afbryder
Naturloven, og denne kan nu een Gang ikke afbrydes. \textit{A.} Du vil dog vel
ikke udgive denne Paastand for et videnskabeligt Beviis? \textit{B.} Den er
imidlertid at ansee som den uomstødelige Grundsætning, paa hvilken
Culturbevidstheden og den hele moderne Videnskabelighed hviler. \textit{A.}
Hvorfor er Underet altsaa en Umulighed? \textit{B.} Fordi Naturloven ikke kan
afbrydes. \textit{A.} Hvorfor kan Naturloven ikke afbrydes? \textit{B.} Fordi
det vilde være unaturligt. \textit{A.} Nu glemmer du jo, hvad du selv nylig
sagde. \textit{B.} Nei, ikke just fordi det er ligefrem unaturligt; men fordi
... fordi det Overnaturlige ikke kan gjennembryde det Naturlige. \textit{A.}
Altsaa fordi Underet er umuligt? \textit{B.} Ja! \textit{A.} Summa Summarum:
Naturloven kan ikke afbrydes, og hvorfor ikke? Fordi Underet er umuligt. Underet
er umuligt, hvorfor? Fordi
% PRINTER'S DEFECT: „Naturloveu“ for „Naturloven“ --- a turned n. Verified at
% 900 dpi (PDF p. 113) against the four correct „Naturloven“ on the same page:
% the final sort here has the u's upward exit at the top right and lacks the n's
% bottom-right foot. Both OCR witnesses read u. Transcribed as printed. This is
% the first of a cluster of turned/wrong sorts in this gathering: see also
% „udeu“ (pp. 97, 107), „Paastaud“ (p. 97), „uok“ (p. 98), „Universnm“ (p. 99),
% „kuu“ (p. 104), „hgr“ and „ajør“ (p. 106).
Naturloveu ikke kan afbrydes. \textit{Ergo:} Underet er umuligt, fordi det er
umuligt. Finder du nu ikke, at det er en særdeles tilforladelig Grundsætning,
hvorpaa Culturbevidstheden og den hele moderne Videnskabelighed hviler?
\textit{B.} Ligemeget! Jeg troer paa dens Ufeilbarhed, og Ingen ... \textit{A.}
Tys! Lad os ikke begynde paa Troesbekjendelsen; jeg kunde jo ellers let gribe
den samme Udflugt, og paastaae, at Underet er muligt, fordi det er muligt.
\textit{B.} Men man kan jo dog saa levende fornemme paa sig selv og paa hele
Verden, at Underet maa være en absolut Umulighed! \textit{A.} Husk nu bare paa
vort „almeenvidenskabelige Standpunkt“, og hav saa den Godhed at svare mig
ganske bestemt paa følgende Spørgsmaale: Kan man bevise Underets Umulighed,
eller kan man ikke? \textit{B.} Man kan ikke. \textit{A.} Kan man af Underets
saakaldte Usandsynlighed bevise dets
% --- p. 97 ---
Uvirkelighed? \textit{B.} Vi bleve jo desværre enige om, at Underet maatte være
usandsynligt. \textit{A.} Maa dette nu nødvendigviis indrømmes, saa falder det
vel heller ikke saa vanskeligt at bevise, at den moderne Bibelkritik, der ved
sine „almeenvidenskabelige“ Sandsynlighedsberegninger deels angriber, deels
forsvarer Evangeliet, i sit inderste Væsen er meningsløs. \textit{B.}
Hvorledes? \textit{A.} Den negative Kritik, der gaaer ud paa at tilintetgjøre
den hellige Histories Troværdighed ved at paavise
% Rettelser: printed „Undergjerningens“; corrected to „Undergjerningernes“
% (Side 97, Linie 8 fra oven). Printed line 8 verified on the image at 900 dpi
% (PDF p. 114) and reads „hellige Histories Troværdighed ved at paavise
% Undergjerningens“, exactly as the errata leaf reports. A contemporary reader
% has struck out the final „ens“ in pencil in the margin of this copy.
Undergjerningernes allerhøieste Usandsynlighed, kommer idelig,
% PRINTER'S DEFECT: „udeu“ for „uden“ --- the same turned n as „Naturloveu“
% (p. 96); verified at 900 dpi, the final sort matching the u of „udeu“'s own
% first letter. Both OCR witnesses read u. Transcribed as printed.
udeu selv at ane det, til at fremhæve et af deres umiskjendeligste Kriterier; og
den apologetiske Middelveiskritik, der anstrenger sig for at godtgjøre
Sandsynligheden, kommer idelig tvertimod sin Hensigt til at berøve dem deres
alleruundværligste Kjendemærke. Er dette ikke en Bagvendthed, der i enhver anden
Videnskab forgjæves søger sin Mage? \textit{B.} Maaskee! Dog --- derom ville vi
ikke stride.

Evangelietroen grunder i Aabenbaringens Under. Den har Intet at befrygte af
Philosophien, efterdi denne ligesaa lidt er istand til at bevise Underets
Umulighed, som dets Nødvendighed. Den hellige Historie har heller Intet at
befrygte af den almindelige Sandsynlighedskritik, hvad enten saa denne iøvrigt
bereder sig til Forsvar eller til Angreb. Spørger Naturvidenskaben, hvad der har
bevæget Gud til at aabenbare sig paa denne Klode, da viser den Troende trøstig
hen til --- Kjærlighedens Under. Rykker Kritiken frem med den suffisante
% PRINTER'S DEFECT: „Paastaud“ for „Paastand“ --- turned n again; verified at
% 900 dpi against the n of „man“ two words later on the same line. Both OCR
% witnesses read u. Transcribed as printed.
Paastaud, at man ikke maa bedømme det Gl. og Nye Test.s Mirakelfortællinger
efter anden Regel end de hedenske Myther: da holder den Troende fast ved ---
Kjærlighedens Under. Vil verdslig Kløgt gjøre sig vigtig med de Tilfældigheder
og Hændelser, der i den virkelige Verden maa have truet Frelserens Liv: da
minder den Troende blot om --- Kjærlighedens Under. En saadan eensformig
Gjentagelse
% --- p. 98 ---
af een og samme „ubeviislige“ Forudsætning (og det i en Verden, der tør sætte
sin Salighed i Pant paa, at Gud ikke er istand til at gjøre et eneste Under)
lyder jo rigtignok saa mat og klangløs, at mangen Evangelieforsvarer vel vilde
undsee sig ved den, og
% „strax“, not „strar“: verified at 900 dpi (PDF p. 115) by the test used on
% pp. 84 and 86 --- the final sort drops a diagonal below the baseline, unlike
% the r two letters earlier in the same word, and is identical to the true x of
% „strax“ on p. 6. Tesseract reads r here; the image does not.
strax ile med at opsøge sig nogle videnskabelige Beviser. Men just af den Grund
kan Methoden være ret brugbar til praktisk Øvelse; med den videnskabelige
Formidling har det jo slet ingen Hast. Undseer den Troende sig derimod ved at
føre en saa eenfoldig Tale ligeover for de stærke Tænkere og skarpsindige
Kritikere, der saa vel forstaae at mestre den hellige Skrift; er han bange for
at rykke aabenlyst frem med sin „uvidenskabelige Forudsætning“: da kan han
ligesaa gjerne overgive sig først som sidst; thi Modstanderne have allerede
Seieren ihænde. Er Underets Mulighed derimod først tilfulde betrygget mod
ethvert almeenvidenskabeligt Angreb; er dets historiske Virkelighedsform
tilstrækkelig sikkret imod Sandsynlighedskritikens Insinuationer; vaagner
Evangelietroen endelig igjen af sin lange Bedøvelse, og lærer at indsee, at den
begaaer et Selvmord, hvis den ikke aabenlyst og uforbeholdent bekjender sig for
al Verden at være en Under-Tro, en Tro paa den forløsende Kjærligheds ufattelige
Under: da kommer den kritiske Retfærdiggjørelse som af sig selv, da har det slet
ingen Nød med den videnskabelige Formidling.

Og nu tilbage til den hellige Familie! --- En Opvakt, der har taget Under-Troen
forfængelig og beruset sig i Syner og Aabenbarelser, vilde maaskee her forvente
en aldeles overmenneskelig Forherligelse af de Udvalgte, der samle sig om
Frelserens Vugge. At Evangeliets eenfoldige Simpelhed ogsaa i denne Henseende
danner en priselig Modsætning til Legendernes svulstige Overdrivelser, er
imidlertid ofte
% PRINTER'S DEFECT: „uok“ for „nok“ --- turned n; verified at 900 dpi against
% the n of „nogle“ and „nu“ on the same page. Both OCR witnesses read u.
% Transcribed as printed.
uok blevet paaviist og anerkjendt af besindige Kritikere. Men, forat denne Dom
kan være mere end en religiøs
% --- p. 99 ---
Smagsdom, maa det Overnaturlige paa den ene Side skarpt fremhæves, og dets
virkelige, skjøndt sparsommere, Tilsynekommen i den kanoniske Fremstilling paa
den anden Side vise sig begrundet i Sagens Natur.

Den hellige Familie lever som paa Grændsen mellem tvende Verdner: Evigheden
kaster et forklarende Skjær hen over dens enkelte Skikkelser; og dog tilhører
den det Nærværende saa ganske, at end ikke et eneste af de Baand, hvorved den er
forenet med det borgerlige Samfund, kan siges at briste. Himmelens Kræfter røre
sig i dens Midte; og dog foregaaer alt det Vidunderlige saa stille, saa skjult,
at den travle Omverden egentlig slet Intet mærker. Ja, skjult og hemmeligt skeer
det Kjærlighedens Under, hvorved det Hellige fødes i en troende Sjæl; skjult og
hemmeligt var ogsaa det Kjærlighedens Under, hvorved det hellige Ord blev Kjød i
Slægten. --- Betragte vi nu
% „Vexelvirkningen“, not „Verelvirkningen“: verified at 900 dpi (PDF p. 116).
% The sort after „Ve“ drops a diagonal below the baseline; the r of „virkningen“
% in the same word does not. Both OCR witnesses read r. The same holds of the
% recurrence on p. 101.
Vexelvirkningen imellem de tvende Verdner lidt nærmere, da tager Forholdet sig
ved første Øiekast ud, som om den overnaturlige i Grunden var den svagere, fordi
dens Skikkelser ere saa „luftige“, og dens Virkninger saa skjulte; medens den
naturlige derimod synes at have Virkelighedens afgjørende Overvægt paa sin Side.
Men, er Underet ikke et Blændværk, saa forholder Sagen sig jo omvendt, saa har
det overnaturlige Rige jo den væsentlige Overmagt, saa er det jo denne
nærværende Verden, der maatte opløse sig og svinde som et Luftsyn, hvis hine
Evighedsmagter brøde frem med Vælde. Hvad er det da, der holder de Vældige
tilbage, at de, siden de dog i deres Usynlighed ere os saa nær, ikke ryste, ikke
rokke dette naturlige
% PRINTER'S DEFECT: „Universnm“ for „Universum“ --- the converse of the turned
% n, an n standing where the u belongs. Verified at 900 dpi; both OCR witnesses
% read n. Transcribed as printed.
Universnm, ikke forrykke denne Tingenes Orden? Den Troende har kun eet Svar:
Guds uendelige Kjærlighed! Det er denne Kjærlighed, der giver Underet sit sande
evangeliske Grundpræg. Ja at de himmelske Syner kunne vise sig for et
menneske\-%
% --- p. 100 ---
ligt Øie, uden at det naturlige Syn med det Samme for evigt slukkes; at
himmelske Røster kunne lyde i et menneskeligt Øre, uden at dette naturlige Øre
med det Samme for evigt lukkes; at et Gjennembrud af Skaberaandens
overnaturlige Kraft kan foregaae i en menneskelig Sjæl uden at forstyrre dens
Natur og forvirre alle Følelser og Tanker: er dette ikke et Kjærlighedens Under?

% SUBSECTION HEAD, part-way down p. 100, about two fifths of the way from the
% top. As with the chapter heads, the numeral „1.“ stands centred on a line of
% its own above the title; the title runs over two centred lines and ends with a
% full stop, and a scripture reference in a type smaller than the body is
% centred beneath it. NO RULE accompanies this head: the page was scanned at
% 600 dpi and every row of the head region tested for a long horizontal run of
% ink; the only one found (114 px at y 1538--1541) is the em dash inside the
% title itself. Compare the head of chapter II (p. 42), which has rules of
% 0.38 and 0.12, and of chapter III (p. 81), which has one of 0.08.
%
% INDHOLD vs. BODY: the Indhold prints „1. Marias Besøg hos Elisabeth ---
% Johannes Døberens Fødsel ...... 100--128“, with no point after „Elisabeth“.
% The printed head has one: „Marias Besøg hos Elisabeth. --- Johannes Døberens
% Fødsel.“ Both readings stand; the head follows the page, the TOC entry follows
% the Indhold, and neither is normalised toward the other.
%
% The book prints „Luk.“ here (it printed „Luk.“ at p. 42 and „Luc.“ at p. 53);
% transcribed as printed.
\subsection*{1.\\[4pt] Marias Besøg hos Elisabeth. --- Johannes Døberens Fødsel.}
\addcontentsline{toc}{subsection}{1. Marias Besøg hos Elisabeth --- Johannes Døberens Fødsel}

\begin{center}
  {\footnotesize (Luk.\ 1, 39--80.)}
\end{center}

Aabenbaringen er intet Skyggespil, der spotter Tilværelsens Love; den er ingen
Leg, hvori de Himmelske letsindig blande sig med de Jordiske: i alle Yttringer
den samme Strenghed, fra alle Sider det samme Grundpræg af guddommeligt Alvor!
% „følelige“: the e between „føl“ and „lige“ is very weakly inked and all but
% invisible in the scan; both OCR witnesses drop it and read „sol“/„sol-“. At
% 900 dpi under autocontrast the ghost of the sort is there, and the initial is
% f, not long s. Read „følelige Slag“.
Det er følelige Slag, hvormed det Overnaturlige rammer de Udvalgte.
% Letterspaced (Sperrsatz) from „Zacharias“ to „Joseph.“, confirmed on the image
% and by spacing.py, which returns four RUNs across these four lines. The
% surrounding sentences are ordinary body Fraktur.
\emph{Zacharias er bleven maalløs, Engelen har bundet hans Tunge; Elisabeth
holder sig skjult i sit eget Huus, for at undvige Menneskenes Blikke; selv Maria
er i sit Hjem saa taus og stille, at hun end ikke kan sige et Ord til Joseph.}
Der er i denne Forstummen, i denne Skjulthed og Indesluttethed Noget, der
uvilkaarligt ængster os, og, naar vi dvæle derved, næsten bringer Blodet til at
isne. Det er, som om det Menneskelige med Eet var blevet overvældet af det
Guddommelige, som vare de Enkelte saaledes bortrykkede fra denne Tingenes Orden
og fra hinanden indbyrdes, at de ikke mere kunde finde Sammenhæng og Mening i
det Nærværende. Men denne Fortabthed i det betyngende Dunkle viser paa den anden
Side hen til en naturlig Grundvirken, hvorunder Tilegnelsen foregaaer. Den
bundne Tilstand gaaer over i en høiere Befrielse: et saligt Liv
gjennemstrøm\-%
% --- p. 101 ---
mer Sjælene med Lys og Glæde. Det er denne Befrielsens og Forstaaelsens Jubel,
der toner i Marias Hymne, i Zacharias's prophetiske Lovsang. Begge
Lovpriisninger ere tillige saa eiendommelig bestemte ved de virkelige
Omgivelser, og svare, hver især, saa aldeles til deres Anledning, at
Vexelvirkningen imellem det Overnaturlige og det Naturlige derved end mere
klarer sig for et troende Øie. --- Hvor naturligt, hvor reent menneskeligt lader
det sig ikke forklare, at Maria, under de givne Omstændigheder, besøger sin
Frænke Elisabeth, og ved sin Indtræden i hendes Huus modtages med aabne Arme,
med kjærlig Hilsen; hvor naturligt og reent menneskeligt gaaer det ikke til, at
de to Veninder, der begge medbringe et Haab, en glæderig Velsignelse, hvorpaa
den hele Betydning af deres Liv beroer, nu mødes, den Unge hos den Gamle, udøse
deres Hjerter for hinanden, glæde sig, og takke Gud. Men i deres Hilsen og
gjensidige Meddelelse forekomme strax saadanne Udbrud og Yttringer, der slet
ikke tillade nogen naturlig Forklaring.
% PRINTER'S DEFECT: no full stop closes this quotation. The line reads
% „... fyldt med den Hellig Aand“ and is followed by ordinary word space and
% „Skal denne Indledning“; verified at 900 dpi (PDF p. 118). Both OCR witnesses
% show the same. The quotation marks themselves are balanced. Transcribed as
% printed.
„Og det begav sig, der Elisabeth hørte Marias Hilsen, sprang Fosteret i hendes
Liv, og Elisabeth blev fyldt med den Hellig Aand“ Skal denne Indledning
forstaaes som en senere Opdigtelse, fordi den berører „det Usandsynlige“; da ere
de følgende Glædesyttringer ligeledes enten sværmeriske Udbrud af Qvinderne
selv, eller sindrige Opfindelser af en digtende Erindring. Opdage vi derimod,
hvorledes det Overnaturlige endnu vedbliver at røre sig som i en dæmrende
Baggrund; betragte vi dets øieblikkelige Indvirkning paa Elisabeth som en
betydningsfuld Eftervirkning af de stærkere Frembrud, der skete i begge
Bebudelser: da har Elisabeths Hilsen til Maria: „Velsignet er du iblandt
Qvinderne, og velsignet er dit Livs Frugt“, sin gode, tilstrækkelige Grund; den
frydefulde Overraskelse, hvori hun tilkjendegiver, at den Besøgende staaer i
Værdighed høit over den, der
% --- p. 102 ---
modtager Besøget: „Hveden kommer mig det, at min Herres Moder kommer til mig?“
viser da hen til en umiddelbar Indflydelse af den Hellig Aand. Til Elisabeths
Stemning svarer igjen Marias Henrykkelse; de benaadede Qvinder forstaae
hinanden: deres Sjæle smelte sammen i den ene, salige Grundfølelse, at Frelsens
Tid nu endelig er kommen.

Paa lignende Maade forholder det sig med den følgende Begivenhed. Elisabeths Tid
fuldkommedes, og hun fødte sin Søn. „Og Naboer og Slægtninge hørte, at Herren
havde gjort sin Barmhjertighed stor imod hende; og de glædede sig med hende. Og
det begav sig paa den ottende Dag, da kom de at omskjære Barnet, og de kaldte
det efter hans Faders Navn Zacharias.“ En mere naturtro og gemytlig Fremstilling
af et jødisk Familie-Optrin, end denne, vil man under de foreliggende
Omstændigheder dog vel ikke forlange. Men denne de Paarørendes velvillige og
naturlige Deeltagelse giver just Anledning til, at det Overnaturlige desto
bestemtere kommer tilsyne. De gode Slægtninge studse, da de erfare, at Elisabeth
ikke vil lade Barnet kaldes med Faderens Navn. „Der er jo slet Ingen i Familien,
der hedder Johannes“. Valget af dette Navn forekommer dem derfor saa
besynderligt, at de nikke til den Stumme, for at erfare hans Mening; og, da
Faderen nu til yderligere Bekræftelse skriver paa Tavlen:
% Letterspacing of a single short word: „e r“ inside the quotation, set with the
% wide letterspace of Sperrsatz against the ordinary setting of the words on
% either side. Verified at 900 dpi (PDF p. 119). spacing.py does not flag a
% two-letter word, so this one is found only on the image.
„Hans Navn \emph{er} Johannes“, forundre de sig saare. Neppe ere de Deeltagende
ved deres Forundring og de Anelser, denne maa have fremkaldt, forberedte paa det
Overordentlige, før der ogsaa skeer et nyt Gjennembrud. Zacharias befries med
Eet fra sin Stumhed: hans Mund oplades, hans Tunges Baand løses, han taler og
priser Gud. Imidlertid gaaer dog det Overnaturlige saaledes op i det Naturlige,
at Underet selv derved næsten forvandles til en psychologisk Kjendsgjerning. Ved
Underets Ind\-%
% --- p. 103 ---
træden er Stemningen hævet til det Punkt, at Lovsangen kan begynde. Dens Indhold
lader sig opfatte som et høitideligt Svar paa de Deeltagendes forventningsfulde
Yttring: „Hvad mon der skal blive af dette Barn?“ Thi ogsaa dette hører jo med
til Aabenbaringens Bestemmelse, at ikke blot det Vidunderlige skeer, men at
Forstaaelsen af det Skete med det Samme indledes og forberedes. „Og der kom en
Frygt over Alle, som boede omkring dem, og alle disse Ting rygtedes paa alle
Judæas Bjerge. Og Alle, som det hørte, lagde det paa deres Hjerte.“

Førend vi imidlertid gaae videre i den religiøse Udvikling af Fortællingens
Indhold, ville vi her standse nogle Øieblikke, for at høre Kritikens
Indvendinger.

% Letterspaced paragraph opening, confirmed on the image and by spacing.py
% (RUN). The same device introduces „Besvarelse.“ (pp. 104, 107) and „Anden
% Indvending.“ (p. 105).
\emph{Første Indvending.} „Om man end til en vis Grad vil indrømme Underets
religiøse Betydning, saa følger dog endnu ingenlunde heraf, at Evangelisterne
have gjengivet de enkelte vidunderlige Tildragelser med alle Biomstændigheder
aldeles overeensstemmende med den virkelige Sandhed. En flygtig Prøvelse vil
strax overbevise os om det Modsatte. Forat Joseph ikke skal forlade Maria, maa
Engelen jo aabenbare ham Hemmeligheden. Men af hvis Mund kan nu Joseph, før
Drømmen indtraf, have faaet hiin ominøse Underretning om Marias Svangerskab?
% „forvexler“: the sort after „forve“ drops a diagonal below the baseline,
% unlike the r of „for“ at the head of the same word; verified at 900 dpi
% (PDF p. 120). Tesseract reads r, the ABBYY layer reads r; the image does not.
Matthæus lader som Ingenting, og forraader derved, at han forvexler sin egen
Forestilling om Joseph med Josephs virkelige Stilling. Endvidere: Elisabeths
Hilsen forudsætter jo, at hun allerede forud veed, hvorledes det staaer til med
den indtrædende Maria. Men hvorfra har hun denne Kundskab? Det springende Foster
har dog vel ikke kunnet fortælle hende Noget; en særlig Aabenbaring omtales
ikke, og ligesaalidt er der Spor af, at Efterretningen bliver hende meddeelt ad
den naturlige Vei. Lukas ændser slet ikke dette Punkt, og viser derved, at han
% „forvexler“ broken across the leaf, „for=“ at the foot of p. 103 and „vexler“
% at the head of p. 104; the x verified at 900 dpi on PDF p. 121 by the same
% test. Joined with \-% , not a bare hyphen.
for\-%
% --- p. 104 ---
vexler sin egen Forestilling med den virkelige Begivenhed. Dog, den største
Vanskelighed staaer endnu tilbage. Af hvilke Kilder have Evangelisterne øst
disse Beretninger? Havde Zacharias maaskee optegnet, hvad der tildrog sig med
ham og hans Hustru? Havde Maria maaskee efterladt Noget skriftligt, eller havde
hun maaskee senere hen fortalt, hvad hun vidste, til En og Anden iblandt
Apostlene? Til hvem skulde Maria i dette Tilfælde vel hellere have betroet Sligt
end til Johannes? Men hvorledes forklares det da, at Johannes i sit Evangelium
slet ikke ændser disse Fødselsmirakler? Disse Vanskeligheder kunne endnu forøges
med mange lignende. Enhver, der ikke vil forhærde sig imod Sandheden, og
forsætlig tillukke Øinene for Kritikens allertydeligste Vink, vil idetmindste
see sig nødsaget til at indrømme dette, at den hele Cyclus af Familiescener, der
gaae forud for Døberens Optræden, præsenterer sig som Billeder i Erindringens
% ANTIQUA against the Fraktur, and a PRINTER'S DEFECT within it: the book prints
% „magiea“ for „magica“. Verified at 900 dpi (PDF p. 121) --- the fourth letter
% is an unmistakable antiqua e, with its horizontal bar, not a c. Both OCR
% witnesses read e. The antiqua here is of a lighter weight than the bold
% antiqua of the speaker letters and of „Ergo:“ (p. 96); all are rendered
% \textit{}. Transcribed as printed.
\textit{Laterna magiea}: flygtige, ubestemte Familiesagn, der, til Forherligelse
af Jesu Fødsel og Barndom, vel i en sildigere Tid ere blevne omarbeidede, men
dog ikke tilstrækkelig sammenarbeidede!“

\emph{Besvarelse.} Det kan ikke have været Evangelisternes Hensigt at opregne
Alt, hvad der har tildraget sig i den hellige Familie;
% PRINTER'S DEFECT: „kuu“ for „kun“ --- turned n; verified at 900 dpi against
% the n of „Optrin“ on the same line. Both OCR witnesses read u. Transcribed as
% printed.
kuu de mærkeligste Optrin og disses meest charakteristiske Hovedtræk ere
fremhævede. Hvad de manglende Mellemled angaaer, da falder det den Troende, der
tør tage de overnaturlige Forkyndelser med i Beregningen, ligesaa let at digte
sig en Sammenhæng efter sin Synsmaade, som det falder Troesfornegteren, hvem det
Overnaturlige overalt staaer i Veien, let at digte sig saa mange Urimeligheder,
han selv finder for godt. Slige subjective Hypotheser ere som Sæbebobler: de
glimre og briste. Istedetfor at indlade sig paa Spørgsmaal, der ikke kunne
besvares, bør den Troende helst lade det Skjulte være skjult, det Ubestemte
ubestemt
% --- p. 105 ---
og kun sige ganske ligefrem: „Man veed det ikke!“ Altsaa: Hvo har givet Joseph
den første Efterretning om Marias Svangerskab? „Man veed det ikke!“ I hvilket
Sprog har den Hellig Aand talt til Elisabeth om Marias Udkaarelse? „Man veed det
ikke!“ Af hvis Mund have Evangelisterne hørt de Familiehistorier, de i dette
Afsnit fortælle os? „Man veed det sandelig ikke!“ Dette kan da end ydermere tjene
til Beviis for, at Evangelietroen er og maa være en Frihedens Sag, der umulig
kan fremtvinges ved historiske Beviser. Er jeg stemt for Troesfornegtelsen, da
skal det ikke mangle mig paa Midler til at opløse den hellige Historie i lutter
Modsigelser; er jeg derimod tilfulde greben af Evangeliets Aand, og har lært at
indsee, hvad Underet vil sige, da vil jeg ogsaa i Skriftens Antydninger finde
Midler nok til at paavise den hellige Histories indre Overeensstemmelse med sig
selv, saavel i det Hele som i alle væsentlige Led og Dele. Det er ingenlunde
Kritiken, der gjør Udslaget, naar det gjælder om Tro eller Vantro; men det er
omvendt Forskjellen imellem Tro og Vantro, der bestemmer og leder Kritiken.
Apologetens Anstrengelse for at gjendrive Troesfornegterens historisk-kritiske
Indvendinger, er spildt Umage. Kun naar den negative Kritik giver sig Mine af at
gaae ind paa den troende Opfattelse, og da med alskens Fordreielser søger at
bringe denne i Modsigelse med sig selv, bør Evangelietroen tage til Gjenmæle;
thi det Umulige og det i sig selv Urimelige kan jeg ikke troe.

\emph{Anden Indvending.} „Det er Troens umiskjendelige Opgave at gjøre den
hellige Historie saa virkelig, saa historisk, som muligt. Hvis den vantroe
Kritik ikke vedblev med sine drillende Modsigelser, vilde den Troende sikkerlig
aldrig falde paa, at underkaste de bibelske Fortællinger nogen streng
videnskabelig Prøvelse. Forelsket i de overnaturlige Kjendsgjerninger, faaer han
% --- p. 106 ---
naturligviis saa meget at bestille med de himmelske Syner, de uudtømmelige
Spaadomme og den guddommelige Viisdoms mangehaande Hemmeligheder, at der neppe
bliver Tid til at indlade sig paa Tvivlens møisommelige Undersøgelser. Naar nu
desuagtet den nyere Bibelkritik omsider
% PRINTER'S DEFECT: „hgr“ for „har“. The middle sort is a g with a full
% descender; verified at 900 dpi (PDF p. 123). Both OCR witnesses read g.
% Transcribed as printed.
hgr drevet det dertil, at endog den kirkeligste Theolog ikke tør fragaae, at
navnlig Fortællingerne i dette første Afsnit hvile paa en svævende Overlevering,
og at Begivenhederne følgelig alle tilhobe ere hyllede i Sagnets Slør; da maa
denne Indrømmelse ligefrem betragtes som aftvungen. At Troen nu ikke desto
mindre
% PRINTER'S DEFECT: „ajør“ for „gjør“ --- an a standing where the g belongs;
% the sort has no descender, unlike the g of „gjøre“ elsewhere on these pages.
% Verified at 900 dpi. Both OCR witnesses read a. Transcribed as printed.
ajør en Dyd af Nødvendigheden, og beraaber sig paa den frie Tilegnelse, viser
sig tydeligt nok som et Forsøg paa at skjule Forlegenheden. Det ligger jo i
Aabenbaringens Natur, at Evangelietroen maa staae og falde med de hellige
Kjendsgjerninger: kun Den, der antager Kjendsgjerningen, kan anerkjende Ideen,
kun Den, der fastholder Bogstaven, kan være sikker paa Aanden. Indrømmes det nu,
at de paagjældende Fortællinger allesammen ligge i Sagnets Tvetydighed; saa har
Troen tabt sit egentlige Fodfæste. Nærer man endog den allersvageste Formodning
om, at de enkelte Træk muligviis ikke ere i fuldkommen ordlydende
Overeensstemmelse med de virkelige Begivenheder; saa har man med det Samme
opgivet, hvad der ikke tør opgives, den evangeliske Virkelighedsform, og maa da
ogsaa beqvemme sig til at aflede Kjendsgjerningens Betydning af Ideen. Men ved
denne Vending er den Troende, inden han selv veed deraf, midt inde i
Philosophiens og Kritikens Labyrinther. Saalænge Evangelietroen kan lade være at
tænke paa nogensomhelst Uovereensstemmelse imellem Begivenhedens Virkelighed og
Fremstillingens Form, kan den nok holde Tvivlen ude; men, saasnart den indlader
sig med Sagnet, er Uovereensstemmelsens Mulighed sat. Man gjøre nu iøvrigt denne
Mulighed saa snæver, man vil;
% --- p. 107 ---
den er altid en Sprække, hvorigjennem Tvivlen lister sig ind i den troende
Bevidsthed. Kort og godt: Vil Evangelietroen undgaae Philosophiens Skjær ved at
klamre sig til Bogstaven, saa strander den paa den historiske Kritik; forsøger
den at undvige Kritiken ved at tye ind i Ideen og Aanden, da drukner den i den
philosophiske Skepsis. Den, der nu i denne Snævring mener, at kunne hjælpe paa
Troen ved at anbefale den Middelveien, han maa tilvisse være en snild Mand, saa
snild, at han ikke engang ændser, hvad det er, hvorom her handles. At følge
Middelveien er nemlig det samme, som at forudsætte en blot \emph{omtrentlig}
Overeensstemmelse imellem det Virkelige og Formen; men saa kommer jo
Philosophien strax med sine Ideer, og Kritiken med sine Sagn, og saa --- ja saa
faaer Evangelietroen en brat Opløsning. Evangelietroen maa vel have en Anelse om
denne skjulte Selvmodsigelse; thi vistnok philosopherer den til en vis Grad, og
taaler Kritiken indtil en vis Grad, men bliver dog snart lidenskabelig og
trodsig, naar Videnskaben trænger ind paa den, for at forhjælpe den til høiere
Selvbevidsthed.“

\emph{Besvarelse.} Paa Sagnets Enemærker øver Kritiken sit frieste Spil; som
Drengen ryster sit Kaleidoskop, behøver den blot at ryste sine Hypotheser, og
strax danne sig de allerbesynderligste Figurer. Den ordinære Forstandskritik
udvisker Begivenhedernes overnaturlige Træk, og viser os de enkelte Brudstykker
i Virkelighedens sædvanlige Træsnit
% PRINTER'S DEFECT: „udeu“ for „uden“ again, as on p. 97; verified at 900 dpi
% (PDF p. 124). Both OCR witnesses read u. Transcribed as printed.
udeu Ideens Glands. Den philosophiske Ideekritik forvandler saavel de naturlige
som overnaturlige Bestanddele til poetiske Momenter i Fremstillingens
Beskuelighed, og foreviser en Række af mythiske Luftbilleder, der Intet have med
Virkeligheden at skaffe. Begge Methoder gjøre naturligviis Fordring paa den
allerstrengeste Videnskabelighed; men den „strenge Videnskab“ er desuagtet ikke
istand til at forlige
% --- p. 108 ---
dem indbyrdes. Forstandskritiken, der forsøger sig i den „naturlige Forklaring“,
indseer ret godt, at man ikke uden videre kan opgive en historisk Overlevering,
fordi denne indeholder sagnagtige Træk, og kalder desaarsag den gjennemførte
Ideekritik \emph{phantastisk}. Denne Beskyldning for Uvidenskabelighed
tilbagebetaler nu den philosophiske Kritiker fra sit
% PRINTER'S DEFECT: „Standpnnkt“ for „Standpunkt“ --- n for u, verified at 600 dpi
% (the glyph carries no breve). Same wrong-sort cluster as pp. 96--104.
Standpnnkt med rigeligt Overskud. Han beviser nemlig, hvor ukritisk den naturlige
Fortolkning gaaer tilværks, idet den indbilder sig „at have udskilt det Mythiske
og Poetiske af en Fortolkning, naar den afplukker nogle Qviste og Blomster af
dette Skud, men lader den mythiske Rod ligge
% PRINTER'S DEFECT: „nantastet“ for „uantastet“ --- n for u, verified at 600 dpi.
nantastet ved det Reenthistoriske“. Derfor kalder han da reentud en saadan
Halvkritik \emph{ideeløs og borneret}. Denne Uenighed er ingenlunde tilfældig
eller blot foreløbig; den har sin dybere Grund i Sagen selv. Den gjennemførte
Ideekritik er virkelig phantastisk, efterdi den i Principet undergraver al
historisk Tro: skulde en Begivenhed ansees for uhistorisk, blot fordi den i en
paafaldende Grad anskueliggjør en stor Idee; hvormange virkelige Begivenheder
maatte da ikke udslettes af Historien! Den ordinære Forstandskritik er i Sandhed
borneret, eftersom den ikke engang kan mærke paa sig selv, at den maa slaae Ideen
ihjel, for at redde Kjendsgjerningen. Med det Overnaturliges Forsvinden forsvinder
Aabenbaringen selv. Fradrages altsaa de overnaturlige Bestanddele som formeentlige
Tilsætninger eller Opdigtelser, saa maa det naturlige Skelet falde bort af den
hellige Historie; thi Aabenbaringshistorier, der Intet aabenbare, høre blandt de
aabenbare Modsigelser. Før den strenge Videnskab overtager sig at forhjælpe Troen
til „Selvbevidsthed“, vilde det vist være ret hensigtsmæssigt, hvis den forsøgte
paa at forhjælpe sig selv til lidt klarere Bevidsthed om Betydningen af de hellige
Sagn. --- Denne Bevidsthed har Troen med sand og oprindelig Vished i sig selv. Den
Troende veed, at Sagnets
% --- p. 109 ---
Aand i Evangeliet ikke er en Skinnets og Blændværkets Aand, men Sandhedens hellige
og ufeilbare Aand; hermed er enhver Tvivl om Beretningernes Troværdighed absolut
udelukket. Der kunne vel opkomme mange Formodninger om den Gang, Overleveringen er
gaaet, før den omsider slog sig til Ro i vor foreliggende Fremstilling; men har
Sandhedens Aand fulgt den paa alle Omveie, da kunne vi ogsaa forlade os paa, at
den er bleven sikkret mod enhver Udskeielse. Der kunne vel danne sig mange
Meninger om Aandens Indvirkning paa dens udvalgte Redskaber; men, er Underet ikke
en Indbildning, saa har Sandhedens Aand tilfulde forstaaet at udtrykke det
guddommelige Indhold i den menneskelige Form, saa at enhver Forvanskning af de
aabenbarede Kjendsgjerninger er og forbliver
% PRINTER'S DEFECT: „ntænkelig“ for „utænkelig“ --- n for u, verified at 600 dpi.
ntænkelig. Der kunne vel reise sig mange Stridigheder om det endelige Forhold
imellem de umiddelbare Begivenheder og deres evangeliske Fremstilling, men staaer
det urokkelig fast, at den samme Aand, der virkede i Begivenheden, ogsaa har været
medvirkende i Fremstillingen: saa har jo Læseren den selvsamme Forvisning, som et
Øievidne; thi intet Øievidne vilde været istand til at finde Mening i hine
Begivenheder, hvis han ikke havde seet Alt i Aandens Lys, og hørt Alt med et
aandeligt Øre. Den Troende kan derfor, i denne Henseende, anstille hvilken kritisk
Reflexion, han vil: han kommer dog altid til det Resultat, at Evangeliets Form er
den eneste sande og absolut tilfredsstillende Form, den Form, hvori Ordet har
forklaret Underet, og Aanden gjennemtrængt Bogstaven.

\emph{Tredie Indvending.} „Naar et Menneske først har sat sig det i Hovedet, at
han Intet vil see, høre, føle, tænke eller tale uden i Troens Navn, saa er det
rigtignok en ganske naturlig Selvfølge, at ingen Kritik kan bringe ham i
Forlegenhed. Man kan jo ikke paa een Gang være baade i Troen og udenfor Troen, ikke
% --- p. 110 ---
i eet og samme Aandedrag bekræfte og benegte Troen. Det nytter derfor heller ikke
stort at indlade sig med det Slags Mennesker, hvis hele Skarpsindighed er traadt
saa udelukkende i Troens Tjeneste. Da de altid ligge i Strid med den sunde
Fornuft, ere de stundom endog ret opfindsomme, og vide at vende og dreie deres
Paastande saalænge, indtil disse omsider endog faae et vist Skin af
Videnskabelighed. Den Indvending, der nu skal fremsættes, er derfor heller ikke
beregnet paa at bringe Troen til Fornuft, men har kun til Hensigt at oplyse de
Uhildede om, hvad det er for Trivialiteter, de Troende paadutte Sandhedens Aand.
Ifølge Evangelietroens Forudsætning maa man jo antage, at de Hymner, der lægges
Maria og Zacharias i Munden, ere umiddelbare Udstrømninger af den hellige Aand.
Men under en saadan Forudsætning vil man da ogsaa undskylde Kritikeren, naar han
finder det paafaldende, at en Begeistring, der flød saa umiddelbart af en
guddommelig Kilde, ikke er falden mere original ud, men, for at minde om et Ord af
% „Strauß“ is set in Fraktur here, not antiqua (checked at 600 dpi); no \textit{}.
Strauß, viser sig saa udpyntet med Reminiscenser af det Gl. Tst.“

\emph{Besvarelse.} Altsaa: Hymnerne mangle poetisk Originalitet! Ganske rigtigt.
Maria er ingen Skjønaand, og Zacharias ingen Digter. Er det maaskee
Theaterkritikens Maalestok, der skal lægges an paa de evangeliske Optrin; da maae
% PRINTER'S DEFECT: „strar“ for „strax“ --- as on pp. 56, 86. At 600 dpi the final
% glyph throws no descender below the baseline, so it is r, not the Fraktur x.
vi strar give Kritikeren Ret, og kun forhjælpe ham til at udtale sig med lidt
større Bestemthed. „De tvende Scener i Zacharias's Huus, nemlig: Marias Besøg og
Naboernes Lykønskninger i Anledning af Sønnens Fødsel, ere naturligviis anlagte
paa at motivere et Par religiøs-lyriske Impromtuer. Men i denne Henseende kan den
Hellig Aand neppe siges at have været sin Opgave voren. Elisabeth driver det
% PRINTER'S DEFECT: „knn“ for „kun“ --- n for u, verified at 900 dpi.
knn til nogle temmelig vidt drevne Glædesudbrud, og hele Marias Foredrag gaaer paa
anden Haand. Den
% --- p. 111 ---
i og for sig ret smukke Begyndelse: „Min Sjæl ophøier Herren o. s. v.“, røber
tydeligt nok, at Ideen er laant af Hannas Lovsang (1 Sam.\ 2, 1). Ubetydelige
Afvigelser fra Originalen, der forekomme hist og her, ere Compilationer af
Propheter og Psalmer. Ja endog den ene Strophe, hvori den Benaadede skildrer sin
egen Ophøielse (V. 48): „See! Nu herefter skulle alle Slægter prise mig salig“, er
en reen Efterklang af Elisabeths Udraab (V. 45): „Og salig hun, som troede“. Det
hele Forsøg falder saaledes temmelig fattigt ud: ikke en original Tanke, ikke en
eneste Vending, der røber noget egentligt Talent! Det Samme gjælder om
Zacharias's Parti. Med Undtagelse af nogle meget dunkle og smagløse Hentydninger
% PRINTER'S DEFECT: the parenthesis that closes after „Elisabeth“ below has no
% opening bracket in the printing --- verified at 600 dpi that no „(“ stands
% before „nemlig“. Transcribed as printed; the sense wants „(nemlig ved Ordet …“.
til de Paagjældendes Navne, nemlig ved Ordet: „Barmhjertighed“ (V. 72) alluderes
til Navnet: „Johannes“; ved Udtrykket „tænke paa“ ɔ: ihukomme (V. 72) til ---
„Zacharias“; og ved Ordene: „den Eed, han svor“ (V. 73) til --- „Elisabeth“), er
hele den lange, almindelige Indledning (V. 68--75) holdt i en saa løs Forbindelse
med Situationen, at den ligesaa godt kunde passe sig til enhver anden opbyggelig
Anledning. Da Lovsangen omsider naaer det Vendepunkt, hvor Faderglæden skulde
% PRINTER'S DEFECT: the quotation „Og du Barnlille o. s. v. is opened but never
% closed --- verified at 900 dpi. Transcribed as printed; this leaves the batch's
% „/“ count one opener in surplus.
bryde frem i hele sin Styrke: „Og du Barnlille o. s. v. (V. 76), driver
Forfatteren det ikke videre end til en Gjentagelse af, hvad vi allerede forlængst
have hørt af Propheten og Engelen.“ --- Lad dette saa være det sørgelige Udfald af
den negative Kritik.

Have vi saaledes efter bedste Evne forsøgt at sætte os ind i den moderne
Bevidstheds Tankegang, saa vil „den strenge Videnskab“ sikkerlig heller ikke
indvende Noget imod, at Evangelietroen nu ogsaa anlægger sin Maalestok, og
forsøger en Opfattelse af de foreliggende Hymner, svarende til dens Begreb om den
Hellig Aand. Herregud! Verden holder nu altid saa meget af Talen\-%
% --- p. 112 ---
terne og de originale Ideer, og tænker saa lidt paa de Tusinder, der, uden at være
Talenter, og uden at kunne fatte de originale Ideer, dog ogsaa skulle leve et
menneskeligt Liv ved Aandens daglige Brød. Vi tvivle ingenlunde om, at den moderne
Bevidsthed jo nok vilde blive troende, naar den blot med det Første kunde faae en
ganske ny original Aabenbaring; men hvad vi derimod høilig betvivle, er, at disse
Hymners Originalitet eller Ikke-Originalitet skulde i saa Henseende kunne gjøre
det Allermindste enten til eller fra. Hvad hjalp det vel, om den Hellig Aand havde
været betænkt paa at indsætte et Par originale Lovsange flere i det N. Tste, naar
% The book abbreviates the New Testament „N. Tste“ here (twice) but „N. Tst.“ on
% pp. 45 and 48, and „Gl. Tst.“ on pp. 110 and 117. Each transcribed as printed.
det N. Tste dog nu selv allerede er saa gammelt, at dets Ideer alene af den Grund
ikke længere kunne overraske ved deres Nyhed eller gjøre Lykke ved deres
Originalitet? Men er det dog ikke en egen Sag med disse mange originale
Aabenbaringer, hvori Tidens Viisdom idelig fødes, for at døe paany, og Ordet
blomstrer, for igjen at falme. Den Overraskelsesleg, Verdensaanden under alle
geniale Forklædninger leger med sine opvakte Børn, grunder sig øiensynligt paa en
forfængelig Brug af Livets og Sandhedens Ord. Verden kræver, at den nye Aand maa
tale med nye Tunger, og heri har den heller ikke Uret: Verden er ligesaa klog og
% PRINTER'S DEFECT: „fordie“ where the sense wants „fordre“. At 900 dpi the fifth
% glyph is a plain i, with none of the r's shoulder. Transcribed as printed.
snild, naar der handles om at kræve og fordie, som den er uforstandig og daarlig,
% PRINTER'S DEFECT: „brnge“ for „bruge“ --- n for u, verified at 900 dpi.
naar det gjælder om at modtage og brnge. I den letfærdige Brug af det aandelige
Ord ligger Forfængeligheden. Geniet, der bringer den nye Viisdom for Lyset, er
aandelig fortabt i sin originale Frembringelse: istedetfor at fremstille, hvad han
i Virkeligheden selv har oplevet, oplever han, egentlig talt, kun Fremstillingens
ideelle Smerte og Glæde. Tidens Opvakte, der føle sig kaldede til at anerkjende og
udbrede den originale Viisdom, sætte sig da i Bevægelse, og begynde nu paa deres
Viis at virke og arbeide; forsaavidt bliver der jo stridt
% --- p. 113 ---
og arbeidet nok i Ideens Tjeneste. Men det gaaer med den udbredende og tilegnende
Virksomhed, som med den frembringende: Alle ville de fryde sig ved den opgaaende
Idee, Alle ere betænkte paa Verdensforandringer i det Store; men ingen Enkelt
faaer Tid at gjøre Forsøg i det Smaae, som f.\ Ex.\ ved først at indrette sit eget
Liv efter Viisdommens Regel. Derfor gaaer det saa hurtigt med Omskiftelsen af de
originale Aabenbaringsformer, hvis Løsen er den overraskende Genialitet, og hvis
væsentlige Maal er Ideenydelse. Jo mere overraskende Formen er, desto større er
den øieblikkelige Nydelse; og jo mere øieblikkelig Nydelsen er, desto
% PRINTER'S DEFECT: „bnrtigere“ for „hurtigere“ --- two wrong sorts in one word.
% At 900 dpi the first glyph closes its bowl on the baseline and throws no
% descender (contrast the h of „hur=tigt“ two lines above, which does), and the
% third glyph carries no breve. Rejected reading: „hurtigere“.
bnrtigere indtræder igjen den Slaphed, der forventer en ny Overraskelse. Tilsidst
kommer nok den geniale Frembringelse i et saa uhyre Misforhold til Nydelsen, at
Verdensbevidstheden fortærer flere originale Ideer i et Døgn, end alle Verdens
Genier kunne frembringe i en Menneskealder. Forholder dette sig nu saaledes, da er
det jo saare godt, at den Hellig Aand ikke har anvendt sin uudtømmelige
Skaberkraft paa at overbyde det formende og digtende Genie, men er gaaet den
omvendte Vei, og har nedlagt sin guddommelige Viisdomsfylde i saadanne store,
eenfoldige Grundord, som med sund Tarvelighed kunne høres atter og atter uden
overraskende Formforandringer, efterdi de i deres første oprindelige Form ere
bestemte til at være --- alle Sjæles daglige Brød.

% Letterspaced, NOT larger and NOT bold: at 600 dpi the fount matches the body of
% „Saamange, som have Øre …“ that follows it on the same line in cap height and
% weight. The same line recurs, again letterspaced, as its own paragraph on p. 117.
\emph{Min Sjæl ophøier Herren, og min Aand fryder sig i Gud min Frelser.}
Saamange, som have Øre for den Benaadedes Høisang, høre i disse Ord den rene,
klare Gjenlyd af Aabenbaringens eget uforanderlige Grundord, af det Evighedsord,
om hvilket Forjættelsen stadfæster, at det skal bestaae, naar Himmel og Jord
forgaae. Eller skulde disse Ord maaskee ogsaa være forgængelige, som et Menneskes
Ord, eller skulde de maaskee no\-%
% --- p. 114 ---
gensinde trænge til en væsentlig Forandring? Hvilken Forandring vilde man da vel
foreslaae? Ere de ikke ved sig selv forstaaelige nok? Ethvert Barn kan jo forstaae
dem, saa klar og tydelig er deres Mening! Eller ere de maaskee ikke i sig selv
uforstaaelige nok? Vil Nogen ret lægge an paa at udforske deres Hemmelighed, see,
da er det jo, som om en levende Aand rørte sig i dem. De holde da ligesom Øie med
Forskeren, fra hvilken Side han kommer, at han ikke skal overraske dem; de passe
nøie paa, hvilke Forudsætninger og hvilket Sindelag han bringer med, at han ikke
skal tage dem forfængelig; de slutte sig, saa at sige, tættere sammen,
efterhaanden som Fortolkeren rykker nærmere, og danne tilsidst en uoverstigelig
Skillemuur imellem Tro og Vantro. Den Ikke-Troende kan desaarsag med al sin Kløgt
og Snildhed ikke aflokke dem en eneste Forklaring, endnu mindre trænge sig ind i
deres Allerinderste. Forgjæves opbyder han sin hele Tænkekraft: Ordenes
umiddelbare Forstaaelighed afspærrer ham enhver dybere Forstaaelse. Her er i
Grunden slet Intet at tænke, og derfor maa Tanken staae stille. Her nytter ingen
Sammenligning; thi selv den bittreste Smerte er i et naturligt Gemyt end ikke saa
forskjellig fra den høieste Glæde, som denne igjen er forskjellig fra den
Benaadedes Glæde i Gud. Her er ingen anden Udvei: enten lyde disse Ord paa en tom
Indbildning, og da er al videre Undersøgelse overflødig; eller de forudsætte et
Under (et sjæleligt Under), og da er enhver Forklaring umulig. Med Troen derimod
er Forstaaelsen given. Dette er saa uimodsigelig vist, at det Ene endogsaa
ligefrem kan maales med det Andet: Troen med Glæden, og Glæden med Troen. Er Troen
falsk, saa er Glæden i Gud ogsaa falsk; er Troen svag, saa er ogsaa Glæden svag;
er Troen stærk og levende, saa er Glæden i Gud det ikke mindre. Saalænge et
Menneske lever i Troen, erfarer han ogsaa Glæden; om han end straffes af Gud
% --- p. 115 ---
i Angerens Smerte, saa vidner Smerten dog om, at den Helliges Ild brænder i ham,
at Gud selv staaer ham ganske nær, og holder ham fast, og denne Vished er jo en
høiere Glæde. Saa hvile da Lovsangens Ord uforanderlige i deres Grundord, indtil
denne Verden selv gaaer under i den sidste Forandring. --- Men er det nu ogsaa
vist, at disse Ord kunne bestaae, naar Himmel og Jord forgaae? See, Evigheden
trænger paa; den varsler i Ordet, og minder os om, at alt det Endelige faaer en
Ende. Men Evigheden er derfor ikke utidig paatrængende; den veed at give Tid, at
see Tiden an; den giver Forandringen Tid, at det naturlige Hjerte dog omsider maa
forandre sig, og gribe det Uforanderlige. Saalænge Foranderligheden varer, er Alt
forblommet og blandet: Sorgen blander sig i Glæden, og Glæden i Sorgen;
Bedrøvelsen efter Verden blander sig i Bedrøvelsen efter Gud, og med Glæden i Gud
forblandes saa let, ja saa saare let den underfundige Selvkjærligheds egen
forfængelige Glæde. Men Evigheden trænger paa, at det Blandede maa klare sig, og
% PRINTER'S DEFECT: „Knn“ for „Kun“ --- n for u, verified at 600 dpi against the
% breved u of „Ueensartede“ on the same line.
det Ueensartede skille sig ad. Knn Sjælens reneste Glæde, Glæden i Gud, indeholder
jo ingen Blanding. Hvorledes skulde da Lovsangens Ord kunne opløses og forgaae,
naar de slet ingen Blanding indeholde? --- Forandringens timelige Tid er for den
% spacing.py reports a RUN at „en Betænkningens Tid“; the image at 600 dpi shows a
% tightly set line, not letterspacing. No \emph{} here.
vælgende Frihed tillige en Betænkningens Tid. Evigheden giver sine Vælgere Tid til
Betænkning; derfor falde dens Ord forskjelligt efter Tidernes Leilighed. Der er et
Ord: „Hvo, som ikke er med mig, er imod mig“, det
% PRINTER'S DEFECT: „verler“ for „vexler“ --- the same r-for-x as „forverle“
% (p. 73) and „forverler“ (p. 119); no descender at 900 dpi.
verler med et andet Ord: „Hvo, som ikke er imod mig, er med mig“. Hiint, det
strenge Ord, peger mod det Evige, dette, det mildere Ord, tager Hensyn paa det
Timelige; thi den timelige Tid er Betænkningens Tid. Betænkningen hører
Endeligheden til, og maa have en Ende (en Betænkningstid, der aldrig hører op,
fordi man med al sin Betænken dog aldrig kan faae sig betænkt, en evig
Betænkningstid er en Modsigelse). Men, naar
% --- p. 116 ---
Betænkningens Tid hører op, saa maa jo ogsaa den blandede Følelse, den blandede
Villie, den blandede Sorg og den blandede Glæde høre op. Hiint skaansomme Ord:
„Hvo, som ikke er imod mig, er med mig“, forsvinder da i det sidste strenge Ord,
der lyder ind i Evigheden som et Dommedagsord: „Hvo, som ikke er med mig, er imod
mig“. Men nu er det jo dog umuligt, at nogen Sjæl skulde kunne være med Gud, uden
tillige at ophøie Gud, og glæde sig i Gud. Det sidste Vilkaar er altsaa dette:
Hvo, som ikke er med Gud i Evighedens fuldkomne Glæde, maa stride en
Evighedsstrid mod Gud, og have Glæden til Fiende. --- Det er en Rædsel at tænke
sig en evig Fordømmelse og en evig Qval; men det er en uendelig Trøst, at den
dømmende Retfærdighed aabenbarer den guddommelige Naade. Det store Spørgsmaal,
der paa Dommens Dag skal gaae igjennem alle Samvittigheder, kan da vel udtydes
saaledes: „Fatter du nu, o Menneskesjæl, at du oprindelig skabtes til Glæde;
forstaaer du nu, at det var dette Maal, hvortil al Endelighedens Uro, alle Savn og
Sorger, alle Smerter og alle Lidelser og alle Længselens Taarer viste hen, som til
det sidste Endemaal: har du nu endelig lært, hvad det vil sige, at glæde sig i
Gud? Her gjælder ingen Halvhed, ingen undvigende Tvetydighed; her spørges paa Ja
eller Nei! Saligheden ligger i et Ja, Usaligheden i et Nei!“ --- Forholder dette
sig nu saaledes, og betænke vi det ret, at det virkelig maa være saaledes: da kan
der aldeles ikke være Tale om at forandre eller forbedre Lovsangens Ord, efterdi
disse Ord jo indeholde det klareste Svar paa det høieste Spørgsmaal. Ved
Indgangen til Salighedens Rige skal Den, „der fuldendte Løbet og bevarede Troen“,
da heller ikke have Behov at bede Aanden om et nyt Ord, som om det gamle alt var
forældet, og ikke ret vilde passe sig til en saa høitidelig Anledning; men det vil
være nok, ja til Fuldkommenhed fuldkommen nok, naar
% --- p. 117 ---
Striden er endt, naar Dommen er endt, og Himlen aabner sine Porte, at istemme
Marias Sang:

\emph{Min Sjæl ophøier Herren, og min Aand fryder sig i Gud min Frelser.}

Vi dvæle ved Lovsangens Begyndelse; thi denne Begyndelse er ikke blot en Deel af
et Hele, men selv et Hele; er ikke blot en enkelt Lovsang blandt flere andre, men
Lovsangen selv i alle evangeliske Lovsange, de Forløstes Lovsang i Tid og
Evighed. Vel sandt! den er kort, meget kort, denne Lovsang, den har kun faa Ord;
men i disse faa Ord er Salighedens Evangelium indesluttet som i sit Grundord.
Hvo, der forstaaer dette Ord, forstaaer alle Troens Lovsange; hvo, der
misforstaaer det, misforstaaer dem alle. Men, forat Nogen skal kunne forstaae
det, maa Grundordet bevise sin Kraft i et sjæleligt Under; thi den sande Glæde i
Gud er et sandt, et ubeskriveligt Under. Den, der miskjender Underet, miskjender
ogsaa Ordet, og naar Grundordet forsvinder, da er det jo intet Under, at de
enkelte Ord synke sammen i døde Sætninger og tomme Talemaader. Kritikens Anke
over Mangel paa Originalitet falder desaarsag som af sig selv. Man hænger sig i
Sætninger og Talemaader. Finder man lignende Sætninger i det Gl. Tst., saa maa
det naturligviis være derfra, de ere laante. Nu gaaer det i en Fart med at afsige
% The opening „ before „Ere“ is very lightly inked --- both OCR witnesses miss it,
% but it is there at 900 dpi, and the closing “ after „Efterligning.“ demands it.
Dommen. „Ere Lovsangens Talemaader laante fra det Gl. Tst., saa ere de jo ikke
originale; ere de ikke originale, saa ere de jo heller ikke oprindelige; ere de
ikke oprindelige, nu, saa ere de jo aandløse, og det Hele løber ud paa en mat
Efterligning.“ --- Om den formelle Kritik dog imidlertid ikke forløber sig ved
saaledes at sammenblande Begreberne: Originalitet og Oprindelighed, ville vi
forsøge at oplyse fra et andet Synspunkt. Originalitet er et Kunstbegreb,
Oprindelighed et Naturbegreb. Lad os engang kaste et Blik paa Naturen. Uden
% --- p. 118 ---
Oprindelighed ingen Natur, være sig guddommelig eller verdslig: enhver
Naturfrembringelse maa nødvendigviis bære Oprindelighedens Præg; men deraf følger
endnu ikke, at den er original. Det Originale taaler ingen Gjentagelse, men
Naturen indlader sig paa Gjentagelser i det Uendelige.

Vi erindre meget vel, hvad der for Resten kan siges om det Ulige i alt det Lige,
om Forskjelligheden i det Eensartede, hvorledes det ene Blad ikke er som det
andet, den ene Vanddraabe ikke ganske som den anden. Men denne Forskjellighed har
egentligt slet Intet med Originaliteten at gjøre. Ethvert Blad paa det samme Træ,
% PRINTER'S DEFECT: „Grnndform“ for „Grundform“ --- n for u, verified at 600 dpi
% against the breved u of „uden“ on the same line.
ja paa alle Træer af samme Art, gjentager den samme Grnndform, uden at denne
Eensformighed i mindste Maade forvansker Oprindelighedens Indtryk. Selv om Løvet i
et følgende Foraar antog aldeles de selvsamme Former, som i et foregaaende, skulde
Foraaret dog ikke tabe det Allerringeste af sin oprindelige Skjønhed. Ja man vilde
vel neppe engang blive Eensformigheden vaer (thi der hører dog unegtelig en stærk
Hukommelse til, for at kunne huske, hvorledes Træernes Blade sidde og danne sig
det ene Aar efter det andet). Med Kunstfrembringelser er det en anden Sag. Her
skyer man den eensformige Gjentagelse, fordi den altid gaaer paa Oprindelighedens
% PRINTER'S DEFECT: „strar“ for „strax“ again (cf. p. 110 above).
% PRINTER'S DEFECT: the stop after „Vildfarelsen“ prints as a period with a short
% raised stroke over it --- a damaged or wrong sort. Read as a full stop, since
% „Thi“ is capitalised.
Bekostning. --- Falder Kritiken nu paa, at bedømme Naturens Værker efter samme
Forudsætninger, som Kunstens, mærker man strar Vildfarelsen. Thi saalænge
Naturinstinctet virker lige umiddelbart i alle Frembringelser, har man
Oprindeligheden i enhver Gjentagelse, og spørger slet ikke om det Originale. Eller
vilde man ikke smile, hvis en Kritiker lod haant om Nattergalens Sang, fordi den
nu slet ikke „faldt original ud, men viste sig saa udpyntet med Reminiscenser af
hiin gamle længst bekjendte Nattergalesang, man endelig maatte være kjed af, fordi
man havde hørt den saa tidt.“ Kun naar Gjen\-%
% --- p. 119 ---
tagelsen beroer paa en fri Efterligning, maa Originalitetsbestemmelsen træde til,
for at godtgjøre Oprindeligheden. Originalitet er afledet Oprindelighed,
Oprindelighed i skaberisk Efterligning, men ikke Skaberevnens egen oprindelige
Grundbestemmelse. Geniet maa arbeide sig igjennem til original Fremstilling; den
sande Originalitet er dets Maal, Formen for dets høieste Aabenbarelse. Men
Skaberens Aand, den Uoprundne, hvori, og hvoraf Alt oprinder, er høit ophøiet over
Originalitetens Begrændsning.

For at levere et ret anseeligt Bidrag til Kritikens almindelige Sprogforvirring,
behøvede man blot at holde en Lovtale over Skaberens Originalitet: Gud Fader
maatte da i en saa aandrig Taleform blive et reent Genie, ja et uhyre Genie,
% ANTIQUA against the Fraktur --- „par excellence“ stands in roman type, plain and
% unspaced (checked at 600 dpi). Per the book's convention this is \textit{}.
Geniet \textit{par excellence}, Sønnen hans genial-originale Selvaabenbarelse, og
den gudmenneskelige Christus lige genial og lige original i Ord og Gjerninger. ---
% PRINTER'S DEFECT: „forverler“ for „forvexler“ --- r for x, exactly as „forverle“
% on p. 73; no descender at 600 dpi.
Den Forvirring, som herfra udgaaer, er alsidig nok. Man forverler det Hellige med
det Profane, det Overnaturlige med det Naturlige, Naturfrembringelse med
Kunstfrembringelse, Sagen selv med Ideen, Virkelighed med Mulighed, det Oplevede
med dets Fremstilling. Naar alle Grændser saaledes gjøres flydende, da bliver
enhver bestemt Forudsætning jo ogsaa flydende. En Begrebsforvirring er en smidig
Forudsætning. Den moderne Kritik har altsaa den store Fordeel, at kunne tilberede
sig de allerbeqvemmeste Forudsætninger, og dog rose sig af at beherske alle
Forudsætninger, og være „forudsætningsløs“. At være forudsætningsløs er her dog
nok omtrent det Samme, som at begynde med løse Forudsætninger. Skulde denne
Paastand maaskee forekomme En og Anden noget overdreven? Ved Overdrivelse vindes
Intet! Vi ville derfor prøve den paaankede Begrebsforvirring i det Exempel, der
ligger vor evangeliske Opgave allernærmest.

Naar der tales om at glæde sig, kan Glæden vel være af
% --- p. 120 ---
meget forskjellig Art, sandselig, aandelig, verdslig, religiøs o. s. v.; men
hvilken Glæde man end vil nævne, saa ligger det i dens Natur, deels at være
oprindelig, deels at yttre sig paa en eiendommelig Maade. Betragte vi nu for det
% PRINTER'S DEFECT: „strar“ for „strax“ once more. Note that the SAME page prints
% a true Fraktur x in „strax“ twelve lines below (the glyph throws a clear
% descender there and none here); both are transcribed as printed.
Første den reent naturlige eller verdslige Glæde, saa opdage vi strar en Forskjel
imellem Glæden selv og Glædens ideelle Fremstilling. Eet er det, at være glad og
yttre sin Glæde uvilkaarlig i Ord og Gebærder, et Andet er det, at fremstille
Glædens Yttringer i ideel Efterligning: Eet er det, at synge af Glæde, et Andet,
at besynge Glæden. Den Glade, der uvilkaarlig søger et Udtryk for sin Glæde, er
ikke altid Digter; og selv om han var det, saa er han dog nu saa glad, saa
sjæleglad, at han slet ikke tænker paa at digte, men kun paa at opleve Glæden. Paa
den anden Side er Digteren just heller ikke altid glad, idet han stræber at
% PRINTER'S DEFECT: a stray inked mark shaped like an apostrophe stands between
% „hvor“ and „mangen“ --- a risen space or foul sort, not a letter. Not transcribed.
fremtrylle Billedet af den Glade (hvor mangen Digter har ikke vædet sine
Glædesbilleder med Savnets Taarer). Skulle de nu, hver for sig, den Glade og
Digteren, underkastes en Kritik, hvor forskjellig maa da ikke Forudsætningen være?
Den Lykkelige, der oplever Glæden, giver sit Hjerte Luft, og søger et Udtryk for
sine Følelser. Den naturlige Følelse kræver et naturligt Udtryk; det forstaaer
sig. Men Eftertrykket falder imidlertid paa Stemningen, ikke paa Udtrykket; paa
Glæden selv, ikke paa den særegne Fremstilling af Glædens Idee; paa
Oprindeligheden, ikke paa Originaliteten. Ordet strømmer fra det overstrømmende
Hjerte, og Ingen tænker paa, at fastholde det flygtige Ord, saalænge han hører den
Glade. Om Den, der opfyldes af Glæden, ikke strax kan finde det rette Ord, om han
er ubehjælpsom, og gjentager de samme Udbrud, gjør slet Intet til Sagen. Jo flere
Gjentagelser, desto inderligere Glæde. Den glade Sjæl er altfor glad, til at tænke
paa det Originale; derfor skal han ogsaa have Lov til at vælge sine Ord, som de
bedst kunne falde. Er
% --- p. 121 ---
% p. 121 opens FLUSH at the body margin, continuing the paragraph that runs from
% p. 120 („Er“ is the last word there); verified on the image (PDF 138). No word
% is broken across the leaf.
det kun almindelige Vendinger? De blive jo udtryksfulde ved den blotte Betoning!
Er det kun døde Sætninger? Han gjør dem jo levende i sin Følelses Inderlighed! Er
det maaskee længst bekjendte og fra Andre laante Talemaader? Hvad bryder han sig
derom, de blive jo hans sande Eiendom i Glædens Oprindelighed! --- Anderledes med
Digteren. Han giver os ikke Glæden selv; men han gjengiver Glædens Idee: Ideen er
Alt ved sin Form, og derfor er den originale Fremstilling det Høieste. Digteren
er ikke bunden til Virkeligheden, og just derfor skal han gjøre den kunstneriske
% TURNED SORT: printed „Brnser“ (n set for u) where the sense wants „Bruser“.
% Confirmed at 600 dpi --- the first minim carries no breve, unlike the u of
% „Bruger“-type words on the same page; both OCR witnesses read „Brnser“.
% Transcribed as printed. First of a heavy cluster of turned sorts in this
% gathering (see also pp. 121, 122, 124, 126, 130, 131, 133, 134 below).
Brug af Muligheden. Brnser Følelsen, røres Sprogets Vande, strømme Billeder og
Tanker ind paa ham; da skal Begeistringen bevise sin Kraft i Besindigheden. Midt
i Stemningernes Oprør maa Digteraanden sidde rolig, og bevare Overblikket; thi
der skal gjøres Regnskab for ethvert Ord, endog det ubetydeligste. Taber Digteren
den ideelle Ligevægt, mislykkes ogsaa Fremstillingen: man kan jo ikke vel paa een
% Letterspaced: „spille“ alone. „den Berusede.“ that follows is ordinary body
% Fraktur --- checked at 600 dpi, where s-p-i-l-l-e stand visibly apart and
% B-e-r-u-s-e-d-e do not. spacing.py reported the RUN „ſpille Beruſede, For-“;
% only the first word is really spaced.
Gang selv baade være beruset, og dog tillige \emph{spille} den Berusede.
% PRINTER'S DEFECT: printed „Forverles“ (a genuine r where the sense wants x),
% the same fault as „strar“ pp. 56, 86, 110, 118, 120, „verler“ p. 115 and
% „forverler“ p. 119. Confirmed at 600 dpi: the letter sits on the baseline like
% the r of „Forudsætninger“ two words later, with no descender. As printed.
--- Forverles nu de Forudsætninger, der passe paa Kunstfrembringelsen, med dem,
der gjælde den naturlige Yttring, hvor uklar og meningsløs maa da ikke Dommen
blive.

Men Begrebsforvirringen gaaer endnu videre. Saalænge der nemlig kun spørges om
verdslig Glæde, kan det i mange Tilfælde være tvivlsomt, hvorvidt den umiddelbare
Livsnydelse er høiere end Ideenydelsen, Sagen selv høiere end Fremstillingen,
eller omvendt; men, naar Talen er om den religiøse Glæde, er Tvivlen hævet. Ingen
skal kunne negte, at Sjælens reneste Glæde, Glæden i Gud (forudsat, at
% TURNED SORT: printed „eu“ for „en“ (u set for n) --- the second minim carries
% the breve at 900 dpi. Both OCR witnesses read „eu“. As printed.
eu saadan virkelig gives), maa være uendelig mere værd, end alle dens ideelle
% TURNED SORT: printed „uok“ for „nok“ --- the first minim carries the breve,
% unlike the n of „nu“ at the head of the same line. As printed.
Fremstillinger, disse være nu iøvrigt uok saa geniale og originale. Kritikens
Angreb er desaarsag aldeles ubeføiet. At Maria har kjendt Hannas Hymne,
% --- p. 122 ---
og maaskee mange Gange opbygget sig ved den, at Zacharias har gransket i Loven og
Propheterne, at begge saaledes have de væsentligste Former og Udtryk for deres
fromme Følelser og Forestillinger fra det Gl.\ Tst., hvo vil betvivle dette? Det
Modsatte vilde jo netop være os aldeles uforklarligt. Men Sagen er, at her
hverken kan eller skal tænkes paa nogen original Fremstilling, saa lidt som paa
nogen vilkaarlig Efterligning. Det hele Eftertryk hviler derimod paa den
umiddelbare Oprindelighed. Aanden aander paa Hjertet, og det rette Udtryk giver
% TURNED SORT: printed „eu“ for „en“ again, twelve lines after the p. 121
% instance. At 900 dpi the breve is unmistakable, and the „en“ that opens the
% same line is set with a plain n. As printed.
sig af sig selv. --- Kun en fortumlet Kritik kan her ville see eu mekanisk
Sammensætning af laante Sætninger. Den Sjæl, der „fryder sig i Gud sin Frelser“,
% Both OCR witnesses read „ſor at“ here; at 900 dpi the sort carries a crossbar
% reaching to the right, which the long s of „Hukommelſen“ on the same line does
% not. It is f: „for at“.
behøver i Sandhed ikke at anstrenge Hukommelsen, for at oplede en Talemaade: den
har i sig selv, hvad den søger, eller rettere, den søger slet ikke, den finder,
finder, hvad Hjertet giver, finder sin Glædes eenfoldige Ord, som Fuglen finder
sin Sang, naar Solen gaaer op.

Betragte vi nu de foreliggende Hymner, hver for sig, i deres Heelhed, da viser
det sig klart, at deres indre Eiendommelighed egentlig fremgaaer af de Talendes
særegne Stilling og sjælelige Tilstand.

Maria, der har hørt Engelens Budskab, erfarer dets Kraft, fatter dets Betydning.
Den ydre Aabenbarelse bliver herved, saa at sige, oversat i en indre Aabenbaring.
Alt, hvad den Benaadede indvortes oplever, udtaler hun i Hymnen. Kun en saadan
Meddelelsesform egner sig for et saadant Indhold. Hymnen selv faaer derved en
% „Marie“ as printed, though the book has „Maria“ throughout the surrounding
% pages; not normalised.
dobbeltsidig Betydning, en ubetinget nemlig, forsaavidt Marie taler af Aandens
Drift, som for sig selv alene; og en betinget, forsaavidt hun med det Samme
besvarer Elisabeths Hilsen, og meddeler sig til hende. Hvad Bebudelsens
Hemmelighed angaaer, kan her umuligt tænkes paa nogen privat Underholdning
% --- p. 123 ---
imellem de to Veninder. Kun i den høieste religiøse Stemning kunne de udsige det
Uudsigelige. Elisabeth er ingen nysgjerrig Qvinde, at hun skulde tillade sig at
udfritte „sin Herres Moder“, for nærmere at erfare, hvorledes Engelen tog sig ud,
da han traadte ind, hvor tydeligt han talte, hvor stærk hans Stemme klang, og paa
hvilken Maade eller ad hvilken Vei den Himmelske igjen forsvandt. Havde
Zacharias's Hustru begyndt med slige Spørgsmaal, var der neppe fulgt noget Svar.
Nu derimod, da den prøvede Arons Datter, „selv opfyldt af den Hellig Aand“,
træffer Grundordet i sin udtryksfulde Hilsen: „Salig hun, som troede; thi det
skal fuldkommes, som hende er sagt af Herren“, --- kan Maria svare i det samme
Sprog, og udøse sit Hjerte i Lovsangen. --- Men endnu fra en anden Side træder
denne Hymne i et mærkeligt Forhold til Aabenbarelsen.

Engelen havde jo udtrykkeligen talt til Maria om Sønnens kongelige Værdighed,
sigende: „Gud Herren skal give ham Davids, sin Faders Stol. Og han skal regjere
over Jakobs Huus evindelig, og der skal ikke være Ende paa hans Kongerige
% The book prints „Luk.“ here (it prints „Luk.“ at pp. 42 and 100, „Luc.“ at
% p. 53); transcribed as printed.
(Luk.\ 1, 32--33).“ At blive Moder til en saadan Søn, maa dog i verdslig Forstand
vel ansees for en ubeskrivelig Lykke. „Ja, den Qvinde, der faaer en saadan
Udvælgelse, har sandelig Noget at glæde sig over!“ Og dog, hvor mærkværdigt,
disse Engelens Ord gjenlyde slet ikke med noget særligt Eftertryk i Lovsangen.
Nei! Maria tænker aldeles ikke paa at gjøre Lykke i Verden. Om hun skal see sin
Søn paa Davids-Thronen, og selv sidde ved hans høire eller venstre Side, eller om
hun maaskee skal staae under hans Kors, og et Smertens Sværd trænge sig igjennem
hendes Sjæl, derom veed den Benaadede endnu Intet, slet Intet. --- Intet
Sideblik forstyrrer de fromme Tanker, ingen verdslige Forhaabninger faae Lov at
blande sig i hendes rene Glæde. Kun dette veed hun
% --- p. 124 ---
med Troens Vished, at Israel ved hende nu faaer en Frelser af Guds
Barmhjertighed, at Herren har seet til sin Tjenerindes Ringhed, og at derfor alle
Slægter skulle prise hende salig. Kunde det nu nytte, at føre Beviser imod dem,
der altid rose sig af at føre Beviser imod Evangeliet, da behøvede vi blot at
% Letterspaced, confirmed at 600 dpi and by spacing.py (two RUNs). Same weight
% as the body: this is Sperrsatz, not the bold Fraktur of pp. 63--64.
pege paa disse vidunderligt gribende Ord: \emph{See, nu herefter skulle alle
Slægter prise mig salig.} --- Eller have disse Ord maaskee heller ingen
Oprindelighed? Saa siger os dog, hvorfra de ere laante. Hvilken anden Qvinde har
% TURNED SORT: printed „nkjendt“ for „ukjendt“ --- the first minim carries no
% breve at 1200 dpi. As printed.
da nogensinde tilladt sig at føre et saadant Sprog? Vel sandt! Sværmeriets
% Rettelser: printed „at han tør“; corrected to „at han her før“.
% The printed line 11 from the top of p. 124 does end „... at han tør“, exactly
% as the errata leaf reports, and the leaf's correction has been checked at
% 1800 dpi: the last word begins with f (ascender, crossbar AND descender), not
% with the t of „tør“ two words earlier. So the leaf really does read
% „at han her før“, and the correction is applied here per this edition's policy.
% BUT the resulting sentence is not Danish: „hvo er saa formastelig, at han her
% før mene at høre ...“ has no finite verb, whereas the printed „at han tør mene
% at høre“ reads perfectly. The leaf's own „før“ is almost certainly a misprint
% for „tør“, the substantive correction being only the insertion of „her“
% („at han her tør mene at høre“). This copy carries a contemporary pencil
% annotation in the right margin at exactly this line: „her“ with a caret
% between „han“ and „tør“ --- i.e. an early reader applied the erratum that way.
% Flagged for the editor; not resolved here.
Formastelighed er desværre ikke nkjendt i Verden; men hvo er saa formastelig, at
han her før mene at høre en Sværmerindes Brusen? Den Benaadede er i
Henrykkelsens Udbrud som en Forklaret, og hendes Røst som en Aanderøst fra de
Saliges Hjem; men hun er reen for alle overspændte Forestillinger, fra alle
bedaarende Indbildninger. Hun har i sit Hjertes Ydmyghed ikke forglemt
Tjenerindens Ringhed, ikke forglemt, at hun endnu er i det Forkrænkelige, ikke
forglemt, hvorledes den virkelige Verden seer ud. Verdens-Billedet viser sig jo
% TURNED SORT: printed „Throuer“ for „Throner“ --- the fifth sort carries the
% breve at 1200 dpi, where the n of „undertrykke“ on the same line does not.
% Both OCR witnesses read „Throuer“. As printed.
for hende med alle de velbekjendte Træk: „Her ere de Hovmodige, der spotte Gud i
deres Hjertes Tanker; her sidde de Mægtige paa deres Throuer og undertrykke de
Ringe; her svælge de Rige i Overdaadighed, medens de Nøgne fryse, og de Hungrige
sukke for Brød.“ Saaledes, just saaledes seer Verden jo ud; og Maria tager ikke
Feil. Men hvorfra skal nu Hjælpen komme? Nedbeder hun maaskee den guddommelige
Hevn over de Ugudelige, truer hun i en ophidset Sindsstemning med Sønnens
sønderknusende Vrede, skuer hun allerede i Aanden, hvorledes den vældige Hersker
skal drage ud til Strid, at udøse sine Fjenders Blod, og træde paa de Overvundnes
Nakke? Nei langtfra! End ikke ved et Ord mindes vi om de Gjengjældelsens
Sonoffre, hvormed det selvsyge
% --- p. 125 ---
Sværmeri kun altfor gjerne vil mætte den graadige Lidenskab. --- I al
Sagtmodighed udtaler Maria den Forvisning, at Frelsen kommer fra ham, „som er
% TYPE DEFECT: the second „Slægt“ is set from a damaged t whose ascender has
% broken off; at 2000 dpi the crossbar and foot are present but the upper stem
% is gone, and both OCR witnesses therefore read „Slægr“. Read „Slægt“.
mægtig, hvis Navn er helligt, og hvis Barmhjertighed varer fra Slægt til Slægt
mod dem, som frygte ham.“ Det er den frelsende og forløsende Gud, der
gjennemtrænger Verden med Kraftens Ord, og selv vender alle verdslige Forhold om,
% Letterspaced, confirmed at 600 dpi. The emphasis begins at „øver“: „thi han“
% at the end of the line is set close, and spacing.py's RUN starts at „Magt“.
uden at lade sig hindre af nogen endelig Skranke. De Trodsige kunne ikke
modstaae; thi han \emph{øver Magt med sin Arm, og adspreder dem, som ere
hovmodige i deres Hjerters Tanke.} Voldsherskerne kunne ikke modstaae; det er
forgjæves, at de ruste sig; thi den Hellige taaler ingen Uretfærdighed: han
\emph{drager de Mægtige ned af Stolene, og ophøier de Ringe.} Men de Nødlidende
og Hjælpeløse kunne trygt haabe paa ham, hvis Retfærdighed er idel
Barmhjertighed: \emph{de Hungrige opfylder han med gode Gaver, og de Rige afviser
han tomme.} Dette er, naar vi betragte Sagen i Almindelighed, slet ingen ny
Opdagelse: enhver Troende veed, hvorledes det er at forstaae, og Maria vidste det
forud i Troen. Men derved, at Troens Lys opgaaer i Glæden, bliver det Altsammen
nyt. I Forstaaelsen af Frelsens Under forstaaer Maria Alt; thi hun forstaaer med
Hjertet. Derfor fryder hendes Aand sig i Gud, derfor er hun nu salig i sin Glæde,
derfor veed hun med Ydmyghedens ubedragelige Vished, at alle Slægter skulle prise
hende salig.

Naar vi i den følgende Hymne finde Tonen noget nedstemt, saa mene vi ikke
dermed, at dens Indhold, i og for sig, skulde lide af nogen væsentlig Mangel, men
kun, at Indtrykket er særlig betinget ved de Forudsætninger, under hvilke
Zacharias udtaler sin Glæde. Her er ikke denne sjælfulde Inderlighed, denne
Henrykkelsens Salighed, som i den foregaaende Hymne. Begyndelsen selv: „Lovet
være Herren, Israels Gud, at han haver besøgt og
% --- p. 126 ---
forløst sit Folk“, antyder allerede Forskjellen. Zacharias er „bleven fyldt med
den Hellig Aand“, at han maa vidne om Forjættelsernes Opfyldelse; men han er ikke
saaledes gudopfyldt, at han kan leve, røre sig og aande i Gud, som Maria. Hans
Stilling er, i Forhold til hendes, Betragterens og Tilskuerens Stilling. Hans
% PRINTER'S DEFECT: printed „strar“ (r for x) --- confirmed at 1200 dpi, where
% the final sort sits on the baseline exactly like the r of „griber“ that
% follows, with no descender. Contrast the true „strax“ of pp. 128, 131, 132 and
% 134 in this same batch, where the descender is plain. As printed.
Hymne begynder egentlig der, hvor Marias ender. Medens hun, den Benaadede, strar
griber det Nærværende, og finder Opfyldelsen i sit eget Hjerte, søger den Gamles
Blik derimod tilbage til Folkets historiske Fortid, og hviler sig paa Davids
Huus, paa Propheterne, som have været fra fordums Tid. Zacharias er ikke saaledes
% TURNED SORT: printed „sknlde“ for „skulde“ --- the third minim carries no
% breve at 1200 dpi, where the u of „Middelpunkt“ on the same line does. As
% printed.
et Frelsens Middelpunkt, at han nærmest sknlde tænke paa sig selv; men det er
hans inderlige Glæde, at han frelses ved sit Folks Frelse. „Betryggelsen mod
Fjendernes Vold, Gudsdyrkelsen i Hellighed og Retfærdighed“ ere de Velgjerninger,
der fylde den Bedagede med Taknemmelighed, og lader ham glemme sit Livs Korthed
og Dødens Bitterhed. For sit eget Huus forlanger han Intet; han tænker ikke paa,
at den nyfødte Søn engang skal gjøre Familien Ære, og bevare hans Navn i sildige
Slægters Ihukommelse. Det lille Barn er en Fryd for Oldingens Hjerte, hans Sjæl
--- hænger ved Drengen; men han tænker slet ikke paa at beholde ham for sig selv;
% TURNED SORT: printed „dn“ for „du“ --- no breve on the second minim at
% 1200 dpi. Both OCR witnesses read „dn“. As printed.
%
% LETTERSPACING: the scripture runs „Og du, Barnlille, skal kaldes ... at berede
% hans Veie.“, but only from „Barnlille“ on is it letterspaced. „Og dn,“ closes
% the preceding line in ordinary body Fraktur (O-g and d-n stand close at
% 1200 dpi); the wide gap spacing.py flagged is the word space, not Sperrsatz.
% Emphasis marked where it is actually set.
han giver ham hen, idet han, hengiven i hans Beskuelse, giver ham sin
Velsignelse: Og dn, \emph{Barnlille, skal kaldes den Høiestes Prophet; thi du
skal gaae frem for Herrens Aasyn, at berede hans Veie.} Barnet bliver saaledes
ikke udseet til at være Herren selv, men den Tjener, der skal give Herrens Folk
„Kundskab om Saliggjørelsen ved deres Synders Forladelse.“ I Hymnens Slutning,
saavelsom i Begyndelsen (V.\ 66), høre vi en Gjenlyd af Marias Lovsang. „Lyset
fra det Høie“ har allerede besøgt Israel (V.\ 78), „for at skinne for dem, som
sidde i
% --- p. 127 ---
Mørke, og i Dødens Skygge, at styre vore Fødder paa Fredens Vei“ (V.\ 79).

Saaledes staaer Zacharias i sit Livs Aften ved den Grændse, hvor det Gamle ender,
og det Nye begynder; saaledes er den Familieaand, der i sit Skjød modtager den
Høiestes Prophet, selv forklaret i den Hellig Aand, og saaledes forstaae vi da
% PRINTER'S DEFECT: printed „vorte“ (r for x) where the sense wants „voxte“ ---
% the same r/x fault as „strar“ on p. 126 above; at 1200 dpi the sort sits on
% the baseline like the r of „Barnet“ four words earlier. As printed.
ogsaa Slutningens Ord (V.\ 80): „Barnet vorte, og blev styrket i Aanden, og var i
Ørkenerne indtil den Dag, han fremstillede sig for Israel.“

See vi nu tilbage paa disse tvende Optrin i den hellige Familie, da kunne vi
endnu fra en anden Side overbevise os om, at Aabenbaringens Kjendsgjerninger kun
ere til for de Troende. Deres hele Betydning beroer jo, som vi have seet,
udelukkende derpaa, at den hellige Glæde ikke blot er virkelig oplevet, men ogsaa
i sig selv har indre Sandhed, ja uendelig Gyldighed. Men dette kan ikke
godtgjøres ved nogensomhelst historisk Kritik. Sæt endog, at en Kritiker kunde
tage Maria og Zacharias selv i Forhør, og finde dem villige til at afgive enhver
mulig Forklaring, saa vare vi endda lige nær, eftersom jo Kritikeren ikke blot
vilde spørge dem, „hvad de troe at have erfaret“, men --- „om de nu ogsaa
virkelig have erfaret det“; og en saadan Adskillelse er i dette Tilfælde for de
Paagjældende en reen Umulighed. Ingen kritisk Inqvisition vilde, selv om
Begivenhederne foregik lige for dens Øine, være istand til at indhente andre
Oplysninger end dem, der indeholdes i de evangeliske Fortællinger. Kun Aanden
selv kan give det fuldgyldige Vidnesbyrd om deres indre Sandhed. Vil Kritiken
ikke troe disse Fortællinger for den Aands Skyld, der taler igjennem dem, saa
nytter det heller ikke, om Maria og Zacharias opstode fra de Døde, for at bevidne
deres historiske Virkelighed. Maria glædede sig, fordi hun troede; Zacharias
glædede sig, fordi han troede: hvo
% --- p. 128 ---
der vil dele deres inderlige Glæde i Gud, han opgive alle Kritikens Tvivl, og
komme hid, at dele deres Tro.

%% ============================================================
%%  SUBSECTION HEAD, printed part-way down p. 128 (PDF 145).
%%
%%  Set exactly as the head of subsection 1 on p. 100: the numeral „2.“ alone on
%%  its own line, then the title in a larger display Fraktur, then a smaller
%%  scripture-reference line. NO RULE anywhere in the head region --- checked at
%%  600 dpi over the whole block; the only long horizontal run of ink is the em
%%  dash inside the title itself. (Compare chapter II, p. 42, rules 0.38 and
%%  0.12 of the measure, and chapter III, p. 81, rule 0.08.)
%%
%%  The title is NOT letterspaced --- the display fount is simply larger, and
%%  its letters stand close (B-e-t-h-l-e-h-e-m at 600 dpi). spacing.py's RUN for
%%  this line is the usual display-heading artefact.
%%
%%  INDHOLD vs. BODY: the Indhold prints „2. Jesu Fødsel --- Hyrderne i
%%  Bethlehem ...... 128--166“, with no point after „Bethlehem“. The printed
%%  head has one: „Jesu Fødsel. --- Hyrderne i Bethlehem.“ --- and a further
%%  point after „Fødsel“ that the Indhold also lacks. Both readings stand; the
%%  head follows the page, the TOC entry follows the Indhold. Exactly the same
%%  disagreement as at the head of subsection 1 on p. 100.
%%
%%  The reference line differs in FORM from the p. 100 one, which is
%%  parenthesised („(Luk. 1, 39--80.)“). Here it is unparenthesised, and the
%%  final „V“ carries NO POINT --- verified at 1800 dpi. A dropped point;
%%  transcribed as printed.
%%
%%  The book prints „Luk.“ here (as at pp. 42, 100, 123); „Luc.“ at p. 53.
%% ============================================================
\subsection*{2.\\[4pt] Jesu Fødsel. --- Hyrderne i Bethlehem.}
\addcontentsline{toc}{subsection}{2. Jesu Fødsel --- Hyrderne i Bethlehem}

\begin{center}
  {\footnotesize Evangelium Luk.\ 2 K.\ 1--20 V}
\end{center}

% The Gospel is set at FULL MEASURE in a fount a size LARGER than the body ---
% cap heights and x-heights are visibly greater than the two body lines at the
% head of p. 128, and the weight is the same. This is the \large display
% quotation of this edition's convention, not the bold Fraktur of pp. 63--64.
\begingroup\large
Og det begav sig i de samme Dage, at der udgik en Befaling fra Keiser Augustus,
at al Verden skulde beskrives til Skat. Denne Beskrivelse var den allerførste,
som skete, der Cyrenius var Landsherre i Syrien; og de gik alle at lade sig
beskrive til Skat, hver til sin Stad. Da drog Joseph op af Galilæa af den Stad
Nazaret til Jødeland, til Davids Stad, som kaldes Bethlehem (fordi han var af
Davids Huus og Slægt), at han skulde lade sig beskrive til Skat med Maria, sin
trolovede Hustru, som var frugtsommelig. Og det skete, medens de vare der, da
blev Tiden fuldkommen, at hun skulde føde. Og hun fødte sin førstefødte Søn, og
svøbte ham, og lagde ham i en Krybbe; thi de havde ikke Rum i Herberget. Og der
vare Hyrder i den samme Egn, som laae paa Marken og holdt Nattevagt over deres
Hjord. Og see, Herrens Engel stod for dem, og Herrens Herlighed skinnede over
dem, og de frygtede saare. Men Engelen sagde til dem: frygter ikke! thi see, jeg
forkynder Eder en stor Glæde, som skal vederfares alt Folket; thi Eder er idag en
Frelser født, som er den Herre Christus i Davids Stad. Og dette skal være Eder
til et Tegn: I skulle finde et Barn svøbt, liggende i en Krybbe. Og strax var der
% „strax“ with a true x here (descender plain at 1200 dpi), twelve pages after
% the „strar“ of p. 126 --- the book alternates, exactly as on p. 120.
hos Engelen en himmelsk Hærskares Mangfoldighed, som lovede Gud og sagde: Ære
% --- p. 129 ---
være Gud i det Høie! og Fred paa Jorden! og i Menneskene en Velbehagelighed!

Og det skete, der Englene fore bort fra Hyrderne til Himmelen, da sagde Hyrderne
til hverandre: saa lader os nu gaae lige hen til Bethlehem, at vi og maae see
Det, som der er skeet, hvilket Herren nu haver aabenbaret for os. Og de kom
hasteligen, og fandt baade Maria og Joseph, og Barnet liggende i Krybben. Og der
de havde seet det, da forkyndte de overalt det Ord, som var sagt til dem om dette
Barn, og Alle, som dette hørte, forundrede sig over det, som Hyrderne sagde til
dem. Men Maria bevarede alle disse Ord, og overveiede dem i sit Hjerte. Og
Hyrderne vendte tilbage igjen, prisede og lovede Gud for alt Det, som de havde
hørt og seet, ligesom det var sagt til dem.
\par\endgroup

% The verse quotation is set as an indented block in a fount a size SMALLER than
% the body (x-height visibly less than the body line that follows it), not
% letterspaced and not bold.
\begingroup\small
\begin{verse}
„Enhver, der er vorden af Verden forladt,\\
Bringe vi Trøst i den eensomme Nat.\\
Enhver, der Dagen over maa lide,\\
Tør vi husvale ved Midnatstide.\\
Den Stund, da Frelseren fødtes paa Jord,\\
Fødtes en Trøst for den bittreste Smerte.\\
En Trøst for den, der haaber og troer,\\
En Gjenlyd af de Saliges Chor,\\
Stiger da ned i Menneskets Hjerte.“
\end{verse}
\endgroup

I disse skjønne, rørende Toner har Digteren tolket os „Englenes Lovsang“, og
derved stemt os til en inderlig Tilegnelse af de Saliges Evangelium, der
forkyndes for Hyrderne i den stille Nat. Men det hele Evangelium har Digteren
ikke forklaret. Sligt et Forsøg har han vist ikke turdet vove. Digteren er jo en
Kunstner, og den sande Kunstner har en dyb Ærefrygt for enhver af det
% --- p. 130 ---
Guddommeliges Aabenbarelser, og tager derfor aldrig Mere, end Det, der bliver ham
givet, at han ikke ved selvraadig Anmasselse skal krænke det Hellige, og vorde
straffet af den fortørnede Aand. De forskjellige Kunster maae desaarsag i Forhold
til Ideens mangesidige Aabenbarelse række hinanden Haanden: eet og samme
% TURNED SORT: printed „mauge“ for „mange“ --- the second minim carries the
% breve at 900 dpi. Both OCR witnesses read „mauge“. As printed.
guddommeligt Baand omslutter og forbinder dem indbyrdes, saa mauge som de ere.
--- Saa kommer da nu ogsaa Maleren til, og hører Evangeliet, og faaer
Beaandelsen, og maler os et Billede, et deiligt Billede. Det er Billedet med det
lille Barn i Krybben, hvorpaa Alles Blikke hvile. „Jomfru Maria, den fattige og
dog saa rige Moder, sidder nær hos Krybben, og nærer sin Sjæl --- af Barnets
Beskuelse. Hos hende staaer Joseph, den trofaste Ven, saa inderlig stille: hans
Øie vogter paa Barnet. Og Hyrderne nærme sig, den Ene med den Anden, den Ene
efter den Anden, og deres Øine søge hen til Krybben at finde Barnet. Og der er et
hensvindende Mørke, der viger tilbage i Billedet, og der er et seirende Lys, der
forklarer Billedet. Det er Lysets Seier over Mørket, ja Lyset maa seire: see, det
straaler om Barnet i Krybben!“ --- Saaledes faaer da ogsaa Farvernes skjønne
Kunst sin Andeel i Evangeliets Fremstilling; men den faaer ikke det hele
Evangelium. Her er endnu Noget, Kunstneren ikke tør røre, det er Keiser August,
% spacing.py reported a RUN here („han ike“); the image at 900 dpi shows an
% ordinary, rather tightly set line with no letterspacing at all. Not emphasis.
som beskriver al Verden til Skat. Ikke, at Kunstneren jo nok tør male en Keiser
Augustus; men --- han tør ikke male den store Keiser i det samme Billede, hvori
han maler os det lille Barn. Ja, om han endog malede Keiseren i al hans
Herlighed paa et Billede for sig; han skulde dog ikke kunne male os al den
Travlhed, der opstod i Byer og Stæder, da Befalingen udgik at beskrive al Verden
til Skat. Al Verden! Hvilken Kunstner er vel istand til at male al Verden?
Derfor, hvis Nogen forlanger at vide bedre Besked om Keiser Augustus, hvor stor
hans Magt har
% --- p. 131 ---
været, hvor vidt hans Rige strakte sig, og hvor mange Folkeslag han kunde
beskrive til Skat, da maae vi henvise ham til en anden Kunstner, og bede ham
raadføre sig med Verdenshistorien. --- Ogsaa den historiske Fremstilling er jo en
Kunst. Men Verdenshistorien, der veed saa Meget at fortælle om Keiseren i Rom,
har kun saare Lidet, fast Intet at fortælle os om Barnet i Bethlehem. Den lange
Strid om Skattebeskrivelsen selv er saa indviklet og lærd, at de Eenfoldige ikke
kunne følge dens Gang; men Evangeliet er jo for Troens Eenfoldige. Med faa og
eenfoldige Ord berører Evangelisten det historiske Vendepunkt, og strax fornemme
vi et Indtryk af Verdensrigets Storhed og Magt. Thi det kan et Barn forstaae, at
% TWO TURNED SORTS in one line: printed „kau“ for „kan“ (breve present on the
% third minim) and „ndgaae“ for „udgaae“ (breve absent on the first). Both
% confirmed at 900 dpi; both OCR witnesses read them so. As printed.
denne Augustus maa være en stor og mægtig Keiser, siden han kau lade en saadan
Befaling ndgaae, at al Verden beskrives til Skat. Med faa og eenfoldige Ord
beretter Evangelisten, at Joseph tilligemed Maria, sin trolovede Hustru, drog fra
Nazaret til Bethlehem, fordi han var en ringe Mand, der maatte adlyde Keiserens
strenge Befaling. Men det kan den Eenfoldige forstaae, at her foruden Keiser
August endnu maa være en anden Herre, som har at befale over al Verden: det
burde jo Christus, Davids Søn, at fødes i Bethlehem, Davids Stad. Med rene Ord
siger Evangelisten os, at Marias Førstefødte blev lagt i en Krybbe, fordi der
ikke var Rum i Herberget; men --- de himmelske Hærskarer vidne jo for Hyrderne
paa Marken, at Jordens Pragt og Verdens Herlighed kun er Tomhed og Tant, et Intet
for Ham, den Ærens Konge, der hviler i Krybben. --- Som nu alle disse Hovedtræk
samle sig for Beskueren til eet Billede, der ener det Jordiske med det
Himmelske, som nu alle disse Tilværelsens stærke Modsætninger bryde sig paa
% DEFECT: printed „ndre“ where the sense wants „indre“ --- at 1200 dpi the word
% is n-d-r-e, the first sort a plain n with no breve, and no i before it. A
% dropped or wrong sort; both OCR witnesses read „ndre“. As printed.
hinanden, og gjennemlyse hinanden i Billedets ndre Uendelighed, hvor maa da ikke
al menneskelig Fremstillings\-%
% --- p. 132 ---
kunst føle sig et fattigt, forfængeligt Stykværk imod dette Rigdoms-Dyb af
aabenbaret Sandhed og Naade! ---

„Sandhed og Naade?“ hvisker en Tvivlens Aand, og den negative Kritik begynder
igjen at hæve sin Stemme. „Historisk Sandhed er her i denne Fortælling kun Lidet
af; men, det skal ikke negtes, desto Mere er her af den opbyggelige Illusion. ---
Forat Jesusbarnet kan blive født i Bethlehem, udkræves intet Mindre, end at
Keiser August maa beskrive den hele Verden til Skat! Skulde Nogen nu mene, at der
til Opnaaelsen af et saa rimeligt Øiemed dog forlangtes altfor urimelig Meget
baade af Keiser August og den hele Verden; saa behøver han blot at besinde sig et
Øieblik og see, at den evangeliske Skattebeskrivelse egentlig først foregaaer paa
den Tid, da Cyrenius (Qvirinius) var Landsherre i Syrien. Denne Opdagelse vil
strax forsone den kritiske Kjender med Evangelistens „troværdige Beretning“; thi
heraf er det jo aabenbart nok, at det i slige troværdige Beretninger, egentlig
talt, slet ikke kommer an paa, hvad der nu virkelig skeer her paa Jorden, naar
man kun faaer en desto nøiagtigere Beskrivelse over de Mærkværdigheder, der
tildrage sig i Himlen. Ifølge Historiens Vidnesbyrd er det nemlig ikke Qvirinius,
der i Herodes den Stores sidste Regjeringsaar var Syriens Præses, men det er
Sentius Saturninus. Efter ham fulgte Qvintilius Varus, og først længere Tid efter
Herodes's Død tiltraadte Qvirinius Bestyrelsen af Syrien.“

--- Saa stor, saa uhyre stor er denne kritiske Anstødssteen, naar Vantroens Øie
ret faaer den at see: jo længere man seer paa den, desto større bliver den! ---
Saa lille, saa ubetydelig og lille er denne kritiske Anstødssteen, naar den sees
med et troende Øie: jo længere man seer paa den, desto mindre bliver den! --- I
% PRINTER'S DEFECT: printed „Grundtert“ (r for x) where the sense wants
% „Grundtext“ --- at 1200 dpi the sort sits on the baseline like the r of
% „Grund-“, with no descender. The same fault as „strar“ p. 126 and „vorte“
% p. 127. As printed.
den evangeliske Grundtert staaer det jo ikke skrevet, at den hele Verden virkelig
blev beskattet, da Christus fødtes; men der staaer saaledes: \emph{en Befa\-%
% --- p. 133 ---
% The word „Befaling“ is broken across the leaf (p. 132 ends „en Befa=“, p. 133
% opens flush with „ling“) and is joined with \- above. The letterspacing begins
% on p. 132 at „en Befa-“ and runs to „Mandtal“.
ling udgik fra Keiser Augustus, at al Verden skulde skrives i Mandtal} (for
% PRINTER'S DEFECT: „var“ is set twice --- „hvis det saa var var muligt!“.
% Confirmed at 1200 dpi; both OCR witnesses read the doubling. As printed.
derefter nemlig at foretage en Census, og, hvis det saa var var muligt! inddrive
Skatten). Havde det nu været Lukas's Hensigt, at sætte os omstændelig ind i
Datidens mange Skattehistorier, vilde han vistnok have udtrykt sig nøiagtigere
saaledes: „Denne Folketælling var det første Skridt til den første romerske
% TURNED SORT: printed „Cyrenins“ for „Cyrenius“ (the word is broken „Cy-“ /
% „renins“ across the line) --- no breve on the fourth minim at 1200 dpi. Both
% OCR witnesses read „renins“. As printed; contrast the correct „Cyrenius“ on
% pp. 128 and 132.
Skatte-Udskrivning i Judæa, en Skatte-Udskrivning, der dog først senere kom
istand, nemlig ti Aar efter, da Cyrenins var Landsherre i Syrien.“ Nu derimod, da
Evangelisten ikke vil anføre os til at studere Beskatningssystemet i Rom, men
derimod forkynde al Verden det glade Budskab om Frelserens Fødsel i Bethlehem, er
han, hvad Skatte-Udskrivningen selv angaaer, vel berettiget til at vælge en
kortere Udtryksmaade; og dette saa meget mere, som det historiske Eftertryk her
% Letterspaced: „Befaling“ ALONE. „den keiserlige“ before it and „denne
% Befalings nærmere Udførelse“ after it are ordinary body Fraktur --- checked at
% 900 dpi. The comma stands outside the emphasis, separated by the usual
% letterspacing gap.
nødvendigviis maa falde paa den keiserlige \emph{Befaling}, ikke paa denne
Befalings nærmere Udførelse.

% QUOTE STRUCTURE: three openings („Men“ ... „om man nu ogsaa ... „Først lader
% der sig ...) and, at the close on p. 134, a single „Historie.“ --- so the
% nested quotation of Strauß is closed once for two openings. Transcribed as
% printed; the imbalance is the compositor's, not ours.
%
% „David Strauß“ is Fraktur (with ß), not antiqua --- checked at 1200 dpi.
% „Schleiermacher“ below IS set in a contrasting fount, but it is a lighter,
% narrower FRAKTUR (the ch ligature and its descender are unmistakable at
% 1800 dpi), not antiqua; so neither \textit{} nor \textbf{} applies and it is
% left plain. Logged because it is the only such contrast in this batch.
„Men“ --- vedbliver den moderne Kritik --- „om man nu ogsaa vil lade det gaae med
Beskatningen, som det bedst kan; hvad bliver der saa af Englene og de himmelske
Hærskarer? Lad os en Gang høre et Alvorsord af David Strauß: „Først lader der sig
med Billighed spørge, hvortil Engle-Aabenbarelsen skal tjene? Det nærmeste Svar
er: til at gjøre Jesu Fødsel bekjendt. Men den bliver jo derved saalidet
bekjendt, at Magerne siden efter bringe den første Efterretning om den nylig
fødte Jødernes Konge til det saa nær liggende Jerusalem, og at der overhovedet i
den senere Historie ikke findes noget Spor af en saadan Begivenhed ved Jesu
Fødsel. Kan altsaa Hensigten af denne overordentlige Begivenhed ikke have været,
at gjøre Jesu Fødsel bekjendt i videre Kredse, maa man vel antage med
Schleiermacher, at dens Formaal har været, at virke umiddelbart paa Hyrderne
selv. Men derhos
% --- p. 134 ---
maatte man da forudsætte med samme Forfatter, at Hyrderne havde, som hiin Simeon,
været særdeles opfyldte af messianske Forventninger, og at Gud har villet belønne
og bestyrke denne deres fromme Tro ved hiint Syn. Men hverken beretter
Fortællingen Noget om en saadan Stemning hos Hyrderne, ei heller bliver der
bemærket nogen blivende Virkning paa dem; men især synes, ifølge den hele
Fremstilling, Intet Hyrderne angaaende at have været Formaalet for Synet, men
blot Forherligelsen og Bekjendtgjørelsen af Jesu Fødsel som Messias. Men da det
Sidste, som vi allerede have bemærket, ikke blev opnaaet, og det Første, som et
tomt Gjøgleri, ikke var noget Gud værdigt Formaal, stiller denne Omstændighed,
allerede for sig betragtet, en ikke ubetydelig Hindring i Veien for den
supranaturalistiske Opfattelse af denne Historie.“

Uden strax at paakalde Evangelietroen, kunne vi i nærværende Tilfælde frit
henvende os til ethvert Menneske, der har nogen Sands og Følelse for det
Guddommeliges Tilsyneladelser i denne Sandsernes Verden, og spørge ham, om dog
ikke dette Evangelium med Keiser Augustus og Barnet i Krybben, med Hyrderne paa
Marken og Englenes Chor, om ikke dette Evangelium, sige vi, maa i den Grad
% QUOTE STRUCTURE / PRINTER'S DEFECT: the sigh is printed with ONE opening „ and
% TWO closings --- after „Sandhed;“ and again after „sandt!“. At 1800 dpi under
% autocontrast the mark after „Sandhed;“ is a heavily over-inked closing “ (two
% merged commas at cap height), not a blot; both OCR witnesses record a
% quotation mark there. The compositor evidently dropped the opening „ before
% „men“. Transcribed as printed; the surplus “ is the only unpaired closing
% quotation mark in this batch.
%
% „ak desværre“, not „af“: tesseract reads f, but at 1200 dpi the sort has the
% arm-and-leg of Fraktur k and sits on the baseline, where the f of
% „forunderlig“ on the line above carries a plain descender. The ABBYY witness
% also reads k.
tiltale hans bedre Væsen, og bevæge Hjertet saa forunderlig saligt, at han
uvilkaarligt kunde sukke: „Give Gud, at det var Sandhed;“ men --- ak desværre,
dette Syn er vist altfor saligt til at være sandt!“ --- Det Salige har ingen
Sandhed, og det Sande ingen Salighed: her er en Følelse af Tilværelsens hemmelige
Smerte. Derfor har Mennesket en naturlig Angst for den kolde Sandheds evige
% TURNED SORT: printed „iud“ for „ind“ --- the second minim carries the breve at
% 1200 dpi. „sneg“ on the same line is correct (plain n). As printed.
Gaade: derfor tænker man med Gru paa hiin forvovne Yngling, der ved Maanens Skin
sneg sig iud i Templet at løfte Sløret fra Gudindens Aasyn. „Det Salige har ingen
Sandhed, og hvad der i sig selv ingen Sandhed har, kan
% --- p. 135 ---
% p. 135 opens FLUSH at the body margin (PDF 152), continuing the paragraph and the
% quotation that run from p. 134; the quotation opened there closes here at
% „virkeligt.“ and the paragraph runs on in the same line. No word broken across the leaf.
aldrig blive virkeligt.“ Dette er en af de trøsteløse Grundsætninger, hvortil den
moderne Bevidsthed støtter sig; derfor kan den negative Kritik neppe dølge sin
uvidenskabelige Harme, skjøndt den med tilsyneladende Upartiskhed fremsætter det
„videnskabelige“ Spørgsmaal: „hvortil saa det hele syngende Engle-Chor vel skulde
tjene?“ Det Virkelige har ingen Salighed, og det Salige ingen Virkelighed: ikke sandt?
denne Tanke er som en nagende Orm, der ikke lader Sjælen Ro, før Mennesket beslutter sig
til, enten ganske at fornegte Aandens dybere Krav, eller ogsaa at bøie sig i Ydmyghed
under den Almægtiges Viisdom, og antage Underet. Før det imidlertid kommer saa vidt, er
det naturligt nok, at Den, hvis bedre Væsen ikke tillader ham at trodse Salighedens
Aand, men hvis Verdensdannelse dog paa saa mange Maader sætter sig op imod Underets
umiddelbare Virkelighed, gjør mange Forsøg paa at affinde sig med Verdslighedens Lov, og
% „Exempel“ here has a TRUE x --- the sort descends plainly below the baseline at
% 1200 dpi, unlike the r of „Bar-“ on p. 141. Compare „eregetiske“ eleven lines below,
% which has none.
forsone sig med Aabenbaringens „Idee“. --- Som Exempel paa et saadant Forsøg kunne vi her
anføre den æsthetiske Opfattelse af de evangeliske Beretninger.

„Hvor er den heftige Strid om Aabenbaringens Troværdighed dog ikke en gold og ufrugtbar
Strid! Der staae nu disse „Rettroende“, og ivre for det bogstavelige Udtryk, og kunne
% TURNED SORT: printed „uu“ for „nu“ (u set for n), confirmed at 1200 dpi --- both
% minim pairs carry the u-bowl, and both OCR witnesses read „uu“. As printed. First of
% a heavy cluster of turned sorts in this gathering: see pp. 138, 140, 142, 143, 147.
ikke see den simple Sandhed for lutter Vidundere. Det skal uu altsammen have samme
Virkelighed, baade det Naturlige og det Overnaturlige; som om en vis Sum af mirakuløse
Phænomener, trods al deres Flygtighed, nu virkelig kunde udgjøre Troens væsentlige og
blivende Grundlag. Denne egensindige og borneerte Vidhængen ved det Iøinefaldende og
Tilfældige ægger saa igjen de eensidige Kritikere, der med al deres Skarpsynethed dog
% WRONG SORT: printed „eregetiske“ for „exegetiske“ --- the sort after the first e is a
% plain r with no descender at 1200 dpi (compare the g of „-getiske“ immediately after).
% The book alternates x and r within a page in this gathering; transcribed as printed.
ikke kunne faae Øie paa Sagen selv, fordi deres Forstand er saa beskjæftiget med at
opdage Supranaturalisternes eregetiske Vilkaarligheder, og
% --- p. 136 ---
efterjage historiske og metaphysiske Modsigelser i alle Retninger. Denne Strid er
ligesaa utaalelig for den sunde Forstand, som oprørende for den fromme Følelse.
Vildfarelsen stikker deri, at begge Parter stride om en Virkelighed, for hvis egentlige
Betydning de have forsømt at gjøre sig Rede. --- De Engle, der viste sig for Hyrderne,
have neppe været saadanne udvortes Skikkelser, som enhver Anden ogsaa vilde kunnet see
med sine naturlige Øine, hvis han til rette Tid var gaaet over Markerne udenfor
Bethlehem; den Klarhed, der omskinnede Hyrderne, vilde sikkert ligesaa lidt være bleven
synlig, som de Toner, de hørte, vilde været hørlige for nogen Uvedkommende, om han end
tilfældigviis havde befundet sig i deres Nærhed. Derimod maa den hele Aabenbarelse
antages at være foregaaet i de fromme Mænds eget Indre; og, naar dette antages og
% TYPE DEFECT: this „Kjendsgjerning“ is set from an i whose dot is gone --- at 1800 dpi
% the stem stands bare between the two n's, and both OCR witnesses therefore read
% „Kjendsgjernmg“. The next line's „Kjendsgjerning?“ carries its dot plainly. Read
% „Kjendsgjerning“.
forstaaes, som det bør forstaaes, kan den negative Kritik ikke længere forarge sig over
denne Kjendsgjerning.“

„Denne Kjendsgjerning? Saa er her jo slet ingen Kjendsgjerning!“ --- „Indvendinger af
denne Natur bevise just, paa hvor lavt et Trin Striden holder sig. Alt, hvad der ikke
har udvortes haandgribelig Virkelighed, anseer man for uvirkeligt; ret som om der slet
ingen indvortes Virkelighed gaves. Er da ikke Sandhedens Væsen usynligt? Maa ikke
Phantasiens Skikkelser sees med det indvortes Øie? Tør Nogen vel paastaae, at de Syner,
der i Begeistringens lykkelige Øieblik opgaae for Digterens indre Sands, kun ere tomme
Indbildninger, fordi vi ikke forefinde saadanne Skikkelser i den udvortes Verden?
Hvorfra komme de høiere Indskydelser, uden fra den uudgrundelige Skaber-Aand, der selv
giver alle Ting Liv og Aande og Alt? I Sandhed, ogsaa de salige Drømme ere fra Gud! Hvad
skulde det vel ellers betyde, at Drømme og Syner i det Gl. Tst. (hvad det N. T. angaaer,
tænke man paa Josephs Drøm, Matth.\ 1, 20, paa Peters Syn,
% --- p. 137 ---
% p. 137 opens flush; the citation continues the parenthesis begun at the foot of p. 136.
% The book prints „Ap. G.“ with a point after Ap (not the comma tesseract reads), and
% „Joh.“, not „Ioh.“, here.
Ap.\ G.\ 10, 10---11, paa Joh.\ Aabenb.\ 1, 10 o.\ fl.\ St.) fremsættes som
guddommelige Aabenbarelser? Hvorledes vilde man i modsat Fald kunne forstaae en
Ezechiels eller en Daniels prophetiske Syner? Eller høre slige Syner maaskee slet ikke
med blandt de aabenbarede Kjendsgjerninger? --- Den Indvending, at alle Visioner i
Grunden kun ere Fostre af Menneskets egen Indbildning, falder bort, naar vi nærmere
overveie Sagen. Hvorpaa skal jeg kjende, at en udvortes Gjenstand har Selvstændighed,
uden derpaa, at jeg kan fornemme den Modstand, den gjør, og føle den Nødvendighed, der
nøder mig til at forestille mig den, som den nu engang er given. Men skulde den indre
Virkelighed nu ikke have et lignende Kjendemærke? --- Naar En siger, at han elsker,
eller, at han hader: hvorledes skulde han da kunne vide med sig selv, at han siger
% NOT emphasis: the four lines from „Det Samme gjælder“ to „Følelse.“ are loosely
% justified --- the word gaps are wide but the letters inside each word stand close at
% 600 dpi, and spacing.py reports no run. Ordinary body setting.
Sandhed, hvis han ikke umiddelbart følte enten Kjærlighedens eller Hadets uvilkaarlige
Indskydelser? Det Samme gjælder om Sandhedens Erkjendelse: den levende Overbeviisning er
altid begrundet i en umiddelbar Følelse. Kjendetegnet paa den indre Virkelighed er
altsaa den Nødvendighed, hvormed Tilværelsens Magter paavirke den indre Sands. Saaledes
kunne vi da forstaae Aabenbaringens Syner: naar det Guddommelige vinder Skikkelse, og
fremstiller sig for Sjælens Øie med en saadan Bestemthed og Nødvendighed, at Seeren
% WRONG SORT: printed „forverle“ for „forvexle“ --- plain r, no descender, at 1200 dpi.
% „t. Ex.“ four lines below has a true x on the same page.
umulig kan forverle Synet selv med Phantasiens øvrige, vilkaarlige Frembringelser, da
maae vi tillægge den aabenbarede Skikkelse baade væsentlig Selvstændighed og indre
Virkelighed i Kraft af en overnaturlig Indflydelse. --- Den Indvending, at det paa denne
Maade bliver saare vanskeligt at adskille det Sande fra det Falske, det Guddommelige fra
det blot Menneskelige, maa ligeledes falde bort ved nærmere Prøvelse. Der kan t.\ Ex.\
føres megen Strid om den oprindelige og afledede, den sande og falske Poesi;
% --- p. 138 ---
men det er dog indlysende, at denne Strid ikke kan afgjøres ved historisk Kritik. Hvis
% TURNED SORT: printed „omspnrgte“ for „omspurgte“ (n set for u), confirmed at 600 dpi;
% both OCR witnesses read „omspnrgte“. As printed.
dette var Tilfældet, maatte man jo i Forveien vide, hvo der var den sande Digter, og hele
Undersøgelsen dreie sig om, hvorvidt det omspnrgte Digt var af den foregivne Forfatter,
eller ikke. Denne Strid er meningsløs. Hvo er den sande Digter? Svar: Den, hvem Musen
bringer den digteriske Aabenbarelse. Hvo har vel faaet den digteriske Aabenbarelse?
Svar: Den, der producerer den sande Digtning. Paa Digtningens indre Sandhed skal jeg
altsaa kjende den sande Digter, den Lykkelige, hvem Musen besøger; og ikke omvendt.“

„Paa lignende Maade forholder det sig ogsaa med den vidtløftige Strid, der føres om
Adskillelsen imellem de kanoniske og apokryphiske Evangelier. Denne Strid kan maaskee
være en gavnlig Motion for de Lærde indbyrdes; men den giver kun et ringe Udbytte for
den religiøse Erkjendelse. Den historiske Kritik kan aldrig faae Ende; jo længere man
strider, desto større bliver Forvirringen, og hvorfor? Fordi Striden i sig selv er
meningsløs. Hvo er den sande Evangelist? Den, der forkynder os det sande Evangelium! Men
hvorpaa skal jeg saa kjende det sande Evangelium, uden derpaa, at det tilfredsstiller
min dybeste religiøse Trang? Guds Aand vidner med min Aand, at dette maa være det sande,
uforfalskede Evangelium. --- Hvad her er fremsat i Almindelighed, finder en særegen
Anvendelse paa den foreliggende Beretning. Hvilken Forskjel imellem dette kanoniske
Evangelium og de apokryphiske Opdigtelser, hvormed en from, men overtroisk Phantasie har
villet udsmykke Jesu Fødsel! Da Joseph, hedder det i Jakobs Protevangelium, havde ført
Maria paa et Esel til Bethlehem, for at indskrives til Skat, begyndte hun i Nærheden af
Staden at tee sig snart sorrigfuld, snart munter. Da man spørger hende om Aarsagen,
% The mark between the two numerals sits on the baseline and is read as a comma here
% (the book's form for chapter and verse throughout); the ABBYY witness reads a point.
giver hun det Svar, at hun (efter 1 Moseb.\ 25, 23) saae to
% --- p. 139 ---
% The book prints „Luc.“ on this page (as at p. 53); the head of the subsection on
% p. 128 prints „Luk.“
Folk for sig, det ene grædende, det andet leende; det vil sige, efter den ene Forklaring,
de to Dele af Israel, hvoraf Jesu Fremtræden skulde blive (Luc.\ 2, 34) den Ene til
Fald, den Anden til Opreisning; men, ifølge den anden Udlægning, Jødernes Folk, der siden
% „strax“ with a TRUE x --- the sort descends plainly below the baseline at 1200 dpi,
% quite unlike the r of „vorer“ on p. 141. Same on p. 143.
efter forkastede Jesus, og Hedningernes, der antog ham. Da Maria strax derpaa bliver
overfalden af Fødselssmerterne, bringer Joseph hende til en ved Veien liggende Hule, og
beskriver nu Fødselstimens Høitidelighed i følgende Udtryk: Jeg (Joseph) saae Himlens
Pol staae stille, Luften studse, og Himlens Fugle skjælve. Jeg saae hen over Jorden:
Faarene vare spredte omkring; de gik hverken frem eller bleve staaende. Hyrden opløftede
Haanden, for at drive dem frem med sin Stav; men hans Haand blev staaende oprakt. Jeg
saae hen til Flodens Strøm, saae Kiddene strække Hovederne over Vandet, men uden at
drikke o.\ s.\ v.“

„Den Indvending, at Adskillelsen imellem den kanoniske og den apokryphiske Fremstilling
ifølge denne Anviisning dog i det Høieste kun bliver en Adskillelse imellem den ægte og
den forfeilede Mythe, beroer paa en reen Misforstaaelse. --- Kirkeligsindede Theologer
have paa en vis Maade Grund til at forarge sig over den hyppige Anvendelse af Ordet
Mythe, naar Talen er om Fortællinger i det Gl.\ og N.\ Tst. Men, for at modarbeide den
mythographiske Omtydning, nytter det kun lidet, at ville oversætte de indre
Aabenbarelser i den ydre Virkelighed, eller beraabe sig paa tørre Kjendsgjerninger og
historiske Vidnesbyrd. Istedetfor at stille det Virkelige uden videre op imod det
Mythiske, gjorde man vist bedre i at eftervise den indre Modsætning imellem Mythe og
Aabenbaring. Disse tvende Begreber forholde sig ikke ligefrem til hinanden, som det
Falske til det Sande (ogsaa Mythen udtrykker paa sin Viis en væsentlig Idee, og
forsaavidt en Sandhed); men de forholde sig
% --- p. 140 ---
% This page carries the densest cluster of turned sorts in the batch --- „Aabenbarigen“,
% „deu“, „Standpuukt“, „Dauuelsen“, „Cnlturtriu“, „udeu“ --- each confirmed at 1200 or
% 1800 dpi and read alike by both OCR witnesses. All transcribed as printed.
som det Profane til det Hellige, som det Naturlige til det Overnaturlige. Ideen, der
kommer til Gjennembrud i Mythens Form, viser tilbage til Verdens egen gudbeaandede
% DROPPED SORT: printed „Aabenbarigen“ for „Aabenbaringen“ (the n of -ring- is missing);
% confirmed at 1200 dpi, comma following. As printed.
Naturgrund; den Idee derimod, der vinder Skikkelse i Aabenbarigen, er gjennemtrængt af
Guddommens eget overnaturlige Væsen, og bestemt til at forkynde Guds hellige Villie.
Herpaa beroer vistnok den egentlige Forskjel imellem Mythologiens Guder og Evangeliets
Engle. Med denne Forskjel er det da tillige givet, at Mythen maa forholde sig til
Aabenbaringen, som det Forsvindende til det Blivende. Da nemlig det Naturlige og blot
Menneskelige har en eensidig Overvægt i den mythiske Forestilling, da det
% WRONG SORT: printed „forverlet“ for „forvexlet“ --- plain r, no descender, at 1800 dpi.
Ikke-Guddommelige herved paa en uklar Maade bliver forverlet med det Guddommelige: saa
indeholder Mythen forsaavidt en Usandhed; det ligger da i Sagens Natur, at deu opvorende
% TURNED SORT „deu“ for „den“, and WRONG SORT „opvorende“ for „opvoxende“ (plain r) ---
% both in the same line, both at 1800 dpi. „Standpuukt“ on the next line is a further
% turned sort (u for n).
Tænkning efterhaanden maa opdage denne Usandhed, og føre ud over det mythiske
Standpuukt. Anderledes med Aabenbaringen. I den hellige Idee har det guddommelige
Indhold saaledes gjennemvirket Formen, at her ikke er Mere af det Endelige og
% WRONG SORT: printed „sor“ (long s) for „for“ --- the stem carries no crossbar at
% 1800 dpi, and both OCR witnesses read s. As printed.
Menneskelige tilbage, end hvad der netop udkræves, forat den rene, overnaturlige Sandhed
kan anskueliggjøres sor Mennesket, og paavirke det hele Menneske. Eftersom nu Dauuelsen
% „Dauuelsen“ for „Dannelsen“ and „Cnlturtriu“ for „Culturtrin“ --- two turned sorts in
% each word, confirmed at 1200 dpi. „Reflerionen“ is the plain-r form again (no
% descender); the book prints no true x anywhere on this page.
skrider frem, kan Reflerionen vel bidrage Sit til at inderliggjøre Aabenbaringens
uendelige Gehalt; men aldrig vil den menneskelige Tænkning paa nogetsomhelst Cnlturtriu
kunne opløse de evangeliske Forestillinger, uden at geraade i en umenneskelig Modsigelse
med den guddommelige Sandhed.“

% NOT emphasis and NOT a larger fount: this paragraph is very loosely justified --- the
% word gaps are wide --- but the line pitch (123 px at 600 dpi) and the x-height are
% exactly those of the body above it, and the letters inside each word stand close.
% Compare the genuine \large display Gospel on p. 128, whose x-height is a third greater
% at the same pitch. Ordinary body setting.
„Et slaaende Beviis for denne Sætnings Rigtighed have vi her i Evangeliet om Christi
Fødsel. Dette Evangelium er os overleveret i en saa ædel, reen og fuldkommen Form, at
% TURNED SORT: „udeu“ for „uden“.
det, udeu at forandre sin indre Charakteer, ledsager det religiøse Gemyt paa Dannelsens,
Erfarin\-%
% --- p. 141 ---
% The running head is 141; the 4 is worn and both OCR witnesses read it as 1.
gens og Selverkjendelsens forskjellige Livstrin, ret ligesom en Stjerne, der, endskjøndt
den staaer ubevægelig paa Himlen, dog synes at bevæge sig, for at ledsage Vandreren.
Dette Evangelium, saa fatteligt, eenfoldigt og elskeligt, er jo, som var det kun skrevet
for Børn og barnlige Sjæle. Den barnlige Phantasie bliver ligesaa let fortrolig med
Englene, som med Hyrderne. Adskillelsen mellem den ydre og indre Virkelighed gjælder
slet ikke for Barnet; her kan altsaa ikke være Tale om nogen Mythe: „det er Altsammen
% x AND r WITHIN TWO LINES: „voxer“ here has a true x (plain descender at 1200 dpi);
% „vorer“ on the next line, and both occurrences twelve lines below, have a plain r with
% none. Each transcribed as the glyph stands.
baade vist og sandt i alle Maader“. Saa voxer Barnet op til Yngling; men Evangeliet
vorer med. En smertelig Strid imellem den haarde, mangekantede Virkelighed og de
Længsler og Fordringer, der reise sig i den Unges Bryst, begynder. „Det Nærværende vil
ikke tilfredsstille; dog, Tilfredsstillelsen maa vel komme, Haabet kan ikke skuffe: hist
i det Fjerne vinker Idealet“. Har den ungdommelige Phantasie nu Kræfter nok til at
fastholde sit Ideal (og det er jo Ungdommens herlige Fortrin at have saadanne Kræfter);
har det religiøse Gemyt baade Styrke og Mod til at opløfte Øiet mod Salighedens
Forbilleder, og anvende disse Forbilleder paa det virkelige Livs arbeidende Higen (og
det er jo Ungdommens skjønne Forret at have en saadan Styrke og et saadant Mod): maa da
% Letterspaced, confirmed at 600 dpi. The emphasis begins at „en“ --- „er“ at the end of
% the line is set close --- and ends with the point after „Folket“; „Saa fortsættes“
% resumes in the ordinary tracking. The words are Luk. 2, 10, quoted from the Gospel set
% at the head of the subsection on p. 128.
ikke Fremtiden lyse for den begeistrede Yngling med Evangeliets Klarhed, og Englene selv
kræves til Vidner, at der dog virkelig er \emph{en stor Glæde, som skal vederfares alt
Folket.} Saa fortsættes Striden: Ynglingen vorer op til Mand; men Evangeliet vorer med,
og forklarer sine Forjættelser i Striden. „Paa denne Jord er Saligheden vel ikke at
finde; men Jorden er dog ogsaa velsignet af Gud. Ja, Velsignelsen har taget Forbandelsen
% WRONG SORT: printed „Løu“ for „Lov“ (u set for v) --- confirmed at 1200 dpi, where the
% final sort carries the u-bowl and matches the u of „Frugt“ on the same line. Both OCR
% witnesses read u. The Rettelser page records the same fault at Side 165 („Loø“ ->
% „Lov“) but not this one; transcribed as printed.
bort; den redelige Stræben bærer sin Frugt, Arbeidet har sin Løu, det Hele skrider frem
mod Fuldendelsen hisset“. Vistnok indtræder nu her den skarpe Adskillelse
% --- p. 142 ---
imellem det Virkelige og det Ideale; men saa har ogsaa Reflerionen naaet den Modenhed,
at den kan gjøre tilbørlig Forskjel imellem indre og ydre Virkelighed. Inderliggjørelsens
Kraft vorer under de mange udadgaaende Bestræbelser. Istedetfor at forflygtige sig til
% PRINTER'S ERROR, twice in one line: printed „Manddomntens“ for „Manddommens“ and
% „Barndontmens“ for „Barndommens“ --- the mm of each word is set as m+n+t and n+t+m
% respectively. Confirmed at 1200 dpi; both OCR witnesses read the same. As printed.
en Mythe, aabner Aabenbaringen nu først sin indre og sande Kjerne: Manddomntens
Overbeviisning bekræfter Barndontmens Tro. --- Med god Samvittighed kan da den Ældre
læse dette Evangelium igjen for et Barn, og besvare ethvert Troskyldighedens Spørgsmaal:
„Er det virkelig sandt?“ Det bør Du troe, mit Barn! „Du troer det jo ogsaa?“ „Ja dette
bør vi Alle troe, thi det er jo Guds Ord!“ --- Saa ældes Manden, og bliver gammel, men
Evangeliet følger den Gamle, skjøndt det ikke kan ældes. Livet er et Kredsløb, siger
man, dette Udsagn har ogsaa en dybere Mening. Oldingen bliver et Barn igjen, og kommer
det Evige nærmere. Den ydre Virkelighed med dens Larm og Strid fjerner sig længere bort:
det bliver stille omkring ham. Men, idet det bliver stille som i Natten trindt omkring,
stige Synerne frem af Sjælens Grund: den indre Verden oplader sig for Øiet. „Er det
Haabet, der vinker fra et bedre Hjem? Er det en Barndomserindring, der hilser til
Afsked? Det er jo Evangeliets Engle, de prise Gud i det Høie, de forkynde Mennesket hans
% SHARED CLOSER: the single “ after „Jorden.“ closes both the inner quotation opened at
% „Er det Haabet“ three lines above AND the long quotation opened at „Et slaaende Beviis“
% on p. 140. The compositor sets no second mark, so that opener has no closer of its own
% and this fragment is one „ ahead on that account. As printed.
Velbehag; de lyse Fred over Jorden.“

Den æsthetiske Opfattelse kan paa mange Maader godtgjøre sin Berettigelse ligeoverfor
den blot historiske Forstandskritik; den kan i mange Vendinger bidrage til at hævde
% The point after „Reflerion“ is badly inked; at 1200 dpi it carries a dot above and a
% tail below, i.e.\ a semicolon. Second plain-r „Refler-“ on this page.
Følelsens og Indbildningskraftens Ret ligeoverfor den prosaisk-endelige Reflerion; den
kan, især naar den gaaer igjennem Speculationen, med stor Overlegenhed forsvare den
Hovedsætning, at Aabenbaringens Gud ligesaavel maa være Idealets som Virkelighedens,
ligesaavel Phantasiens som Tankens Gud; med andre Ord: den æsthetiske
% --- p. 143 ---
Opfattelse hævder den Fordring, at der, saa sandt der gives en aabenbaret Religion,
ogsaa maa gives en religiøs Phantasie-Anskuelse, for hvis Gehalt og Gyldighed
Sandhedens Gud selv vil indestaae. --- Hvorledes den æsthetiserende Speculation nu
iøvrigt bærer sig ad med en saadan Beviisførelse, vedkommer slet ikke denne Undersøgelse.
Den ovenfor antydede Fortolkning skal kun tjene til at vise, hvorvidt Evangelietroen kan
% „strax“ with a true x again on this page (plain descender at 1200 dpi), four pages
% after the „vorer“ of p. 141.
være tilfreds med en saadan Forklaring, eller ikke. Paa dette Spørgsmaal maae vi strax
% TURNED SORT: printed „Gruudsvaghed“ for „Grundsvaghed“ (u set for n), confirmed at
% 1200 dpi; both OCR witnesses read „Gruud-“. As printed.
uden videre give et benegtende Svar. Den æsthetisk-speculative Fremgangsmaade lider af
en Gruudsvaghed, der er saa meget farligere, som den med sine dulmende Omsvøb ellers let
kan bedaare den religiøse Følelse. Den naturlige Fromhed benytter kun altfor gjerne den
Lighed, den har med Troen, til at udgive sig for Troen selv, og sætte sig fast i dennes
% PRINTER'S ERROR: printed „det vælter“ where the sense wants „der vælter“ --- the sort
% is a plain t at 1800 dpi, and both OCR witnesses read „det“. As printed.
Sted. --- At forvandle enhver Engle-Aabenbarelse uden videre til et magisk Aande-Syn,
gaaer aldeles ikke an; thi da maatte jo t.\ Ex.\ den Engel, det vælter Stenen fra
Christi Grav, ogsaa være et saadant Aande-Syn. Begivenheden selv fortæller Matthæus
% BOLD FRAKTUR, not the \large display fount. Measured at 600 dpi against the body lines
% immediately above and below on this same page: x-height 40--43 px against the body's
% 37--39, but the median stem width goes 7 -> 8--9 px, a disproportionate thickening.
% The \large Gospel on p. 128 behaves quite differently (x-height 46--55 at stem 8), and
% the bold verses of Jes. 7, 14 and 9, 6 on p. 64 measure exactly as this one does. The
% citation „(K. 28, 2--6 V.)“ stands outside, in the ordinary weight.
(K.\ 28, 2---6 V.) saaledes: \textbf{Og see, der skete et stort Jordskjælv; thi Herrens
Engel nedfoer af Himmelen, traadte til, og væltede Stenen fra Indgangen, og satte sig
der. Men hans Skikkelse var ligesom Lynet, og hans Klædebon hvidt som Sneen. Men
Vægterne skjælvede af Frygt for ham, og bleve som døde. Men Engelen tiltalede
Qvinderne, og sagde: frygter I ikke! thi jeg veed, at I lede efter Jesum, den
korsfæstede. Han er ikke her; thi han er opstanden, som han haver sagt. Kommer hid, seer
Stedet, hvor Herren laae.} --- At her i denne Fortælling handles om noget Mere end et
reent Aande-Syn, forstaaes af sig selv. --- Vil nu den æsthetiserende Fortolkning
adskille baade „Vægternes“ og „Qvindernes“ Syn fra de i Jordskjælvet og Lynet vir\-%
% --- p. 144 ---
% „Textens“ has a true x --- the sort descends plainly at 1200 dpi.
kende Naturkræfter, da strider en saadan Adskillelse ikke blot imod Textens rene Ord; men
den hele Fortolkning priisgiver med det Samme sig selv til den Kritik, den just lægger an
paa at bekjæmpe. Den foreslaaede Adskillelse mellem Mythe og Aabenbaring er, naar den
besees lidt nærmere, selv mythisk: det Mythiske ligger i Tvetydigheden. Skulde nemlig
Aabenbaringens overnaturlige Gjennembrud i den Troendes Samvittighed lade sig forklare
af Ideens geniale Frembrud i Digterens Phantasie, da maatte disse jo ligesaavel være
overnaturlige, som hine. Forskjellen imellem de evangeliske og de digteriske
Aabenbarelser nedsættes herved til en blot Gradforskjel. Denne Gradforskjel være nu for
Resten saa stor, den vil: Begrebet om det Overnaturlige bliver dog alligevel slaaet
sammen med det Naturliges Begreb; og saaledes have vi da igjen den mythiske Opfattelse.
Skal Meningen derimod være denne, at de digteriske Aabenbarelser skee ifølge naturlige
Indflydelser, medens de evangeliske derimod foregaae ved et for os aldeles uforklarligt
Under: hvad opnaaer man da ved den paaberaabte Adskillelse imellem den indre og ydre
Virkelighed? --- Imod Adskillelsen selv have vi, forsaavidt den holder sig i reen
Almindelighed, Intet at indvende. Indvendingen gjælder kun det Øiemed, hvori den er
fremsat; thi den er øiensynlig lagt an paa, at holde Aabenbarelsen i en saa svævende
Form, som muligt, at Miraklet ikke skal vække Forargelse. Imod det Slags Forsøg maa
% PRINTER'S ERROR: printed „den ei en Livssag“ where the sense wants „den er en Livssag“
% --- the second sort is an i with its dot, confirmed at 1800 dpi; both OCR witnesses
% read i. As printed.
Evangelietroen af al Magt sætte sig til Modværge. Troen er i allerhøieste Forstand en
Livssag; men just fordi den ei en Livssag, just fordi Den, der vil leve i Troen, maa
sætte sit Haab til noget ganske Andet, end Den, der Intet har med Troen at gjøre, just
derfor maa det med uomtvistelig Bestemthed kunne vides, hvad det er, vi her have at
holde os til. --- Er det Høieste, hvorefter et Menneske bestemmer sin hele Livsretning,
Noget, der baade kan
% --- p. 145 ---
og skal begribes: da kan dette Høieste i alt Fald ikke være Gud (thi Gud selv kan jo
intet Menneske begribe); men, naar dette Høieste derimod er Noget, som hverken kan eller
skal begribes, saa røber det jo Mangel paa sundt Begreb, hvis Nogen spilder Tiden paa at
begribe Det, hvorefter han i ethvert Øieblik skal bestemme sit Liv. --- Ligesom det
overhovedet er en stor Anbefaling for en veldannet Mand, at han veed, naar han skal tie,
og naar han skal tale; saaledes bør jo Enhver i Forhold til Aabenbaringen fremfor Alt
lægge sig efter at vide, for Hvem og paa hvilket Sted han skal tie, og med Hvem og paa
hvilket Sted han har Lov at raisonnere. --- Da nu Troen paa det Allerstrengeste byder
Forstanden at tie stille for den Almægtiges Under, saa falder hiin sindrige Adskillelse
bort af sig selv. Om den Maade, hvorpaa Englene bleve istand til at aabenbare sig for
Hyrderne, veed den meest udviklede Mand ikke Mere, end det uskyldige Barn. Ikke desto
mindre er der dog en stor Forskjel imellem begges Viden. Barnet veed nemlig slet ikke
af, at her paa Begribningens Vegne er Noget i Veien. Dette veed Manden derimod meget
godt; men han veed ogsaa, hvorfor han hverken kan eller skal begribe. Er dette givet, da
bliver Engle-Aabenbarelsen ganske ligefrem at opfatte som umiddelbar Kjendsgjerning.

% NOT a larger fount, despite the look of the page: this paragraph is very loosely
% justified, but measured at 600 dpi against the two body lines directly above it, the
% x-height (38--44 px) and the median stem width (8 px) are identical, and the line
% pitch is the body's 123 px throughout. No \large, no \emph{}. --- „Vexelvirkning“ and
% „existentielle“ both carry a true x (plain descenders at 1200 dpi).
Evangelietroen er uendelig interesseret i de overnaturlige Begivenheders kritiske
Behandling. Kun forsaavidt den har forstaaet denne sin Interesse, har den forstaaet sig
selv. Paa denne Selvforstaaelse concentrerer den al Begriben. Troen vil --- vel at
mærke --- slet ikke begribe Aabenbaringens Vexelvirkning med Naturlovene og den
% „Lessing“ is set in the ordinary FRAKTUR, not antiqua --- verified at 1200 dpi, where
% the long-s pair and the Fraktur g are plain. Compare „Schleiermacher“ (p. 133) and
% „Strauß“ (p. 119). No \textit{}.
sandselige Virkelighed; men den vil forstaae sit eget existentielle Forhold til
Evangeliets Under. --- Det er ganske sandt, hvad Lessing har sagt, at „den Overgang,
hvorved man paa en historisk Efterretning vil bygge
% --- p. 146 ---
en evig Sandhed, er et Spring.“ Men denne Indvending kan slet ikke anfægte
Evangelietroen, der jo netop forudsætter, at Enhver, der kalder sig en Troende, maa have
gjort dette „Spring“. Lad os her lidt omhyggeligere undersøge, hvad Meningen vel kan
være med denne Tale. --- Ingen endelig Erfaring og ingen Fornuftvidenskab er istand til
at afnøde det tænkende Menneske den Tilstaaelse, at der nødvendigviis maa gives en
overnaturlig Aabenbaring. Naar nu Mennesket desuagtet siger: „jeg veed med urokkelig
Vished, at Gud har aabenbaret sig i Christo til Verdens Frelse, jeg veed det lige saa
vist, som jeg veed, at der er en Gud til, og at jeg selv lever i dette Nu“: saa gjør han
jo et Spring, idet han med Eet springer af fra alle vidtløftige Raisonnementer, og
griber den ubetingede Vished i Troen. --- Ingen historisk Forskning og intet
Fornuftbeviis formaaer at paatvinge det tænkende Menneske den Overbeviisning, at Guds
Aabenbaring just skulde være at finde i disse Evangelieberetninger. Naar nu Mennesket
desuagtet føler sig overbeviist om, at disse Beretninger med alle derhen hørende
Kjendsgjerninger ere det rette uforfalskede Udtryk for Aabenbaringens positive Indhold;
da kan han jo ikke have naaet denne Erkjendelse ad den historisk-kritiske Vei, men maa
% The dash closing this paragraph is only half inked --- about a third of the length of
% the em dash before „i Troen“ on the same line at 1200 dpi. Read as the book's em dash.
være kommen til den ved et umiddelbart Spring --- i Troen. ---

Endelig gjælder det her, at den Overgang, hvorved den Troende bryder alle
Sandsynlighedsberegninger af, og til Trods for Verdslighedens modstridende Analogier
finder Indgangen i de evangeliske Kjendsgjerningers overnaturlige Verden, ikke kan
adskilles fra den umiddelbare Overgang, hvorved han overhovedet kommer til Vished om, at
der gives en virkelig Aabenbaring: begge Erkjendelser opnaaes ved eet og samme „Spring“.

Ikke desto mindre gives der dog en Forstandskritik, der raisonnerer saaledes: „den
positive Aabenbaring og den subjective
% --- p. 147 ---
Tro forudsætte vistnok hinanden; men, for at de rettelig kunne enes til gjensidig
Understøttelse, maae de begge, hver for sig, underkastes en sindig Drøftelse, en
omhyggelig Prøvelse. Af en saadan sindig og videnskabelig Prøvelse vil den fuldstændige
og alsidig begrundede Evangelietro da først kunne resultere.“ --- Denne Sindighed kan
imidlertid let afstedkomme en Deel Forstyrrelse. Hvordan er det nu man pleier at bære
sig ad? --- Med høitidelig slæbende Halvhed, det vil sige, med en vis ærbar
Anerkjendelse af Oldtidens historiske Vidnesbyrd i Almindelighed, og en salvelsesfuld
% TURNED SORT: printed „Cursns“ for „Cursus“ (n set for u), confirmed at 900 dpi; both
% OCR witnesses read „Cursns“. As printed.
Ærefrygt for vore evangeliske Overleveringer i Særdeleshed begynder man det
historisk-kritiske Cursns. Men hvad bliver Resultatet? Kommer man virkelig ved
Anvendelse af den høiere og lavere, den indvortes og udvortes Kritik omsider til en
afgjørende Dom om de hellige Bøgers Authentie? „Ikke saa ganske!“ --- Faaer man maaskee,
ved at lægge den almindelige Indledning sammen med den specielle, tilsidst den lange
Undersøgelse over det Nye Testamentes Kanon afsluttet? „Ikke saa ganske!“ Seer man sig
da istand til at beseire alle Tvivl om de fire Evangeliers Ægthed og indbyrdes
Overeensstemmelse? „Ikke saa ganske!“ --- Men de overnaturlige Kjendsgjerninger, dem maa
man da vel see at sikkre mod ethvert videnskabeligt Angreb, som man sikkrer sin
Øiesteen, eftersom det jo dog fornemmelig er dem, til hvilke den subjective
Aabenbaringstro skal støtte sig, naar den lider Anfægtelse? „Nei, ikke saa ganske! Disse
Kjendsgjerninger have tvertimod de allersvageste Vidnesbyrd blandt alle. Det forstaaer
sig, Noget maa der naturligviis være seet, noget Overordentligt; men hvad dette
Overordentlige er, hvori det nu dog egentlig har bestaaet, derom tør den besindige
Kritik, reent ud sagt, endnu ikke have nogen bestemt Mening.“

Er det nu ved disse sindige Indrømmelser blevet vitterligt
% --- p. 148 ---
for Alle, at Evangeliernes Troværdighed staaer paa saa svage Fødder, saa maae vi da med
desto større Forventning see os om til den anden Side, om vi maaskee ved den „sindige
Drøftelse“ kunne opdage Troens faste, urokkelige Holdningspunkt i selve det religiøse
Gemyt. --- Hvorledes staaer det sig da med den subjective Tro? Faaer den Troende maaskee
nu og da saadanne indre Aabenbarelser og umiddelbare Syner, at han med Sikkerhed deraf
kan slutte sig til de evangeliske Kjendsgjerningers Virkelighed? „Ingenlunde!
Miraklernes Tid er længst forbi; lad os for al Ting see, at holde Sværmeriet ude.“ ---
Har den subjective Aabenbaringstro da en saadan Klarhed og Tilforladelighed i sig selv,
at den paa egen Haand kunde værge sig mod enhver Tvivl om det Overnaturliges Realitet,
% PRINTER'S DEFECT: the mark before „Nei“ is the HIGH, closing form “, set where the
% opening low „ belongs --- the answer „Nei, ingenlunde! ... Livet.“ is a fresh
% quotation, and no quotation is open at this point. Verified at 900 dpi against the
% unmistakable low „ of „Ingenlunde!“ six lines above. Transcribed as printed. The mark
% is thus BOTH a spurious closer here and a missing opener for the answer that follows,
% so it puts this fragment two “ ahead by itself.
endogsaa i det Tilfælde, at det virkelig skulde lykkes den negative Kritik, at opløse
Evangeliernes Indhold i phantastiske Myther og vilkaarlige Sagn?“ Nei, ingenlunde! Vor
religiøse Bevidsthed vil altid behøve et fast, opklarende Ord; den troende Anskuelse,
der trænger til Exempler og Forbilleder, maa idelig styrke sig ved at see hen til hans
Forbillede, der er Veien og Sandheden og Livet.“ --- Et meget opbyggeligt Svar! Men vil
dette nu med andre Ord sige, at den Troende blot behøver disse Exempler for
Anskuelighedens og Klarhedens Skyld; saa kunde man vel, uden at bekymre sig stort om den
historiske Virkelighed, gjerne forvandle den hele Evangeliebog til en folkelig
% Letterspaced, confirmed at 600 dpi: „Indledningen“ alone is spaced, and the rest of
% the line („at vor svage Tro behøver ...“) stands in the ordinary tracking.
Billedbog, og den religiøse Overbeviisning vilde endda være lige betrygget? „Ingenlunde!
Derfor bliver det jo udtrykkelig sagt i \emph{Indledningen}, at vor svage Tro behøver et
historisk Støttepunkt i den positive Aabenbaring.“ --- Sammenfatte vi nu, hvad der paa
begge Sider bringes ud af denne Halvkritiks sindige Drøftelse, da maae vi rigtignok
tilstaae, at den Troskab, hvormed den sub\-%
% --- p. 149 ---
jective Tro og den positive Aabenbaring i dette Tilfælde understøtte hinanden, ikke
just er saa ganske paalidelig.

Et Exempel: Et Menneske søger Vished i en Sag, der naturligviis er ham af
allerhøieste Vigtighed, efterdi det er en Livssag. Man raader ham at henvende sig til
A. A trøster ham paa det Bedste, og yttrer den Formodning, at han nok skal faae den
fuldkomne Vished, naar han kun vil raadføre sig med B. B giver ham ligeledes det
bedste Haab, naar han nu bare vil vende tilbage til A, og trænge ind paa ham, og nøde
ham til at erklære sig med et bestemt Ja. Men A forsikkrer høitidelig, at han paa
ingen Maade tør afgive en saadan Erklæring, før B har sagt Ja. Da begynder B at
sværge ved Alt hvad helligt er, at han for sin Samvittigheds Skyld ikke kan sige Ja,
før A rykker ud med Sproget. Imidlertid er Mennesket i megen Ængstelse og stor
Forlegenhed; han har kun disse Tvende at holde sig til, og Sagen er af yderste
Vigtighed, efterdi det er en Livssag. Da møder han til al Lykke en klog Mand paa
Veien, et kritisk Hoved, der giver ham følgende Raad: „Det er jo rigtignok sandt, at
Ingen af Vidnerne tør beedige Sagen med et fuldstændigt Ja; men Enhver af dem har dog
alligevel den bedste Formodning. En god Formodning er, naar vi ret overveie det, jo
dog omtrent det samme som et halvt Ja. Tvende Halve udgjøre en Heel, det er en
afgjort Sag. Vær du derfor ganske rolig: læg de to halve Ja'er sammen, saa har du den
fuldkomne Bekræftelse“. Maa nu Mennesket ikke være en saa besindig Mand høilig
forbunden for hans alsidige Drøftelse? --- I Sandhed, hvis Halvkritikens Besindighed
virkelig vidste, hvad den selv gjør, kunde man næsten fristes til at antage, at den
saa smaat lagde an paa at have baade Evangeliet og Troen til Bedste.

Hvilken Afstand imellem det blot Sandsynlige og det Sande!
% --- p. 150 ---
Da det altsaa er umuligt at tilveiebringe en heel Sandhed ved at sammenlægge to
Sandsynligheder, saa ligger det i Sagens Natur, at Evangelietroen kun ved et
afgjørende Spring kan opnaae den fulde Forvisning. --- For, saa vidt muligt, at
forebygge enhver Mistydning, ville vi betragte Sagen endnu fra en anden Side.
Meningen med det anførte „Spring“ er nemlig ikke den, at en Troende skulde paa een
% Printer's defect, p. 150: „kuune“ for „kunne“ --- a turned n, verified at 1200 dpi
% (the sort's connecting stroke lies at the foot, where the neighbouring n has it at
% the head). Transcribed as printed.
Gang og, saa at sige, med eet Sæt kuune sætte over fra Uvisheden til den i alle
% Printer's defect, p. 150: „fnldkomne“ for „fuldkomne“ --- the same fault the other
% way about, an n where the u belongs. Transcribed as printed.
Henseender fuldendte Forvisning og fnldkomne Tro. Tvertimod, den sande og levende Tro
lever kun i en fortsat Tilegnelse. Under denne Tilegnelse maa den Troende idelig
modarbeide Tvivlen, idet han under stadig Paavirkning af den Aand, der rører sig og
virker i Menighedens umiddelbare Vidnesbyrd (den apostoliske Troesbekjendelse) deels
„randsager Skrifterne“, deels prøver sig selv og sine indre Erfaringer som for Guds
Aasyn; forsaavidt befinder han sig da i en uafladelig Forsken og Stræben. „I skulle,
siger Christus (Joh. 8, 32) erkjende Sandheden, og Sandheden skal gjøre Eder frie.“
Men, forat denne Stræben efter Sandhedserkjendelse kan føre til Maalet, maa den
Troende med Bestemthed vide, hvad det er, der foraarsager Tvivlen og svækker
Overbeviisningen. --- Ifølge Kritikens Paastand, er det naturligviis den
ufuldstændige og lidet troværdige Aabenbaring, der maa bære al Skylden; men ifølge
den troende Opfattelse falder den hele Skyld paa Mennesket, der ikke veed at tilegne
sig den aabenbarede Sandhed. De evangeliske Kjendsgjerninger og Aandens Vidnesbyrd i
den vaagne Samvittighed kunne umulig være Tvivl underkastede, hvis den Troende ikke
skal have Lov til at gaae i Rette med Gud og modsige sig selv.

% Printer's defect, pp. 150--152: this long speech in the voice of Evangelietroen
% opens with „ here and is never closed. It plainly ends at „upartisk.“ on p. 152,
% where the compositor set a plain full stop and no “. Verified at 1200 dpi on both
% leaves. Transcribed as printed; the opener is left standing.
„Evangelietroen skyer ikke Lyset, den vil ikke krybe i Skjul for Videnskaben, den
skal nok holde ud med Halvkritiken, den er slet ikke bange for den moderne
Bevidsthed. Bliver den
% --- p. 151 ---
udfordret efter Ærens og Ridderlighedens Love, er den strax beredt at tage sin
Rustning paa og træde i Skranken med sin Modstander. Ingen Rigmand kan stride med
større Iver for at bevare sine Eiendomme for Tyve og Røvere, end den Troende strider
for at hævde de evangeliske Kjendsgjerninger imod Kritikens Angreb. Men er han først
bleven saa fortrolig med Aabenbaringens Indhold, at han lever og aander deri, saa
skal det ogsaa falde Kritikeren vanskeligt nok at fravriste ham endog blot en Tøddel
af de evangeliske Beretninger. --- Den negative Kritik kan paapege enkelte Træk af
Profanhistorien, der ikke synes at stemme med Evangelistens Udsagn; men hvad beviser
saa dette? Er den profane Historie maaskee useilbar? Er det afgjort, at den paaankede
Modsigelse er meer end tilsyneladende? Kan en Evangelist ikke have havt Grund til at
fremhæve en enkelt Begivenhed, der er forbigaaet af Historikeren, og omvendt? --- Den
negative Kritik kan fremdeles udpege en enkelt tilsyneladende Uovereensstemmelse
imellem Evangelierne indbyrdes. Men alle Beretningerne bære dog det samme Grundpræg.
Forskjelligheden lader sig forklare deraf, at Enhver af Vidnerne har seet Sagen med
sit Øie, og gjort sit Udvalg efter sit særegne Øiemed. Afvigelsen, hvor en saadan
virkelig findes, er desuden saa uvæsentlig, at den i ethvert Tilfælde kun minder os om
Fremstillingens reent menneskelige Side. --- Den negative Kritik kan endvidere gjøre
Vanskeligheder med Hensyn til de ældste Vidnesbyrd, der findes hos Fædrene, hos
Kjætterne og Christendommens Fiender, om vore kanoniske Evangelier. Men, naar disse
Vidnesbyrd veies paa Troens Vægtskaal, er det saa langt fra, at de bekræftende skulde
opveies af de benegtende, at disse endog indsvinde til et reent Intet i Sammenligning
med Kirkens store, eensstemmige Vidnesbyrd. --- Den negative Kritik kan endelig gjøre
den Indvending, at den Troendes
% --- p. 152 ---
% The quotation opened on p. 150 ends here, at „upartisk.“, with no closing “ --- see
% the note at its opening. Transcribed as printed.
Dom ikke er upartisk. Denne Indvending maae vi udtrykkelig fremhæve; thi den er fast
den eneste, som har nogen Vægt. Den Troende er, det maae vi tilstaae, i disse
Anliggender ingenlunde, hvad man ellers i historisk-kritisk Forstand kalder upartisk.
Men, før vi bryde Staven over Partiskheden, maae vi betænke Sagen lidt nøiere. Det
forholder sig ikke med Aabenbaringens, som med Verdenshistoriens Kjendsgjerninger.
Disse sidste kunne og maae bedømmes fra Upartiskhedens, det vil sige, fra
Interesseløshedens Standpunkt; thi, da deres Indhold er af endelig Natur, kan det i
høiere Forstand være Forskeren ligegyldigt, om de nu virkelig have tildraget sig
saaledes eller anderledes, naar han kun kan udfinde hvorledes. Med Aabenbaringens
Kjendsgjerninger er det en anden Sag; deres Indhold har uendelig Betydning: de gjøre
Fordring paa at troes til Salighed. Vil man nu kalde Den, der troer dem, partisk, da
maa man ogsaa kalde Den, der ikke troer dem, partisk. Vantroen er altsaa i dette
Punkt ligesaa partisk, som Troen. Men hvor Upartiskheden er en Umulighed, ophører
Partiskheden at være en Udyd; dette veed den resolute Fritænker lige saa godt, som
den oprigtige Troende. Kun Halvkritiken, der paa een Gang vil være baade i Troen og
% A reader's pen-stroke stands after „ikke.“ in the scan; it is manuscript, not type,
% and is ignored. The printed point is a full stop (so read by the ABBYY layer; the
% raised wedge tesseract takes for the upper dot of a colon belongs to the pen-stroke).
udenfor Troen, veed det ikke. Saaledes kommer den da til at staae ikke blot paa
Sindighedens og Alsidighedens og Upartiskhedens, men ogsaa paa Ligegyldighedens og
Middelmaadighedens og Tankeløshedens høiere Standpunkt. Den middelmaadige Halvkritik
er i visse Maader uovervindelig. Det er forgjæves her at anbefale den eenfoldige,
barnlige Tro: Middelmaadigheden paatager sig en videnskabelig Mine; den fordrer
Indsigt i de guddommelige Ting; den vil nu selv til at tænke. Saa kommer da
% Printer's defect, p. 152: „Speculatiouen“ for „Speculationen“ --- a turned n,
% unmistakable at 1200 dpi beside the true n two sorts later. Transcribed as printed.
Speculatiouen, og lægger sine tunge Tanker paa dens trinde Hoved; men neppe fornemmer
Middelmaadigheden, hvor det værker
% --- p. 153 ---
i Hjernen, før den sætter et eenfoldigt Ansigt op: „nei, intet Menneske er dog istand
til at begribe Guds Hemmeligheder;“ nu vil den da allerhelst holde sig til „den
positive Aabenbaring“, og saa til --- „den kritiske Indledning“. --- Imod to Magter
skal det imidlertid falde Middelmaadigheden vanskeligt nok at holde Stand, og disse
% Printer's defect, p. 153: „Frelsereus“ for „Frelserens“ --- a turned n, confirmed at
% 1200 dpi against the true n of „Evangelium“ on the same line. Transcribed as printed.
to Magter ere: Verdensaanden med de titaniske Naturkræfter og Frelsereus Evangelium
med de salige Engle.

Hvorvidt Middelveislæren, i det Hele taget, kan finde sig til Rette med den
almindelige Verdensaand, bliver et Spørgsmaal for sig. Her skulle vi især fremhæve
dette, at Halvkritiken ikke formaaer at tilegne sig Evangeliet. Fremfor alt maae vi
da vel vogte os for at forvexle den raisonnerende Middelmaadighed med den umiddelbare
Eenfoldighed. Forskjellen viser sig strax, naar Forholdet imellem det Ideale og det
Ufuldkomne, det Salige og det Virkelige nærmere skal bestemmes. Middelmaadigheden,
der klamrer sig fast til det Endelige, begrændser Idealet efter de naturlige
% Printer's defect, p. 153: „Eeufoldigheden“ for „Eenfoldigheden“ --- a turned n, and
% the word stands correctly three lines below. Transcribed as printed.
Betingelser, og opgiver det Salige, forsaavidt det kommer i Strid med det anerkjendte
Virkelige. Eeufoldigheden derimod seer op til Idealet, og glemmer Endeligheden;
tilegner sig det Salige, og takker Gud. Middelmaadigheden er selvklog, Eenfoldigheden
er ydmyg; Middelmaadigheden er gjenstridig, Eenfoldigheden er lydig;
Middelmaadigheden tvivler og raisonnerer, Eenfoldigheden tier og troer. Paa Grund af
denne Forskjellighed kan den ene ikke engang forstaae den anden. At denne Paastand
ikke lider af nogen vilkaarlig Overdrivelse, vil Evangeliet kunne bevise. --- Joseph
og Maria ere eenfoldige Sjæle; men Middelmaadighedens Tankegang kjende de ikke. Havde
de i de Dage, da den keiserlige Befaling udgik, slaaet ind paa den kritiske
Middelvei, da var Reisen fra Nazaret til Bethlehem tilvisse bleven dem begge en høist
besværlig Reise. Ja, hvis deres endelige Forstand
% --- p. 154 ---
var i dette Tidspunkt kommen til at beregne Modsigelsen imellem det Salige og det
Virkelige, var Striden i deres Indre vel bleven saa fortvivlet, at de sikkerlig aldrig
havde naaet --- „Krybben“, men vare omkomne paa Veien. Imidlertid kan det maaskee
falde Middelmaadighedskritiken vanskeligt nok at opdage, hvori denne Vanskelighed
bestaaer. Man vil vel snarere antage, at den Strid, hvis Mulighed her antydes, langt
bedre lader sig udlede af en dialektisk Grille, end af den evangeliske Fortælling.

Idet Halvkritiken besinder sig, og sammenholder de forskjellige Optrin i den hellige
Familiehistorie med hinanden indbyrdes: møder den en saadan Overflødighed af
Aabenbarelser, at den maa finde det meget tvivlsomt, at det Salige nu virkelig kan
have yttret sig i saa mange og saa umiddelbare Meddelelser. „Ikke nok, at Josephs
Drømmesyn staaer saa lidet motiveret ligeoverfor Bebudelsen til Maria; men der
indtræder fast ikke et eneste Vendepunkt, udenat der jo strax skeer et Under. Jesu
Fødsel forkyndes for Hyrderne ved et Under; Simeon møder i Templet som ved et Under;
ved et Under faaer Joseph den Paamindelse, at han skal flygte med Moder og Barn til
Ægypten, for at undgaae Herodes's Efterstræbelse; ved et Under bestemmes igjen
Familiens Tilbagevenden fra Ægypten (Matth. 2, 19); ved et nyt Under advares Joseph
mod Archelaus, og finder saa endelig Veien til Nazaret (Matth. 2, 23). Ja, naar man
overalt finder sig omgiven og veiledet af saa mange Undere, maa det just ikke være
saa farligt, at foretage endog en længere Reise. Det manglede blot endnu, at en Engel
i Drømme ogsaa skulde have forklaret Joseph, hvorfor den keiserlige Befaling absolut
burde adlydes. Istedetfor at gjøre Ophævelser over den korte Reise fra Nazaret til
Bethlehem, maa man snarere forundre sig over, at Joseph og Maria paa deres
forskjellige Reiser deelagtiggjøres i saa mange Aabenbarelser, at de endog synes at
glemme den ene over
% --- p. 155 ---
% Three scripture passages on this page are letterspaced (Sperrsatz), not bold and not
% of a larger fount: measured at 600 dpi the x-height and stem match the body exactly.
% „(Jesus)“ after the second one stands in ordinary tracking and is left outside the
% \emph{}; the letterspacing of the first begins at „Men“, after the colon.
den anden. Da (Luc. 2, 18) Alle forundrede sig over det, som Hyrderne sagde til dem,
tilføier Evangelisten (V. 19): \emph{Men Maria bevarede alle disse Ord, og overveiede
dem i sit Hjerte.} Desuagtet hedder det dog siden efter i Anledning af Simeons Ord
(Luc. 2, 33): \emph{Og Joseph og hans Moder forundrede sig over de Ting, som bleve
sagte om ham} (Jesus). Skulde denne Forundring, som Nogle mene, være at henføre
dertil, hvorledes den hellige Hemmelighed kunde være bleven Simeon aabenbaret: da
maatte Ordene jo ikke, saaledes som nu, vise hen til Simeons Tiltale, som Gjenstanden
for Forundring, men til den Tiltalendes Personlighed. Endelig, og det er endnu det
allerforunderligste, erfare vi (Luc. 2, 50), at Forældrene ikke engang forstode de
Ord, den tolvaarige Jesus talede til dem, der de havde fundet ham i Templet siddende
midt iblandt Lærerne: \emph{Hvi ledte I efter mig? vidste I ikke, at mig bør at være
i min Faders Gjerning?}“

Denne Opfattelse er nu slet ikke saa forunderlig, naar man kun vil stille sig paa
Middelmaadighedens Standpunkt, og betragte de anførte Optrin i behørig Afstand.
Betragteren veed i Forveien, at Maria, ifølge Beretningen, nu virkelig er bleven Guds
Søns Moder; deraf slutter han, at Joseph ligesaa let maatte kunne vide det.
Betragteren veed forud, at, om Barnet end maa lægges i Krybben, fordi der ikke var
Rum i Herberget, saa ville Hyrdernes Beretning om de himmelske Hærskarer snart trøste
baade Maria og Joseph over deres Fornedrelse; deraf slutter han da, at Joseph og
% Printer's defect, p. 155: „kunnet“ where the sense wants „kunne“. The sort is
% unambiguous at 1200 dpi. Transcribed as printed.
Maria, endog uden Hyrdernes Budskab, let maatte kunnet trøste sig selv. Kort sagt,
Betragteren er med sine Tanker altid et Skridt forud for Begivenheden: deraf slutter
han da, at Maria og Joseph ligeledes maae være det. Betragteren er saa fremsynet,
fordi han kan den hele Historie udenad; deraf formoder han, at Joseph
% --- p. 156 ---
og Maria maae have været ligesaa fremsynede; og glemmer, at de just Skridt for Skridt
maatte opleve „det Utrolige“. --- At dette nu er en bagvendt Opfattelse, forstaaes vel
af sig selv; men her er endnu en anden Modsigelse, hvori Middelmaadighedskritiken gjør
sig skyldig. Den antager nemlig uden videre, at en Aabenbarelse er en Aabenbarelse, og
at saa Sagen dermed kan være afgjort. „Har man først faaet en virkelig Aabenbarelse,
og dermed en tydelig Besked om Guds udtrykkelige Villie, da veed man ogsaa, hvad man i
Hovedsagen har at rette sig efter: alt det Øvrige kan Gud da gjerne overlade til
Menneskets egen Conduite.“ --- Ja, naar man først har faaet en virkelig Aabenbarelse!
Men hvoraf kommer det da, at Middelmaadigheden selv har saa vanskeligt ved at overtyde
sig om, at de her omhandlede Aabenbarelser ere virkelige?

Middelmaadighedskritiken tænker nu som saa: „Havde jeg blot været i Josephs eller
Marias Sted, jeg skulde nok have overbeviist mig om, hvorvidt den første foregivne
Aabenbarelse var sand eller falsk; og, naar jeg da ved en sindig Drøftelse var bleven
enig med mig selv, skulde jeg ogsaa have vidst, at tage mine Forholdsregler
derefter.“ --- Et saadant Raisonnement beviser just, hvor lidt Middelmaadigheden er
% Letterspaced, confirmed at 1200 dpi: the run opens at „at“ and closes with the full
% stop after „Naturlige“, „Na-“ / „turlige“ being broken across the line.
istand til at sætte sig enten i Josephs eller Marias Sted. Vanskeligheden ligger just
i, \emph{at det Overnaturlige er ueensartet med det Naturlige.} Denne Vanskelighed kan
% Printer's defect, p. 156: „Meuneske“ for „Menneske“ --- a turned n standing directly
% beside a true n; the two joins are opposite at 1200 dpi. Transcribed as printed.
umulig ophæves derved, at et Meuneske selv oplever Aabenbarelsen. Tvertimod, et
saadant Menneske maa netop føle, hvor stærk Modsætningen imellem det Salige og det
Virkelige tynger og trykker med al sin Haardhed. Lad os blot et Øieblik forsøge at
lægge Middelmaadighedens egen Reflexion ind i Joseph og Maria.

Nogle Maaneder før Keiser Augusts Befaling skulde efter\-%
% --- p. 157 ---
kommes, have altsaa de Trolovede, hver for sig, i en særegen Aabenbarelse erfaret den
guddommelige Raadslutning. Den Tilstand, hvori Maria nu virkelig befinder sig, synes
% Printer's defect, p. 157: Maria's speech opens with „ here and is never closed ---
% it plainly ends at „Skikkelse?“, where the leaf carries the em-dash and no “.
% Verified at 1200 dpi. Transcribed as printed.
at bekræfte Aabenbarelsens Sandhed; da reiser Tvivlen sig i hendes Sjæl. --- „Det var
et saligt Øieblik: Engelen stod for mig, og hilste mig som den Benaadede. Men er jeg
da virkelig den Benaadede? Kan Joseph overbevise sig herom? Hvis Verden anede min
Hemmelighed, den vilde staae op imod mig, at myrde mig for min Gudsbespottelse, eller
at haane mig for min Afsindighed. Det var et saligt Syn; men Synet er jo svundet:
denne Tilstand er usalig! Er det i Virkeligheden saaledes at være den Benaadede? Jeg
er mig ingen Brøde bevidst; men er det Hele ikke som en Trolddom? Kan ikke ogsaa
Mørkets Aand iføre sig en Lysets Engels Skikkelse? --- Vi behøve ikke at gaae videre,
for at indsee, hvad Middelmaadighedsreflexionen vilde virke under saadanne Vilkaar:
en eneste Tvivl, og Maria er fortabt --- i Vanvid. Det Samme gjælder om Joseph. Ogsaa
i hans Indre maa det Salige stride en forfærdelig Strid med det Virkelige, hvis den
raisonnerende Kritik finder Indgang. --- „En salig Drøm, den løste mig den mørke
Gaade! Engelen aandede sin Fred dybt i min Sjæl: saa kunde jeg igjen forstaae Maria.
--- Turde jeg nu stole paa en Drøm, da var min Lod den skjønneste paa Jorden. Som
Den, der er af Davids Huus og Slægt, skulde jeg da staae ved Messias's Vugge. Under
% Printer's defect, p. 157: „opvore“ for „opvoxe“ --- a plain r where the sense wants
% x. The r sits flat on the baseline; the true x on this same page (in
% „Middelmaadighedsreflexionen“, nine lines above) drops a tail below it. This is the
% book's habitual fault; transcribed as printed.
mit faderlige Tilsyn skulde da Guds Søn, Immanuel, opvore og forfremmes i Alder og i
Viisdom. Men det er jo umuligt. --- Kun en Drøm! Hvad er jeg i den virkelige Verden?
En ringe Haandværksmand i Nazaret; saa dybt er Davids Slægt da sunken. --- Urimeligt!
Hvad er jeg i mit Forhold til Maria? Som en bedragen Mand. Hvad vilde Menneskene
dømme, hvis jeg stod frem, og talte til dem om min
% --- p. 158 ---
Gaade? De vilde lee, og have Ret. --- Saa er jeg da en Daare! Den virkelige Verden
seer paa mig, og gjækker mig med sine kloge Øine: den spotter mig, den leer jo ad min
Drøm.“ ---

% Printer's defect, p. 158: „sindige Drøftelse opens with „ and is never closed --- a
% comma follows „Drøftelse“ and no “. Verified at 1200 dpi. Transcribed as printed.
Det gaaer ikke saa let med den „sindige Drøftelse, som Halvkritiken synes at mene; jo
længere man drøfter, desto tungere bliver Forstaaelsen. --- Hvorledes er Modtagelsen i
Bethlehem? Maria trænger til et Herberge; hendes Time er kommen. Men der er ikke Rum i
Herberget! hun finder kun en liden Plads ved Krybben. --- „Skal dette være Opfyldelsen
af Prophetens Spaadom (Micha 5, 1)? Kan dette kaldes Tidens Fylde? Det synes dog i
Virkeligheden at være en saare ubeleilig Tid. At kastes saaledes hen, som En, der er
given i Tilfældets Vold, see, det skal være Tegnet for den Benaadede!“ --- Nei, med
Middelmaadighedsreflexionen er det her aldeles umuligt at komme af Stedet. Modsigelsen
imellem det Salige og det Virkelige er altfor skrigende. --- Men i den rene
Eenfoldighed er Modsigelsen overvunden; thi den eenfoldige Tro holder Tvivlen ude.
Tømmermanden Joseph er visselig ikke noget kritisk Hoved. Halvkritiken maa nok smile
over hans Ubetydelighed. Thi --- ikke sandt? --- det er dog en høist eenfoldig Mand,
denne Joseph? eenfoldig, at han troer paa en Drøm; eenfoldig, at han igjen antager sig
den forladte Jomfru; eenfoldig, at han saa ubetinget adlyder Keiserens Befaling;
eenfoldig, at han i et saa kritisk Tidspunkt fører Maria til Bethlehem, uden i
Forveien at have bestilt Herberge. --- Men er det nu sandt, at Troens Eenfoldighed
alene formaaer at ene det Salige med det Virkelige, saa maae vi jo prise Joseph og
Maria salige, at de vare eenfoldige Sjæle: Held dem, at de ikke kjendte
Middelmaadighedens Tankegang.

Forholder det sig nu saaledes med Tilegnelsen af de overnaturlige Kjendsgjerninger,
da vil en Troende ogsaa forstaae,
% --- p. 159 ---
% Three turned sorts on this page --- „persoulige“, „Tiugenes“, „gjøreude“ --- all
% verified at 1200 dpi against true n's in the same lines, and all transcribed as
% printed. This is the heavy cluster of turned n's that runs from p. 150 to p. 160.
hvad de mange forskjellige Aabenbarelser i den hellige Familiekreds vel kunne betyde.
--- Under Verdslighedens og den blot naturlige Fromheds almindelige Forudsætninger er
Menneskelivet i en vis Forstand langt simplere anlagt, end det derimod kan være det
efter Troens Maalestok. Naturmennesket veed, hvad det har at holde sig til,
forsaavidt det nemlig indseer, at det i denne sandselige Verden maa bruge sine Sandser
og sin Forstand, og finde sig til Rette med Omstændighederne. Her er kun een
Grundlov, Verdensloven: her er kun een Verden. Anderledes med Livet i Troen.
Formedelst det persoulige Gudsforhold kommer Mennesket i Forhold til en overnaturlig
Tiugenes Orden. Den Troende maa altsaa paa een Gang leve i tvende Verdener, den ideale
og den virkelige, den salige og den forkrænkelige. For en saadan Grændsebeboer maa
Tilværelsen nødvendigviis være langt mere forviklet, end for Den, der har opslaaet sin
Bolig midt i Naturlighedens velbebyggede, velbetryggede Egne. Troens Dobbeltverden er
overalt gjennemkrydset af modsatte Veie, modsatte Bevægelser, modsatte Love. Idet den
Troende vil underordne sig det Evige, maa han ligesom udfordre Endeligheden; thi
Evighedens Lov skal fuldbyrdes i det Endelige. Idet han vil adlyde Gud, maa han
tillige vogte paa sin Forstand; thi Daarskab og Vanvid behage ikke Gud. Paa en saadan
Livstraad danner sig da en Mangfoldighed af Knudepunkter, der ikke kunne løses uden
Guds særegne Bistand. Jo mere afgjøreude Mennesket handler, desto mere ubetinget maa
han henstille Sagen til Gud, og jo mere ubetinget han overgiver sin Sag i Guds Haand,
desto stærkere er han igjen opfordret til at anstrenge sine Evner og Kræfter. Den
Troende er en Høitbetroet, der skal staae til Ansvar. --- Kan dette nu gjælde om
Troeslivet i Almindelighed, hvor meget mere maa det da ikke kunne gjælde om den
hellige Families Liv i Særdeleshed. See, Joseph og Maria
% --- p. 160 ---
% Printer's defect, p. 160: the first line is weakly inked --- the final e of
% „Høitbetroede“ prints only as a raised tick and the e of „dem“ is half wanting.
% Both OCR witnesses read an apostrophe for the first. The sorts are the right ones;
% read as e.
ere Høitbetroede: Guds egen Søn er dem betroet. Middelmaadighedskritiken behøver ikke
at frygte for, at de mange Aabenbarelser skulle virke bedøvende paa deres Sind: de
trænge ved ethvert afgjørende Skridt til Guds underfulde Bistand. Det guddommelige
Barn er jo selv et Under, en overordentlig Gave fra Himmelen: hvad Under da, at
Himmelen holder Øie med det, at Englene vaage over det, at Guds Aand drager de
Fromme, og samler dem omkring dette Barn, at de maae glæde sig i Tilbedelsen og
aflægge Vidnesbyrdet. --- Forsaavidt Josephs og Marias Forundring forudsætter en
saadan almindelig Forstaaelse, er denne Forundring nærmest at opfatte som en inderlig
salig Beundring, der vidner om, at disse eenfoldige Sjæle med enhver ny Aabenbarelse
modtage et nyt, oprindeligt Indtryk af det Guddommelige. Men paa den anden Side
kommer dog ethvert Gjennembrud tillige som noget Uventet og Uforudseet.
Sandsynlighedsberegningen er afbrudt: Ingen veed, hvad Gud i næste Øieblik vil gjøre;
Ingen begriber, hvad der nu i denne virkelige Verden skal skee med dette Barn,
hvorledes det skal udvikle sig, hvilke Stadier det skal gjennemgaae. Just af den
Grund maa hvert enkelt Optrin med dets forskjellige Biomstændigheder faae en uendelig
Betydning for Joseph og Maria. I denne Henseende kan deres Forundring betragtes som en
religiøs Studsen, en hellig Overraskelse, der vækker den dybe, halvforstaaede Anelse.
--- Det kan dog umulig gaae til med den inderlige Forstaaelse af Guds Søns Indtræden i
Verden, som med den endelige Forstaaelse af en sædvanlig Begivenhed.
Middelmaadighedskritiken gjør Maria Uret, naar den forestiller sig Sagen saaledes:
„Det skulde have været mig, der af Bebudelsens Engel forud havde faaet den Forklaring,
at Barnet virkelig var Guds eenbaarne Søn; jeg skulde saavist ikke være falden i
synderlig Forundring over Hyrdernes Beretning, men have tænkt i mit stille
% --- p. 161 ---
Sind: ja, det forstaaer sig, Englene maae naturligviis love og prise min Søn, thi
denne min Søn er jo Messias selv, den Høiestes Søn. Men lad det nu ogsaa være, hvad
det er, med den første Overraskelse og den første Forundring: over Simeons Ord i
Templet var der da nu aldeles ingen Grund til at blive forundret. Alt, hvad den fromme
Olding i spaadomsfuld Begeistring taler, idet han holder Barnet paa sine Arme, maatte
Maria og Joseph jo kunne saa godt som udenad i Forveien. Istedetfor at lade sig
velsigne af Simeon, burde Maria snarere selv have velsignet den gamle Mand, fordi han
talte saa klart og saa correct. Jo længere Historien skred frem, desto sikkrere maatte
man jo blive i sin Sag. Da derfor Jesus i Templet svarede den beængstede Moder: „Mig
bør at være i min Faders Gjerning“; burde Maria aabenbart, ifølge det hidtil
Passerede, have vexlet et betydningsfuldt Øiekast med Joseph, istedetfor at de nu
begge staae der, som et Par --- sandselige Mennesker (1 Cor. 2, 14), der ikke fatte de
Ting, som høre Guds Rige til.“

Hvis Middelmaadigheden imidlertid ikke med slet saa megen Sikkerhed paatog sig at
dømme om Ting, den ikke forstaaer, vilde en saadan Misforstaaelse aldrig kunne
opkomme. Joseph og Maria kunde i al deres Eenfoldighed jo vel forstaae, at det
guddommelige Barn maatte være Messias, Guds egen Søn; men de kunde ogsaa forstaae, at
dette Barn nu var betroet til deres Omsorg, at de derfor ikke blot skulde glæde sig
over denne vidunderlige Guds Gave, men ogsaa lære at fatte Guds Villie og Raad. Derfor
kan her ikke være Tale om overfladisk Sikkerhed og vanemæssig Tryghed. Maria, der seer
paa sin Søn, seer ind i et uudgrundeligt Dyb af Viisdom og Naade; derfor er hun saa
aarvaagen, og dog saa inderlig stille: „hun bevarer alle disse Ord, og overveier dem i
sit Hjerte.“

% p. 162 opens a new paragraph with the printed indent („Men, vedbliver
% Middelveislæren, om Joseph og Maria med samt de Øvrige...“), so the blank line above
% carries the paragraph break across the fragment joint. Do not close it up.
% --- p. 162 ---
Men, vedbliver Middelveislæren, om Joseph og Maria med samt de Øvrige i den
hellige Familiekreds nu ogsaa paa denne Maade have kunnet forbinde det Salige
med det Virkelige i deres fromme Eenfoldighed: saa er det dog alligevel ikke den
Side, hvorfra man maa opfatte Sagen, hvis Evangeliet skal trænge igjennem i
Nutiden. Det Salige vil ikke længere nærme sig Bevidstheden i saa naive
Tilsyneladelser. Ligesom det enkelte Menneske umulig kan blive Barn igjen, om det
end nok saa gjerne vilde, naar det har naaet den modnere Alder; saaledes er det
ogsaa umuligt for Menneskeheden, paa dens nærværende Culturtrin, at skrue sig
selv tilbage paa hiint Oldtidens barnlige Standpunkt. Hvoraf kommer det, at den
Dannede, selv med den bedste Villie, ikke kan troe disse Barndomsmirakler, uden
deraf, at Tidens mægtige Aand stemmer imod? Nutiden er desto værre snarere vantro
end troende. At byde en vantro Tidsalder en saa pointeret Fremstilling af
mirakuløse Optrin, er kun at forøge Vantroen og gjøre Ondt værre. Vil man altsaa
tage Sagen fra den rette praktiske Side: da lade man de evangeliske Billeder
gjøre deres æsthetiske Indtryk, og see for Resten, at fremhæve den
religiøs-sædelige Substans. Menneskeslægten kan ikke frelse sig selv; Frelseren
blev født til Verden, at det kunde blive vitterligt for Alle, at Gud ikke
forlader Menneskene, om de end i deres Forvildelse stundom synes at forlade ham:
„\emph{Ære være Gud i det Høie!}“ Vel er der Strid i Verden, og jo friere
Kræfterne virke, desto større bliver Anstrengelsen, desto heftigere Striden. Men
stride vi kun efter bedste Evne for Sandhed og Ret, da kunne vi rolig see op til
ham, der fødtes til at stride den store Strid, og i vor egen Samvittighed høre
den himmelske Røst: „\emph{Fred paa Jorden!}“ Vel er det Onde i Verden:
Ugudeligheden kan ikke behage Gud. Men hverken den hemmelige Skyld eller den
aabenbare Forbrydelse lader sig udrydde
% --- p. 163 ---
ved Mirakler og Mirakeltro. Lader os derfor med mandigt Mod holde ud i Frihedens
og Sædelighedens ærefulde Kamp; vi kunne da trøste os med, at Gud ved Christus
har aabenbaret sig, ikke som Vredens, men som Kjærlighedens Gud:
„\emph{I Menneskene en Velbehagelighed?}“
% Printer's defect: the closing mark of the third clause of the Gloria is a
% QUESTION MARK, verified at 600 dpi (curved head, point below). The two
% parallel clauses on p. 162 („Ære være Gud i det Høie!“, „Fred paa Jorden!“)
% and the full verse at pp. 128--129 all print „!“. Transcribed as printed.
% The Fraktur I/J is one sort; the sense („i Menneskene“ at the head of the
% quotation) settles it as I.

Om Middelveislæren forøvrigt ved slige opbyggelige Trøstegrunde og moralske
Omtydninger seer sig istand til at stille den urolige Verden tilfreds, lades her
uafgjort; kun dette bemærke vi, at slige Trøstegrunde ogsaa kunne haves uden
særligt Hensyn til Evangeliet. Vi indrømme gjerne, at Middelmaadigheden er en
Magt i Verden; vi ere heller ikke uskjønsomme imod den oekonomiske
Forstandighed, der saa smukt forstaaer at give det Timelige en passende Bismag
af det Evige, det Profane en sømmelig Tilsætning af det Hellige, og derved seer
sig istand til at uddele til de Trængende, om ikke just Livets Brød, saa dog
nogle brugbare Leveregler. Ja, Middelmaadigheden er en Magt i Verden: naar
Horizonten mørknes, og Tidens Vande røres, da er det gjerne Middelmaadigheden
selv, der i egen Person griber efter Roret, og tager Commandoordet. Det gaaer
saa an i uroligt Veir, naar Stormen dog ikke er altfor stærk. Men bryder Orkanen
løs, bliver Menneskeheden opbragt, fordi den føler sig skuffet og bedragen for
det Høieste; da kunne Lidenskaberne let bruse op, og slaae Middelmaadigheden
omkuld, og slaae over i Fortvivlelse.

--- Det er dog en underlig Tale, naar det hedder, at den dannede Verden umulig
kan vende tilbage til hiint Barnlighedens og Eenfoldighedens Trin, paa hvilket
Sjælene ret tør glæde sig ved Evangeliet. Som om der for os Mennesker kunde
gives andet Forhold til Gud, end det barnlige Eenfoldighedsforhold. --- Den
raisonnerende Halvkritik synes her at glemme, hvad den dog saa ofte selv siger,
og hvad den gjennemførte Speculation tilstrækkelig
% --- p. 164 ---
bekræfter, at Viisdommens Dybder ere uudgrundelige. Men om nu ogsaa de Vise
begyndte at hovmode sig, og de Skriftkloge at forhærde sig: skulde disse Tidens
Tegn da virkelig være et Beviis for, at Menneskeheden ikke mere kan vende
tilbage til hiint Barnlighedens og Eenfoldighedens Trin, hvor den kan finde sin
Trøst og Glæde i det Saliges Aabenbarelse? Er da maaskee hvert andet Menneske
nuomstunder blevet en Tænker, en Kritiker eller en Skriftklog? Vare der maaskee
paa Keiser Augusts og Koug
% Printer's defect: „Koug“ for „Kong“ --- a turned n, verified at 1200 dpi (two
% stems with square foot-serifs and no top diagonal, unlike the n of „ingen“ on
% the same line). Transcribed as printed.
Herodes's Tid ingen selvkloge Vise og ingen forhærdede Skriftkloge?

Men er det saaledes, at de ædleste og bedst begavede Mennesker altid maae glæde
sig ved at forstaae det Salige med Troen i al Eenfoldighed: skulde da
Evangeliet ikke endnu, som forhen, kunne prædikes for de Fattige. Den Engel, der
vidner for Hyrderne om Frelserens Fødsel, siger jo ikke: jeg forkynder Eder en
Glæde, som dog kun kan vederfares de Vise og Skriftkloge; men han siger
udtrykkelig: \emph{see, jeg forkynder Eder en stor Glæde, som skal vederfares alt
Folket.}---Det er dog en besynderlig Tale, naar det hedder, at Evangelietroen
skulde kræve for Meget af Nutidens Mennesker, at Tilegnelsen af de aabenbarede
Kjendsgjerninger nu skulde overgaae menneskelige Kræfter. Hvor Meget blev der da
krævet og fordret af hine Tiders Fromme? Aabenbaringen kom som en Gave fra
Himmelen! Der meldes Intet om, at Hyrderne havde arbeidet for nogen Opvækkelse,
eller deeltaget i nogen storartet Befrielseskamp: de besørgede deres ringe
Gjerning, de laae paa Marken og holdt Nattevagt over deres Hjord. Men, da
Herrens Engel stod for dem, og Herrens Klarhed skinnede om dem, da det glade
Budskab forkyndtes dem; da gave de sig ikke til at kritisere og raisonnere, men
de troede Aabenbarelsen, og sagde til hverandre: \emph{lader os dog}
% --- p. 165 ---
\emph{gaae til Bethlehem, og see det, som Herren haver kundgjort os.} Og de bleve
jo heller ikke skuffede; det var en sanddru Engel, der havde bragt dem
Evangeliet; \emph{og de kom hastelig, og fandt baade Maria og Joseph, og Barnet
liggende i Krybben.} --- Er det da at forlange for Meget af Menneskene, naar der
kun kræves, at de skulle troe Aabenbaringens Ord, og søge Frelsen, hvor den er at
finde? Kan den kloge Forstandskritik da ikke forstaae, at den sande Eenfoldighed
ligesaa vel maa være den menneskelige Udviklings Fuldendelse som dens
Begyndelse?--- Ville de Vise og Skriftkloge ikke indsee dette, men misbruge den
aandige Dannelse, og sætte deres egne Vilkaarligheder i de evige Sandheders Sted:
ligemeget! saa vil Tilværelsen selv med Omstændighedernes Magt vel lære Folket at
see, hvor Faren er, og hos hvem Frelsen er at finde. Vel sandt, man bliver just
ikke til noget Stort i Verden ved at antage Christi Evangelium. Hyrderne fik vel
heller ingen Forfremmelse i deres ydre Stilling; men --- de saae dog Barnet i
Krybben, og de kundgjorde, hvad der var talet til dem om dette Barn. Og der var
jo talet til dem om, at dette Evangelium skulde bringe Fred paa Jorden. Det er en
sanddru Tale: naar alt Folket først kunde fatte og tilegne sig den store Glæde,
som skal vederfares det formedelst Evangeliet, da vilde Broder forsone sig med
Broder, da vilde Menneskene glemme deres Kiv og Strid, da vilde der være „Fred
paa Jorden“. Og der blev talet til Hyrderne om, at Gud
% Printer's spacing defect at the foot of the line: the word-space between „at“
% and „Gud“ is omitted and a space stands inside „Gu d“ instead (600 dpi). The
% next line is not letterspaced, so this is loose setting, not emphasis.
for Frelserens Skyld fandt Velbehag i Menneskene. Det er en sanddru Tale: kunde
alt Folket kun fatte og tilegne sig den store Glæde, som skal vederfares det
formedelst Evangeliet: da vilde Hjerterne bøie sig i Lydighed under den hellige
Lov;
% Rettelser: printed „Loø“; corrected to „Lov“
% (Side 165, Linie 3 f. n. --- confirmed on the image: the third line from the
% foot reads „under den hellige Lo?“, the last sort a heavily inked ø with its
% diagonal visible at 600 dpi. Compare the correctly set „Guds Lov“ on the very
% next line. Cf. the „Løu“ for „Lov“ at p. 141.)
da vilde Frihedens Tider oprinde, thi Ingen, der lyder Guds Lov, kan
være Vilkaarligheders Træl: da skulde Sandhedens Aand
% --- p. 166 ---
vidne, at der var „i Mennesker en Velbehagelighed“. --- At det nu holder saa
haardt, inden dette Evangelium kan finde Indgang hos alt Folket, at det maa gaae
gjennem Trængsel og Nød, at Barnligheden først gjenvindes i Tugtelsen, og
Eenfoldigheden kjøbes med Smerte, derpaa tænker Kundgjørelsens Engel slet ikke.
De Salige ændse ikke de kjødelige Hjerters Gjenstridighed, det er for dem, som
forstod det sig af sig selv, at en falden Slægt ret gjerne maa ville lade sig
frelse. De himmelske Hærskarer, der lyse for Hyrderne i Nattens Fred, see slet
ikke paa Verdens Strid Forvildelse og Forhærdelse; men de see paa Frelsens,
Forsoningens og Salighedens Aabenbarelse i Verden; de see paa Menneskehedens
evige Væsen, paa det Guds Billede, hvori vi alle ere skabte, paa den Guds
Eenbaarne, der hviler i Krybben, og ved hvem dette Væsensbillede igjen skal vinde
Skikkelse. Derfor glæde de sig paa Menneskenes Vegne, derfor lyder deres Lovsang
i saligt Chor: {\large Ære være Gud i det Høieste, og Fred paa Jorden, og i
Mennesker en Velbehagelighed.}
% The Gloria here is set in the \large display fount of this edition: at 600 dpi
% the x-height of „nne“ in „Mennesker“ runs ~45 px against ~37 px for the „nne“
% of „Menneskenes“ two lines above, at the same stem weight. Not bold Fraktur.
% It begins mid-line after „Chor:“ and ends mid-line, so it is set inline rather
% than as a \begingroup\large block.

% The verse is set as an indented block in a fount a size SMALLER than the body,
% not letterspaced and not bold --- the same five closing lines as the quotation
% at p. 129, but with „Smerte,“ for „Smerte.“ and no comma after „Chor“.
\begingroup\small
\begin{verse}
„Den Stund, da Frelseren fødtes paa Jord,\\
Fødtes en Trøst for den bittreste Smerte,\\
En Trøst for den, der haaber og troer,\\
En Gjenlyd af de Saliges Chor\\
Stiger da ned i Menneskets Hjerte.“
\end{verse}
\endgroup

%% ============================================================
%%  SUBSECTION HEAD, printed part-way down p. 166 (PDF 183).
%%
%%  Set as the heads of subsections 1 (p. 100) and 2 (p. 128): the numeral „3.“
%%  alone on its own line, then the title in the larger display Fraktur. NO RULE
%%  anywhere in the head region --- swept at 600 dpi over the whole block, and
%%  there is no long horizontal run of ink at all.
%%
%%  UNLIKE the two earlier subsection heads, this one carries NO SCRIPTURE
%%  REFERENCE LINE. Subsection 1 prints „(Luk. 1, 39--80.)“ parenthesised;
%%  subsection 2 prints „Evangelium Luk. 2 K. 1--20 V“ unparenthesised. Here
%%  the body text („For at kunne opfatte nogen...“) follows the title directly.
%%
%%  The title is NOT letterspaced --- the display fount is simply larger.
%%
%%  INDHOLD vs. BODY: the Indhold prints „3. De evangeliske Forvarsler ......
%%  166--183“, without a point. The printed head has one: „De evangeliske
%%  Forvarsler.“ Both readings stand; the head follows the page, the TOC entry
%%  follows the Indhold. The disagreement is confined to the full stop, unlike
%%  the heads of subsections 1 and 2, which also differ internally.
%% ============================================================
\subsection*{3.\\[4pt] De evangeliske Forvarsler.}
\addcontentsline{toc}{subsection}{3. De evangeliske Forvarsler}

For at kunne opfatte nogen, det være nu hvilkensomhelst, af de Omvæltninger, der
foregaae deels i Naturens deels i Aandens Verden, maa man især see at rette sin
Opmærksomhed paa de indledende Bevægelser, paa Forvarslerne og de første
Frembrud, d. e. man maa fremfor Alt see at opleve Begyndelsen.

% --- p. 167 ---
Indtræffer en Omvæltning derimod saa uforudseet, at den forrige Tingenes Orden
alt er kuldkastet, og alle bundne Elementer løsladte, før Nogen aner det; da
bliver man let overvældet, og er ikke istand til at fatte Noget, efterdi man selv
taber al Fatning. Det hele Optrin staaer da for En, som et uforklarligt og
forfærdeligt Særsyn, indtil man omsider naaer saa vidt, at man igjen kan gaae
tilbage, og besinde sig paa Begyndelsen.

At Spørgsmaalet om Begyndelsen altid maa udøve en egen Tiltrækning paa det
contemplative Gemyt, er ligefrem grundet i Sagens Natur. Den umiddelbare
Begyndelse er som et springende Punkt, der danner Grændsen imellem det, der
forgaaer, og det, der opstaaer; et underligt tvetydigt Noget, der rører sig i
Overgangen fra det Mulige til det Virkelige.

Forend vi imidlertid gjøre nogen nærmere Anvendelse heraf paa den evangeliske
Historie, ville vi først forvisse os om Begyndelsens fleersidige Betydning ved at
vælge et Par Exempler, af Naturen og af Verdenshistorien.
% „Exempler“ with a true x --- the descender hook is plain at 1200 dpi, unlike
% the r of „Par“ on the same line. So too „t. Ex.“ two lines below.

Hvilket Naturoptrin er vel mere skikket til at vække det Aandige i Mennesket, end
Uveiret med Lyn og Torden? Og dog, hvis et saadant Uveir pludselig (som t. Ex. i
en Bjergegn) overfalder Vandreren med al sin Heftighed, bliver der neppe Tid til
beskuende Nydelse. Det blændende Lyn, det bedøvende Bulder, og da især den
nedskyllende Regn gjør jo En aldeles uskikket til ophøiede Betragtninger og
aandrige Iagttagelser. Skulde den Reisende maaskee tage Mod til sig, og trodse de
fire Elementer, for ret at nyde det store Skuespil, og tabe sig i Beundringen af
Naturaandens Majestæt? Det gaaer vist ikke an. Majestæten er i Oprør: han taaler
ingen Betragtning, og vil slet ikke vide af nogen Beundring at sige. Betragte
Ham, den Frygtelige? En anden Gang! Nu er han vred, hans Pande mørk: han truer.
Beundre Ham?
% --- p. 168 ---
Han er saa heed, saa vild, hans Øine brænde: hvor det dog flammer i de sorte
Bryn! --- Og disse tunge Drøn, denne vilde Latter! ---
% Printer's defect: the em dash after „Latter!“ prints broken and under-inked
% --- at 1200 dpi only fragments survive, one of them solid, all at em-dash
% height rather than on the baseline. Read as the dash the parallel
% constructions on this page and on p. 169 („Rasen. --- Hvad“) require.
Fly Vandrer, om du kan, til Bjergets Hule! afsted i Skjul! thi hvis han mærker
dig, da er du Dødsens.

Vælge vi derimod det Tidspunkt da Uveiret drager op, og nærmer sig sin
Begyndelse: saa faaer Naturbegivenheden et ganske andet Udseende, saa bliver her
Noget at betragte. En forunderlig Ro, en næsten beængstende Stilhed er udbredt
overalt. Luften er saa trykkende, Solstraalen stikker, alle Kræfter slappes: det
er, som fik Naturen Qvalme, og vilde til at besvime. Denne Tilstand er i
Virkeligheden høist ubehagelig: man døier, og venter paa en Forandring. Da viser
sig en lille Sky i Horizontens Rand; deu voxer, den formerer sig lidt efter lidt,
% Printer's defect: „deu“ for „den“ --- a turned n, verified at 1200 dpi against
% the n of „Rand“ on the same line. Transcribed as printed. „voxer“ in the same
% breath carries a true x (descender plain at 1200 dpi).
den taarner sig, endogsaa lige imod Vinden: det er et Forvarsel. Snart er
Himmelen overtrukken med Skymasser i mangfoldige Lag af de allerforskjelligste
Former og Farver. En Uerfaren, der ikke vidste, hvad alt Dette skulde betyde,
kunde let falde paa den Tanke, at det Hele maaskee var et Slags Gjøgleri, at
Naturen under sit Ildebefindende var falden i Slummer, og nu drømte om --- Drager
og Kjæmper, Patriarcher og Engle, Oversvømmelser og Ildebrand, og foreviste sine
Drømmebilleder i Luften. Men vi vide jo, at det dog alligevel ikke er noget
Gjøglespil; vi vide, at disse underlige Skymasser ere Værksteder for sunde og
stærke Kræfter. Der arbeides derinde, skjøndt alt det Ydre tyder hen paa
vilkaarlige Indfald; der arbeides derinde: dæmrende Blus, hurtig forsvindende
Strimer skimtes igjennem de mørke Lag. Der arbeides derinde: Uveirets Aander,
som boe i de luftige Telte, færdes i Travlhed frem og tilbage med Ild og Lys, som
skulde de tænde Lynet. Og dog er Alt endnu i tvetydig Ro, i ængstende Stilhed:
det Hele er endnu kun en Mulighed. Forvarslerne have viist
% --- p. 169 ---
sig: Forberedelserne ere til Ende. Da skeer det i et Nu; som paa et Vink fra En
deroppe, der fører an, skeer det i et Nu: Lynet blinker, og Luften zittrer, og
Vinden rører sig; og Regnskyen luder, og sænker sig dybt, og dens skrøbelige Bund
brister, og store, tunge Draaber falde paa Træernes Blade, paa den støvede Jord.
Det er et vederqvægende Øieblik, vi opleve Begyndelsen.

Og nu i Historien, hvilken Omvæltning kan vel her være mere storartet og for
Betragteren mere tiltrækkende, end den Omvæltning, vi hos os (betegnende uok!)
% Printer's defect: „uok“ for „nok“ --- a turned n, verified at 1200 dpi against
% the n's of „betegnende“ in the same parenthesis. Transcribed as printed.
betegne med det fremmede Ord: en Revolution? Men ogsaa her gjælder det, at Den,
der vil opfatte det aandige Indtryk, ikke maa styrte sig ubesindig i
Begivenhedernes Malstrøm. Naturkræfterne ere frygtelige, naar de komme i Oprør;
men --- det være sagt baade til Ære og Skam for Menneskeheden --- naar
Lidenskaberne blusse op, og tænde Statsbygningen i Brand: da er al Elementernes
Strid kuu som en uskyldig Spøg imod et ophidset Folks vanvittige Rasen.
% Printer's defect: „kuu“ for „kun“ --- a turned n, verified at 1200 dpi.
--- Hvad Historien angaaer, forstaaer det sig vel af sig selv, at Betragteren,
for at finde det rette Udgangspunkt, maa gaae tilbage til Begyndelsen. Men det
gjælder nu ogsaa om at \emph{opleve} Begyndelsen; thi den Strid, Muligheden her
paa de første Overgangstrin fører med den givne Virkelighed, er uendelig
betydningsfuld. Lad os, for at oplyse dette nærmere, give de tvende stridende
Magter, Muligheden og Virkeligheden, en bestemtere Form, at vi desto bedre kunne
opfatte Indtrykket.

Den givne Virkelighed er da for det Første et forældet Statslegeme, der hviler
sig paa sin historiske Grund, og desaarsag gjør Fordring paa at ansees for
Retsideens naturlige Skabning, for Retfærdighedens udtrykte Billede. Men det er
desværre en sørgelig Selvmodsigelse; den virkelige Ret er paa alle Punkter
forvandlet til væsentlig Uret. Fyrsten paa sin Throne kalder sig Statens Hoved og
% --- p. 170 ---
Folkets Fader. Han har virkelig Ret; han hævder det Nedarvede, han gjør det
Samme, hans Fædre gjorde før ham. Og dog har han alligevel væsentlig Uret. Som
ved en syvdobbelt Ringmuur af Privilegier og vilkaarlige Begunstigelser er
Statens Hoved adskilt fra sit Legeme, og Folkefaderen, der holdes som indesluttet
i sin Borg, omringet af de Begunstigede, der kalde sig Thronens Støtter, og vogte
Palladsets Indgange, kan, om han end nok saa gjerne vilde, dog ikke høre Folkets
Stemme, der kalder som fra det Fjerne. Hver Gang nu de Store benytte deres
Forret, gjøre de under disse Vilkaar væsentlig Uret; og, naar de Undertrykte
prøve paa at tage sig selv til Rette, begaae de en virkelig Uret. Saa trykkende,
saa utaalelig kan Tilstanden være. Og dog er det Hele endnu, som var det i sin
Orden: Eensformighedens Ro og Sædvanens Stilhed er udbredt overalt. Men det er en
ubyggelig Ro, en ængstende Stilhed, thi det ulmer i Grunden, og Modsigelsen
arbeider i det Skjulte. Da komme Forvarslerne, det ene efter det andet i al
Tvetydighed. --- Man byder Loven Trods. Er det fortvivlet Nødværge eller
Raahedens Anmasselse? „Siger Intet. Kun en Opbrusen af urolige Hoveder, løse,
magtesløse Angreb paa det Bestaaende“. --- Men --- der danne sig nye Talemaader,
nye Slagord: underlige Tankelyn glimte i Masserne. See, Folket rører sig som En,
der plages af urolige Drømme, og anstrenger sig, for at vaagne. --- „Skulde der
virkelig være Fare paa Færde. Lad da Øvrigheden vise sig med Trusler og Løfter,
og minde Undersaatterne om deres Troskabsed og deres Pligter.“ Det hjælper! Hvor
Øvrigheden træder frem, viger Folket tilbage: Individerne føle dog en hemmelig
Skye, en uvilkaarlig, dunkel Ærefrygt for den høiere Magt. --- „Ja, sagde jeg det
ikke nok: det Hele er i Grunden et Gjøglespil, det velsindede Folk har endnu ikke
glemt Lydigheden.
% --- p. 171 ---
Man skal kun træde dristig op, og imponere det med Loven; saa hjælper det strax:
nu kommer det snart til Ro.“ --- Dog nei, det hjælper ikke! Det brænder i
Grunden, Modsigelsen arbeider i det Skjulte. Urolige Skarer bevæge sig hid og
did. Den uhyre Folkemasse er som et bølgende Hav; man hører den fjerne Brusen:
det tyder paa Storm og Uveir. „Skulde det blive til Alvor, da maa Allerhøieste
selv sætte sig i Bevægelse: der maa skee noget Overordentligt!“ --- Saa stiger da
Fyrsten ned fra sin Throne, træder ud af Palladsets Dør, og taler trøstende,
advarende og formanende Ord til sit kjære Folk, som en Fader til sine Børn. ---
„See, det hjælper. Den vrimlende Mængde udbryder i Jubel. Hvor det er en rørende
Scene! Det Hele er vist et gemytligt Sværmeri. Det troskyldige Folk har længtes
efter at see Landets Fader. Nu har det seet hans Aasyn, og hørt hans Løfter: nu
har det ingen Fare, nu gaaer det nok over.“ --- Men det brænder i Grunden,
Modsigelsen arbeider i det Skjulte. Samfundets Piller ryste. Hvert Stød bringer
nye Forviklinger, nye offentlige Ulykker. Det er som Udbrud af en underjordisk
Ild, hvilken Ingen kan slukke. „Ville Underverdenens Magter røre sig, da er det
paa Tide, at adspørge Oraklet: man maa dog høre den offentlige Stemme!“ --- Saa
træde da Folkets Talsmænd op. I lange Rækker stille de sig ligeoverfor de
Mægtige: Tvetydigheden stiger til det Høieste. Undres ikke paa, at disse Folkets
Mænd ere saa tankefulde, saa alvorlige, saa ordknappe: de staae jo paa Grændsen
af tvende Verdener, Mulighedens og Virkelighedens, Fremtidens og Fortidens
Verden. Undres ikke paa, at disse fremmede Skikkelser ere saa stive og kolde, som
vare de støbte i Malm: Retfærdighedens Sværd er over deres Isse, og Sværdet er i
Uro, thi --- Mulighed strider mod Virkelighed. Undres ikke paa, at disse nye
Talsmænd ere saa be\-%
% --- p. 172 ---
vægede, at deres Tunger gløde, at de hviske imellem hinanden, lig Sammensvorne,
der have viet sig til en uforglemmelig Daad. I ville forstaae dem, naar I kun
erindre, hvorfra, og hvorfor de ere komne; naar I see, hvad det er, her foregaaer
trindt omkring. --- „En talløs Vrimmel bedækker Omkredsen i vilde Grupper:
phantastiske Væsner, uhyggelige Gestalter, fordreiede Ansigter, som Ingen før har
seet ved Dagens Lys! Det tager sig jo ud som et Carneval, en gyselig
Folkeforlystelse, et Maskespil.“ --- Nei, det er intet Maskespil og intet
Carneval; men det er den arbeidende Mulighed, der nu nærmer sig sin virkelige
Begyndelse. --- Man dvæler, og venter, og aander saa tungt. --- „Hvad venter man
paa?“ --- Et afgjørende Ord (thi Ordet er og bliver dog altid Menneskehedens
Løsen.) Hør, hvor der strides derinde om Ret og Uret! --- Man dvæler, og venter,
og aander saa tuugt.
% Printer's defect: „tuugt“ for „tungt“ --- a turned n, verified at 1200 dpi.
% The same phrase four lines above prints „tungt“ correctly. Transcribed as
% printed.
Da skeer det i et Nu, som paa et Vink fra En dernede, der fører an, skeer det i
et Nu: et Ord! et Raab! et skrigende Raab! det er \emph{Frihed og Frelse!} Og
Trylleriet svinder, og Folkestrømmen svulmer, og bryder alle Dæmninger, og
bortskyller Thronen, og begraver en Fortids-Verden med al dens Jammer og al dens
Herlighed i sit Svælg: her har Betragteren fundet det rette Øieblik, nu opleve vi
Begyndelsen.

Men, befinde vi os ret, har da ikke Jesu Komme til Verden frembragt en
Forandring, en aandelig Omvæltning, som aldrig før eller siden er seet? Er der
noget Vendepunkt i Historien, om hvilket det med Sandhed kan siges, at en hidtil
ukjendt Mulighed brød igjennem, og virkeliggjorde sig i en ny Begyndelse, da er
det vel det Vendepunkt, Apostelen kalder \emph{Tidens Fylde.}

Man dømme nu iøvrigt om Forløsningsværket, som man vil: det maa dog i Sandhed
have været et stort Øieblik, da disse Christus-Kræfter rørte sig for første Gang,
og dette Værk tog sin
% --- p. 173 ---
virkelige Begyndelse. Det maa dog have været et stort Øieblik, da Røsten
begyndte at raabe i Ørkenen: \emph{Bereder Herrens Vei, gjører hans Stier
jævne. Hver Dal skal opfyldes, og hvert Bjerg og Høi skal nedtrykkes, og det
Krumme skal gjøres lige, og de ujævne Veie skulle slettes; og alt Kjød skal see
Guds Frelse.} Det maa dog i Sandhed have været et stort Øieblik, da Han, der
vidste sig som Guds eenbaarne Søn, traadte op som \emph{Menneskens Søn} i
Menneskenes Forsamling, og gjentog Forgjængerens strenge Ord i Kjærlighedens
milde Formaning: \emph{Omvender Eder, thi Himmeriges Rige er nær.}

% The Isaiah/Luke verses and „Omvender Eder …“ are LETTERSPACED, not bold: at
% 600 dpi the strokes match the body exactly and only the fit is opened. The
% Thomas verse below is the other device --- bold Fraktur, body x-height with a
% visibly heavier stem --- so both stand on this one page and must not be
% conflated.
Ere vi nu ogsaa paa nogen Maade istand til at gjenkalde os disse Begivenheders
oprindelige Indtryk, at opleve denne Begyndelse? --- Der er en Paamindelse, vi
først maae høre, og det er den Opstandnes Paamindelse til den tvivlende
Discipel: \textbf{Fordi du saae mig, Thomas, troede du; men salige ere de, som
ikke have seet, og dog troe.} „Hemmelighedsfuld, som Hans Bortgang fra Verden,
er ogsaa Hans Komme til Verden“, siger Kritiken, og den har Ret; men
Hemmeligheden er --- aabenbaret de Troende.

Er Nogen opsat paa at \emph{begribe} denne Begyndelses Hemmelighed, da kan han
frit beskylde Evangelisterne for aandelig Indskrænkethed, indbyrdes
Uovereensstemmelse og Modsigelse; de svare ham ikke, de forsvare sig ikke, de
sidde som adspredte, hver for sig, saa stumme, som Billeder af Steen. Sige vi
derimod: Vi forlange ikke at \emph{begribe} (dette er os her baade for Meget og
for Lidet); men vor Attraa er, at \emph{opleve} Begyndelsen (som jo de Troende i
Tilegnelsens Inderlighed gjerne ville opleve hele Frelserens Liv); da rører det
sig i de fire Vidner, da er det, som nærmede Evangelisterne sig til os og derved
til hinanden indbyrdes
% --- p. 174 ---
hver med sit Evangelium, og pegede paa det Første, og sagde: \emph{Kommer og
seer!}

Er der noget Punkt i den hellige Historie, om hvilket man med Sandhed kan sige,
at det er alsidig opfattet, da er det vel Begyndelsen.

Vende vi os til Johannes, Evangelisten Johannes, --- Kirken seer ham forklaret i
et Billede: Ørnen ved hans Fod, Himlen i hans Øie ---, da tyder han os jo den
evige, den oprindelige Begyndelse, Begyndelsen herovenfra. \emph{I Begyndelsen
var Ordet, og Ordet var hos Gud, og Ordet var Gud.} Men, naar han nu dernæst,
som svævede han hen over Endelighedens Mellemled, tilføier: \emph{Og Ordet blev
Kjød og boede iblandt os (og vi saae hans Herlighed, en Herlighed, som den
Eenbaarnes af Faderen), fuld af Naade og Sandhed;} da smelter det Evige i det
Timelige, og Lyset blænder vort Øie. Føle vi nu, at vi ere paa Jorden, og mangle
„Ørnens Vinger“ og Synets Klarhed, og trænge til den endelige Forestilling; da
tilbyde Matthæus og Lukas deres Veiledning. De berette os da, hvorledes det gik
til, at den Eenbaarne, som var i Faderens Skjød, sænkede sig i Menneskehedens
Skjød, og blev holdt for en Søn af Joseph og Maria: vi møde hos dem den hellige
Familie i For-Evangeliets mærkeligste Optrin. Medens Johannes da især dvæler ved
den \emph{overhistoriske} Begyndelse, meddele Matthæus og Lukas derimod de
\emph{for-historiske} Begivenheder, og nu forstaae vi tillige Markus, der
umiddelbart begynder med det offentlige Liv, og derved strax frembringer det
stærke \emph{historiske} Indtryk.
% „stærke“ and „Indtryk“ are NOT letterspaced --- only „historiske“; verified at
% 600 dpi.

Ikke destomindre have Theologien og Kritiken dog smukt forenet sig om at
misforstaae denne Begyndelse, og det i dobbelt Henseende. De mange mislykkede
Forsøg, paa at udfinde Forbindelsen
% --- p. 175 ---
imellem For-Evangeliet og den virkelige Historie, vise det tydeligt nok. Først
antager man, at de for-historiske Beretninger maae i alle Maader stilles paa
lige Trin med de historiske. Men, naar man saa (hvad der for Resten giver sig af
sig selv) bliver opmærksom paa, at det dog giver en Forskjel, om det er en
hemmelighedsfuld Bebudelse af en overordentlig Mands Fødsel, eller denne Mands
egen offentlige Fremtræden der beskrives, at man t. Ex. ikke kan tale om Engelen
Gabriel i samme Sprog, som om en Johannes den Døber, saa mener man at have gjort
en stor kritisk Opdagelse. Istedetfor at forstaaeliggjøre sig Aabenbaringens
ideale Kjendsgjerninger, og besinde sig paa, hvad det dog egentlig er,
Evangelietroen kræver, farer man nu lige over i den modsatte Yderlighed, slaaer
en Streg over det Guddommelige, og reducerer saa det hele For-Evangelium til et
mythologisk Forspil. --- Det gaaer nu saa temmelig let, ikke just fordi det er
„Videnskaben selv“, der gjør det, og Alt, hvad „Videnskaben“ godtgjør,
naturligviis maa være godt gjort, men især, fordi den hele Sag, hvorom
Forhandlingerne dreie sig, jo ikke er Andet end en eenfoldig Troessag. Det
skulde blot være en anden og betydeligere Sag, t. Ex. en almindelig Æressag, en
Nationalitetssag, eller en høist alvorlig Pengesag, hvis Opgjørelse var
afhængig af disse Begivenheders rette Fortolkning: vi skulde da snart see den
sunde Sands reise sig og gribe „Videnskaben“ i Brystet, og ryste den lange
„kritiske Indledning“ saa længe, indtil Theologien omsider lærte at indsee, hvad
en overordentlig Begyndelse har at betyde. Men --- „en Troessag? --- Ikke
Andet!“

% No extra leading here: the baseline pitch across this break was measured at
% 600 dpi (118 px against a page mean of 125) --- an ordinary paragraph, not a
% section break, although both OCR witnesses insert a blank line.
I Theologiens kritiske Misforstaaelser have Evangelisterne imidlertid ingen
Skyld. Evangelisterne ere Aandens Mænd, og Aandens Mænd maae jo dog vide, at den
overordentlige Begyndelse kun opleves og fremstilles, idet man begynder med at
% --- p. 176 ---
see Mulighedens Forvarsler bag Virkelighedens Dække, som bag et halvt
sønderrevet Forhæng.

Lader os altsaa for nogle Øieblikke sige Farvel til den moderne Bevidsthed, og
vandre i Apostlenes Spor tilbage til Palæstina, til det hellige Land, til det
Tidspunkt, da Jesu offentlige Liv tog sin virkelige Begyndelse. See,
Evangelisten Lukas vil jo selv lede os paa Veien, idet han fortæller Følgende:

% Luk. 3, 1--3 stands at full measure in the larger display fount with the
% ordinary paragraph indent, like the pericopes of pp. 1--2 and Joh. 1, 1--5 on
% p. 80; hence \begingroup\large … \par\endgroup.
\begingroup\large
Men i Keiser Tiberii femtende Regjeringsaar, der Pontius Pilatus var
Landshøvding i Judæa, og Herodes var Fierdings-Fyrste i Galilæa, men hans Broder
Philippus var Fierdings-Fyrste i Ituræa og Trachonitis Land, og Lysanias
Fierdings-Fyrste i Abilene; der Annas og Caiphas vare Ypperste-Præster: skete
Guds Ord til Johannes, Zacharias' Søn i Ørken. Og han kom i den hele Egn
omkring Jordan og prædikede Omvendelsens Daab til Syndernes Forladelse
(Luk. 3, 1--3).
\par\endgroup

Lige før Begyndelsen indtræder, opfordres vi altsaa til at kaste et Blik paa
Jødefolkets daværende Tilstand, en sørgelig, en høist beklagelig Tilstand. Deelt
og udstykket imellem Fjerdingsfyrsterne og den romerske Landshøvding, er det
lille Rige, der igjennem saamange Aarhundreder har stridt for sin
Selvstændighed, og igjennem saamange Farer og Ødelæggelser er paa vidunderlige
Maader frelst og igjenopreist ved Forsynets Haand, dog nu alligevel yderlig nær
ved sin Afgrund. Alt tyder hen paa en forestaaende Opløsning. --- Den gamle Lov
bestaaer vel endnu: Herodes har smykket Jehovas Tempel. Men Aanden er vegen
bort, og Kraften er svunden. Hoffærdige Pharisæer, med lange Bønner og brede
Bræmme, sidde paa Mose Stol, og ivre for den fædrene Tro, og lære Folket de
mange Bud. Det er et Skyggespil: de mange Bud ere Menneskebud, og de tomme Ord
kalde ikke det
% Printer's defect: the second „mange“ in „de mange Bud ere Menneskebud“ is set
% with a TURNED u for the n (verified at 1200 dpi against the n of „Menneskebud“
% on the same line). The word reads „mange“ and is transcribed so.
% --- p. 177 ---
Døde tillive. Den kolde Sadducæer foragter Religionens Trøst og Mængdens Tro,
han har Nok i sin egen Dyd: ingen Overtro, ingen Fordom skal udelukke ham fra
Nydelsen af Culturens Glæder. Men Folket selv, hvorledes ere dets Kaar? Det
lider Vold, og øver Vold: dets Hjerte bløder. Uden Lærere og uden Ledere, er den
ustadige Mængde, \emph{som Faar paa Bjergene, der ikke have Hyrde.} Det er en
sørgelig, en fast utaalelig Tilstand. Men hvad skal der da skee med dette Folk?
Ja --- „hvad vil der dog skee?“ Dette Spørgsmaal vanker som en urolig Tanke
omkring i Jerusalem og det ganske Land, ængster de Alvorlige og underholder de
Letsindige. Hvad vil der dog skee; thi det gjælder om Frihed og Frelse?

Vilde Fædrenes Gud nu blot som Fædrenes Sønner! --- „da skulde Frihedshelten
komme som Lynet, thi et Oprør ulmer i Grunden.“

Vilde Fædrenes Gud nu blot som Fædrenes Sønner! --- „da skulde Zion hæve sig
over alle Bjerge, og et Rige reise sig, som aldrig før blev seet paa Jorden.“

Vilde Fædrenes Gud dog blot som Fædrenes Sønner! --- „da kom Messias selv før
Hanegal, i næste Nattevagt: det er jo den yderste Tid, det gjælder om Frihed og
Frelse.“ --- Han maa vel komme. „Ja, han skal komme! Om han kommer fraoven, som
\emph{Menneskesønnen i Himmelens Skyer}, eller han skyder sig op af \emph{Jesai
Rod}, som en \emph{Qvist paa den gamle Stamme}: ligemeget, naar Han kun kommer.
Ligemeget, naar Han kun kan \emph{tale til Jordens Konger i sin Vrede, og
forfærde Hedningene i sin Grumhed.} Ligemeget, naar Han kun \emph{slaaer} os
disse Fiender \emph{med sit Jernspiir, og sønderbryder dem som en Pottemagers
Kar!}“
% The letterspacing here follows the SOURCE, not the sentence: the psalm and
% prophet words are spaced and Nielsen's own connectives („som en“, „Han kun
% kan“, „os disse Fiender“) are not. Verified word by word at 1200 dpi; the
% breaks are real and are reproduced.

Og see! Han er kommen. --- Fædrenes Gud har holdt sine
% --- p. 178 ---
Løfter; men Sønnerne forstaae det ikke. Befrieren er kommen: men de vide det
ikke. \emph{Han staaer midt iblandt dem; men --- de kjende ham ikke.}

Men hvad vil der dog skee, nu det gjælder om Frihed og Frelse?

Det Forfærdelige vil skee. Dette Folk vil korsfæste Messias. Dette Folk vil i
Lidenskabernes Oprør forsøge Befrielsen ved egen Kraft; det vil ende med
Fortvivlelse: Jerusalems Ødelæggelse vil blive Historien et Rædselsminde, til
hvilket Menneskene knytte Forestillingen om unævnelig Jammer.

„Men --- før dette skeer --- vil Underet skee! Før Ødelæggelsen skeer, vil
Herlighedens Herre aabenbare sig midt i sit Folk. Før dette gamle Tempel, som er
bygget med Hænder, nedbrydes, opreiser Han det nye Tempel, som ikke er bygget
med Hænder. Men --- det Gamle nedbrydes: Forhænget brister, Offer-Alteret
brænder, Zions Mure falde: det er Lovens Opfyldelse, det er Verden et Varsel, at
Guds Rige er kommet med Frihed og Frelse.“

Staae vi nu her fast paa Virkelighedens Grund, og see med Eftertidens Øie
tilbage paa den første Begyndelse, paa de underfulde Varsler, hvorved det
betegnes de Troende, at Himmeriges Rige nærmer sig: da faae ogsaa vi et
umiddelbart Indtryk af For-Evangeliets Betydning i den hellige Historie.

Jødefolkets daværende Tilstand have vi betragtet. Denne sørgelige Virkelighed er
nu, saa at sige, det Forhæng, igjennem hvis Rifter Evangelisterne Lukas og
Matthæus lade os see enkelte Aabenbarelser af Salighedens stille Rige; ligesom
man ved Nattetider stundom seer --- en enkelt Stjerne gjennem Drivskyens Rifter.

% --- p. 179 ---
% The refrain „Varslerne komme og svinde …“ is BOLD FRAKTUR, not the larger
% display fount: at 600 dpi its x-height matches the body and only the stem is
% heavier. It recurs on pp. 179 (twice), 180, 181 and 182, and is bold every
% time. The visionary paragraphs it introduces are LETTERSPACED instead, so the
% two devices alternate down these four pages.
\textbf{Varslerne komme og svinde; thi Himmeriges Rige er nær.}

\emph{I Jerusalem, i Templets Helligdom staaer en gammel Præst foran Røgelsens
Alter. Han er saa underlig tilmode, saa dybt bevæget, --- alene med Gud i det
Hellige. Da kommer Synet, et guddommeligt Syn! Det forfærder ham, Frygt falder
over ham: see, Herrens Engel staaer ved den høire Side af Røgelsens Alter! Det
er Forvarslet for den Høiestes Prophet, den anden Elias.} --- Men Folkets
Mangfoldighed, der venter udenfor i Røgelsens Time, bliver urolig, og begynder
at spørge sig selv, hvi dog Præsten tøver derinde. „Der skulde dog aldrig være
skeet ham et Vanheld? --- Den gamle Mand skulde dog vel ikke være falden i
Afmagt? --- Nei, nu træder han ud af Helligdommen. Men hvilken besynderlig Mine!
Han giver et Vink: han er stum! --- Der maa være hændet ham noget
Overordentligt, han maa have seet et Syn i Templet.“ Saa taler Mængden, og
skilles ad, og --- glemmer Hændelsen.

\textbf{Varslerne komme og svinde; thi Himmeriges Rige er nær.}

\emph{Er der endnu et Sted paa Jorden, hvorhen Forbandelsen ikke har naaet? See,
det er Hende, den rene Jomfru med det rene Hjerte! Og Engelen kommer ind til
hende og hilser hende; og de tale sammen, de Tvende, som Sjæl til Sjæl,
beaandede af Gud. --- Det seer jo ud, som i Paradisets Egne, i Edens Have, før
Cheruben drog sit Sværd!} --- Men det maa dog nok være en Phantasi; vi ere jo
ikke i Paradis. Dette er Nazaret, Flækken Nazaret ved Tiberias' Sø.
% --- p. 180 ---
Her kan vist ikke være skeet noget Overordentligt; ikke et Menneske i Byen har
erfaret det Allermindste. --- „Skulde det maaskee være det Rygte, der gaaer om
Hende, som er trolovet med Tømmermanden Joseph? Ganske rigtigt! Man har en
Mistanke. Dog, det jævner sig vel nok.“ Og see! det jævner sig virkelig. Joseph
annammer Maria, sin trolovede Hustru, den borgerlige Sædelighed tilfredsstilles;
nei see dog, hvor det jævnede sig! „Slige Tilfælde bør man ikke regne saa nøie.“

\textbf{Varslerne komme og svinde; thi Himmeriges Rige er nær.}

Det er en stor Glæde, der vederfares Zacharias's Huus i Judæ Stad. Naboer og
Slægtninge høre, at Herren har gjort sin Barmhjertighed stor imod Elisabeth; og
de forsamle sig, og glæde sig med hende. --- „Men det er dog besynderligt, at
Barnet ikke skal kaldes med sin Faders Navn!“ --- Zacharias skriver paa sin
Tavle: \emph{Hans Navn er Johannes;} --- og de forundre sig Alle. See, da
kommer Varslet: \emph{Den Stummes Mund oplades, hans Tunges Baand løses, han
taler og priser Gud. Saa bliver det da omsider vitterligt for Alle, at her er
skeet et virkeligt Under. Og det rygtes paa alle Judæas Bjerge, og Alle, som det
høre, lægge det paa deres Hjerte.} --- Ja, det rygtes! --- „Den gamle Mand var
som forklaret i Glæden: han saae paa sin lille Søn, og var saa glad. Skulde
Glæden ikke kunne løse den bundne Tunge? Skulde det dog ikke være gaaet naturlig
til, at Faderkjærlighedens Glæde oplod den Stummes Mund?“ --- Dog, Verden faaer
Andet at tænke paa.

Der er udgaaet en Befaling fra Keiser Augustus til Landsherren i Syrien, at hele
Landet skal beskrives til Skat, og hver Mand møde i sin egen Stad at lade sig
beskrive. --- I Bethlehem,
% --- p. 181 ---
Davids Stad, er Tilstrømningen saa stærk, at fattige Folk ikke finde Rum i
Herberget. --- „En ung Moder er bleven forløst med en Søn, hun svøber ham og
lægger ham i en Krybbe.“ --- Forunderligt! --- „Er det da Noget at undre sig
over, at den stakkels Moder maa lægge Barnet i en Krybbe, naar der ikke er Rum i
Herberget?“ --- O hvilket Haab, o hvilken Glæde! --- „Er det da nu Noget at
glæde sig over, at en Moder i Folket er forløst med en Søn: den Fattiges Arving
pleier just ikke at vække nogen stor almindelig Glæde.“ --- Men bag hiin travle
Dag opgik den stille Nat. \emph{Det er dog et Under, det er dog en Glæde. See,
de himmelske Magter undre sig: dem er det et Saligheds Under. Guds Engle i
Himmelen glæde sig; de ville ei dølge deres Glæde. --- „Det maa forkyndes den
faldne Slægt, at Verdens Frelser er født.“ Men Menneskens Børn ændse det ei: de
sove og drømme --- om Dagens Møie. --- See, disse Hyrder vaage dog; de holde
Nattevagt ved deres Hjord. Det maa da forkyndes til Hyrdernes Trøst, at
Frelseren fødtes i denne Nat; det maa forkyndes de fattige Mænd, at Glæden
vederfares alt Folket.“
% Printer's defect: this second angelic speech is CLOSED but never OPENED ---
% there is no „ before „Det maa da forkyndes til Hyrdernes Trøst“ (verified at
% 600 dpi). The stray closing “ is transcribed as printed; it is the only
% unbalanced mark in pp. 173--183.
--- Og Hyrderne paa Marken høre Engelens Budskab, og --- Synet straaler i den
dybe Nat: en himmelsk Hærskares Mangfoldighed! Hellige Toner i den tanse Nat:
Ære være Gud i det Høie!}
% Printer's defect: „tanse“ for „tause“ --- a TURNED u, set as n. At 1200 dpi the
% glyph carries the n's right-hand tail and is identical to the n of „den“ two
% words earlier, and unlike the u of „undre“ higher on the page. Transcribed as
% printed.

\textbf{Varslerne komme og svinde; thi Himmeriges Rige er nær.}

\emph{„Og Fred paa Jorden?“} --- Paa den virkelige Jord er der ingen Fred. Hvad
er vel Jordens Fred? Et Solglimt paa en Uveirsdag! og der er slet ingen
\emph{„Velbehagelighed i Menneskene.“} Hvo kan nu finde Behag i en Keiser
August, der
% --- p. 182 ---
vil beskrive al Verden til Skat? Og \emph{„en himmelsk Hærskares
Mangfoldighed?“} Ja en velsignet frugtbar Overtroens Mangfoldighed. Opbyggelige
Historier af nogle Hyrder, der ikke vide, hvad de selv sige! --- „Et skinnende
Luftsyn; nogle Trækfugles Skrig, see det er nok det Hele.“

\textbf{Men Varslerne komme og svinde; thi Himmeriges Rige er nær.}

Tolv Aar hengaae; de svinde snart (og dog kan i den Tid Meget høres, og Meget
glemmes).

Det er Paaskehøitiden, der holdes i Jerusalem. Templet er omgivet af Haller. I
disse Haller have de Skriftkloge Skoler, som kaldes Tempelskoler. I en af disse
Tempelskoler træffer det sig nu, at et halvvoxent Barn træder ind midt iblandt
Lærerne, og hører dem, og begynder at adspørge dem. --- „Hvilke Spørgsmaal og
hvilke Svar af et tolvaars Barn?“ \emph{Alle, som høre Ham, forundre sig saare
paa hans Forstand og Gjensvar.} --- „Men hvo kan dog have tilskyndet dette Barn
til her at søge Lærere i Guds Huus?“ --- Man veed det ikke. Det er vel hans egen
Drift; han er nok gaaet bort fra sine Forældre. „Ogsaa dette bliver vel omsider
oplyst.“ --- Ja vist bliver det oplyst. See, der komme jo hans Forældre. „Hvor
Moderen har været bekymret for den kjære Søn! Hun henter ham med en mild
Bebreidelse.“ --- Det er en Tømmermands Søn fra Nazaret. --- Herlige Gaver! „Kom
han i de Lærdes Skole: det kunde blive en stor Rabbi.“

Men han kommer ikke i de Lærdes Skole. \emph{Nu drager han ned med sine
Forældre, og kommer til Nazaret, og er dem underdanig. Saa voxer han op, og
forfremmes i Viisdom og Alder og Naade hos Gud og Mennesker.} --- „Et smukt
Vidnesbyrd,“ sige de Lærde til alle
% „halvvoxent“ and „voxer“ carry a genuine Fraktur x on this page, checked at
% 1200 dpi against the r of the same lines; no r-for-x defect here.
% --- p. 183 ---
dem, de møde fra Nazaret, „et meget smukt Vidnesbyrd. Men vil man være Folkets
Lærer, maa man først selv lade sig belære; vil man virke paa Verden, maa man
kjende dens Tankegang og benytte de Kyndiges Veiledning.“ --- Nu er Han snart
tredive Aar. --- „Den bedste Læretid er altsaa forbi: han bliver dog aldrig
nogen sand Rabbi.“ --- „Nei, Menneskeheden veed nu lidt Mere, end denne
selvlærte Rabbi!“ gjentager Kritiken til Ære for den moderne
Bevidsthed\footnote{%
% Printed mark „*)“. The note is German set in the LIGHTER Fraktur this book
% uses for German, not in antiqua, and so stays plain --- except „Dr.“, which is
% genuine antiqua and is italicised.
Die geistig weltliche Bildung Jesu überschreitet dasjenige nicht, was bei
reichen Anlageu in den Volksschulen Galiläas und auf den Volksversammlungen zu
Jerusalem erworben werden konnte; andre, unter gebildeten Völkern und in
glücklicheren Zeiten geboren, waren ihm darin überlegen. \textit{Dr.} Karl Hase.
Kirchengeschichte. Leipz. 1834. S. 25 \S\ 83.%
% Printer's defect: „Anlageu“ for „Anlagen“ --- a TURNED n, set as u. At 1200 dpi
% the glyph is the round-footed u of „unter“, not the tailed n of „den“ two words
% later. Transcribed as printed.
}. Men Evangelisterne løfte Forhænget, pege paa Varslerne, og tale til os i
Aandens Sprog: \emph{Kommer og seer!}

Ja kommer og seer, hvi denne Lærer, denne Ene, kan lære Alle, uden selv at
„belæres“. Kommer og seer, hvi denne Mester, denne Ene, kan bevæge en Verdens
Tanker, uden at rette sig efter Verdens Tankegang. Hører det dog og forstaaer!
Naar denne Ene opløfter sin Røst, og begynder at tale, naar Han udrækker sin
Haand, og begynder at handle: da er Himmeriges Rige nær, da skeer der en sand,
en virkelig Begyndelse.

% CLOSE OF FØRSTE AFSNIT, at the foot of p. 183. There is NO three-asterisk
% ornament here and NO closing scripture in the display fount --- unlike the end
% of chapter II on p. 80. The only apparatus is the footnote rule and, well below
% the note, a single centred rule that closes the Afsnit; the rest of the page is
% blank. Measured at 600 dpi on PDF p. 200: the closing rule runs x 941--1401
% (461 px) against a body measure of 2102 px (left 129, right 2230) = 0.219,
% i.e. c. 19.5 mm, and is centred on the measure (rule centre 1171, body centre
% 1180). It is thicker than a hairline (c. 6 px solid at 600 dpi) and is probably
% the same thick-thin double rule used at pp. 42 and 80, where it measured 0.38;
% as elsewhere in this file it is reproduced as a single \rule. The footnote rule
% above it runs x 126--913 (788 px) = 0.375 of the measure.
\begin{center}
  \rule{0.22\textwidth}{0.4pt}
\end{center}

%% ============================================================
%%  ANDET AFSNIT  (printed pp. 184--530 = PDF 201--547)
%%
%%  Chapter ranges per the Indhold (PDF 16), under the standing head
%%  „Forløseren“:
%%    I.   Forgjængeren ................................... 184--201
%%    II.  Christus iblandt sine Egne ..................... 201--227
%%    III. Christus er et Tegn, som imodsiges ............. 228--289
%%    IV.  Christus aabenbarer Hjerternes Tanker .......... 289--334
%%    V.   Christus og Mørkets Aand ....................... 334--393
%%    VI.  Christus er mægtig i Ord og Gjerning ........... 393--458
%%    VII. Jesus Christus er Verdens Frelser .............. 458--530
%%
%%  NOT YET BEGUN. Markers will be added when this Afsnit is batched; the head
%%  below is set from the Indhold and must be checked against printed p. 184.
%% ============================================================

\end{document}
